\documentclass{article}
\usepackage[english]{babel}
\usepackage{geometry,amsmath,amssymb,stmaryrd,enumerate,hyperref,latexsym,xcolor,theorem,slashed,centernot,cancel}
\geometry{a4paper,
%total
left=25mm,
right=25mm,
top=25mm,
%bottom=20mm
}
\usepackage[T1]{fontenc}

%%%%%%%%%% Start TeXmacs macros
\catcode`\<=\active \def<{
\fontencoding{T1}\selectfont\symbol{60}\fontencoding{\encodingdefault}}
\catcode`\>=\active \def>{
\fontencoding{T1}\selectfont\symbol{62}\fontencoding{\encodingdefault}}
\newcommand{\Tau}{\mathrm{T}}
\newcommand{\assign}{:=}
\newcommand{\asterisk}{\mathord{*}}
\newcommand{\backassign}{=:}
\newcommand{\cdummy}{\cdot}
\newcommand{\llangle}{{\langle\!\langle}}
\newcommand{\mathd}{\mathrm{d}}
\newcommand{\mathi}{\mathrm{i}}
\newcommand{\nobracket}{}
\newcommand{\rrangle}{{\rangle\!\rangle}}
\newcommand{\tmcolor}[2]{{\color{#1}{#2}}}
\newcommand{\tmop}[1]{\ensuremath{\operatorname{#1}}}
\newcommand{\tmrsub}[1]{\ensuremath{_{\textrm{#1}}}}
\newcommand{\tmscript}[1]{\text{\scriptsize{$#1$}}}
\newcommand{\tmsep}{, }
\newcommand{\tmtextbf}[1]{{\bfseries{#1}}}
\newcommand{\tmtextit}[1]{{\itshape{#1}}}
\newenvironment{enumerateroman}{\begin{enumerate}[i.] }{\end{enumerate}}
\newenvironment{itemizedot}{\begin{itemize} \renewcommand{\labelitemi}{$\bullet$}\renewcommand{\labelitemii}{$\bullet$}\renewcommand{\labelitemiii}{$\bullet$}\renewcommand{\labelitemiv}{$\bullet$}}{\end{itemize}}
\newenvironment{proof}{\noindent\textbf{Proof\ }}{\hspace*{\fill}$\Box$\medskip}
\newenvironment{proof*}[1]{\noindent\textbf{#1\ }}{\hspace*{\fill}$\Box$\medskip}
\newtheorem{theorem}{Theorem}
\newtheorem{corollary}[theorem]{Corollary}
\newtheorem{definition}[theorem]{Definition}
\newtheorem{lemma}[theorem]{Lemma}
\newtheorem{proposition}[theorem]{Proposition}
{\theorembodyfont{\rmfamily}\newtheorem{remark}[theorem]{Remark}}
\newcounter{nnremark}
\def\thennremark{\unskip}
{\theorembodyfont{\rmfamily}\newtheorem{remark*}[nnremark]{Remark}}
\newcommand{\tmkeywords}{\textbf{Keywords:} }
\newcommand{\tmmsc}{\textbf{A.M.S. subject classification:} }
%%%%%%%%%% End TeXmacs macros
\newcommand{\tmaffiliation}[1]{\\ #1}

\newcommand{\slashnabla}{\slashed{\nabla}}
\newcommand{\slashk}{\centernot{k}}
\newcommand{\slashgen}[1]{\cancel{#1}}
\newcommand{\AMSclasstext}{MSC(2020)}

% replace \not with \centernot

\begin{document}

\title{
  Grassmannian stochastic analysis
  and the stochastic quantization
  of Euclidean Fermions
}



\author{Sergio Albeverio, Luigi Borasi,\\ Francesco C. De Vecchi, Massimiliano Gubinelli\\[.5em]
  \tmaffiliation{Hausdorff Center of Mathematics \&\\
  Institute of Applied Mathematics\\
  University of Bonn, Germany}\\[.2em]
  \tt{\small albeverio@iam.uni-bonn.de, borasi@iam.uni-bonn.de,}\\
  \tt{\small fdevecchi@.uni-bonn.de, gubinelli@iam.uni-bonn.de}
}

\date{January 2020}


\maketitle

\begin{abstract}
  We introduce a stochastic analysis of Grassmann random variables suitable
  for the stochastic quantization of Euclidean fermionic quantum field
  theories. Analysis on Grassmann algebras is developed here from the point of
  view of quantum probability: a Grassmann random variable is an homomorphism
  of an abstract Grassmann algebra into a quantum probability space, i.e. a
  $C^{\ast}$-algebra endowed with a suitable state. We define the notion of
  Gaussian processes, Brownian motion and stochastic (partial) differential
  equations taking values in Grassmann algebras. We use them to study the long
  time behavior of finite and infinite dimensional Langevin Grassmann
  stochastic differential equations driven by Gaussian space-time white noise
  and to describe their invariant measures. As an application we give a proof
  of the stochastic quantization and of the removal of the space cut-off for
  the Euclidean Yukawa model, indicating also how this method can be extended
  to other models of quantum fields.
\end{abstract}

\tmmsc{81S20}{\tmsep}{81T08}{\tmsep}{60H15}

\tmkeywords{Grassmann algebras, Euclidean Fermion fields, stochastic
quantization, non-commutative probability, non-commutative stochastic partial
differential equations, constructive quantum field theory, Yukawa model.}

{\tableofcontents}

\section{Introduction}

Euclidean quantum field theories
(QFTs)~{\cite{simon_po2_1974,glimm_quantum_1987,rivasseau_perturbative_1991}}
are a tool to construct and analyze mathematical models of relativistic
quantum fields, that is the quantum theory of elementary particle in
interaction satisfying the basic requirements of special relativity, i.e.
Poincar{\'e} covariance and
locality~{\cite{MR0277208,streater_pct_2000,jost_general_1965,bogolubov_general_1990,Haag1,Segalbook,MR1799198,strocchi_introduction_2013,friedrichs_mathematical_1953}}.
In the case of Bose--Einstein particles, the Euclidean theory is given by a
probability measure on Schwartz distributions over the Euclidean space
$\mathbb{R}^d$. In other words, bosonic Euclidean fields $\varphi$ are random
distributions~{\cite{symanzik,nelson_construction_1973,MR1340626}}. By a
well--known result of Osterwalder and Schrader~{\cite{glimm_quantum_1987}}
(see also~{\cite{zinoviev_equivalence_1995}}), a basic set of properties,
essentially Euclidean covariance, reflection positivity and regularity, is
sufficient to reconstruct a well-behaved relativistic QFT from these
probabilistic data.

In the case of Fermi--Dirac particles (e.g. the electron), the Euclidean
theory is expressed in terms of a linear functional over an infinite
dimensional Grassmann
algebra~{\cite{schwinger_euclidean_1958,schwinger_euclidean_1959,nakano_quantum_1959,osterwalder_euclidean_1973,osterwalder_euclidean_1973-2}}.
The value of this functional on Grassmann monomials gives the correlation
functions (Schwinger functions) of the Euclidean theory, which are represented
by Berezin integrals~{\cite{Berezin1,Berezin2}} of the type
\begin{equation}
  \langle O (\psi, \bar{\psi}) \rangle = \frac{\int \mathd \psi \mathd
  \bar{\psi} O (\psi, \bar{\psi}) e^{- S_E (\psi, \bar{\psi})}}{\int \mathd
  \psi \mathd \bar{\psi} e^{- S_E (\psi, \bar{\psi})}},
  \label{eq:eqft-berezin}
\end{equation}
where the fields $\psi, \bar{\psi}$ are the generators of the Grassmann
algebra, $O$ is a functional of $\psi$ and $\bar{\psi}$ generating the
Schwinger functions, and $S_E (\psi, \bar{\psi})$ is the Euclidean action
function, usually the sum of a quadratic part and a polynomial interaction
(possibly involving also bosonic fields, which then will have to be averaged
according to an appropriate probability
measure)~{\cite{feldman_renormalizable_1986,disertori_continuous_1998,mastropietro_non-perturbative_2008}}.
The formula~{\eqref{eq:eqft-berezin}} is an algebraic expression, devoid of
analytic meaning and is used as a convenient bookkeeping for various kinds of
expansions used to define and study the properties of the Schwinger
functions~{\cite{rivasseau_perturbative_1991,salmhofer_renormalization,mastropietro_non-perturbative_2008}}.
This is essentially due to the fact that Grassmann algebras do not have a
natural notion of positivity, nor a norm, associated with them. For early
discussions of related issues involving Euclidean fermions see
e.g.~{\cite{frohlich_is_1975,janson_gaussian_1997}}.

\

Stochastic quantization
(SQ)~{\cite{jona_lasinio_stochastic_1987,damgaard_stochastic_1988}} is an
approach to the construction of the correlation functions of an Euclidean QFT
introduced by Parisi and Wu~{\cite{parisi_perturbation_1981}}. The basic idea
is to consider an additional variable (usually a fictitious time) and
interpret the Euclidean fields $\varphi$ as the stationary solution to a
stochastic partial differential equations (SPDEs) involving this additional
variable and an external Gaussian source of noise $\xi$. This strategy has
been employed to rigorously study bosonic Euclidean theories starting with the
work by Jona-Lasinio and
Mitter~{\cite{jona_lasinio_stochastic_1985,jona_lasinio_large_1990,bertini_stochastic_1993}}
on the stochastic quantization of the $\Phi^4$ model for a scalar particle in
two (Euclidean) dimensions and with quartic interaction. More recently the
work of Hairer on regularity structures~{\cite{hairer_theory_2014}} opened the
way to the study of the three dimensional $\Phi^4$ model, see
also~{\cite{catellier_paracontrolled_2013,kupiainen_renormalization_2016}}.
While the original implementation of stochastic quantization gives rise to
parabolic SPDEs, variants can be constructed involving
elliptic~{\cite{albeverio_elliptic_2018,albeverio_elliptic_2019}} or
hyperbolic
equations~{\cite{gubinelli_renormalization_2018,gubinelli_paracontrolled_2018}}.
See the introductions to~{\cite{albeverio_invariant_2017}} and
to~{\cite{gubinelli_pde_2018}} for further references.

\

From physicists' point of view, stochastic quantization gives an alternative
approach to define and regularize quantum theories (especially theories with
gauge invariance)~{\cite{damgaard_stochastic_1988,bertini_stochastic_1993}}.
However, in the past few years, it has been realized that stochastic
quantization has also interesting properties from the mathematical point of
view. By solving the stochastic evolution it is possible to express the random
fields $\varphi$ of a (bosonic) Euclidean theory as well behaved functionals
of the external noise, i.e. $\varphi = \varphi (\xi)$, and this allows to
study the non-perturbative features of the former by leveraging the Gaussian
structure of the latter. Another conceptual advantage of stochastic
quantization is that it is essentially insensitive to questions of absolute
continuity with respect to a reference measure. On the other hand, lack of a
suitable reference measure is the main difficulty of the Gibbsian point of
view as expressed in equations of the form~{\eqref{eq:eqft-berezin}}. This new
perspective led to a series of results on the global space-time control of the
stochastic
dynamics~{\cite{MWcomedown,gubinelli_global_2019,albeverio_invariant_2017,moinat_space_time_2018}}
and to a novel proof of the constructions of non-Gaussian bosonic Euclidean
quantum field theories in three dimensions~{\cite{gubinelli_pde_2018}}.

\

The main motivation of the present work is to develop a stochastic
quantization of Euclidean fermionic quantum field theories (QFTs). In
consideration of the Grassmannian nature of the fermionic Euclidean fields as
expressed by~{\eqref{eq:eqft-berezin}}, this must involve stochastic partial
differential equations taking values in Grassmann algebras, including the
study of their long time behavior, invariant ``measures'' and regularity
properties. Additionally we aim for a theory which is both simple and
effective and which can be used in the construction and analysis of Euclidean
QFTs that involve fermions.

The main differences between the bosonic and the fermionic cases is the need
of extending the notion of stochastic processes to a non-commutative
framework. The approach we follow is to frame this problem in the context of
\tmtextit{non-commutative} probability (sometimes referred to as
\tmtextit{quantum} or \tmtextit{algebraic} probability)~
{\cite{parthasarathy_introduction_1992,biane_calcul_1995,Meyerbook,streater_classical_2000,accardi_outline_2015}},
namely we define the relevant Grassmannian objects as non-commutative random
variables.

A non-commutative probability space $(\mathcal{A}, \omega)$ is given by a
$C^{\ast}$-algebra $\mathcal{A}$ and a state $\omega$, which is a linear
normalized positive functional on $\mathcal{A}$. Inspired by the general
approach of Accardi et al.~{\cite{accardi_quantum_1982}} we will define random
variables with values in a Grassmann algebra $\Lambda$ as homomorphisms of
$\Lambda$ into $\mathcal{A}$. We will not require the homomorphisms to respect
the natural involution present in any $C^{\ast}$-algebra, since there is no
canonical candidate for that in Grassmann algebras.

It is useful to keep in mind how classical commutative random variables fit
in this algebraic approach. Consider a random variable $X : \Omega \rightarrow
\mathcal{M}$ on a probability space $(\Omega, \mathcal{F}, \mathbb{P})$,
taking values in a manifold $\mathcal{M}$ and let $\mathcal{B}$ be the algebra
of all measurable, complex bounded functions on $\mathcal{M}$. All the
relevant probabilistic information about $X$ is encoded in the homomorphisms
of algebras $f \in \mathcal{B} \mapsto f (X) \in \mathcal{A}$ where
$\mathcal{A}$ is the algebra of bounded measurable complex random variables on
the basic space $\Omega$ endowed with the linear functional given by the
expectation $\mathbb{E}$ associated to the probability measure $\mathbb{P}$
defined on a $\sigma$-algebra $\mathcal{F}$. In particular, this is the point
of view with which SDEs on manifolds are defined~{\cite{hsu_stochastic_2002}}.

\

The embedding of $\Lambda$ into $\mathcal{A}$ allows to use the topology of
$\mathcal{A}$ to do analysis on Grassmann algebras. Of course the associated
analysis will not be canonical from the point of view of the Grassmann algebra
itself, but it will turn out to be powerful enough to allow us to obtain a
satisfactory theory of stochastic quantization for fermions. Let us mention
another analogy which can help the reader to understand this point of view.
The standard approach to study Gaussian processes in Hilbert
spaces~{\cite{Gross1967,janson_gaussian_1997,hida_white_2013}} requires to
have a non-canonical embedding of the Hilbert space into a larger Banach space
$B$ which supports the Gaussian measure. This Banach space is not canonical
and there are various possible choices for any given Hilbert space $H$. For
example Brownian motion is associated with the Hilbert space $H$ of \
functions on $\mathbb{R}_+$ with square integrable derivative, but Wiener
measure is supported on the Banach space $C^{\gamma}$ of $\gamma$ H{\"o}lder
functions for any $\gamma < 1 / 2$. In the case of random variables taking
values in a Grassmann algebra $\Lambda$, the role of this ``bigger'' Banach
space is played by a suitable Clifford sub-algebra of $\mathcal{A}$ which
provides us with enough Grassmann generators to realize our homomorphisms and
which has a canonical notion of norm. It will be for us just ``a convenient
place where to hang our (analytic) hat on''.

\

The main finding of our work is that, once one accepts this point of view,
the rest of the analysis falls in place quite naturally. Objects like
Grassmann white noise, Brownian motion and free fields are relatively simple
to define and control. On top of them a novel theory of stochastic
differential equations (SDEs) with values in Grassmann algebras can be
initiated and carried to the point to be able to discuss their long time
evolution and invariant ``measures''. Moreover many considerations carry over
to SDE in infinite dimensions and allow to define and solve, in particular,
the non-commutative nonlinear stochastic partial differential equations
(SPDEs) appearing in stochastic quantization, again in relatively simple
terms. \ As a test of the effectiveness of our approach we give a proof of
existence of the infinite volume limit of the massive Euclidean Yukawa model
for small
coupling~{\cite{osterwalder_euclidean_1973,osterwalder_euclidean_1973-2,lesniewski_effective_1987}}
(in Appendix \ref{appendix_models} it is shown how the results of this paper,
in particular in Section \ref{sec:yukawa}, hold also for several other models
of quantum fields involving fermions, including \ quantum electrodynamics in
two dimensions with a positive photonic mass, pseudo-scalar Yukawa model, as
well as the Gross--Neveu model in two dimensions and some simplified
supersymmetric models ,including the Wess--Zumino model). At the end, we find
that analysis of Grassmann SPDEs is in some respect \tmtextit{much simpler}
than its commutative counterpart: the discussion of certain parts of the
theory, like the existence of global space-time dynamics, are made relatively
trivial by the boundedness of Grassmann Gaussian variables.

\

As we already remarked, our approach follows Osterwalder and
Schrader~{\cite{osterwalder_euclidean_1973,osterwalder_euclidean_1973-2}} who
were the first to do rigorous analysis of Euclidean Fermions. See the review
of Palmer~{\cite{palmer_euclidean_1980}} for a relatively recent account of
their theory. They were using the language of Fock space operators and not of
Berezin calculus, and this allowed them to consider rigorously expressions
involving infinite dimensional Grassmann algebras. The Fock space approach has
also been recently used to discuss supersymmetry in the context of
Parisi--Sourlas dimensional
reduction~{\cite{klein_supersymmetry_1983,klein_supersymmetry_1985}} and
Euclidean stochastic
quantization~{\cite{albeverio_elliptic_2018,albeverio_elliptic_2019,de_vecchi_note_2019}}.
The results of these papers could also be framed in the language of Grassmann
random variables introduced in the present work.

Our work provides also non-trivial examples of non-commutative non-linear
SPDEs and of their qualitative analysis. Analysis of partial differential
equations in non-commutative algebras is not well developed in general.
Minkowski QFTs provide examples of PDEs in non-commutative
algebras~{\cite{Segalbook}}. Other relevant examples we are aware of are the
work of Rosenberg~{\cite{rosenberg_noncommutative_2008}} on the (linear)
Laplace equation on the non-commutative torus and the subsequent work on
linear operators and functional spaces on non-commutative
spaces~(e.g.~{\cite{savin_noncommutative_2010,xiong_sobolev_2018}}), works on
the algebraic and geometric features in PDEs (see
e.g.~{\cite{prastaro_geometry_1994,connes_noncommutative_1994}}), the work of
Khrennikov on linear equations and differential operators on
superspace~{\cite{khrennikov_equations_1991}} and the work of
Osipov~{\cite{osipov_quantum_2000}} on solutions of quantum field equations
via Wick kernels. However none of these works concern SPDEs and all rather use
indirect methods to find solutions.

We have left open the problem of removing the small-scale regularization in
the stochastic quantization of the Yukawa model. In order to make progress in
this direction one will have to understand the renormalization problem for
such singular SPDEs in our setting, along the line of the recent work of
Hairer and
others~{\cite{hairer_theory_2014,gubinelli_paracontrolled_2015,kupiainen_renormalization_2016}}.
Apart from the more delicate analysis required by the low regularity of the
random fields, the removal of these divergences will make appear renormalized
Wick products of Euclidean fermion fields. These operators are unbounded in
general and their appearance make more delicate to close the estimates, but we
are confident that this can be achieved in our framework.

It is our hope that this study of Grassmann stochastic equations could pave
the way to a deeper understanding of non-commutative stochastic analysis and
PDE theory in general.

\

{\noindent}\tmtextbf{Historical remarks.}

\

The rest of this introduction (before we describe the structure of our work)
is dedicated to discuss the connections of our work with previous research.

The interest in studying time evolutions and stochastic processes with values
in non commutative algebras arose in connection with the development of
quantum mechanics in the 20s--30s of last century and later in quantum field
theory, particularly around the 50s. In these theories one associates physical
observables with operators evolving in time and acting on a, in general
infinite dimensional, complex Hilbert space. Mathematically, this lead to the
representation theory of certain operator algebras (CCR-algebras resp.
CAR-algebras in the case of fermions).

In this line, quantum fields require the study of certain infinite dimensional
(non-commutative) algebras, the study of which has been particularly
stimulated by developments in the period 1940--1960, where the mathematical
description of systems with an infinite number of degrees of freedom, in
particular for quantum fields and statistical mechanics has been worked out.
Let us mention the first works on von~Neumann algebras by Murray and
von~Neumann~(see~e.g.~{\cite{VonNeumann}}) and the work on
$C^{\asterisk}$-algebras developed originally and quite independently by
I.~M.~Gelfand and his school (see, e.g.,~{\cite{Gelfand}}), since the
beginnings of the 40s, and by I.~E.~Segal and his school
(see~e.g.~{\cite{Segal1,Segalbook}}). Both types of algebras are associative
and equipped with an involution (the adjoint operation). Whereas
$C^{\ast}$-algebras are looked upon as abstract normed algebras, without
necessarily being realized as algebras of operators acting in Hilbert spaces,
the study of von Neumann algebras (also called $W^{\ast}$-algebras) rather
involves looking at them as self-adjoint algebras of bounded operators in
Hilbert spaces, closed in the weak operator topology. Both algebras can be
looked upon as particular cases of normed algebras, that in the case of
closedness are called Banach algebras. More precisely, $C^{\ast}$-algebras are
normed algebras with a unit $\mathbb{I}$ and an involution $a \mapsto
a^{\ast}$, any element $a$ of which satisfies $\| a a^{\ast} \| = \| a \|^2$.
Although in principle every von~Neumann algebra can be looked upon as a
$C^{\ast}$-algebra, structural and conceptual developments in the study of von
Neumann algebras involve particularly their realizations in Hilbert space. The
mathematical literature on normed, Banach, $C^{\ast}$ and $W^{\ast}$ algebras
is very rich and there are a number of excellent monographs on them, see e.g.,
Naimark~{\cite{Naimark}} (see also, e.g. {\cite{bonsall_complete_1973}}), and
Zelazko~{\cite{MR0448079}} for normed spaces and algebras, and Kadison and
Ringrose~{\cite{Kadison1,Kadison2}} (and references therein) for the other
type of algebras. For applications in mathematical physics, in particular
quantum mechanics, quantum field theory and statistical mechanics, there are
monographs by Bratteli and Robinson~{\cite{Bratteli1,Bratteli2}},
Emch~{\cite{Emch}}, Baumg{\"a}rtel~{\cite{Baumgartel1}}, as well as
Baumg{\"a}rtel and Wollenberg~{\cite{Baumgartel2}}, Holevo {\cite{Holevo}},
Segal~{\cite{Segalbook2}},~Schm{\"u}dgen~{\cite{schmudgen_unbounded_1990}};
for relativistic quantum fields, we refer particularly to the books by Araki
{\cite{MR1799198}}, by Baez, Segal and Zhou~{\cite{Segalbook}} and
Haag~{\cite{Haag1}}. The articles by Summers~{\cite{summers2012perspective}}
and Doplicher~{\cite{Doplicher}} survey particularly some recent developments
related to application of such algebras in quantum field theory.

Osterwalder and
Schrader~{\cite{osterwalder_euclidean_1973,osterwalder_euclidean_1973-2}}
discovered that Euclidean fermions satisfy the relations defining a Grassmann
algebra. However see the papers of Schwinger and Nakano for the first
discussion of Wick rotation for
fermions~{\cite{schwinger_euclidean_1958,schwinger_euclidean_1959}} and also
the
papers~{\cite{van_nieuwenhuizen_euclidean_1996,williams_euclidean_1974,MR0452281,lehmann_probabilistic_1991,borasi}}.
Grassmann algebras constitute a class of associative non commutative algebras
that were introduced by Grassmann in the second half of the 19th century. They
are connected with Clifford algebras, that in turn are the underlying
structure to the study of CCR and CAR algebras of quantum mechanics. All such
algebras are nowadays part of the general theory of algebra as pioneered by
E.~Noether in the 20s-40s, see e.g.
Van~der~Waerden~{\cite{van2003algebra1,van2003algebra2}},
Corry~{\cite{corry_modern_2004}}.

While the original paper of Osterwalder and
Schrader~{\cite{osterwalder_euclidean_1973-2}} uses the same point of view
which is developed here, the majority of the subsequent work with Euclidean
fermions involved the formalism of Berezin's integral~{\cite{Berezin2}}. This
formalism has been subsequently developed into an analysis on Grassmann
algebras, in the form of super-analysis~{\cite{Berezin2}}. For super-algebras
and super-analysis see also, e.g., Rogers~{\cite{Rogers2007}} and
Leites~{\cite{Leites}}, on a more heuristic level also DeWitt~{\cite{DeWitt}}.
For a review see Pestov~{\cite{pestov_analysis_1994}}. Berezin integrals have
been useful in calculating fermionic systems, see e.g.
Zinn-Justin~{\cite{ZinnJustin}}, Izkyson and Zuber~{\cite{Itzykson}}. Such
methods are also quite essential in most of the mathematical literature on
fermionic quantum fields, including renormalization theory, see, e.g.,
Feldman, Kn{\"o}rrer and Trubowitz~{\cite{MR1905424}},
Salmhofer~{\cite{salmhofer_renormalization}}. Euclidean quantum field theories
involving fermions have surprisingly simpler behavior than the bosonic ones
and have been very much studied also from the rigorous point of
view~{\cite{gawedzki_gross_neveu_1985,feldman_renormalizable_1986,lesniewski_effective_1987,rivasseau_perturbative_1991,disertori_continuous_1998,benfatto_functional_2007,mastropietro_non-perturbative_2008}},
mainly in the formalism of Berezin integration. In this context let us mention
also the probabilistic representation introduced by de~Angelis et
al.~{\cite{angelis_berezin_1998}} for finite-dimensional Grassmann-valued
Markov processes using Poisson processes.

\

After Parisi and Wu seminal paper, the stochastic quantization of fermions
has been discussed heuristically in the physics literature, starting with
Kakudo et al.~{\cite{kakudo1983gauge}}, Fukai et
al.~{\cite{fukai_stochastic_1983}}, Damgaard and
Tsokos~{\cite{damgaard_stochastic_1984}} and She--Sheng and
Ting--Chang~{\cite{she-sheng_stochastic_1985}}, \ up to the recent paper of
Efremov~{\cite{efremov_stochastic_2019}}. These authors stress particularly
the relations with ``Grassmann valued random variables'', solving formally
partial differential equations with Grassmann valued Gaussian white noise. The
invariant measures of physical interest are here described by averages with
respect to Berezin integrals. However all these papers never properly discuss
the analytical difficulties of considering (stochastic) differential equations
in Grassmann algebras and the related stochastic analysis.

From the point of view of probability theory, pioneering work in the theory of
non-commutative processes has been done by Accardi, Streater and Hudson and
Parthasarathy and others in the 70s--80s, see
e.g.~{\cite{accardi_outline_2015,parthasarathy_introduction_1992,streater_classical_2000,biane_calcul_1995}}.
This includes the study of non-commutative Markov semigroups, non-commutative
Brownian motion and stochastic analysis based on on it, in particular also
SDEs but with no space variable dependence, hence no SPDEs. Let us also
mention the work of
Gross~{\cite{gross_noncommutative_1971,gross_existence_1972,gross_hypercontractivity_1975,gross_formula_1977}}
applying Segal's non-commutative integration
theory~{\cite{segal_non_commutative_1953,nelson_notes_1974}} to Clifford
algebras. Work extending SDEs with values in algebras other than Grassmann
algebras, like $C^{\ast}$ and von Neumann algebras are by authors like
Streater and coworkers~{\cite{Streater1,Streater2,Streater3}}, Applebaum and
Hudson~{\cite{Applebaum1,Applebaum2}}, Belavkin~{\cite{Belavkin}}, Carlen and
Kree~{\cite{Carlen1}}, Gordina~{\cite{Gordina}}, Sinha and
Goswami~{\cite{Sinha}}, K{\"u}mmer {\cite{Kummer}}, see also the monographs by
Meyer~{\cite{Meyerbook}}, Holevo~{\cite{Holevo}}, and the recent survey
article by Cipriani~{\cite{Cipriani}} and the references therein.

Grassmann variables are in a certain sense \tmtextit{completely
non-commutative numbers}: as such they appear also in non-commutative central
limit theorems~{\cite{von_waldenfels_non_commutative_1986}} and in the setting
of non-commutative processes with independent
increments~{\cite{schurman_quantum_1991,bozejko_interpolations_1996}}. A
rigorous approach to infinite dimensional Grassmann algebras and the related
stochastic analysis is that of
Rogers~{\cite{rogers_graded_1986,rogers_stochastic_1992_I,rogers_stochastic_1992_II,rogers_path_1995,leppard_feynman_kac_2001}}
which essentially consists in looking at all the finite dimensional
sub-algebras, use the Berezin integral to compute averages and require certain
natural consistency conditions in order to obtain a projective system. A
similar line of research was also carried on by Kupsch and
Haba~{\cite{kupsch_fermionic_1993,haba_supersymmetry_1995}}. These last works
follows the observation of Hudson and
Parthasarathy~{\cite{hudson_unification_1986}} and
Le~Jan~{\cite{le_jan_temps_1987}} (see also Meyer's book~{\cite{Meyerbook}})
which developed a unified representation of the Bosonic and Fermionic Brownian
motion (and related stochastic calculus) using Fock space techniques. None of
these lines of research seems to have reached a stage where the theory is
powerful enough to easily and naturally accommodate SPDEs (or even SDEs) of
the kind needed in stochastic quantization. One of the technical problem
involved is that (infinite dimensional) Grassmann algebras, as we already
remarked, do not come with natural analytic structured suitable for non-linear
analysis. Recent discussions of this problem can be found in works of
Ivashchuk~{\cite{ivashchuk_infinite_dimensional_2000}} and Alpay
et~al.~{\cite{alpay_distribution_2019}}.

\

More successful rigorous attempts to define Grassmann stochastic quantization
were undertaken in a series of papers by Scherbakov et
al.~{\cite{shcherbakov_elements_1993,shcherbakov_elements_1993-1,shcherbakov_elements_1993-2}}
using the locally convex topology obtained by considering all the possible
correlation functions of the fields and the noise as the family of semi-norms.
Convergence in this setting therefore corresponds to convergence of all the
correlation functions. This approach has the disadvantage of not allowing the
introduction of norms strong enough to study differential equations via
standard approaches. Despite this technical shortcoming, which makes all the
proofs quite involuted, it should be said that these papers define and
successfully solve equations with finitely many degrees of freedom (i.e. with
small scale and large scale cut-offs) and then address the convergence when
the large scale cut-off is removed (infinite volume limit) using cluster
expansion techniques. Other results in this direction were obtained later on
by Ignatyuk et al.~{\cite{ignatyuk_convergence_1993}}, with applications in
statistical mechanics.

\

The main aim of our work is to provide a framework suitable to discuss
singular Grassmann-valued PDEs and in particular SPDEs. Besides proving
methods for existence, uniqueness of local and global solutions, and of
invariant measures, that can in principle be used for several applications, we
go in details in one of the simplest but physically relevant models of quantum
fields involving fermions. This is the Yukawa model, that couples a scalar (or
pseudo-scalar) massive field with a pair of Fermi fields. Historically it was
introduced by Yukawa in 1935 and further studied
e.g.~in~{\cite{matthews_propagators_1955}}. Its mathematical construction was
achieved, in space-time dimension $2$, as part of the constructive approach to
relativistic quantum fields, see
e.g.~{\cite{lanford_iii_construction_1966,mcbryan_volume_1975,glimm_quantum_1985,dimock_estimates_1972,McBryan1975,Magnen1976,Magnen1978,Federbush1971,Jaffe1971}}.
The Euclidean approach was developed
in~{\cite{osterwalder_euclidean_1973-2,seiler_schwinger_1975,seiler_bounds_1975,Gross_Euclidean_1976,williams_euclidean_1974}}.
Our work show that the Euclidean model can be obtained independently as a
stationary state from a Grassmann-valued stochastic quantization equation. It
is expected that also the high energy regularization can be treated in this
framework.

\

{\noindent}\tmtextbf{Structure of the paper.} In
Section~\ref{sec:grassmann-rv} we first define Grassmann random variables and
present various related notions like compatibility and independence. A
distance and its corresponding topology on a space of Grassmann-valued random
variables is introduced in Section~\ref{sec:prob} and Section~\ref{sec:top}
and its completion is taken. Moreover, functions on the Grassmann algebra
$\Lambda V$ over a pre-Hilbert space $V$ are introduced together with their
right, resp. left, derivative as maps from $\Lambda V$ to $\Lambda V \otimes
V$ (resp. $V \otimes \Lambda V$), and a Taylor formula is proven. Then basic
notions needed for formulating SDEs for Grassmann-valued processes are
introduced. \ In particular, in Section~\ref{sec:gaussians} Grassmann Gaussian
variables are discussed and constructed together with corresponding notions of
Gaussian white noise and Brownian motions.

Section~\ref{sec:finite-dim} studies finite dimensional Grassmann
differential equations with additive Gaussian white noise. We first study the
linear case (Grassmann Ornstein--Uhlenbeck motion) and exhibit an explicit
formula expressing its solution, in particular its stationary solution, that
remains bounded when the linear drift is given by a strictly negative operator
(Proposition~\ref{proposition_linearlimit}). In Theorem~\ref{theorem_picard}
and Theorem~\ref{th:global-sol} we present a unique solution for the
non-linear case given by an odd polynomial. Here we use an iteration of
Picard's type and fix-point arguments in a Banach space. Global existence is
established, in our setting, essentially by means of a Gronwall inequality. An
(average) It{\^o} formula for solutions of finite dimensional Grassmann SDEs
(Theorem~\ref{th:ito}) is proven in details. In Section~\ref{sec:inv-meas}
invariant measures are defined for the general case. For this we prove a
criterium that can be looked upon as a Grassmannian analogue of the concept of
infinitesimal invariance of measures (expressed via a Fokker--Planck equation)
in the commutative case (Lemma~\ref{lemma_invariant}). The existence of an
invariant measure is proven in Theorem~\ref{th:sq-finite}, in the case of
total drift consisting of a linear term plus a non-linear odd-degree term. The
long time behavior is studied in Section~\ref{sec:long-time}. Here the ``lack
of positivity'' in Grassmann algebras is bypassed by assuming a smallness
condition on the non-linearity and strict negativity of the linear part of the
SDE. The main results here are Theorem~\ref{theorem_approximation}, giving the
convergence of an approximation to the unique stationary solution, and
Theorem~\ref{theorem_invariant2}, that exhibits the invariant measure (state).

Section~\ref{sec:infinite-dim} is devoted to extend the previous results to
SDE in infinite dimensions of the kind relevant for stochastic quantization.
This allows us to use these processes to represent the correlation functions
of Euclidean fermionic theories like those described
by~{\eqref{eq:eqft-berezin}}. To achieve these results, the infinite
dimensional setting is made precise, by taking as underlying space for the
generalization of the Grassmann algebra, a suitable infinite dimensional
Fr{\'e}chet nuclear space $V$. The main problem to overcome is the choice of a
suitable space $\mathcal{G} (V)$ of homomorphisms for the Grassmann algebra
over $V$ to a $C^{\ast}$-algebra $\mathcal{A}$. We present a choice in terms
of operator-valued Besov--H{\"o}lder spaces $C^s (\mathbb{R}^d, \mathcal{A})$.
The appropriateness of this setting, permitting to extend finite dimensional
methods to the infinite dimensional case, is demonstrated by
Theorem~\ref{theorem_existence2} and Theorem~\ref{theorem_longtime2}, that
essentially show that smallness of the non-linear part and a generator of a
strictly contractive semigroup on $C^s (\mathbb{R}^d)$ for the linear part,
suffice for having a unique solution of the stochastic Grassmann infinite
dimensional equation for all times, as well as its convergence to a stationary
solution for long times.

In Section~\ref{sec:yukawa}, as an illustration of the use of our results for
Grassmann SPDEs, we present an application to the study of the Yukawa model of
Euclidean quantum fields. In Section~\ref{sec:os-fields} the Euclidean
construction of the model, in two Euclidean dimensions and with ultraviolet
(UV) and infrared (IR) cut-offs, is presented in our setting, following the
construction in~{\cite{osterwalder_euclidean_1973-2}}. In
Section~\ref{sec:yukawa-sq} the corresponding stochastic quantization equation
(SQE) is formulated and studied. Two main theorems are formulated and proved
here. In the first one (Theorem~\ref{theorem_main1}) we obtain the
representation of the relevant correlation functions for solutions of the SQE
as averages with respect to the probability state, for small values of the
coupling constant $\lambda$ in the non-linear part of the equation. The second
theorem, Theorem~\ref{theorem_main2}, achieves the removal of the space
cut-off (infrared, or infinite volume limit, required for the model to be
euclidean invariant). The proof is quite technical and on its way several
results are established that have an interest also for other applications. In
Section~\ref{sec:finit-dim-yukawa} finite dimensional approximations are
constructed, involving the introduction of the model first on a lattice
approximation of a finite torus and then passing to the full torus. The
passage to the full space is obtained by introducing weighted Besov spaces.
The convergence result is presented in Theorem~\ref{theorem_torusplane}. This
then yields, in particular, a construction of the fermionic sector of the
massive Yukawa Euclidean quantum field model in two dimensions, in the weak
coupling region and with an ultraviolet cutoff, by our method of
Grassmann-valued SPDEs.

In Appendix~\ref{app:GP} we briefly describe the tracial representation of
Euclidean Fermi fields of Gross~{\cite{Gross_Euclidean_1976}} and
Palmer~{\cite{palmer_euclidean_1980}}, its relation with our description of
Gaussian Grassmann random variables (see Section \ref{sec:gaussians}) \ and
another formulation based on Segal's ``real wave''
representation~{\cite{Segalbook}}. In Appendix~\ref{app:series} we include
some side results about the convergence of series expansions for the solutions
to finite dimensional Grassmann SDEs, while In Appendix~\ref{appendix_Besov}
we collect various technical results about functional spaces of Banach-valued
functions. Appendix \ref{appendix:convergenceinfinitedim} contains the proof
of convergence of perturbative series in the infinite dimensional case of
Yukawa$_2$ model and, finally, in Appendix~\ref{appendix_models} we review
some further quantum field theory models that can be approached using the
ideas presented in the present paper.

\

{\noindent}\tmtextbf{Acknowledgements.} This work is supported by DFG via the
grant AL 214/50-1 \ ``Invariant measures for SPDEs and Asymptotics''. The
authors are funded by the DFG under Germany's Excellence Strategy - GZ 2047/1,
project-id 390685813. M.~G. would like to thank Hao Shen for comments on an
earlier version of the paper.

\section{Grassmann random variables}\label{sec:grassmann-rv}

We denote by $\mathcal{L} (A, B)$ the space of linear maps between vector
spaces $A$ and $B$, if both $A, B$ have topologies we consider all the maps to
be continuous. We let $\mathcal{L} (A) =\mathcal{L} (A, A)$. With $\tmop{Hom}
(A, B)$ we denote homomorphisms between algebras $A$ and $B$.

\subsection{Grassmann probability}\label{sec:prob}

\tmtextit{Algebraic probability.} We consider a complex Hilbert space
$\mathcal{H}$ (to be fixed later) and denote by $\mathcal{A}=\mathcal{L}
(\mathcal{H})$ the $C^{\ast}$-algebra of bounded operators with the operator
norm. Moreover we assume to have a state $\omega$ on $\mathcal{A}$, that is a
positive definite functional normalized by the condition that $\omega
(\mathbb{I}) = 1$, where $\mathbb{I}$ is the identity in $\mathcal{L}
(\mathcal{H})$. The pair $(\mathcal{A}, \omega)$ is a (non-commutative or
algebraic) probability space. We do not require the state $\omega$ to be
either faithful, or tracial.

\

\tmtextit{Grassmann algebras.} Let $V$ be a (finite dimensional or infinite
dimensional and separable) real Hilbert space. Denote by $\Lambda V$ the
Grassmann algebra generated by $V$, i.e. the exterior algebra obtained by
quotienting the tensor algebra $T (V) = \oplus_{n \geqslant 0} V^{\otimes n}$
by the two-sided ideal generated by elements of the form $x \otimes x$ for $x
\in V$. We denote the product of two elements $f, g \in \Lambda V$ by $f
\wedge g$ (or if there is no confusion with other products simply by $f g$).

\

Let $(v_{\alpha})_{\alpha \in A}$ be a fixed basis of $V$. As a consequence
$\Lambda V$ is spanned by elements of the form $v_A \assign v_{\alpha_1}
\wedge \cdots \wedge v_{\alpha_n}$ and $A = (\alpha_1, \ldots, \alpha_n)$ is a
$n$-tuple of elements of $A$ with the convention that $v_{\emptyset} = 1$ is
the unit element in $\Lambda V$. When it is clear from the context we will
denote the product in $\Lambda V$ simply by $v_A = v_{\alpha_1} \ldots
v_{\alpha_n}$. We recall that $\Lambda V$ is $\mathbb{Z}_2$ graded in the
sense that it splits into odd and even parts $\Lambda V = \Lambda_{\tmop{odd}}
V \oplus \Lambda_{\tmop{even}} V$. On $\Lambda V$ there is an super Hopf
algebra structure with coproduct $\Delta : \Lambda V \rightarrow \Lambda V
\otimes \Lambda V$ where $\Lambda V \otimes \Lambda V$ is the
$\mathbb{Z}_2$-graded tensor product algebra which satisfies
\[ (f \otimes g) (h \otimes k) = (- 1)^{| g |  | h |} f h \otimes g k, \qquad
   f, g, h, k \in \Lambda V, \]
with $| \cdot | : \Lambda V \rightarrow \{ 0, 1 \}$ the even/odd grading. The
coproduct $\Delta$ is the algebra homomorphism such that $\Delta v = 1 \otimes
v + v \otimes 1$ for all $v \in V \subseteq \Lambda V$ and the counit
$\varepsilon : \Lambda V \rightarrow \mathbb{C}$ given by $\varepsilon (v_A)
=\mathbb{I}_{A = \emptyset}$.

\

\tmtextit{Random variables.}

\begin{definition}
  \label{def:Grassmann-rv}A $V$-Grassmann random variable $\Psi$ is an algebra
  homomorphism from the Grassmann algebra $\Lambda V$ into $\mathcal{A}$. We
  denote by $\mathcal{G} (V) = \tmop{Hom} (\Lambda V, \mathcal{A})$ the set of
  all such homomorphisms. We call the \tmtextit{law} of $\Psi \in \mathcal{G}
  (V)$ the family of its moments $\omega^{\Psi} (F) : = \omega (\Psi (F))$ for
  all $F \in \Lambda V$, also represented by the linear functional
  $\omega^{\Psi} : \Lambda V \rightarrow \mathbb{R}$.
\end{definition}

Note that $\Psi \in \mathcal{G} (V)$ cannot be assumed to be a
$\ast$-algebra-homomorphism since $\Lambda V$ has no (natural)
$\ast$-operation. If $F \in \Lambda V$ has representation $F = \sum_A F_A v_A$
we shall employ the following dual notation
\[ F (\Psi) \assign \Psi (F) = \sum_A F_A \Psi^A, \]
where $\Psi^A \assign \Psi (v_A)$ and $F_A \in \mathbb{R}$. Since $\Psi$ is
assumed to be an algebra homomorphism, we have e.g. $\Psi^{\alpha}
\Psi^{\beta} = - \Psi^{\beta} \Psi^{\alpha}$, where $\Psi^{\alpha} = \Psi
(v_{\alpha})$ and $(v_{\alpha})_{\alpha}$ is a fixed basis of $V$. Moreover,
$\Psi (\Lambda V)$ is a Grassmann sub-algebra of $\mathcal{A}$ and $\Psi (F)$
has the same degree (even or odd) as $F$. As shown
in~{\cite{van2003algebra1,van2003algebra2}}, even if some arguments are
formulated in a basis dependent way, the definition of $\Lambda V$ and its
characterization by anti-commutation relations is independent of the basis.

\

\tmtextit{Compatibility.} When the context is clear we will abbreviate $X \in
\mathcal{G} (V)$ as $X \in \mathcal{G}$. Let $X \in \mathcal{G} (V)$ and $Y
\in \mathcal{G} (W)$, we say that they are \tmtextit{compatible} if the linear
map $Z : V \oplus W \rightarrow \mathcal{A}$ given by $Z (v) = X (v)$ if $v
\in V$ and $Z (w) = Y (w)$ if $w \in W$, extends to a homomorphism $Z :
\Lambda (V \oplus W) \rightarrow \mathcal{A}$. In this case we write $Z \in
\mathcal{G} (V \oplus W)$ or briefly $(X, Y) \in \mathcal{G}$. Note that we
have a super- algebra isomorphism $\Lambda (V \oplus W) \approx \Lambda (V)
\otimes \Lambda (W)$ where tensor product is in the sense of super-algebras.
From this isomorphism we get in particular $Z (F) = m_{\mathcal{A}} [(X
\otimes Y) (F)]$ for all $F \in \Lambda (V \oplus W)$ where $m_{\mathcal{A}} :
\mathcal{A} \otimes \mathcal{A} \rightarrow \mathcal{A}$ denotes the
multiplication of $\mathcal{A}$. Compatibility can of course be defined for
$X_1, \ldots, X_n \in \mathcal{G}$ in a corresponding way. We shall express
that $X_1, \ldots, X_n \in \mathcal{G}$ are compatible by writing $(X_1,
\ldots, X_n) \in \mathcal{G}$.

\begin{remark}
  The notion of compatibility which we introduce here is not standard in
  algebraic probability. A related concept is that of kinematic independence,
  see e.g. the review of Accardi~{\cite{accardi_outline_2015}}.
\end{remark}

Given $X, Y \in \mathcal{G} (V)$ which are compatible we define $X + Y \in
\mathcal{G} (V)$ as
\[ (X + Y) (F) = F (X + Y) \assign m_{\mathcal{A}}^2 [(X \otimes Y) \Delta F],
   \qquad F \in \Lambda V, \]
where $m_{\mathcal{A}}^k : \mathcal{A} \otimes \mathcal{A} \rightarrow
\mathcal{A}$, $k \in \mathbb{N},$ is defined as
\begin{equation}
  m_{\mathcal{A}}^k (a_1 \otimes \cdots \otimes a_k) = a_1 \cdots a_k
  \label{eq:productA} .
\end{equation}
Note indeed that by compatibility $Z = m_{\mathcal{A}}^2 \circ (X \otimes Y) :
\Lambda (V \oplus V) \approx \Lambda V \otimes \Lambda V \rightarrow
\mathcal{A}$ is an algebra homomorphism and therefore \
\[ (X + Y) (F G) = Z (\Delta (F G)) = Z (\Delta F \Delta G) = Z (\Delta F) Z
   (\Delta G) = (X + Y) (F) (X + Y) (G) . \]
The notation is justified from the fact that $(X + Y) (v) = X (v) + Y (v)$ for
all $v \in V$. Similarly for any $\lambda \in \mathbb{C}$ we can define
$\lambda X$ as the only homomorphism such that $(\lambda X) (v) = \lambda X
(v)$ for all $v \in V$.

\

\tmtextit{Independence.} If $(X_1, \ldots, X_n) \in \mathcal{G} (V_1 \oplus
\cdots \oplus V_n)$ are compatible Grassmann variables with values in the
probability space $(\mathcal{A}, \omega)$, then we say that $X_1, \ldots, X_n$
are (tensor) independent (with respect to the state $\omega$) if, for all $F_j
\in \Lambda V_j$, we have that~{\cite{dirksen_noncommutative_2010}}
\[ \omega \left( \prod_{j = 1}^k \vphantom{\frac{}{}} X_j (F_j) \right) =
   \prod_{j = 1}^k \omega (X_j (F_j)) . \]


For example if $(X, Y) \in \mathcal{G} (V \oplus W)$ are (compatible and)
independent Grassmann random variables with values in $(\mathcal{A}, \omega)$,
then for all $v \in V \approx V \oplus 0, w \in W \approx 0 \oplus W$ we have
$v w \approx ((v \oplus 0) \wedge (0 \oplus w)) \in \Lambda (V \oplus W)
\approx \Lambda V \otimes \Lambda W$ and
\[ \omega ((X \otimes Y) (v w)) = \omega (X (v) Y (w)) = \omega (X (v)) \omega
   (Y (w)) . \]


Let now $X \in \mathcal{G} (V)$ be a Grassmann random variable on the
probability space $(\mathcal{A}, \omega)$. By the GNS construction we can
construct a Hilbert space $\mathcal{K}$ and a vector $\Omega \in \mathcal{K}$
such that there is $\tilde{X} \in \tmop{Hom} (\Lambda V ; \mathcal{L}
(\mathcal{K}))$ such that $\omega (X (F)) = \langle \Omega, \tilde{X} (F)
\Omega \rangle_{\mathcal{K}}$. By restriction, we can always take
$\mathcal{K}= \overline{\tilde{X} (\Lambda V)}$ while keeping this relation.
Here the bar denotes closure with respect to the Hilbert space topology of
$\mathcal{K}$. Then on $\mathcal{K}$ we can introduce by density an involution
$R$ that acts as $R \tilde{X} (F) \Omega = \tilde{X} ((- 1)^{| F |} F) \Omega$
where $| \cdot |$ is the grading on $\Lambda V$. We conclude that, without
loss of generality, as long as we are interested only in the law of a
Grassmann random variable $X \in \mathcal{G} (V)$, we can assume that it is
defined on $\mathcal{L} (\mathcal{H})$ for some Hilbert space $\mathcal{H}$
with a vector state $\omega (\cdot) = (\Omega, \cdot \Omega)$ and that it
comes with an involution $R_X : \mathcal{H} \rightarrow \mathcal{H}$ such that
$R_X X (F) = X ((- 1)^{| F |} F) R_X$ for all $F \in \Lambda V$.

\

As a consequence we can always arrange to realize two Grassmann variables on
the same probability space in such a way that they are compatible and
independent while preserving their law. Indeed for $i = 1, 2$, consider the
Grassmann variable $X_i : \Lambda V_i \rightarrow \mathcal{A}_i$ defined on
the probability space $(\omega_i, \mathcal{A}_i =\mathcal{L} (\mathcal{H}_i))$
with involution $R_i = R_{X_i}$ and vector state $\omega_i = (\Omega_i, \cdot
\Omega_i)$. Let $\mathcal{H}=\mathcal{H}_1 \otimes \mathcal{H}_2$, and $\omega
: \mathcal{L} (\mathcal{H}) \rightarrow \mathbb{C}$ given by $\omega (x)
\assign (\Omega_1 \otimes \Omega_2, x \Omega_1 \otimes
\Omega_2)_{\mathcal{H}}$, $x \in \mathcal{L} (\mathcal{H})$. Moreover we
define the random variables $\tilde{X}_i : \Lambda (V \oplus V) \rightarrow
\mathcal{L} (\mathcal{H})$, obtained by extending the relations $\tilde{X}_1
(v) = X_1 (v) \otimes R_2$ and $\tilde{X}_2 (w) =\mathbb{I}_{\mathcal{H}_1}
\otimes X_2 (w)$ for $v, w \in V$. Then $\tilde{X}_1 (v) \tilde{X}_2 (w) = -
\tilde{X}_2 (w) \tilde{X}_1 (v)$ and since $\Lambda (V \oplus V) \approx
\Lambda V \otimes \Lambda V$ we have also that the map $\tilde{X} (v \oplus w)
= \tilde{X}_1 (v \oplus 0) + \tilde{X}_2 (0 \oplus w)$ with $v, w \in V$ can
be extended to a homomorphism from $\Lambda (V \oplus V)$ to $\mathcal{A}$.
Hence the variables $\tilde{X}_i$, $i = 1, 2$, are compatible. Moreover they
have the same law as $X_i$ because $\omega (\tilde{X}_i (F_i)) = \omega_i (X_i
(F_i))$ for all $F_1, F_2 \in \Lambda V$, $i = 1, 2$. Finally they are
independent: $\omega (\tilde{X}_1 (F_1) \tilde{X}_2 (F_2)) = \omega_1 (X_1
(F_1)) \omega_2 (X_2 (F_2))$.

\subsection{Topology and calculus on Grassmann variables}\label{sec:top}

Before beginning we want to specify the topology we consider on $\mathcal{G}
(V)$. We consider here only the case where $V$ is finite dimensional, since
there is a more or less unique natural topology on $\mathcal{G} (V)$. Some
topologies of $\mathcal{G} (V)$ when $V$ is infinite dimensional are discussed
in Section \ref{sec:infinite-dim}.

\

When $V$ is finite dimensional all norms on $V$ are equivalent. For this
reason we choose the norm induced by the pre-Hilbert space inner product
$\langle \cdot, \cdot \rangle$ (that in the case of $V$ finite dimensional is
a Hilbert space inner product on $V$) related to the construction of
$V$-valued Gaussian random variables $X$ (see
Section~\ref{sec:gaussians}).\quad In this case $\mathcal{G} (V)$ has a
natural metric topology given by the distance
\begin{equation}
  d_{\mathcal{G} (V)} (X, Y) \assign \| X - Y \|_{\mathcal{G} (V)} = \sup_{v
  \in V, | v |_V = 1} \| X (v) - Y (v) \|_{\mathcal{A}},
  \label{eq:distrance-g}
\end{equation}
where $\| \cdot \|_{\mathcal{A}}$ is the natural norm in the
$\asterisk$-algebra $\mathcal{A}$.

\begin{remark}
  \label{remark_infinitenorm}In principle the definition of the
  distance~{\eqref{eq:distrance-g}} and the related norm $\| \cdot
  \|_{\mathcal{G} (V)}$ does not use in any way the fact that $V$ is finite
  dimensional. For this reason in the following, when $V$ is a (in general
  infinite dimensional) per-Hilbert space we will use the notation $\| X
  \|_{\mathcal{G} (V)}$ for the quantity $\sup_{v \in V, | v |_V = 1} \| X (v)
  \|_{\mathcal{A}}$.
\end{remark}

When $V$ is finite dimensional $\mathcal{G} (V)$ is a complete metric space
with respect to $d_{\mathcal{G} (V)}$, in fact we have:

\begin{lemma}
  \label{lemma_completeness}The metric $d_{\mathcal{G} (V)}$ makes
  $\mathcal{G} (V)$ a complete metric space.
\end{lemma}

\begin{proof}
  The quantity $d_{\mathcal{G} (V)} (X, Y)$ as defined
  in~{\eqref{eq:distrance-g}} satisfies the usual properties of a metric. If
  $(X_n)_n$ is a Cauchy sequence with respect to $d_{\mathcal{G} (V)}$, there
  is an element $X \in \mathcal{L} (V, \mathcal{A})$ such that $(X_n)_n$
  converges to $X$ in $\mathcal{L} (V, \mathcal{A})$. The only thing that we
  have to prove is that the linear map $X \in \mathcal{L} (V, \mathcal{A})$
  can be extended to an homomorphism from $\Lambda V$ to $\mathcal{A}$. This
  is equivalent to prove that for any $v, w \in V$ we have $X (v) X (w) = - X
  (w) X (v)$ where the product is the natural one in $\mathcal{A}$. On the
  other hand by the continuity of the product of $\mathcal{A}$ with respect to
  $\| \cdot \|_{\mathcal{A}}$ we have
  \[ X (v) X (w) = \lim_{n \rightarrow + \infty} X_n (v) X_n (w) = - \lim_{n
     \rightarrow + \infty} X_n (w) X_n (v) = - X (w) X (v) . \]
\end{proof}

\begin{remark}
  The distance $d_{\mathcal{G} (V)}$ is not the unique reasonable choice in
  $\mathcal{G} (V)$ and it depends on the topology chosen on $\mathcal{A}$ (in
  this case we choose the topology of the uniform converge of operators). An
  essential choice to preserve the non-linear structure of $\mathcal{G} (V)$,
  exploited in the proof of Lemma~\ref{lemma_completeness}, seems to be that
  the product on $\mathcal{A}$ is continuous with respect to this topology. 
\end{remark}

In the following we will denote $\| Y \|_{\mathcal{G} (V)}$ for $Y \in
\mathcal{G} (V)$ the positive number
\[ \| Y \|_{\mathcal{G} (V)} \assign \| Y \|_{\mathcal{L} (V, \mathcal{A})} =
   \sup_{v \in V, \| v \| = 1} \| Y (v) \|_{\mathcal{A}}, \]
i.e. the norm of the restriction of $Y$ to $\mathcal{L} (V, \mathcal{A})$.
This is useful to formulate a Taylor formula on $\mathcal{G} (V)$. We will use
the simpler notation $m$ for the multiplication $m_{\mathcal{A}}$ in
$\mathcal{A}$.

\

Consider $G \in \Lambda V$ and define the \tmtextit{right derivative}
$\partial_R : \Lambda V \rightarrow \Lambda V \otimes V$ by
\begin{equation}
  \partial_R G \assign (\mathbb{I} \otimes \Pi_V) (\Delta G) : \Lambda V
  \rightarrow \Lambda V \otimes V \label{eq:derivativeR}
\end{equation}
where $\Pi_V : \Lambda V \rightarrow V$ is the projection from the tensor
algebra $\Lambda V$ onto $V$. Then
\[ G (X + Y) - G (X) - m_{\mathcal{A}}^2 [(X \otimes Y) (\partial_R G)] =
   m_{\mathcal{A}}^2 [(X \otimes Y) (\mathbb{I}-\mathbb{I} \otimes \varepsilon
   -\mathbb{I} \otimes \Pi_V) \Delta G] \]
for $G \in \Lambda V$, where we used that $G (X) = m_{\mathcal{A}}^2 [(X
\otimes Y) (\mathbb{I} \otimes \varepsilon) \Delta G]$. Given that
\[ (\mathbb{I}-\mathbb{I} \otimes \varepsilon -\mathbb{I} \otimes \Pi_V)
   \Delta G \in \Lambda V \otimes \Lambda_{\geqslant 2} V \]
where $\Lambda_{\geqslant 2} V$ denotes the subspace of $\Lambda V$ of
elements of degree $\geqslant 2$. We can define recursively the $k$+1-th
derivative as
\[ \partial_R^{k + 1} = (\partial_R \otimes \mathbb{I} \otimes \cdots \otimes
   \mathbb{I}) \partial_R^k : \Lambda V \rightarrow \Lambda V \otimes V
   \otimes \cdots \otimes V \]
where $\partial_R^k$ is the $k$-th order derivative. Note that the right
derivative $\partial_R : \Lambda V \rightarrow \Lambda V \otimes V$ satisfies
\[ \partial_R (f_1 \cdots f_n) = \sum_{k = 1}^n (- 1)^{n - k} \left( f_1
   \cdots \centernot{f_k} \cdots f_n \right) \otimes f_k, \qquad f_1, \ldots, f_n
   \in V. \]
Of course one can define also a left derivative $\partial_L = (\Pi_V \otimes
\mathbb{I}) \circ \Delta : \Lambda V \rightarrow V \otimes \Lambda V$ with
similar properties and
\[ \partial_L (f_1 \cdots f_n) = \sum_{k = 1}^n (- 1)^{k - 1} f_k \otimes
   \left( f_1 \cdots \centernot{f_k} \cdots f_n \right), \qquad f_1, \ldots, f_n \in
   V. \]
We will consider $\Lambda V \otimes V$ as a $\Lambda V$-bimodule and in
particular we define the bilinear form $\langle \cdot, \cdot \rangle :
(\Lambda V \otimes V) \otimes V \rightarrow \Lambda V$ by
\begin{equation}
  \langle f \otimes v, w \rangle = f \langle v, w \rangle, \qquad f \in
  \Lambda V, v, w \in V. \label{eq:bimodule}
\end{equation}
Note that there is no ambiguity on whether $\langle \cdot, \cdot \rangle$
denotes this bilinear form or whether it denotes the scalar product of $V$.

\begin{remark}
  \label{remark_infinitederivative}The definition of derivative $\partial_R^k$
  and the bilinear for {\eqref{eq:bimodule}} do not depend on the fact that
  $V$ is finite dimensional. For this reason we can define $\partial_R^k$ and
  the bilinear for {\eqref{eq:bimodule}} also when $V$ is an infinite
  dimensional pre-Hilbert space.
\end{remark}

\begin{remark}
  Note that if $G_{\tmop{odd}} \in \Lambda_{\tmop{odd}} V$ then $\partial_R
  G_{\tmop{odd}} = \partial_L G_{\tmop{odd}}$. On the other hand, if
  $G_{\tmop{even}} \in \Lambda_{\tmop{even}} V$ then $\partial_R
  G_{\tmop{even}} = - \partial_L F_{\tmop{even}}$. In general, if we denote by
  $\pi_{\tmop{odd}}, \pi_{\tmop{even}}$ respectively the projection onto the
  odd and even parts of $\Lambda V$, then $\partial_R = (\pi_{\tmop{even}} -
  \pi_{\tmop{odd}}) \partial_L$.
\end{remark}

We can now state the version of Taylor formula which we will use below.

\begin{lemma}
  \label{lemma:taylor}Consider $G \in \Lambda V$ and let $X, Y \in \mathcal{G}
  (V)$ be two compatible Grassmann random variables such that $\| Y
  \|_{\mathcal{G} (V)} \leqslant 1$, then
  \begin{equation}
    G (X + Y) = G (X) + \sum_{k = 1}^n \frac{1}{k!} m_{\mathcal{A}}^{k + 1}
    [(X \otimes Y \otimes \cdots \otimes Y) (\partial_R^k G)] + O (\| Y \|^{n
    + 1}_{\mathcal{G} (V)}) \label{eq:taylor}
  \end{equation}
\end{lemma}

\begin{proof}
  We have
  \[ G (X + Y) = m_{\mathcal{A}}^2 [(X \otimes Y) (\Delta G)] = \sum_{k
     \geqslant 0} m_{\mathcal{A}}^2 [(X \otimes Y) (1 \otimes \Pi_{\Lambda^k
     V}) (\Delta G)], \]
  where the sum is finite since $F \in \Lambda V$ is a finite polynomial.
  Equation~{\eqref{eq:taylor}} easily follows as soon as we prove that
  \[ m_{\mathcal{A}}^{k + 1} [(X \otimes Y \otimes \cdots \otimes Y)
     (\partial_R^k G)] = (k!) m_{\mathcal{A}}^2 [(X \otimes Y) (\mathbb{I}
     \otimes \Pi_{\Lambda^k V}) (\Delta G)] . \]
  We have
  \begin{eqnarray*}
    \partial^k_R G & = & (\partial_R \otimes \mathbb{I} \otimes \cdots \otimes
    \mathbb{I}) \cdots (\partial_R \otimes \mathbb{I}) \partial_R G\\
    & = & \left( \vphantom{\tfrac{a}{b}} (1 \otimes \Pi_V) \Delta \otimes
    \mathbb{I} \otimes \cdots \otimes \mathbb{I} \right) \cdots \left(
    \vphantom{\tfrac{a}{b}} (\mathbb{I} \otimes \Pi_V) \Delta \otimes
    \mathbb{I} \right) \left( \vphantom{\tfrac{a}{b}} \mathbb{I} \otimes \Pi_V
    \right) \Delta G\\
    & = & \left( \vphantom{\tfrac{a}{b}} (\mathbb{I} \otimes \Pi_V) \otimes
    \mathbb{I} \otimes \cdots \otimes \mathbb{I} \right) (\Delta \otimes
    \mathbb{I} \otimes \cdots \otimes \mathbb{I}) \cdots \left(
    \vphantom{\tfrac{a}{b}} (\mathbb{I} \otimes \Pi_V) \otimes \mathbb{I}
    \right) (\Delta \otimes \mathbb{I}) \left( \vphantom{\tfrac{a}{b}}
    \mathbb{I} \otimes \Pi_V \right) \Delta G\\
    & = & \left( \vphantom{\tfrac{a}{b}} (\mathbb{I} \otimes \Pi_V) \otimes
    \mathbb{I} \otimes \cdots \otimes \mathbb{I} \right) \cdots \left(
    \vphantom{\tfrac{a}{b}} (\mathbb{I} \otimes \Pi_V) \otimes \mathbb{I}
    \right) \left( \vphantom{\tfrac{a}{b}} \mathbb{I} \otimes \Pi_V \right) 
    (\Delta \otimes \mathbb{I} \otimes \cdots \otimes \mathbb{I}) \cdots
    (\Delta \otimes \mathbb{I}) \Delta G\\
    & = & \left( \vphantom{\tfrac{a}{b}} \mathbb{I} \otimes \Pi_V \otimes
    \Pi_V \otimes \cdots \otimes \Pi_V \right)  (\Delta \otimes \mathbb{I}
    \otimes \cdots \otimes \mathbb{I}) \cdots (\Delta \otimes \mathbb{I})
    \Delta G.
  \end{eqnarray*}
  Let $N \assign \{ 1, \ldots, n \}$, $\#N = n$. For $I = \{ j_1, \ldots, j_k
  \} \subset N$, we put $v_I \assign v_{j_1} \cdots v_{j_k}$ for $v_j \in V, j
  \in N$. Let $\tmop{sgn} (N\backslash I)$ be a sign such that $v_N =
  \tmop{sgn} (N\backslash I) v_{N\backslash I} v_I$ By definition of the
  coproduct $\Delta$ in the exterior algebra $\Lambda V$, we have
  \[ \begin{array}{c}
       (\mathbb{I} \otimes \Pi_V \otimes \cdots \otimes \Pi_V) (\Delta \otimes
       \mathbb{I} \otimes \cdots \otimes \mathbb{I}) \cdots (\Delta \otimes
       \mathbb{I}) (\Delta) (v_1 \cdots v_n) =\\
       = \sum_{j_1, \ldots j_n \in N} \tmop{sgn} (N\backslash I)
       v_{N\backslash \{ j_1, \ldots, j_n \}} \otimes v_{j_1} \otimes \cdots
       \otimes v_{j_n} .
     \end{array} \]
  On the other hand
  \[ \begin{array}{lll}
       (\mathbb{I} \otimes \Pi_{\Lambda^k V}) \Delta v_1 \cdots v_n & = &
       (\mathbb{I} \otimes \Pi_{\Lambda^k V}) \sum_{I \subset N} \tmop{sgn}
       (N\backslash I) v_{N \setminus I} \otimes v_I\\
       & = & \frac{1}{k!} \sum_{j_1, \ldots, j_k} \tmop{sgn} (N\backslash I)
       v_{N\backslash \{ j_1, \ldots, j_k \}} \otimes v_{\{ j_1, \ldots, j_k
       \}} .
     \end{array} \]
  By linearity, we therefore have
  \[ (\mathbb{I} \otimes m^k_{\Lambda V}) (\partial^k_R G) = k! (\mathbb{I}
     \otimes \Pi_{\Lambda^k V}) (\Delta G) . \]
  where $m^k_{\Lambda V} : V^{\otimes k} \subset (\Lambda V)^{\otimes k}
  \rightarrow \Lambda V$ is defined as in equation {\eqref{eq:productA}}.
  Since $X, Y$ are algebra homomorphisms they commute with the product, hence
  we have
  \begin{eqnarray*}
    m_{\mathcal{A}}^{k + 1} [(X \otimes Y \otimes \cdots \otimes Y)
    (\partial_R^k G)] & = & m^2_{\mathcal{A}} [(X \otimes Y) (\mathbb{I}
    \otimes m^k_{\Lambda V}) (\partial_R^k G)]\\
    & = & k!m^2_{\mathcal{A}} [(X \otimes Y) (\mathbb{I} \otimes
    \Pi_{\Lambda^k V}) (\Delta G)] .
  \end{eqnarray*}
  This concludes the proof.
\end{proof}

We want to give now some precise bounds on the norms $\| G (X)
\|_{\mathcal{A}}$ and $\| G (X) - G (Y) \|_{\mathcal{A}}$, where $G \in
\Lambda V$ and $X, Y \in \mathcal{G} (V)$ are two compatible Grassmann random
variables. First we introduce a suitable norm $\| \cdot \|_{\pi}$ on $\Lambda
V$ that is the norm induced by the projective norm on $\bigoplus_{n = 0}^{\dim
(V)} \bigotimes^n V$. Note that there exists an injection $i_{\Lambda^n V} :
\Lambda^n V \rightarrow \bigotimes^n V$ given by the unique linear extension
of the following relation
\[ i_{\Lambda^n V} (v_1 \wedge \cdots \wedge v_n) = \frac{1}{n!} \sum_{\sigma
   \in S_n} (- 1)^{\sigma} v_{\sigma (1)} \otimes v_{\sigma (2)} \otimes
   \cdots \otimes v_{\sigma (n)} . \qquad v_1, \ldots, v_n \in V. \]
We then define the map $i_{\Lambda V} \assign \bigoplus_{n = 0}^{\dim (V)}
i_{\Lambda^n V}$. On $\bigoplus_{n = 0}^{\dim (V)} \bigotimes^n V$ we consider
the projective norm $\| \cdot \|_{\pi}$ defined as follows. If $f \in
\bigotimes^n V$ then
\[ \| f \|_{\pi} \assign \inf \left\{ \sum_{k = 1}^p \| f_1^k \|_V \cdots \|
   f^k_n \|_V, \tmop{where} f = \sum_{k = 1}^p f^k_1 \otimes \cdots \otimes
   f^k_n \right\} . \]
For a general element $g \in \bigoplus_{n = 0}^{\dim (V)} \bigotimes^n V$ we
put $\| g \|_{\pi} \assign \sum_{k = 0}^{\dim (V)} \| \Pi_{\otimes^n V}  (g)
\|_{\pi}$, and for any $w \in \Lambda V$,
\begin{equation}
  \| w \|_{\Lambda_{\pi} V} \assign \| i_{\Lambda V} (w) \|_{\pi} .
  \label{eq:def-lambdapi-norm}
\end{equation}
Since $i_{\Lambda V}$ is an injection, $\| \cdot \|_{\pi}$ defines a norm on
$\Lambda V$. If $G \in \Lambda V$ we define
\[ \deg (G) = \max \left\{ n \in \mathbb{N}, \tmop{such} \tmop{that} \|
   \Pi_{\Lambda^n V} (G) \|_{\pi} \centernot{=} 0 \right\} . \]
Let $W$ a vector space and introduce the symmetrizer $S : T W \rightarrow T W$
as the operator from the tensor algebra generated by $W$ in itself that is the
unique linear extension of
\[ S (w_1 \otimes \cdots \otimes w_n) = \frac{1}{n!} \sum_{\sigma \in S_n}
   w_{\sigma (1)} \otimes \cdots \otimes w_{\sigma (n)}, \qquad w_1, \ldots,
   w_n \in W. \]
It is important to note that for $S ((v + w)^{\otimes n})$ the binomial
formula holds, i.e. we have
\begin{equation}
  S ((v + w)^{\otimes n}) = \sum_{k = 0}^n \left( \begin{array}{c}
    n\\
    k
  \end{array} \right) S (v^{\otimes k} \otimes w^{\otimes (n - k)}) .
  \label{eq:binomialformula}
\end{equation}
\begin{lemma}
  \label{lemma_tensor}If $G_n \in \Lambda^n V$ and $X, Y \in \mathcal{G} (V)$
  are two compatible Grassmann random variables we have that
  \[ m_{\mathcal{A}}^{k + 1} [(X \otimes Y \otimes \cdots \otimes Y)
     (\partial_R^k G_n)] = \frac{n!}{(n - k) !} m_{\mathcal{A}}^n (S
     (X^{\otimes (n - k)} \otimes Y^{\otimes k}) (i_{\Lambda V} (G_n))) . \]
\end{lemma}

\begin{proof}
  For any $\lambda \in \mathbb{R}$, by Lemma~\ref{lemma:taylor} we have
  \begin{equation}
    G_n (X + \lambda Y) = G_n (X) + \sum_{k = 1}^n \frac{\lambda^k}{k!}
    m_{\mathcal{A}}^{k + 1} [(X \otimes Y \otimes \cdots \otimes Y)
    (\partial_R^k G_n)] . \label{eq:expansion1}
  \end{equation}
  On the other hand, by the binomial formula~{\eqref{eq:binomialformula}}, we
  have
  \begin{eqnarray}
    G_n (X + \lambda Y) & = & m_{\mathcal{A}}^n (S ((X + \lambda Y)^{\otimes
    n}) i_{\Lambda V} (G_n)) \nonumber\\
    & = & m_{\mathcal{A}}^n \left( \sum_{k = 0}^n \left( \begin{array}{c}
      n\\
      k
    \end{array} \right) \lambda^k S (X^{\otimes (n - k)} \otimes Y^{\otimes
    k}) i_{\Lambda V} (G_n) \right) .  \label{eq:expansion2}
  \end{eqnarray}
  By comparing the expressions~{\eqref{eq:expansion1}}
  and~{\eqref{eq:expansion2}} as polynomials in $\lambda$ we get the thesis.
\end{proof}

\begin{theorem}
  \label{theorem_function1}Let $X, Y \in \mathcal{G} (V)$ be two compatible
  Grassmann random variables and let $G \in \Lambda V$ then, for any $n
  \leqslant \deg (G) - 1$, we have
  \begin{equation}
    \| G (X) \|_{\mathcal{A}} \leqslant \| G \|_{\Lambda_{\pi} V} (1 + \| X
    \|_{\mathcal{G} (V)})^{\deg (G)} \label{eq:localbound}
  \end{equation}
  and
  \begin{equation}
    \begin{array}{c}
      \left\| G (Y) - G (X) - \sum_{k = 1}^n \frac{1}{k!} m_{\mathcal{A}}^{k +
      1} [(X \otimes (Y - X) \otimes \cdots \otimes (Y - X)) (\partial_R^k G)]
      \right\|_{\mathcal{A}} \leqslant\\
      \leqslant C_{n, \deg (G)} (1 + \max (\| X \|_{\mathcal{G} (V)}, \| Y
      \|_{\mathcal{G} (V)}))^{\deg (G) - n - 1} \| G \|_{\Lambda_{\pi} V} \| Y
      - X \|_{\mathcal{G} (V)}^{n + 1}
    \end{array} \label{eq:locallip}
  \end{equation}
  where $C_{n, \deg (G)} > 0$ is a suitable constant depending only on $n$ and
  $\deg (G)$. In the case $n = 0$ we can choose $C_{0, \deg (G)} = \deg (G)$.
\end{theorem}

\begin{proof}
  First of all we note that if $(G_1, \ldots ., G_n) \in \mathcal{L} (V,
  \mathcal{A})$ and $\tilde{G} \in \bigotimes^n V$ we have
  \begin{equation}
    \| m_{\mathcal{A}}^n ((G_1 \otimes \cdots \otimes G_n) (\tilde{G}))
    \|_{\mathcal{A}} \leqslant \| \tilde{G} \|_{\pi} \| G_1 \|_{\mathcal{L}
    (V, \mathcal{A})} \cdots \| G_n \|_{\mathcal{L} (V, \mathcal{A})} ._{}
    \label{eq:inequalitytensor}
  \end{equation}
  Furthermore, if $X \in \mathcal{G} (V)$ and $G \in \Lambda V$ we have
  \[ G (X) = \sum^{\deg (G)}_{n = 0} m_{\mathcal{A}}^n (X^{\otimes n}
     (i_{\Lambda^n V} (\Pi_{\Lambda^n V} (G)))) . \]
  Using then the definition of $\| X \|_{\mathcal{G} (V)}$ and of $\| G
  \|_{\Lambda_{\pi} V}$ we get
  \[ \| G (X) \|_{\mathcal{A}} \leqslant \sum_{n = 0}^{\deg (G)} \| X
     \|_{\mathcal{G} (V)}^n \| i_{\Lambda^n V} (\Pi_{\Lambda^n V} (G))
     \|_{\pi} \leqslant (1 + \| X \|_{\mathcal{G} (V)})^n \| G \|_{\pi} . \]
  In general, writing $G_h = \Pi_{\Lambda^h V} (G)$, by Lemma
  \ref{lemma_tensor}, we have
  \begin{equation}
    \begin{array}{c}
      G_h (Y) - G_h (X) - \sum_{k = 1}^n \frac{1}{k!} m_{\mathcal{A}}^{k + 1}
      [(X \otimes (Y - X) \otimes \cdots \otimes (Y - X)) (\partial_R^k G_h)]
      =\\
      = m_{\mathcal{A}}^h \left( \left( S (Y^{\otimes h}) - \sum_{k = 0}^n
      \left( \begin{array}{c}
        h\\
        k
      \end{array} \right) S (X^{\otimes (h - k)} \otimes (Y - X)^{\otimes k})
      \right) (i_{\Lambda V} (G_h)) \right) =\\
      = m_{\mathcal{A}}^h \left( \left( \sum_{k = n + 1}^h \left(
      \begin{array}{c}
        h\\
        k
      \end{array} \right) S (X^{\otimes (h - k)} \otimes (Y - X)^{\otimes k})
      \right) (i_{\Lambda V} (G_h)) \right) =\\
      = m_{\mathcal{A}}^h \left. \left. \left. \left( \left( \sum_{k = n +
      1}^h \sum_{\ell = 0}^{k - n - 1} \left( \begin{array}{c}
        h\\
        k
      \end{array} \right) \left( \begin{array}{c}
        k - n - 1\\
        \ell
      \end{array} \right) S (X^{\otimes (h - k)} \otimes (- X)^{\otimes \ell}
      \otimes (Y)^{k - n - 1 - \ell} \nobracket \right. \right. \right.
      \right. \right.\\
      \nobracket \nobracket \nobracket \nobracket \nobracket \nobracket
      \otimes (Y - X)^{\otimes (n + 1)})) (i_{\Lambda V} (G_h))) \nobracket
      \nobracket \nobracket = m_{\mathcal{A}}^h \left. \left. \left( \left(
      \sum_{p = 0}^{h - n - 1} S (X^{h - p} \otimes Y^{\otimes p}) \right.
      \right. \right. \right.\\
      \left. \left. \left. \left. \left( \sum_{\ell = 0}^{h - n - 1 - p} (-
      1)^{\ell}_{} \left( \begin{array}{c}
        h\\
        p + \ell + n + 1
      \end{array} \right) \left( \begin{array}{c}
        \ell + p\\
        \ell
      \end{array} \right) \right) \right) (i_{\Lambda V} (G_h)) \right)
      \right. \right.
    \end{array} \label{eq:equalityexpansion}
  \end{equation}
  Using inequality~{\eqref{eq:inequalitytensor}} in
  relation~{\eqref{eq:equalityexpansion}} we get
  \[ \begin{array}{c}
       \left\| G_h (Y) - G_h (X) - \sum_{k = 1}^n \frac{1}{k!} m^{k +
       1}_{\mathcal{A}} [(X \otimes (Y - X) \otimes \cdots \otimes (Y - X))
       (\partial_R^k G_h)] \right\|_{\mathcal{A}} \leqslant\\
       \leqslant \| i_{\Lambda V} (G_h) \|_{\Lambda_{\pi} V} C_{n, h} (\max
       (\| X \|_{\mathcal{G} (V)}, \| Y \|_{\mathcal{G} (V)}))^{h - n - 1} \|
       Y - X \|_{\mathcal{G} (V)}^{n + 1},
     \end{array} \]
  where we can choose
  \[ C_{n, h} = \sum_{p = 0}^{h - n - 1} \left| \sum_{\ell = 0}^{h - n - 1 -
     p} (- 1)^{\ell}_{} \left( \begin{array}{c}
       h\\
       p + \ell + n + 1
     \end{array} \right) \left( \begin{array}{c}
       \ell + p\\
       \ell
     \end{array} \right) \right| . \]
  In the case $n = 0$ we can get a better constant. Indeed we have
  \[ \| G_h (Y) - G_h (X) \|_{\mathcal{A}} = \left\| \sum_{k = 0}^{h - 1} S
     (X^{\otimes k} \otimes Y^{\otimes h - 1 - k} \otimes (Y - X))
     \right\|_{\mathcal{G} (V)} \| G_h \|_{\Lambda_{\pi} V} \]
  \[ \leqslant h (\max (\| X \|_{\mathcal{G} (V)}, \| Y \|_{\mathcal{G}
     (V)}))^{h - 1} \| X - Y \| . \]
\end{proof}

We now introduce the notion of function ``depending on the space variable'' $v
\in V$ and a suitable norm on the space of such functions. If $F \in
\mathcal{L} (V, \Lambda V)$ we can define the composition
\[ F (X) (v) \assign X (F (v)), \qquad X \in \mathcal{G} (V), v \in V. \]
as a linear (and continuous) map $F (X) : V \rightarrow \mathcal{A}$. We
define $\| F \|_{\Lambda_{\pi} V\mathcal{}}$ as
\begin{equation}
  \| F \|_{\Lambda_{\pi} V\mathcal{}} \assign \| F \|_{\mathcal{L} (V,
  \Lambda_{\pi} V)} = \sup_{v \in V, | v | = 1} \| F (v) \|_{\Lambda_{\pi} V}
  . \label{eq:normfunction}
\end{equation}
where in the r.h.s. we use the norm $\| \cdot \|_{\Lambda_{\pi} V}$ defined
in~{\eqref{eq:def-lambdapi-norm}}.

\begin{remark}
  If $F : V \rightarrow \Lambda_{\tmop{odd}} V$ we have that for any $v_1, v_2
  \in V$ we have
  \[ F (X) (v_1) F (X) (v_2) = - F (X) (v_2) F (X) (v_1) . \]
  This means that $F (X)$ can be extend to an homomorphism $\tmop{Hom}
  (\Lambda V, \mathcal{A})$ (we shall denote this extension by $F (X)$ too).
  Furthermore, since
  \[ F (X) (v_1) X (v_2) = - X (v_2) F (X) (v_1), \]
  we have that $F (X)$ and $X$ are compatible Grassmann random variables.
\end{remark}

We can define also $\partial_R^n F : V \rightarrow \Lambda V \otimes \left(
\bigotimes^n V \right)$ as $\partial_R^n F (v) = \partial^n_R (F (v))$, and
$\deg (F)$ in the obvious way.

\begin{theorem}
  \label{theorem_function2}Suppose that $F \in \mathcal{L} (V, \Lambda V)$ and
  let $X, Y \in \mathcal{G} (V)$ be compatible Grassmann random variables then
  we have that
  \[ \| F (X) \|_{\mathcal{G} (V)} \leqslant \| F \|_{\Lambda_{\pi} V} (1 +
     \| X \|_{\mathcal{G} (V)})^{\deg (F)} \]
  \[ \begin{array}{c}
       \left\| F (Y) - F (X) - \sum_{k = 1}^n \frac{1}{k!} m [(X \otimes (Y -
       X) \otimes \cdots \otimes (Y - X)) (\partial_R^k F)]
       \right\|_{\mathcal{G} (V)} \leqslant\\
       \leqslant C_{n, \deg (F)} (1 + \max (\| X \|_{\mathcal{G} (V)}, \| Y
       \|_{\mathcal{G} (V)}))^{\deg (F) - n - 1} \| F \|_{\Lambda_{\pi} V} \|
       Y - X \|_{\mathcal{G} (V)}^{n + 1}
     \end{array} \]
\end{theorem}

\begin{proof}
  The proof is a simple application of Theorem~\ref{theorem_function1} to $F
  (X) (v)$ and $F (Y) (v)$ for any fixed $v \in V$ and of the definition of
  the norm~{\eqref{eq:normfunction}}.
\end{proof}

\subsection{Grassmann Gaussian variables}\label{sec:gaussians}

Let now $V$ be a real pre-Hilbert space with scalar product $\langle \cdot,
\cdot \rangle$ and with an antisymmetric bounded operator $C : V \rightarrow
V$. We recall that, by Remark~\ref{remark_infinitenorm} and
Remark~\ref{remark_infinitederivative}, we can extend the definitions of $\|
\cdot \|_{\mathcal{G} (V)}$, $\partial_R^k$ and the bilinear
form~{\eqref{eq:bimodule}} from the finite dimensional case to the generic
(infinite) pre-Hilbert space $V$.

\begin{definition}
  A ($V$-)Grassmann (centered) Gaussian variable with correlation $C$ is a
  random variable $X \in \mathcal{G} (V)$ such that
  \begin{equation}
    \omega (X (G) X (f)) = \omega (X (\langle \partial_R G, C f \rangle)),
    \qquad G \in \Lambda V, f \in V. \label{eq:ibp-r}
  \end{equation}
  We also require that $\| X \|_{\mathcal{G} (V)} < \infty$, i.e that the map
  $X : V \rightarrow \mathcal{A}$ must be continuous with respect the topology
  induced on $V$ by the pre-Hilbert product structure and the (norm) topology
  of $\mathcal{A}$.
\end{definition}

If we define $\partial_R X (G) = X (\partial_R G)$, with the understanding
that $X (g \otimes w) = X (g) \otimes w$, $g \in V$ and $w \in \Lambda V$,
then we have the more familiar expression (similar to the bosonic counterpart)
\[ \omega (G (X) X (h)) = \omega (\langle \partial_R G (X), C h \rangle) . \]
Note that the integration by parts formula determines all the moments of the
Gaussians. In particular~{\eqref{eq:ibp-r}} implies that $\omega (X (f_1)
\cdots X (f_n)) = 0$ if $n$ is odd and if $n = 2 k$ is even
\begin{equation}
  \omega (X (f_1) \cdots X (f_{2 k})) = \sum_{\sigma} (- 1)^{\sigma} \prod_{i
  = 1}^k \langle f_{\sigma (2 i - 1)}, C f_{\sigma (2 i)} \rangle_V
  \label{eq:wick}
\end{equation}
where the sum runs over all the pairings $\sigma$ of $\{ 1, \ldots, 2 k \}$
and $(- 1)^{\sigma}$ is an appropriate sign. Equation~{\eqref{eq:wick}} is
often called Wick's rule. The right hand side of~{\eqref{eq:wick}} can be
written as a Pfaffian:
\begin{equation}
  \omega (X (f_1) \cdots X (f_n)) = \underset{1 \leqslant i, j \leqslant
  n}{\tmop{Pf}} \langle f_i, C f_j \rangle_V, \label{eq:pfaff}
\end{equation}
where the Pfaffian is defined, for an anti-symmetric $n \times n$ matrix $M$,
to be zero if $n = 2 k + 1$ and if n=2k as the polynomial in the entries of
$M$ which satisfy the relation
\[ (\underset{1 \leqslant i, j \leqslant n}{\tmop{Pf}} M_{i j})^2 =
   \underset{1 \leqslant i, j \leqslant n}{\det} M_{i j} . \]
The existence of a Gaussian variable, in the above sense, implies the
inequality
\[ \underset{1 \leqslant i, j \leqslant n}{\det} \langle f_i, C f_j \rangle_V
   = [\omega (X (f_1) \cdots X (f_n))]^2 \leqslant \| X (f_1) \|^2 \cdots \| X
   (f_n) \|^2 \leqslant \| X \|^{2 n} \| f_1 \|^2 \cdots \| f_n \|^2, \]
well known in the mathematical physics literature relative to fermionic
expansions e.g. see~{\cite{mastropietro_non-perturbative_2008}}.

\begin{remark}
  The averages of Gaussian variables depend only on the quadratic form $(f, g)
  \mapsto \langle f, C g \rangle_V$, however analysis on the Grassmann algebra
  relies on the scalar product $\langle \cdot, \cdot \rangle_V$ also via the
  requirement that $\| X (f) \|_{\mathcal{A}} \leqslant \| f \|_V$. In
  particular the realization of the Grassmann algebra as a family of bounded
  operators is not canonical. This is the reason we need to require the above
  continuity of the map $X : V \rightarrow \mathcal{A}$.
\end{remark}

In order to construct Grassmann Gaussian variables, we need ``a place to hang
the hat on'', this place will be a (canonical commuting relations) CAR algebra
endowed with its vacuum state. \ To allow for arbitrary correlations $C$ we
can use a standard trick~{\cite{osterwalder_euclidean_1973-2}} which consists
in doubling the generators of the CAR algebra with respect to the generators
of the Grassmann algebra.

\begin{lemma}
  \label{lemma:existence-gaussian}For every antisymmetric and bounded $C : V
  \rightarrow V$ there exists a ($V$-)Grassmann Gaussian variable $X$ with
  correlation $C$ (on a suitable probability space $(\mathcal{A}, \omega)$).
\end{lemma}

\begin{remark}
  \label{remark:commutative}It is important to note that Lemma
  {\eqref{lemma:existence-gaussian}} is peculiar of the Grassmannian setting
  and it has no equivalent in the commutative case in the following sense: if
  $V$ is and pre-Hilbert space and $S$ is a bounded positive, self-adjoint
  operator there is no map $X^{\tmop{com}} : V \rightarrow \mathcal{A}$ such
  that, for any $v_1, \ldots, v_{2 k} \in V$,
  \[ \omega (X^{\tmop{com}} (v_1) \cdots X^{\tmop{com}} (v_{2 k})) =
     \sum_{\mathcal{M} \in \left\{ \tmop{perfect} \; \tmop{matchings} \;
     \tmop{of} \; \{ 1, \ldots, 2 k \} \right\}} \prod_{(i, j) \in
     \mathcal{M}} \langle S v_i, v_j \rangle \]
  for some $C^{\ast}$-algebra $\mathcal{A}$ with a positive state $\omega$.
  Indeed in the commutative case the Gaussian variables are unbounded and so
  we must realize them in a algebra $\widetilde{\mathcal{A}}$ of unbounded
  operators (for example we can consider $\widetilde{\mathcal{A}} = \bigcap_{1
  \leqslant p < + \infty} L^p ((\Omega, \mathcal{F}, \mathbb{P}),
  \mathbb{C})$, for a suitable probability space $(\Omega, \mathcal{F},
  \mathbb{P})$ and $\omega (\cdot) =\mathbb{E} [\cdot]$, which is only a
  $\ast$-algebra and not a normed algebra, thus not a $C^{\ast}$-algebra).
\end{remark}

\begin{proof*}{Proof of Lemma \ref{lemma:existence-gaussian}}
  Let us consider $\Lambda V$ itself as a real pre-Hilbert space with respect
  to the scalar product $\langle \cdot, \cdot \rangle_{\Lambda V}$ on $\Lambda
  V$ given by
  \[ \langle v_1 \wedge \cdots \wedge v_n, w_1 \wedge \cdots \wedge w_m
     \rangle_{\Lambda V} = \delta_{n m} \underset{1 \leqslant j, k \leqslant
     n}{\det} \langle v_j, w_k \rangle, \]
  for $v_j, w_k \in V$, $j = 1, \ldots, n$, $k = 1, \ldots, m$, where $\langle
  \cdot, \cdot \rangle$ denotes as usual the scalar product on $V$. Note that
  $\langle \cdot, \cdot \rangle_{\Lambda V}$ on $\Lambda^n V$ is simply the
  restriction of the Hilbert scalar product on the tensor product $V^{\otimes
  n}$. Let $\mathcal{H}$ be the completion of $\Lambda V$ with respect to
  $\langle \cdot, \cdot \rangle_{\Lambda V}$ and denote by $\Omega$ the
  element in $\mathcal{H}$ which corresponds to $1 \in \Lambda V$, often
  referred to as the \tmtextit{vacuum} vector. Denote by $\lambda : \Lambda V
  \rightarrow \tmop{End} (\Lambda V)$ the left action of $\Lambda V$ on itself
  given by $\lambda (H) G = H \wedge G$, where $H, G \in \Lambda V$. We show
  that this action extends to a representation of $\Lambda V$ on
  $\mathcal{H}$. Indeed $\lambda (v)$, $v \in V$ corresponds to a creation
  operator. Let us denote by $\lambda (v)^{\Tau}$ the adjoint of the operator
  $\lambda (v)$ with respect to the scalar product $\langle \cdot, \cdot
  \rangle_{\Lambda V}$. A standard computation shows that for $v \in V$ we
  have
  \[ \lambda (v)^{\Tau} w_1 \wedge \cdots \wedge w_n = \sum_{\ell = 1}^n (-
     1)^{\ell - 1} \langle v, w_{\ell} \rangle w_1 \wedge \cdots \wedge
     \centernot{w_{\ell}} \wedge \cdots \wedge w_n . \]
  Namely $\lambda^T (v) x$ can be expressed in terms of the left derivative as
  $\lambda^T (v) x = \langle v, \partial_L x \rangle$ (recall the definition
  given in equation~{\eqref{eq:bimodule}}) for all $x \in \Lambda V$. In
  particular we have
  \begin{equation}
    \{ \lambda (v), \lambda (w) \} = \{ \lambda (v)^{\Tau}, \lambda (w)^{\Tau}
    \} = 0, \qquad \{ \lambda (v)^{\Tau}, \lambda (w) \} = \langle v, w
    \rangle, \qquad v, w \in V, \label{eq:car}
  \end{equation}
  where as usual $\{ \cdot, \cdot \}$ denotes the anticommutator. As a
  consequence we obtain, for any $x \in \Lambda V$
  \[ \langle \lambda (v) x, \lambda (v) x \rangle_{\Lambda V} + \langle
     \lambda (v)^{\Tau} x, \lambda (v)^{\Tau} x \rangle_{\Lambda V} = \langle
     x, \{ \lambda (v)^{\Tau}, \lambda (v) \} x \rangle_{\Lambda V} = \langle
     v, v \rangle \langle x, x \rangle_{\Lambda V} . \]
  In particular $\lambda (v), \lambda (v)^{\Tau}, v \in V$, extend to bounded
  operators on $\mathcal{H}$ (still denoted by the same symbol). We now define
  $X : \Lambda V \rightarrow \mathcal{L} (\mathcal{H})$ to be the algebra
  homomorphism given on $V$ by
  \[ X (v) \assign \lambda (C v) + \lambda (v)^{\Tau}, \qquad v \in V. \]
  Then $X$ is extended in a natural way to $\Lambda V$ (still denoted by the
  same symbol). Note that $X (x)$, $x \in \Lambda V$, is indeed a bounded
  operator since it is the sum of products of bounded operators. Finally we
  define $\omega : \mathcal{L} (\mathcal{H}) \rightarrow \mathbb{R}$ to be the
  state on $\mathcal{A} =\mathcal{L} (\mathcal{H})$ defined by $\Omega$, that
  is
  \[ \omega (T) \assign \langle \Omega, T \Omega \rangle_{\Lambda V}, \qquad T
     \in \mathcal{L} (\mathcal{H}) . \]
  We claim that $X$ is a Gaussian Grassmann random variable on $(\mathcal{L}
  (\mathcal{H}), \omega)$. We first show that it is Grassmann:
  \[ \begin{array}{lll}
       \{ X (v), X (w) \} & = & \{ \lambda (C v), \lambda (w)^{\Tau} \} + \{
       \lambda (v)^{\Tau}, \lambda (C w) \}\\
       & = & \langle C v, w \rangle + \langle v, C w \rangle = \langle v,
       C^{\Tau} w \rangle + \langle v, C w \rangle = 0,
     \end{array} \]
  where $v, w \in \Lambda$, because by assumption $C^{\Tau} = - C$. Now we
  have also that
  \begin{eqnarray*}
    \omega (X (v) X (w)) & = & \langle \Omega,^{} \lambda (v)^{\Tau} \lambda
    (C w) \Omega \rangle_{\Lambda V}\\
    & = & \langle \lambda (v) \Omega, \lambda (C w) \Omega \rangle_{\Lambda
    V} = \langle v, C w \rangle_{\Lambda V},
  \end{eqnarray*}
  where the first equality follows from the fact that $\lambda (v)^{\Tau}
  \Omega = 0$. More generally using that $X (v) \Omega = \lambda (C v) \Omega$
  and the commutation relation~{\eqref{eq:car}} we have that, for all $f_1,
  \ldots, f_{2 k} \in V$
  \[ \omega (X (f_1) \cdots X (f_{2 k})) = \omega (X (f_1) \cdots X (f_{2 k -
     1}) \lambda (C f_{2 k})) \]
  \[ = - \omega (X (f_1) \cdots \lambda (C f_{2 k}) X (f_{2 k - 1})) + \omega
     (X (f_1) \cdots X (f_{2 k - 2})) \langle f_{2 k - 1}, C f_{2 k} \rangle
  \]
  \[ = \cdots = \sum_{\ell = 1}^{2 k - 1} (- 1)^{2 k - 1 - \ell} \omega \left(
     X (f_1) \cdots \slashgen{X (f_{\ell})} \cdots X (f_{2 k - 1}) \right)
     \langle f_{\ell}, C f_{2 k} \rangle \]
  Therefore
  \[ \omega (X (G) X (f)) = \omega (X (\langle \partial_R G, C f \rangle)),
     \qquad G \in \Lambda V, f \in V. \]
  Note that, similarly
  \begin{equation}
    \omega (X (f) X (G)) = \omega (X (\langle f, C \partial_L G \rangle)) .
    \label{eq:ibp-l}
  \end{equation}
  
\end{proof*}

\tmtextit{Complex Gaussians.} Later on we will need also complex Grassmann
Gaussian variables, i.e. Gaussian variables taking values in $\Lambda V$ with
a complex pre-Hilbert space $V$ whose Hermitian scalar product we denote by
$(\cdot, \cdot)_V$ and we assume anti-linear in the left variable. As in the
commutative setting, their definition poses no particular problems, however
the interplay of the algebraic and analytic structure is here reflected on the
fact that we need to fix a real structure $\varkappa$ on $V$. We will note by
$\llangle \cdot, \cdot \rrangle_{\varkappa} = (\varkappa \cdot, \cdot)_V$ the
associated bilinear form on $V$ and by $A^{\varkappa} = \varkappa A^{\ast}
\varkappa$ the transposition of the operator $A : V \rightarrow V$ with
respect to $\varkappa$. Please refer to the Appendix~\ref{app:GP} for the
basic notions of real structures on complex Hilbert spaces.

\begin{definition}
  \label{def:complex-gaussians}Let $V$ a complex pre-Hilbert space,
  $\varkappa$ a real structure over it and $C : V \rightarrow V$ a
  $\varkappa$-antisymmetric (i.e. $C^{\varkappa} = - C$) bounded linear
  operator. A \ $(V, \varkappa)$-Grassmann (centered) Gaussian variable with
  correlation $C$ is a random variable $X \in \mathcal{G} (V)$ such that
  \[ \omega (X (G) X (h)) = \omega (X (\llangle \partial_R G, C h
     \rrangle_{\varkappa})), \qquad G \in \Lambda V, h \in V, \]
  where $\Lambda V$ is the Grassmann algebra over $\mathbb{C}$ generated by
  $V$. We require that $\| X \|_{\mathcal{G} (V)} < \infty$, i.e that the map
  $X : V \rightarrow \mathcal{A}$ is continuous.
\end{definition}

\begin{lemma}
  \label{lemma:existence-complex-gaussian}For every $\varkappa$-antisymmetric
  and bounded $C : V \rightarrow V$ there exists a \ $(V,
  \varkappa)$-Grassmann centered Gaussian variable $X \in \mathcal{G} (V)$
  with correlation $C$ (on a suitable probability space $(\mathcal{A},
  \omega)$).
\end{lemma}

\begin{proof}
  The construction of a complex Grassmann Gaussian proceeds as in
  Lemma~\ref{lemma:existence-gaussian} by considering the (complex)
  antisymmetric Fock space $\mathcal{H}$ associated to $V$ with vacuum vector
  $\Omega \in \mathcal{H}$ and creation operators $(a (v))_{v \in V}$ which
  are linear and satisfying the canonical anti-commutation relations (CAR) $\{
  a (w)^{\ast}, a (v) \} = (w, v)_V$ for $v, w \in V$. Let
  \[ X (v) = a (C v)^{\ast} + a (\varkappa v), \qquad v \in V, \]
  and consider the state $\omega (A) \assign \langle \Omega, A \Omega
  \rangle_{\mathcal{H}}$ for any $A \in \mathcal{A}=\mathcal{L}
  (\mathcal{H})$. The verification that the bounded operators $(X (v))_{v \in
  V}$ forms a Grassmann algebra and that
  \[ \omega (X (v) X (w)) = \llangle v, C w \rrangle_{\varkappa}, \qquad v, w
     \in V, \]
  as required, is left to the reader.
\end{proof}

\tmtextit{White noise.} A relevant example of a random variable taking values
in an infinite dimensional Grassmann algebra is $d$-dimensional (Gaussian)
white noise, defined as follows.

\begin{definition}
  \label{definition_Brownian}A $V$-valued $d$-dimensional (Gaussian) white
  noise with correlation $C : V \rightarrow V$ is the (centered) Grassmann
  Gaussian variable $\Xi \in \mathcal{G} (L^2 (\mathbb{R}^d) \otimes V)$ with
  correlation $\tilde{C}$ given by
  \[ (\tilde{C} f) (x) = C f (x) \]
  for all $x \in \mathbb{R}^d$ and $f \in L^2 (\mathbb{R}^d) \otimes V \approx
  L^2 (\mathbb{R}^d ; V)$.
\end{definition}

Take now a one dimensional white noise $\Xi$ with values in $V$ and let $B_t
(v) = \Xi (\mathbb{I}_{[0, t]} \otimes v)$ for $v \in V$. For fixed $t
\geqslant 0$, $B_t$ extends as homomorphism to all $\Lambda V$ and therefore
$B_t \in \mathcal{G} (V)$. Note also that
\begin{equation}
  \| B_t (v) - B_s (v) \|_{\mathcal{A}} \leqslant | t - s |^{1 / 2} \| v \|_V,
  \qquad t, s \geqslant 0, v \in V, \label{eq:regularity-B}
\end{equation}
so $B (v) \in C (\mathbb{R}_+, \mathcal{A})$ (here $C (\mathbb{R}_+,
\mathcal{A})$ is the space of continuous maps from $\mathbb{R}_+$ to
$\mathcal{A}$). Note that $B \in \tmop{Hom} (\Lambda V, C (\mathbb{R}_+,
\mathcal{A}))$: this in particular implies that $B_{t_1}, \ldots, B_{t_n}$ is
a compatible family and that $B = (B_t)_{t \in \mathbb{R}_+}$ is a Gaussian
process with continuous trajectories. We have that $B_0 (v) = 0$,
\[ \omega (B_t (v)) = 0, \qquad \omega (B_t (v) B_s (w)) = \langle v, C w
   \rangle  (t \wedge s), \qquad t, s \geqslant 0, \quad v, w \in V, \]
where $C$ is the correlation of the Grassmann Gaussian random variable $\Xi$.
Increments of $B$ are independent and higher order moments can be computed via
Wick's rule~{\eqref{eq:wick}}. Note in particular that
\begin{equation}
  \sup_{0 \leqslant t \leqslant T} \| B_t \|_{\mathcal{G} (V)} < \infty,
  \label{eq:boundedness-B}
\end{equation}
for all $T > 0$. Let us remark here that properties~{\eqref{eq:regularity-B}}
and~{\eqref{eq:boundedness-B}} are very different from the pathwise properties
of the commutative (bosonic) Brownian motion where the random variable
realizing Brownian motion is unbounded both with respect the probability space
and in time (see Remark \ref{remark:commutative} for other observations on
this topic).

\section{Finite dimensional SDEs}\label{sec:finite-dim}

We want to study simple SDEs taking values in $\Lambda V$ as far as it is
needed for the purpose of stochastic quantization, that is with additive white
noise. We refrain to undertake here a general study of Grassmann SDEs, in
particular no stochastic calculus will be needed below. It seems possible to
devise such a calculus but we leave it for a future work. More precisely in
this section we want to study SDEs driven by an additive Brownian motion $B_t$
when the linear space $V$ is finite dimensional.

\begin{definition}
  Let $F \in \mathcal{L} (V, \Lambda V)$ and assume that $F (v)$ is odd for
  all $v \in V$, and let $\Psi_0 \in \mathcal{G} (V)$ be a random variable
  compatible and independent of the Brownian motion $B \in \mathcal{G} (L^2
  (\mathbb{R}_+, V))$, then $\Psi \in C^0 ([0, T] ; \mathcal{G} (V))$ is a
  solution in $[0, T]$ to the (additive) SDE driven by $B$ with drift $F$ and
  initial condition $\Psi_0$, if,
  \begin{equation}
    \Psi_t (v) = \Psi_0 (v) + \int_0^t \Psi_s (F (v)) \mathd s + B_t (v),
    \qquad t \in [0, T], v \in V, \label{eq:sde}
  \end{equation}
  where the integral with respect to the variable $s$ is understood in
  Bochner's sense (relative to the norm on $\mathcal{G} (V)$ introduced in
  equation {\eqref{eq:distrance-g}}).
\end{definition}

\begin{remark}
  Note that we do not require that $\Psi_t, \Psi_s$ are compatible for $t \neq
  s$, in particular
  \[ C^0 ([0, T] ; \mathcal{G} (V)) \neq \tmop{Hom} (\Lambda V ; C ([0, T],
     \mathcal{A})), \]
  nor we require any compatibility of $\Psi$ and $B$. It will turn out that
  such compatibility holds in fact for the unique solution
  of~{\eqref{eq:sde}}, but it is not necessary to put such restriction to
  formulate the notion of solution.
\end{remark}

\subsection{The Grassmann Ornstein--Uhlenbeck motion}\label{sec:ou}

In this section we introduce the Grassmann analog of the Ornstein--Uhlenbeck
process, that is when $F (v) = A v$ with $A : V \rightarrow V$ is a linear
operator on $V$. In this case we can write down an explicit formula for the
solution to equation~{\eqref{eq:sde}} extended to $t \geqslant 0$.

\begin{proposition}
  \label{proposition_linear}Suppose that $A : V \rightarrow V$ and $F (v) = A
  v$. The unique solution to equation~{\eqref{eq:sde}} is given, in this case,
  by
  \begin{equation}
    \Psi_t (v) = \Psi_0 (e^{A t} v) + \Xi (\mathbb{I}_{[0, t]} (\cdot) e^{A (t
    - \cdot)} v), \qquad t \in \mathbb{R}_+, \label{eq:linear}
  \end{equation}
  where $\Xi$ is the Grassmann Gaussian noise related with the Brownian motion
  $B$ (see Definition~\ref{definition_Brownian} and the discussion that
  follows it).
\end{proposition}

\begin{proof}
  Let $h (t, v) \in L^2 (\mathbb{R}; V)$ be given by
  \[ h (t, v) (s) \assign \mathbb{I}_{[0, t]} (s) e^{A (t - s)} v, \qquad s
     \in \mathbb{R}. \]
  Then~{\eqref{eq:linear}} reads $\Psi_t (v) = \Psi_0 (e^{A t} v) + \Xi (h (t,
  v))$. An explicit computation using that
  \[ h (t, v) (s) -\mathbb{I}_{[0, t]} (s) v = \int_0^t A h (r, v) (s) \mathd
     r = \int_0^t h (r, A v) (s) \mathd r, \qquad s \in \mathbb{R}, \]
  in $L^2 (\mathbb{R}; V)$ gives
  \[ \Xi (h (t, v)) = B_t (v) + \int_0^t \Xi (h (r, A v)) \mathd r. \]
  Moreover, by the definition of exponential of a matrix, $\partial_t \Psi_0
  (e^{A t} v) = \Psi_0 (e^{A t} A v)$ so
  \[ \Psi_t (v) = \Psi_0 (v) + \int_0^t \Psi_0 (e^{A r} A v) \mathd r + B_t
     (v) + \int_0^t \Xi (h (r, A v)) \mathd r \]
  \[ = \Psi_0 (v) + B_t (v) + \int_0^t \Psi_r (A v) \mathd r \]
  as required by equation~{\eqref{eq:linear}} for our assumption on $F$.
\end{proof}

\begin{remark}
  In particular this shows that if $(\Psi_0, \Xi)$ is a Grassmann Gaussian
  process then also the solution $\Psi$ to the SDE~{\eqref{eq:sde}}, is a
  Grassmann Gaussian process compatible with $(\Psi_0, \Xi)$. In order to
  prove this fact it is sufficient to prove that any product of the form
  \begin{equation}
    \omega (\Psi_{t_1} (v_1) \cdots \Psi_{t_{2 k}} (v_2)) \label{eq:product}
  \end{equation}
  can be computed using the products (weighted with a suitable sign due to the
  anti-commutation of $\Psi$) of the covariance $\omega (\Psi_{t_i} (v_i)
  \Psi_{t_j} (v_j))$ for any $i, j = 1, \ldots, 2 k$. This property of the
  expectation~{\eqref{eq:product}} follows from the fact that
  \[ \begin{array}{c}
       \omega ((\Psi_{t_1} (v_1) - \Psi_0) \cdots (\Psi_{t_{2 k}} (v_n) -
       \Psi_0)) = \omega (\Xi (h (t_1, v_1)) \cdots \Xi (h (t_{2 k}, v_{2
       k}))) =\\
       = \sum_{\sigma} (- 1)^{\sigma} \prod_{i = 1}^k \int_0^{\min (t_i, t_j)}
       \langle e^{A (t_i - s)} v_{\sigma (2 i - 1)}, C e^{A (t_j - s)}
       v_{\sigma (2 i)} \rangle_V \mathd s
     \end{array} \]
  where $C$ is the covariance of $\Xi$. Using the fact that $\Psi_0$ is a
  Gaussian random variable independent of $\Xi$, and so $\Psi_0 (e^{A (t - s)}
  v)$ is also a Gaussian random variable independent of $\Xi$, we obtain the
  Gaussian behavior of $\Psi_t$. The compatibility of $\Psi_t$ with $(\Psi_0,
  \Xi)$ follows from the fact that $\Psi$ is a linear function of $(\Psi_0,
  \Xi) .$
\end{remark}

Let us now study the family of random variables
\[ \Psi_t^s (v) = \Xi (e^{A (t - \cdot)} v), \qquad t \in \mathbb{R}. \]
In the next proposition, we will show that (under some specified assumptions
on $A$) this represents the stationary solution to the SDE {\eqref{eq:linear}}
(this motivates our putting the upper index $s$ to $\Psi^s_t$).

\begin{proposition}
  \label{proposition_linearlimit}Assume that all eigenvalues of $A$ have
  strictly negative real part less or equal than $- \lambda_A$, where
  $\lambda_A > 0$. Then, we have $\sup_{t \in \mathbb{R}} \| \Psi^s_t
  \|_{\mathcal{G} (V)} < + \infty$, furthermore for any $G \in \Lambda V$ and
  any $t \in \mathbb{R}$ we get
  \begin{equation}
    \omega (G (\Psi^s_t)) = \omega (G (X^A)), \label{eq:G}
  \end{equation}
  where $X^A \in \mathcal{G} (V)$ is a Gaussian random variable with
  covariance $C_A$ given by
  \begin{equation}
    C_A \assign \int^{\infty}_0 e^{A^{\Tau} s} C e^{A s} \mathd s.
    \label{eq:ca}
  \end{equation}
  where $A^{\Tau}$ denotes the transpose matrix.
\end{proposition}

\begin{remark}
  Expression {\eqref{eq:G}} in particular shows that the independence of $t$
  of the expectation $\omega (G (\Psi^s_t))$, expressing the stationarity of
  $\Psi^s_t$. 
\end{remark}

\begin{proof*}{Proof of Proposition \ref{proposition_linearlimit}}
  We have easily by the definition of $\Psi^s_t$:
  \[ \| \Psi^s_t \|_{\mathcal{G} (V)} \lesssim \| \mathbb{I}_{(- \infty, t]}
     (\cdot) e^{A (t - \cdot)} \|_{L^2 (\mathbb{R})} \lesssim
     \frac{1}{\lambda_A}, \]
  where $\lesssim$ stands for inequality modulo some appropriate positive
  constant. The random variable $\Psi^s_t$ is Gaussian so it is completely
  characterized (in term of moments) by its covariance. Note that the
  covariance of $\Psi_t^s$ can be easily computed as follows
  \[ \omega (\Psi_t^s (v) \Psi_t^s (w)) = \omega (\Xi (\mathbb{I}_{(- \infty,
     t]} (\cdot) e^{A (t - \cdot)} v) \Xi (\mathbb{I}_{(- \infty, t]} (\cdot)
     e^{A (t - \cdot)} w)) = \]
  \[ = \langle \mathbb{I}_{(- \infty, t]} (\cdot) e^{A (t - \cdot)} v,
     C\mathbb{I}_{(- \infty, t]} (\cdot) e^{A (t - \cdot)} w \rangle =
     \int^0_{- \infty} \langle e^{- A^{\Tau} \tau} v, C e^{- A \tau} w \rangle
     \mathd \tau = \langle v, C_A w \rangle \]
  with $C_A$ given by~{\eqref{eq:ca}}. The appearance of the transposition is
  due to the fact that $V$ is a \tmtextit{real} pre-Hilbert space.
\end{proof*}

Let us point out that
\begin{equation}
  A^{\Tau} C_A + C_A A = - C \label{eq:relation-ca}
\end{equation}
by a simple integration by parts applied to~{\eqref{eq:ca}}.

\begin{remark}
  Suppose that $A$ commutes with $C$, then we have that $\Psi_t^s$ is for all
  $t \in \mathbb{R}$ a Grassmann Gaussian process with correlation $C_A = - (A
  + A^{\Tau})^{- 1} C$ (the operator $A$ is invertible since all its
  eigenvalues has strictly negative real part).
\end{remark}

\subsection{Existence and uniqueness for general drift}

\begin{theorem}
  \label{theorem_picard}For any $\Psi_0 \in \mathcal{G} (V)$ compatible with a
  given Brownian motion $B$, there exists $T > 0$ such that
  equation~{\eqref{eq:sde}} (for general odd $F$) admits a unique solution
  $\Psi \in C^0 ([0, T] ; \mathcal{G} (V))$. Moreover $\Psi$ is compatible
  with $\Psi_0 \oplus B$.
\end{theorem}

\begin{proof}
  We are going to construct a solution via Picard's iteration. Let $\Psi^0_t =
  \Psi_0$ for all $t \geqslant 0$ and define the $n + 1$-th Picard's iteration
  by
  \[ \Psi_t^{n + 1} (v_{\alpha}) = \Psi_0 (v_{\alpha}) + \int_0^t \Psi_s^n (F
     (v_{\alpha})) \mathd s + B_t (v_{\alpha}), \qquad t \in [0, T], \]
  for any basis element $v_{\alpha} \in V$. Then $\Psi^0 \in C^0 ([0, T] ;
  \mathcal{G} (V))$ and since $\Psi_0$ is compatible with $B$ we have that,
  for all $t \geqslant 0$, $\Psi^1_t$ belong to the Grassmann algebra
  generated by $\Psi_0$ and $B$ in $\mathcal{A}$ and there $\Psi^1_t
  (v_{\alpha})$ is an odd element for all $\alpha$. Therefore $\Psi^1_t \in
  \mathcal{G} (V)$ for all $t \geqslant 0$ and also $\Psi^1_t$ is compatible
  with $\Psi_0 \oplus B$. Since this is true for all $t$, by approximation in
  the operator norm we see that $\Psi^1 \in C^0 ([0, T] ; \mathcal{G} (V))$
  and that $\Psi^1$ is compatible with $\Psi_0 \oplus B$. By induction we
  prove then that $\Psi^n \in C^0 ([0, T] ; \mathcal{G} (V))$ for all $n
  \geqslant 0$ and that it is compatible with $\Psi_0 \oplus B$.
  
  Now observe that $\Psi_s^n (F (v_{\alpha}))$ is a polynomial function of
  $(\Psi^{n, \alpha}_s)_{\alpha}$ and therefore by a standard contraction
  argument in the Banach space $\mathcal{A}$ we obtain that $(\Psi^n)_n$
  converges, as $n \rightarrow + \infty$, in $C^0 ([0, T] ; \mathcal{A})$ for
  some positive $T$ which depends only on $F$, on $\sup_{\alpha} \|
  \Psi_0^{\alpha} \|_{\mathcal{A}}$ and on $\sup_{0 \leqslant t \leqslant T,
  \alpha} \| B_t^{\alpha} \|_{\mathcal{A}}$. If we call $\Psi$ the limit we
  have that $\Psi \in C^0 ([0, T] ; \mathcal{G} (V))$ by the first part of the
  proof and we deduce easily that $\Psi$ is compatible with $\Psi_0 \oplus B$.
\end{proof}

\begin{theorem}
  \label{th:global-sol}The solution to equation~{\eqref{eq:sde}} exists for
  all times.
\end{theorem}

\begin{proof}
  We will have global existence as soon as we can rule out explosions, that is
  prove that for all $t \geqslant 0$ we have $\sup_{\alpha} \| \Psi_t^{\alpha}
  \|_{\mathcal{A}} < \infty$ Let $\Theta_t (v) = \Psi_t (v) - B_t (v)$ and
  extend this map to an homeomorphism of $\Lambda V$ into $\mathcal{A}$.
  Observe that
  \[ \Theta_t (v) = \Theta_0 (v) + \int_0^t m_{\mathcal{A}} [(\Theta_s \otimes
     B_s) (\Delta F (v_{\alpha}))] \mathd s, \]
  recall that $m_{\mathcal{A}} : \mathcal{A} \otimes \mathcal{A} \rightarrow
  \mathcal{A}$ denotes the multiplication in $\mathcal{A}$. In particular, if
  we consider a (linear) basis of $\Lambda V$ denoted by $(e_A)_A$ and let
  $\Theta^A_t \assign \Theta_t (e_A)$ and $B^A_t \assign B_t (e_A)$ we have
  \[ \Theta_t^A = \Theta_0^A + \int_0^t \sum_B h_{A, B, C} \Theta_s^B B_s^C
     \mathd s, \]
  for real coefficients $(h_{A, B, C})_{A, B, C}$. This is a
  finite-dimensional, linear system of non-autonomous ODEs in $\mathcal{L}
  (H)$ for $(\Theta_t^A)_A$ ($\mathcal{H}$ being the Hilbert space in which
  the elements of $\mathcal{A}$ act, see Section~\ref{sec:grassmann-rv}), and
  \[ \sum_A \| \Theta_t^A \|_{\mathcal{A}} \leqslant \sum_A \| \Theta_0^A
     \|_{\mathcal{A}} + | h | \int_0^t \sum_B \| \Theta_s^B \|_{\mathcal{A}}
     \sum_C \| B_s^C \|_{\mathcal{A}} \mathd s, \]
  with $| h | \assign \sup_{A, B, C} | h_{A, B, C} |$. Therefore by Gronwall
  inequality
  \[ \sum_A \| \Theta_t^A \|_{\mathcal{A}} \leqslant \sum_A \| \Theta_0^A
     \|_{\mathcal{A}} \exp \left( | h | \int_0^t \sum_C \| B_s^C
     \|_{\mathcal{A}} \mathd s \right) < \infty \]
  for all $t \geqslant 0$. In particular $\sup_{\alpha} \| \Psi_t (v_{\alpha})
  \|_{\mathcal{A}} \leqslant \sup_{\alpha} \| \Theta_t (v_{\alpha})
  \|_{\mathcal{A}} + \sup_{\alpha} \| B_t (v_{\alpha}) \|_{\mathcal{A}} <
  \infty$ for all $t \geqslant 0$ and the solution to
  equation~{\eqref{eq:sde}} exists for all times.
\end{proof}

\subsection{An It{\^o} formula for solutions of SDEs}

We want to prove a kind of It{\^o} formula for the solution to
equation~{\eqref{eq:sde}}. In order to do that let $Q_C \in \mathcal{L}
(\Lambda V \otimes V \otimes V, \Lambda V)$ be given by
\begin{equation}
  Q_C (f \otimes v \otimes w) = \langle v, C w \rangle f, \qquad f \in \Lambda
  V, v, w \in V. \label{eq:QC}
\end{equation}
For $F \in \mathcal{L} (V, \Lambda V)$ and $G \in \Lambda V \otimes V$ we
define $G \cdot F \in \Lambda V$ by extending linearly
\[ (f \otimes v) \cdot F = f F (v), \qquad f \in \Lambda V, v \in V, \]
where on both sides we have multiplication of elements in $\Lambda V$.

\begin{theorem}
  \label{th:ito}For the global solution to equation~{\eqref{eq:sde}} we have
  \begin{equation}
    \omega (\Psi_t (G)) = \omega (\Psi_0 (G)) + \int_0^t \omega \left( \Psi_s
    \left( \partial_R G \cdot F + \frac{1}{2} Q_C (\partial_R^2 G) \right)
    \right) \mathd s \label{eq:ito-3}
  \end{equation}
  for all $G \in \Lambda V$ and $t \geqslant 0$.
\end{theorem}

\begin{proof}
  Fix any $T \geqslant 0$. Note that
  \[ \| B_t - B_s \|_{\mathcal{G} (V)} \lesssim | t - s |^{1 / 2}, \]
  for all $0 \leqslant s < t \leqslant T$. From that it follows that, if $F$
  in {\eqref{eq:sde}} is a polynomial of degree $\deg (F)$ by
  Theorem~\ref{theorem_function2}, we have
  \begin{equation}
    \| (\Psi_t - \Psi_s) - (B_t - B_s) \|_{\mathcal{G} (V)} = \left\| \int_s^t
    F (\Psi_s) \mathd s \right\|_{\mathcal{G} (V)} \lesssim (1 + \| \Psi
    \|_{[0, T]})^{\deg (F)} | t - s |, \label{eq:holderreg}
  \end{equation}
  where $\| \Psi \|_{[0, T]} \assign \sup_{t \in [0, T]} \| \Psi_s
  \|_{\mathcal{G} (V)}$ which, by Theorem~\ref{th:global-sol} and the
  continuity of $\Psi$ with respect to the time $t$, is finite for every $T
  \in \mathbb{R}_+$. By Taylor formula (see Lemma \ref{lemma:taylor}, Theorem
  \ref{theorem_function1} and Theorem \ref{theorem_function2}) applied to
  polynomial $G$ of degree $\deg (G)$ we have
  \begin{equation}
    \begin{array}{c}
      \nobracket \| G (\Psi_r) - G (\Psi_s) - (\Psi_s \otimes (\Psi_r -
      \Psi_s)) \partial_R G + \nobracket \nobracket\\
      \left. \left. + \frac{1}{2} (\Psi_s \otimes (\Psi_r - \Psi_s) \otimes
      (\Psi_r - \Psi_s)) (\partial_R \partial_R G) \right\|
      \right._{\mathcal{A}} \lesssim (1 + \| \Psi \|_{[0, T]})^{\deg (F) (\deg
      (G) - 2)} | r - s |^{3 / 2}
    \end{array} \label{eq:holderreg2}
  \end{equation}
  since by {\eqref{eq:holderreg}} $\| \Psi_t - \Psi_s \|_{\mathcal{G} (V)}
  \lesssim (1 + \| \Psi \|_{[0, T]})^{\deg (F)} | t - s |^{1 / 2}$. In a
  similar way it is possible to obtain
  \begin{equation}
    \| \langle \partial_R G (\Psi_u), F (\Psi_u) \rangle - \langle \partial_R
    G (\Psi_s), F (\Psi_u) \rangle \|_{\mathcal{A}} \lesssim (1 + \| \Psi
    \|_{[0, T]})^{\deg (F) (\deg (G) - 2)} | u - s |^{1 / 2}
    \label{eq:holderreg3}
  \end{equation}
  On the other hand, by {\eqref{eq:holderreg}}, we have
  \begin{equation}
    \begin{array}{lll}
      \omega ((\Psi_s \otimes (\Psi_r - \Psi_s)) \partial_R G) & = & \omega
      \left( \int_s^r \langle \partial_R G (\Psi_s), F (\Psi_u) \rangle \mathd
      u \right)\\
      &  & + \omega ((\Psi_s \otimes (B_r - B_s)) (\partial_R G))\\
      & = & \omega \left( \int_s^r \langle \partial_R G (\Psi_u), F (\Psi_u)
      \rangle \mathd u \right)\\
      &  & + \omega (\langle \partial_R G (\Psi_s), (B_r - B_s) \rangle)\\
      &  & + O ((1 + \| \Psi \|_{[0, T]})^{\deg (F) (\deg (G) - 2)} | r - s
      |^{3 / 2})
    \end{array} \label{eq:holderreg4}
  \end{equation}
  In a similar way, using inequality {\eqref{eq:holderreg}} we obtain
  \begin{equation}
    \begin{array}{c}
      \omega ((\Psi_s \otimes (\Psi_r - \Psi_s) \otimes (\Psi_r - \Psi_s))
      (\partial_R^2 G)) = \omega ((\Psi_s \otimes (B_r - B_s) \otimes (B_r -
      B_s)) (\partial_R^2 G))\\
      + O ((1 + \| \Psi \|_{[0, T]})^{\deg (F) (\deg (G) - 2)} | r - s |^{3 /
      2}) .
    \end{array}
  \end{equation}
  Furthermore using the fact that $\partial_R G (\Psi_s)$ is independent of
  $(B_r - B_s)$ and $((B_r - B_s) \otimes (B_r - B_s))$ is independent of
  $\partial_R^2 G (\Psi_s)$ (this is due to the fact that \ $\Psi_s$ is a
  function of $\nobracket \{ \Psi_0, B_k | \nobracket \nobracket k \in [0, s]
  \}$ and, moreover, $B$ is independent of $\Psi_0$ and has independent
  increments, we have:
  \[ \omega (\langle \partial_R G (\Psi_s), (B_r - B_s) \rangle) = \langle
     \omega ((B_r - B_s)), \omega (\partial_R G (\Psi_s)) \rangle = 0 \]
  where on the r.h.s. we understand that averages w.r.t. the state $\omega$
  are taken componentwise, and
  \[ \omega [\langle \partial_R \partial_R G (\Psi_s), (B_r - B_s) \otimes
     (B_r - B_s) \rangle_{V \otimes V}] \]
  \[ = \langle \omega [(\partial_R \partial_R G (\Psi_s))], \omega [((B_r -
     B_s) \otimes (B_r - B_s))] \rangle \]
  \[ = (r - s) \omega [Q_C (\partial_R^2 G (\Psi_s))] . \]
  By taking a partition $\{ t_i \}_{i \in \{ 0, \ldots, n \}}$, of diameter
  $\rho_n \rightarrow 0$ as $n \rightarrow + \infty$, of $[0, t]$ and
  exploiting inequality {\eqref{eq:holderreg2}} we obtain
  \[ \begin{array}{lll}
       \omega (G (\Psi_t)) & = & \omega (G (\Psi_0)) + \sum_{i = 0}^{n + 1}
       \omega (G (\Psi_{t_i})) - \omega (G (\Psi_{t_{i - 1}}))\\
       & = & \omega (G (\Psi_0)) + \sum_{i = 1}^{n + 1} \omega ((\Psi_{t_{i -
       1}} \otimes (\Psi_{t_i} - \Psi_{t_{i - 1}})) \partial_R G)\\
       &  & + \frac{1}{2} \omega ((\Psi_s \otimes (\Psi_r - \Psi_s) \otimes
       (\Psi_r - \Psi_s)) (\partial_R^2 G)) + O (\rho_n^{1 / 2})\\
       & = & \omega (G (\Psi_0)) + \omega \left( \int_0^t \partial_R G
       (\Psi_u) F (\Psi_u) \mathd u \right)\\
       &  & + \frac{1}{2} \sum_{i = 1}^{n + 1} (t_i - t_{i - 1}) \omega [Q_C
       (\partial_R^2 G (\Psi_{t_i}))] + O (\rho_n^{1 / 2}),
     \end{array} \]
  where the constants in $O$ are proportional to $(1 + \| \Psi \|_{[0,
  T]})^{\deg (F) \deg (G)}$ and do not depend on the partition. Taking the
  limit $n \rightarrow + \infty$, and so $\rho_n \rightarrow 0$, and using the
  fact that $\omega [Q_C (\partial_R^2 G (\Psi_s))]$ is continuous in $s$ ($G
  (\Psi_s)$ being a continuous function from $\mathbb{R}_+$ to $\mathcal{A}$
  since $\Psi \in C^0 ([0, T], \mathcal{G} (V))$) we obtain the thesis.
\end{proof}

\subsection{Invariant measures}\label{sec:inv-meas}

Consider an equation of the following kind
\begin{equation}
  \Psi_t (v) = \Psi_0 (v) + \int_0^t (\Psi_s (A v) + \Psi_s (F (v))) \mathd s
  + B_t (v), \qquad t \geqslant 0, v \in V, \label{eq:sde-langevin-0}
\end{equation}
for $\Psi \in C^0 ([0, T] ; \tmop{Hom} (V, \mathcal{A}))$, $F \in \tmop{Hom}
(V, \Lambda V)$, $A \in \mathcal{L} (V)$, $B$ a Grassmann Brownian motion with
correlation $C$ and with initial condition $\Psi_0$ distributed as a Grassmann
Gaussian independent of $B$ with correlation $C_A$ (compare to
equation~{\eqref{eq:ca}}). Fix $U \in \Lambda V$ even, and consider the linear
functional $\tilde{\omega} : \mathcal{A} \rightarrow \mathbb{C}$ given by
\begin{equation}
  \tilde{\omega} (\cdot) \assign \omega (\cdot e^{- 2 U (\Psi_0)}),
  \label{eq:gibbs-U}
\end{equation}
where $\omega$ is, as before, is the chosen positive state on the
$\asterisk$-algebra $\mathcal{A} = \mathcal{L} (\mathcal{H})$.

\begin{definition}
  \label{definition_invariant}We say that the functional $\tilde{\omega}$ is
  an invariant measure for the equation~{\eqref{eq:sde-langevin-0}} if for any
  $G \in \Lambda V$ and any $t \in \mathbb{R}_+$ we have
  \begin{equation}
    \tilde{\omega} (G (\Psi_t)) = \tilde{\omega} (G (\Psi_0)) .
    \label{eq:constant}
  \end{equation}
\end{definition}

\begin{remark}
  The functional is invariant in the sense of
  Definition~\ref{definition_invariant}, only if the expectations of the
  form~{\eqref{eq:constant}} are constant when expressions of the form $G
  (\Psi_t)$ are evaluated. This means in particular that only polynomials in
  $\Psi_t$ are considered.
\end{remark}

There is a simple condition, analogous to the Fokker--Planck equation for
commutative SDEs, for checking the invariance of the functional
$\tilde{\omega}$. Here it is:

\begin{lemma}
  \label{lemma_invariant}If $\Psi$ is a solution
  to~{\eqref{eq:sde-langevin-0}} and
  \begin{equation}
    \tilde{\omega} (\mathcal{L}G (\Psi_0)) = 0, \qquad G \in \Lambda V
    \label{eq:inf-invariance},
  \end{equation}
  for
  \[ \mathcal{L}G \assign \partial_R G \cdot F + \frac{1}{2} Q_C (\partial_R^2
     G), \]
  then $\tilde{\omega}$ is an invariant measure in the sense of
  Definition~\ref{definition_invariant}. In other words,
  eq.~{\eqref{eq:inf-invariance}} is a sufficient condition for having an
  invariant measure.
\end{lemma}

\begin{proof}
  Consider $G = G_A = v_A$ where $(v_A)_A$ is a (finite) linear basis of
  $\Lambda V$. Then we have $\mathcal{L} (G_A) = \sum_B \kappa_A^B v_B$ for a
  suitable family $(\kappa_A^B)_{A, B}$ of constants in $\mathbb{R}$. Write
  $\mathfrak{P}_A (t) \assign \tilde{\omega} (G_A (\Psi_t)) = \omega (G_A
  (\Psi_t) e^{- 2 U (\Psi_0)})$, due to definition {\eqref{eq:gibbs-U}}. By
  It{\^o} formula~{\eqref{eq:ito-3}} we have that $\mathfrak{P}_A$ is the
  unique solution of the system of ODEs
  \begin{eqnarray}
    \partial_t \mathfrak{P}_A (t) & = & \partial_t \omega (G_A (\Psi_t) e^{- 2
    U (\Psi_0)}) = \omega ((\mathcal{L} G_A) (\Psi_t) e^{- 2 U (\Psi_0)})
    \nonumber\\
    & = & \sum_B \kappa_A^B \omega (G_B (\Psi_t) e^{- 2 U (\Psi_0)}) = \sum_B
    \kappa^B_A \mathfrak{P}_B (t), \label{eq:ODE1} 
  \end{eqnarray}
  with initial condition $\mathfrak{P}_A (0) = \omega (G_A (\Psi_0) e^{- 2 U
  (\Psi_0)})$. On the other hand condition~{\eqref{eq:inf-invariance}} implies
  that
  \[ \sum_B \kappa^B_A \mathfrak{P}_B (0) = \sum_B \kappa^B_A \omega (G_B
     (\Psi_0) e^{- 2 U (\Psi_0)}) = \omega (\mathcal{L}G (\Psi_0) e^{- 2 U
     (\Psi_0)}) = 0 \]
  which means that $(\mathfrak{P}_A (t))_A$ is constantly zero by uniqueness
  of solution to the ODEs system~{\eqref{eq:ODE1}}. Since the expectations of
  the form~{\eqref{eq:constant}} are linear combinations of $(\mathfrak{P}_A
  (t))_A$ the proof is complete.
\end{proof}

The following theorem gives a sufficient condition on the
SDE~{\eqref{eq:sde-langevin-0}} to have the invariant
measure~{\eqref{eq:gibbs-U}}.

\begin{theorem}
  \label{th:sq-finite}For any even $U \in \Lambda_{\tmop{even}} V$, the SDE
  \begin{equation}
    \Psi_t (v) = \Psi_0 (v) + \int_0^t (\Psi_s (A v) + \Psi_s (\langle C
    \partial_R U, v \rangle) \nobracket \mathd s + B_t (v) \qquad t \geqslant
    0, v \in V, \label{eq:sde-langevin}
  \end{equation}
  with $B$ a Brownian motion with correlation $C$ and $\Psi_0$ an independent
  Gaussian initial condition with correlation $C_A$ (defined in equation
  {\eqref{eq:ca}}) has $\tilde{\omega}$ as invariant measure provided
  \begin{equation}
    A^{\Tau} C_A - C_A A = 0. \label{eq:gradient}
  \end{equation}
\end{theorem}

\begin{proof}
  The proof is a consequence of Lemma~\ref{lemma_invariant}. Denote
  $\omega^{\Psi_0} (G) \assign \omega (G (\Psi_0))$ the law of $\Psi_0$, and
  observe that for $H \in \Lambda V \otimes V$ we have the integration by
  parts formula
  \[ \omega ((\Psi_0 \otimes \Psi_0) (H)) = \omega (\Psi_0 (Q_{C_A}
     (\partial_R H))), \]
  where $Q_C$ is defined in {\eqref{eq:QC}}. Then for all $G \in \Lambda V$
  and $U \in \Lambda_{\tmop{even}} V$ we have
  \[ \omega ((\Psi_0 \otimes \Psi_0) (A (\partial_R G) e^{- 2 U (\Psi_0)})) =
     \omega^{\Psi_0} (Q_{C_A} (\partial_R (A \partial_R G e^{- 2 U
     (\Psi_0)}))) \]
  \[ = \omega^{\Psi_0} (Q_{C_A} ((A \otimes 1) \partial_R \partial_R G) e^{- 2
     U (\Psi_0)}) - 2 \omega^{\Psi_0} (Q_{C_A} ((A \partial_R G) \otimes
     \partial_R U) e^{- 2 U (\Psi_0)}) . \]
  Observe that using the equality~{\eqref{eq:relation-ca}} we have
  \[ Q_{C_A} ((A \otimes 1) \partial_R \partial_R G) = \tmop{Tr} ((A \otimes
     C_A) \partial_R \partial_R G) \]
  \[ = \frac{1}{2} \tmop{Tr} ((C_A A + A^{\Tau} C_A) \partial_R \partial_R G)
     = - \frac{1}{2} Q_C (\partial_R \partial_R G), \]
  where we have used the trace $\tmop{Tr}$ on $V$ given by $\tmop{Tr} (v
  \otimes w) = \langle v, w \rangle$. Moreover, from~{\eqref{eq:gradient}} we
  have also
  \[ Q_{C_A} ((A \partial_R G) \otimes \partial_R U) = \tmop{Tr} ((A
     \partial_R G) \otimes (C_A \partial_R U)) = \tmop{Tr} (\partial_R G
     \otimes (A^{\Tau} C_A \partial_R U)) \]
  \[ = \tmop{Tr} \left( \partial_R G \otimes \left( \frac{A^{\Tau} C_A + C_A
     A}{2} \partial_R U \right) \right) + \tmop{Tr} \left( \partial_R G
     \otimes \left( \frac{A^{\Tau} C_A - C_A A}{2} \partial_R U \right)
     \right) \]
  \[ = - \frac{1}{2} \tmop{Tr} (\partial_R G \otimes (C \partial_R U)) = -
     \frac{1}{2} \langle \partial_R G, C \partial_R U \rangle, \]
  since we assumed that $A^{\Tau} C_A - C_A A = 0$ and
  used~{\eqref{eq:relation-ca}} again. Therefore
  \[ \omega ((\Psi_0 \otimes \Psi_0) (A (\partial_R G) e^{- 2 U (\Psi_0)})) =
     - \frac{1}{2} \omega^{\Psi_0} (Q_C (\partial_R^2 G) e^{- 2 U (\Psi_0)}) +
     \omega^{\Psi_0} (\langle \partial_R G, C \partial_R U \rangle e^{- 2 U
     (\Psi_0)}) . \]
  Now observing that
  \[ \mathcal{L}G (\Psi_0) = (\Psi_0 \otimes \Psi_0) (A \partial_R G) + \Psi_0
     \left( \langle \partial_R G, C \partial_R U \rangle + \frac{1}{2} Q_C
     (\partial_R \partial_R G) \right) \]
  we deduce that
  \[ \omega (\mathcal{L}G (\Psi_0) e^{- 2 U (\Psi_0)}) = \omega ((\Psi_0
     \otimes \Psi_0) (A \partial_R G) e^{- 2 U (\Psi_0)}) + \]
  \[ + \omega^{\Psi_0} \left( \left( \langle \partial_R G, C \partial_R U
     \rangle + \frac{1}{2} Q_C (\partial_R \partial_R G)  \right) e^{- 2 U
     (\Psi_0)} \right) = 0. \]
\end{proof}

\subsection{Long-time behavior for small non-linearity}\label{sec:long-time}

We will now investigate the existence of solutions which are globally bounded
in time. In the commutative setting, non-linear equations can have globally
bounded solutions only if the non-linearity stays uniformly small or if it
shows some coercivity. Like all notions of positivity, also coercivity however
does not apply well in the Grassmann setting. The only kind of coercive term
we have identified is a linear drift with a negative sign. This is a very mild
coercive term, but it turns out to be enough provided the non-linearity is
small enough.

\

Consider the equation
\begin{equation}
  \Psi_t (v) = \Psi_0 (v) + \int_0^t (\Psi_s (A v) + \lambda \Psi_s (F (v)))
  \mathd s + B_t (v), \qquad t \geqslant 0, v \in V, \label{eq:sde-lambda}
\end{equation}
for $\Psi \in C ([0, T] ; \tmop{Hom} (V, \mathcal{A}))$, with $\lambda
\geqslant 0$ and $F \in \tmop{Hom} (V, \Lambda V)$ where (as in
Section~\ref{sec:ou}) $A$ is an operator with eigenvalues having strictly
negative real part less then $- \lambda_A < 0$. In this setting we introduce
the notion of stationary solution to equation~{\eqref{eq:sde-lambda}},
extending the one we defined for the linear case in Section~\ref{sec:ou}.

\begin{definition}
  We say that $\Psi^s_t \in C^0 (\mathbb{R}, \mathcal{G} (V))$ is a stationary
  solution to equation~{\eqref{eq:sde-lambda}} (extended to all $t \in
  \mathbb{R}$) of norm at most $K \in \mathbb{R}_+$ if the following two
  conditions hold:
  \begin{enumerate}
    \item $\sup_{t \in \mathbb{R}} \| \Psi^s_t \|_{\mathcal{G} (V)} < K$,
    
    \item $(\Psi^s_t)_{t \in \mathbb{R}}$ solves the following integral
    equation
    \begin{equation}
      \Psi^s_t = \lambda \int_{- \infty}^t \Psi^s_{\tau} (F (e^{A \tau} v))
      \mathd \tau + B^A_t \label{eq:stationary} \qquad t \in \mathbb{R}
    \end{equation}
    where \ $B^A_t (v) = \Xi (e^{A (t - \cdot)} (v))$.
  \end{enumerate}
\end{definition}

Hereafter we write
\[ L (A, C) = \sup_{t \in \mathbb{R}} \| B^A_t \|_{\mathcal{G} (V)}, \]
and we denote by $\mathfrak{K}_K \subset C^0 (\mathbb{R}, \mathcal{G} (V))$
the set of $\Psi \in C^0 (\mathbb{R}, \mathcal{G} (V))$ such that
\[ \sup_{t \in \mathbb{R}} \| \Psi_t \|_{\mathcal{G} (V)} \leqslant K_{}, \]
for some constant $K > 0$.

\begin{theorem}
  \label{theorem_stationary}Assume $K \geqslant 3 L (A, C)$ and suppose that
  all the eigenvalues of $A$ have negative real part that are less or equal
  than $- \lambda_A$, then there exists $\lambda_0 (K, L, F)$ depending on
  $K$, $L (A, C)$ and $F$ such that if $| \lambda | \leqslant 2 \lambda_0 (K,
  L, F)$ there exists a unique stationary solution of norm at most $K$ to
  equation~{\eqref{eq:sde-lambda}}.
\end{theorem}

\begin{proof}
  It is a simple application of the Banach fix point theorem by noting that
  the map
  \[ \mathcal{K}_{\lambda} (\Psi)_t \assign \lambda \int_{- \infty}^t
     \Psi_{\tau} (F (e^{A (t - \tau)} v)) \mathd \tau + B^A_t, \qquad t \in
     \mathbb{R}, \]
  is a contraction in $\mathfrak{K}_K$. Indeed, if $\Psi \in \mathfrak{K}_K$
  we have, by Theorem~\ref{theorem_function2},
  \[ \| \mathcal{K}_{\lambda} (\Psi)_t \|_{\mathcal{G} (V)} \leqslant |
     \lambda |  \| F \|_{\Lambda_{\pi} V} (1 + K)^{\deg (F)} \int_{- \infty}^t
     e^{- \lambda_A (t - \tau)} \mathd \tau + L (A, C), \]
  where $\| F \|_{\Lambda_{\pi} V}$ is defined by~{\eqref{eq:normfunction}}.
  
  If we request that
  \begin{equation}
    2 \lambda_0 \frac{\| F \|_{\Lambda_{\pi} V} (1 + K)^{\deg (F)}}{\lambda_A}
    + 2 L < K, \label{eq:contraction2}
  \end{equation}
  (where $\lambda_0 \assign \lambda_0 (K, L, F)$) for $| \lambda | \leqslant 2
  \lambda_0$ we have that $\mathcal{K}_{\lambda}$ maps $\mathfrak{K}_K$ into
  itself. Furthermore
  \[ \| \mathcal{K}_{\lambda} (\Psi_1)_t - \mathcal{K}_{\lambda} (\Psi_2)_t
     \|_{\mathcal{G} (V)} \leqslant \frac{| \lambda |}{\lambda_A} (\deg (F))
     \| F \|_{\Lambda_{\pi} V} (1 + K)^{\deg (F) - 1} \| \Psi_{1, t} -
     \Psi_{2, t} \|_{\mathcal{G} (V)} . \]
  Therefore, if we take
  \begin{equation}
    \frac{2 \lambda_0}{\lambda_A} \| F \|_{\Lambda_{\pi} V} (\deg (F)) (1 +
    K)^{\deg (F) - 1} < 1 \label{eq:contraction}
  \end{equation}
  the map $\mathcal{K}_{\lambda}$ is a strict contraction. The fact that
  $\mathfrak{K}_K \neq \emptyset$ will be a consequence of the Lemma
  (\ref{lemma_bound}) below.
\end{proof}

\begin{remark}
  \label{remark_lambda0}The function $\lambda_0$ in
  Theorem~\ref{theorem_stationary} can be taken to be a decreasing function of
  the parameters $L$ and $K$.
\end{remark}

We introduce an approximation for the solution to
equation~{\eqref{eq:stationary}}. Let $X \in \mathcal{G} (V)$ independent of
$B^A_t$ and consider the element $\Psi_{- T, t}^X \in C^0 (\mathbb{R},
\mathcal{G} (V))$ which is $\Psi_{- T, t}^X = X$ if $t \leqslant - T$, for
some $T > 0$, and is a solution to the SDE
\[ \Psi_{- T, t}^X = X (e^{A (t + T)} v) + \lambda \int_{- T}^t F (\Psi_{- T,
   t}^X (e^{A (t - \tau)} v)) + B^A_t - B^A_{- T} (e^{A (t + T)} v), \]
when $t > - T$.

\begin{lemma}
  \label{lemma_bound}Under the assumptions of
  Theorem~\ref{theorem_stationary}, take $\| X \|_{\mathcal{G} (V)} <
  \frac{1}{6} K$, then, for any $| \lambda | \leqslant \frac{1}{2} \lambda_0
  (K, L, F)$, we have that $\Psi_{- T, t}^X \in \mathfrak{K}_K$ for all $t \in
  \mathbb{R}$.
\end{lemma}

\begin{proof}
  Let $\tau \in \mathbb{R}$, by Theorem \ref{theorem_function2}, then we have
  the bound
  \begin{eqnarray}
    \sup_{t \in (- \infty, \tau]} \| \Psi^X_{- T, t} \|_{\mathcal{G} (V)} &
    \leqslant & e^{- \lambda_A \tau} (\| X \|_{\mathcal{G} (V)} + \| B_{- T}^A
    \|_{\mathcal{G} (V)}) + \nonumber\\
    &  & + \| F \|_{\Lambda_{\pi} V} \frac{| \lambda |}{\lambda_A} (1 +
    \sup_{t \in (- \infty, \tau]} \| \Psi_{- T, t}^X \|_{\mathcal{G}
    (V)})^{\deg (F)} + L (A, C) . \nonumber
  \end{eqnarray}
  On the other hand by definition of $\lambda_0$ we have, for any $\tau > 0$
  \[ e^{- \lambda_A \tau} \left( \frac{1}{3} K + \frac{1}{6} K \right) + \| F
     \|_{\Lambda_{\pi} V} \frac{| \lambda |}{\lambda_A} (1 + K)^{\deg (F)} + L
     (A, C) < K. \]
  Let $R_{\tau} \assign \sup_{t \in (- \infty, \tau]} \| \Psi_{- T, t}^X
  \|_{\mathcal{G} (V)}$ and consider the set $Z = \left\{ \tau \in \mathbb{R}:
  R_{\tau} \leqslant \frac{2 K}{3} \right\}$. The set $Z$ is non-empty, since
  $R_{\tau} < \frac{1}{2} K$ when $\tau \leqslant - T$, open since $\tau
  \mapsto R_{\tau}$ is an increasing continuous function of $\tau$ and closed
  since $\tau \rightarrow R_{\tau}$ is continuous. Then we must have $Z
  =\mathbb{R}$ and
  \[ \sup_{t \in \mathbb{R}} \| \Psi_{- T, t}^X \|_{\mathcal{G} (V)}
     \leqslant \frac{2}{3} K, \]
  for $| \lambda | \leqslant \lambda_0$.
\end{proof}

\begin{theorem}
  \label{theorem_approximation}Under the hypotheses of Lemma~\ref{lemma_bound}
  we have that $\Psi_{- T, t}^X$ converges to $\Psi^s_t$ in $\mathcal{G} (V)$
  as $T \rightarrow + \infty$, uniformly on compact subsets of $\mathbb{R}$.
\end{theorem}

\begin{proof}
  If $t \geqslant - T$, by Theorem \ref{theorem_function2}, we have that
  \begin{eqnarray}
    \| \Psi_{- T, t}^X - \Psi_t^s \|_{\mathcal{G} (V)} & \lesssim & e^{-
    \lambda_A (t + T)} (\| X \|_{\mathcal{G} (V)} + \| B_{- T}^A
    \|_{\mathcal{G} (V)}) + | \lambda | \| F \|_{\Lambda_{\pi} V} (1 +
    K)^{\deg (F)} \times \nonumber\\
    &  & \times \int_{- \infty}^{- T} e^{- \lambda_A (t - \tau)} \mathd \tau
    + | \lambda | (\deg (F)) \| F \|_{\Lambda_{\pi} V} (1 + K)^{\deg (F) - 1}
    \times \nonumber\\
    &  & \times \int_{- T}^t e^{- \lambda_A (t - s)} \| \Psi_{- T, \tau}^X -
    \Psi_{\tau}^s \| \mathd \tau . \nonumber
  \end{eqnarray}
  By Gronwall inequality (see, e.g., Theorem 1.3.2 in {\cite{Pachpatte1998}})
  we obtain
  \begin{eqnarray}
    \| \Psi_{- T, t}^X - \Psi_t^s \|_{\mathcal{G} (V)} & \leqslant & \left( \|
    X \|_{\mathcal{G} (V)} + \frac{K}{3} + \frac{| \lambda | \| F
    \|_{\Lambda_{\pi} V} (1 + K)^{\deg (F)}}{\lambda_A} \right) \times
    \nonumber\\
    &  & \times \left( e^{- \lambda_A (t + T) + \alpha t} \int_{- T}^t e^{-
    \alpha \tau} d \tau \right) \nonumber\\
    & \lesssim & e^{- \lambda_A (t + T)} e^{\alpha t} (e^{\alpha T} - e^{-
    \alpha t}) \lesssim e^{- (\lambda_A - \alpha) (t + T)} \nonumber
  \end{eqnarray}
  where $\alpha = | \lambda | \| F \|_{\Lambda_{\pi} V} (\deg (F)) (1 +
  K)^{\deg (F) - 1}$. Since $\lambda_A - \alpha > 0$, by the
  condition~{\eqref{eq:contraction}}, we have that
  \[ \| \Psi_{- T, t}^X - \Psi_t^s \|_{\mathcal{G} (V)} \lesssim e^{-
     (\lambda_A - \alpha) T} \]
  uniformly in $T$ when $t > P > - T$ for any such fixed $P \in \mathbb{R}$.
  This proves that $\Psi_{- T, t}^X$ converges to $\Psi_t^s$ uniformly on
  compact sets.
\end{proof}

An important consequence of Theorem~\ref{theorem_approximation} is the
following one.

\begin{theorem}
  \label{theorem_invariant2}Under the hypotheses and notations of
  Theorem~\ref{theorem_stationary} and Theorem~\ref{theorem_approximation}, if
  $\| X \|_{\mathcal{G} (V)} < \frac{1}{3} K$, $| \lambda | \leqslant
  \lambda_0 (K, L, F)$ and if $\omega (\cdot e^{- 2 U (X)})$ is an invariant
  measure for equation~{\eqref{eq:sde-lambda}} (in the sense of
  Definition~\ref{definition_invariant}), for any $t \in \mathbb{R}$ and for
  any $G \in \Lambda V$, we have that
  \[ \frac{\omega (G (X) e^{- 2 U (X)})}{\omega (e^{- 2 U (X)})} = \omega (G
     (\Psi^s_t)) . \]
\end{theorem}

\begin{proof}
  Using the previous notation, by Definition~\ref{definition_invariant} with
  $U (X)$ replacing $U (\Psi_0)$, we have that, for any $t \geqslant - T$,
  \[ \omega (G (X) e^{- 2 U (X)}) = \omega (G (\Psi_{- T, - T}^X) e^{- 2 U
     (X)}) = \omega (G (\Psi_{- T, t}^X) e^{- 2 U (X)}) . \]
  By Theorem~\ref{theorem_approximation} we can take the limit at the right
  hand side as $T \rightarrow + \infty$, obtaining
  \[ \omega (G (X) e^{- 2 U (X)}) = \omega (G (\Psi^s_t) e^{- 2 U (X)}) . \]
  Since (by construction) $X$ is independent of $B^A_t$ and since $\Psi^s_t$
  is a function only of $B^A$ we get
  \[ \omega (G (X) e^{- 2 U (X)}) = \omega (G (\Psi^s_t)) \omega (e^{- 2 U
     (X)}), \]
  from which the thesis follows.
\end{proof}

\begin{remark}
  \label{remark_invariant2}By Theorem~\ref{th:sq-finite},
  Theorem~\ref{theorem_invariant2} applies to the case $\lambda F (v) =
  \lambda \langle C \partial_R U, v \rangle$ for $\lambda$ small enough.
\end{remark}

\section{Infinite dimensional SDEs}\label{sec:infinite-dim}

We want to study SDEs of the form
\begin{equation}
  \mathd \Psi_t = (A \Psi_t + \lambda F (\Psi_t)) \mathd t + \mathd B_t
  \label{eq:infinite5}
\end{equation}
in the case where $\Psi$ takes values in $\mathcal{G} (V)$ and $V =\mathcal{S}
(\mathbb{R}^d) \otimes \mathbb{R}^n,$ or $\mathcal{S} (\mathbb{R}^d) \otimes
\mathbb{C}^n$, or $C^{\infty} (\mathbb{T}^d) \otimes \mathbb{R}^n$ for some $n
\in \mathbb{N}_0$ and with $A$ a (possibly unbounded) linear operator on $V$.
We discuss in detail the case $V = \mathcal{S} (\mathbb{R}^d) \otimes
\mathbb{R}^n$, but every reasoning can be easily generalized to the other
cases or more generally to any nuclear space of smooth functions to
Fr{\'e}chet nuclear space (see Remark \ref{remark:tensor} for the reason).

\

The main difference between the finite and infinite dimensional case for
Grassmann SDEs is that in general there is no natural topology on $\mathcal{G}
(V)$. The space $V = \mathcal{S} (\mathbb{R}^d) \otimes \mathbb{R}^n$ is a
Fr{\'e}chet nuclear space. Since not every element of $\mathcal{G} (V)$ is
continuous with respect to this topology we consider here the space
$\mathcal{G}^{- \infty} (V) \subset \mathcal{G} (V)$, the subset of
$\mathcal{G} (V)$ for which $X \in \mathcal{G}^{- \infty} (V)$ if and only if
$X$ restricted to $V \subset \Lambda V$ is a continuous linear map from $V$ to
$\mathcal{A}$. Hereafter we always consider $X \in \mathcal{G}^{- \infty} (V)$
both as an homomorphism from the algebra $\Lambda V$ into $\mathcal{A}$, and
as a continuous linear map from $V$ into $\mathcal{A}$.

\

Since $V$ is a Fr{\'e}chet nuclear space whose dual is nuclear the set of
continuous linear maps from $V$ into $\mathcal{A}$ can be identified with the
tensor product $V^{\asterisk} \hat{\otimes} \mathcal{A} = (\mathcal{S}'
(\mathbb{R}^d) \otimes \mathbb{R}^n) \hat{\otimes} \mathcal{A}$ (where
$V^{\asterisk}$ is the topological dual of $V$, the Schwartz space of tempered
distribution $\mathcal{S}' (\mathbb{R}^d) \otimes \mathbb{R}^n$ and the symbol
$\hat{\otimes}$ in the product $V^{\ast} \hat{\otimes} \mathcal{A}$ means that
we take the closure of $V^{\ast} \otimes \mathcal{A}$ with respect to any
natural topology on $V^{\ast} \otimes \mathcal{A}$ being all of them
equivalent since $V^{\ast}$ is nuclear see, e.g., Theorem 50.1 in
{\cite{Treves}}), see, e.g., Proposition 50.5 in {\cite{Treves}}. In other
words we can look upon $\mathcal{G}^{- \infty} (V)$ as the subset of elements
of $(\mathcal{S}' (\mathbb{R}^d) \otimes \mathbb{R}^n) \otimes \mathcal{A}$
that can be extended to a homomorphism from $\Lambda V$ into $\mathcal{A}$.

\begin{remark}
  \label{remark:tensor}The identification of the set of continuous linear maps
  from $V$ into $\mathcal{A}$ with $V^{\ast} \hat{\otimes} \mathcal{A}$ does
  not hold if $V$ is a generic topological vector space, for example if $V$ is
  a Banach space. Indeed in the case where $V$ is a infinite dimensional
  Banach space the space $V^{\ast} \otimes \mathcal{A}$ has not unique natural
  norm. In particular the completion $V^{\ast} \otimes_{\varepsilon}
  \mathcal{A}$ with respect the injective norm (which is the weaker of them)
  we have $V^{\ast} \otimes_{\varepsilon} \mathcal{A} \subsetneq \mathcal{L}
  (V, \mathcal{A})$ (see, e.g., {\cite{Ryan2002}}). For this reason we prefer
  to work in the case where $\mathcal{L} (V, \mathcal{A}) = V^{\ast}
  \hat{\otimes} \mathcal{A}$ which is for example when $V$ is a Fr{\'e}chet
  nuclear space. 
\end{remark}

Using a reasoning similar to the one exploited in
Lemma~\ref{lemma_completeness}, it is possible to prove that $\mathcal{G}^{-
\infty} (V)$ is a closed subset of $V^{\asterisk} \otimes \mathcal{A}$ with
respect to the strong topology (of both $V^{\asterisk} = \mathcal{S}'
(\mathbb{R}^d) \otimes \mathbb{R}^n$ and $\mathcal{A}$), i.e. if
$(X_{\nu})_{\nu}$ is a net in $\mathcal{G}^{- \infty} (V)$ converging to $X$
in the topology of $V^{\asterisk} \otimes \mathcal{A}$, then $X$ can be
extended to an homomorphism from $\Lambda V$ into $\mathcal{A}$. Working on
$\mathcal{G}^{- \infty} (V)$ is not so easy since $(\mathcal{S}'
(\mathbb{R}^d) \otimes \mathbb{R}^n) \otimes \mathcal{A}$ is not even a
Fr{\'e}chet space. For this reason we introduce now some subsets of
$\mathcal{G}^{- \infty} (V)$ (which are analogous to Besov spaces) which \
admit a stronger metric topology.

\begin{remark}
  In what fallow we use the notation if $X \in \mathcal{G}^{- \infty} (V)$ we
  write, for any $k \in \{ 1, \ldots, n \}$, the map $X^k \in \mathcal{G}^{-
  \infty} (\mathcal{S} (\mathbb{R}^n))$ as follows
  \[ X^k (f) = X (f \otimes e_k) \]
  where $\{ e_k \}_{k = 1, \ldots, n}$ is the standard orthonormal basis of
  $\mathbb{R}^n$.
\end{remark}

We recall here only the definition of the Besov norm for functions taking
values in a Banach space, further references and details are given in Appendix
\ref{appendix_Besov} (see also {\cite{Amann1997,Amann2019}}). Let $K_i \in
\mathcal{S} (\mathbb{R}^d)$, with $i \geqslant - 1$, be the functions
corresponding to the Littlewood--Paley blocks used in the definition of Besov
spaces $B^s_{p, q} \assign B^s_{p, q} (\mathbb{R}^d, \mathbb{R})$, where $s
\in \mathbb{R}$ and $p, q \in [1, + \infty]$. If $X \in \mathcal{G}^{- \infty}
(V)$ we define
\[ \Delta_i X (x) = (\Delta_i X^1 (x), \ldots, \Delta_i X^k (x)) \assign (X
   (K_i (x - \cdot) \otimes e_k))_{k = 1, \ldots, n}, \qquad x \in
   \mathbb{R}^d, i \geqslant - 1 \]
where $\{ e_k \}_{k = 1, \ldots n}$ is the standard basis of $\mathbb{R}^n$.
The function $\Delta_i X^k : \mathbb{R}^d \rightarrow \mathcal{A}$ can be
identified with an element $\Delta_i X^k \in \mathcal{S} (\mathbb{R}^d)
\hat{\otimes} \mathcal{A} = \mathcal{S} (\mathbb{R}^d, \mathcal{A})$ and we
define
\[ \| X \|_{B^s_{p, q, n} (\mathbb{R}^d, \mathcal{A})} \assign \sum_{k = 1}^n
   \left[ \sum_{i \geqslant - 1} 2^{q i s} \left( \int_{\mathbb{R}^d} \|
   \Delta_i X^k (x) \|_{\mathcal{A}}^p \mathd x \right)^{q / p}_{} \right]^{1
   / q} . \]
Given $s \in \mathbb{R}$, we say that $X \in \mathcal{G}^s (V)$ when $X \in
\mathcal{G}^{- \infty} (V)$ and $\| X \|_{B^s_{\infty, \infty} (\mathbb{R}^d,
\mathcal{A})} < + \infty$. Thus we introduce the natural metric
$d_{\mathcal{G}^s (V)} (X, Y) \assign \| X - Y \|_{\mathcal{G}^s (V)} \assign
\| X - Y \|_{B^s_{\infty, \infty, n} (\mathbb{R}^d, \mathcal{A})}$. Hereafter
we always use the notation $\mathcal{C}^s_n (\mathbb{R}^d, \mathcal{A}) =
B^s_{\infty, \infty, n} (\mathbb{R}^d, \mathcal{A})$ and $\mathcal{C}^s =
\mathcal{C}^s_1 (\mathbb{R}^d, \mathbb{R})$. Moreover $\mathcal{C}^{s
\asterisk}$ will denote the topological dual of $\mathcal{C}^s$.

\begin{lemma}
  The set $\mathcal{G}^s (V)$ with the metric $d_{\mathcal{G}^s (V)}$ is a
  complete metric space.
\end{lemma}

\begin{proof}
  The proof is analogous to that of Lemma~\ref{lemma_completeness}.
\end{proof}

The case $s > 0$ can be understood more easily. In this case the element $X
\in \mathcal{G}^s (V)$ can be identified with a vector of $k$ continuous
function $\dot{X}^k : \mathbb{R}^d \rightarrow \mathcal{A}$. Under this
identification if $f = (f_1, \ldots, f_n) \in \mathcal{S} (\mathbb{R}^d)
\otimes \mathbb{R}^n$ we have that
\[ X (f) = \sum_{k = 1}^n \int_{\mathbb{R}^d} f_k (x) \dot{X}^k (x) \mathd x,
\]
where the integral is taken in Bochner sense. If $X \in \mathcal{G}^s (V)$,
with $s > 0$, we can obtain the function $\dot{X}$ in the following way. Let
$\mathfrak{a} : \mathbb{R}^d \rightarrow \mathbb{R}$ be a smooth mollifier and
define $\mathfrak{a}_{x, \varepsilon, k} (y) = \varepsilon^{- d} \mathfrak{a}
\left( \frac{y - x}{\varepsilon} \right) \otimes e_k$, where $x, y \in
\mathbb{R}^d$ and $\varepsilon > 0$. Then we have that
\[ \dot{X}^k (x) = \lim_{\varepsilon \rightarrow 0} X
   (\mathfrak{a}_{\varepsilon, x, k}) = X (\delta_x \otimes e_k), \]
where the limit is taken with respect to the topology of $\mathcal{A}$, and
where $\delta_x$ is the Dirac delta with unit mass at the point $x \in
\mathbb{R}^d$. More generally, if $X \in \mathcal{G}^s (V)$ and $g \in
\mathcal{C}^{s \asterisk} \otimes \mathbb{R}^n$, for any sequence
$(g_{\ell})_{\ell} \subset \mathcal{S} (\mathbb{R}^d) \otimes \mathbb{R}^n$
converging to $g$ with respect to the natural topology of $\mathcal{C}^{s
\asterisk} \otimes \mathbb{R}^n$ we have that $(X (g_{\ell}))_{\ell}$
converges in $\mathcal{A}$ to a unique element
\[ X (g) \assign \lim_{n \rightarrow + \infty} X (g_n) . \]
The map $X$ now defined on all of $\mathcal{C}^{s \asterisk} \otimes
\mathbb{R}^n$ is continuous. In other words, if $X \in \mathcal{G}^s (V)$ then
it can be identified with an element of $\mathcal{L} ((\mathcal{C}^{s
\asterisk} \otimes \mathbb{R}^n), \mathcal{A})$. Furthermore it is simple to
prove, using suitable $\mathcal{S} (\mathbb{R}^d) \otimes \mathbb{R}^n$
approximations, that if $g, h \in \mathcal{C}^{s \asterisk} \otimes
\mathbb{R}^n$, we have
\begin{equation}
  X (g) X (h) = - X (h) X (g) . \label{eq:anticomgeneral}
\end{equation}
In particular we have
\[ \dot{X}^k (x_1) \dot{X}^h (x_2) = - \dot{X}^h (x_2) \dot{X}^k (x_1), \qquad
   x_1, x_2 \in \mathbb{R}^d, h, k \in \{ 1, \ldots, n \} . \]
The relation~{\eqref{eq:anticomgeneral}} has important consequences. For any
$\ell \geqslant 1$, let
\[ \begin{array}{cc}
     \Lambda^{\ell}_{\pi} (\mathcal{C}^{s \asterisk} \otimes \mathbb{R}^n)
     \subset & \underbrace{(\mathcal{C}^{s \asterisk} \otimes \mathbb{R}^n)
     \otimes_{\pi} (\mathcal{C}^{s \asterisk} \otimes \mathbb{R}^n)
     \otimes_{\pi} \cdots \otimes_{\pi} (\mathcal{C}^{s \asterisk} \otimes
     \mathbb{R}^n)}\\
     & \text{$n$ times}
   \end{array} \]
be the closure of $\Lambda^{\ell}_{\pi} (\mathcal{C}^{s \asterisk} \otimes
\mathbb{R}^n)$ in $\otimes^{\ell}_{\pi} (\mathcal{C}^{s \asterisk} \otimes
\mathbb{R}^n)$ using the natural inclusion $i_{\Lambda^{\ell} (\mathcal{C}^{s
\asterisk} \otimes \mathbb{R}^n)}$ of $\Lambda^{\ell} (\mathcal{C}^{s
\asterisk} \otimes \mathbb{R}^n)$ in $\otimes^{\ell} (\mathcal{C}^{s \ast}
\otimes \mathbb{R}^n)$ (see the discussion after Lemma \ref{lemma:taylor} for
more details). Here $\otimes_{\pi}$ is the projective tensor product of Banach
spaces whose norm is defined as follows: if $f \in \otimes^{\ell}
(\mathcal{C}^{s \asterisk} \otimes \mathbb{R}^n)$ then
\[ \| f \|_{\otimes^{\ell}_{\pi} (\mathcal{C}^{s \asterisk} \otimes
   \mathbb{R}^n)} \assign \inf \left\{ \sum_{i = 1}^k \prod_{j = 1}^{\ell} \|
   f^i_j \|_{(\mathcal{C}^{s \asterisk} \otimes \mathbb{R}^n)}, \tmop{where} f
   = \sum_{i = 1}^k f^i_1 \otimes \cdots \otimes f^i_{\ell} \right\} . \]
\begin{lemma}
  \label{remark_regularity}If $X \in \mathcal{G}^s (V)$ then it can be
  extended to a continuous homomorphism from $\bigoplus_{i = 1}^{\ell}
  \Lambda_{\pi}^i (\mathcal{C}^{s \asterisk} \otimes \mathbb{R}^n)$ into
  $\mathcal{A}$ for any $n > 0$. Here $\bigoplus_{i = 1}^{\ell}
  \Lambda_{\pi}^i (\mathcal{C}^{s \asterisk} \otimes \mathbb{R}^n)$ is
  considered as equipped with the product obtained from the standard product
  of $\bigoplus_{i = 1}^{\infty} \Lambda_{\pi}^i (\mathcal{C}^{s \asterisk}
  \otimes \mathbb{R}^n)$ via the projection onto $\bigoplus_{i = 1}^{\ell}
  \Lambda_{\pi}^i (\mathcal{C}^{s \asterisk} \otimes \mathbb{R}^n)$.
\end{lemma}

\begin{proof}
  The proof is based on the fact that $X$ can be identified with an element of
  $\mathcal{L} (\mathcal{C}^{s \asterisk}, \mathcal{A})$ and on the
  relation~{\eqref{eq:anticomgeneral}}.
\end{proof}

In the case $s > 0$, we have for example that
\[ X ((\delta_{x_1} \otimes e_{k_1}) \wedge \cdots \wedge (\delta_{x_{\ell}}
   \otimes e_{k_{\ell}})) = \dot{X}^{k_1} (x_1) \cdots \dot{X}^{k_{\ell}}
   (x_{\ell}), \qquad x_1, \ldots, x_{\ell} \in \mathbb{R}^d, k_1, \ldots,
   k_{\ell} \in \{ 1, \ldots, n \} . \]
We want now to define a natural notion of function $F$ of an element $X \in
\mathcal{G}^s (V)$. For example, if $s > 0$, and we write $Y = F_{\ell} (X)$,
we want to be able to take $F_n$ of the form
\begin{equation}
  \dot{Y}^j (x) = \sum_{k_1, \ldots, k_{\ell} = 1}^n \sum_{(h_1, \ldots,
  h_{\ell}) \in D} K^j_{(h_1, \ldots, h_{\ell}), k_1, \ldots, k_{\ell}}
  \dot{X^{k_1}} (x + h_1) \dot{X}^{k_{\ell}} (x + h_2) \cdots
  \dot{X}^{k_{\ell}} (x + h_{\ell}) \label{eq:examplefinite},
\end{equation}
where $D \subset \mathbb{R}^{d \ell}$ is some finite set and $K_{(h_1, \ldots,
h_{\ell}), k_1, \ldots, k_{\ell}}^j \in \mathbb{R}$ and $j \in \{ 1, \ldots, k
\}$. When $\ell$ is odd, the important property of the function $F_n$ defined
before is that $Y = F (X) \in \mathcal{C}^s_k (\mathbb{R}^d, \mathcal{A})$ and
furthermore for any $g, h \in \mathcal{C}^{s \asterisk}$ and $k, p \in \{ 1,
\ldots, n \}$ we have
\begin{equation}
  \begin{array}{c}
    Y^k (g) Y^p (h) = - Y^p (h) Y^k (g)\\
    Y^k (g) X^p (h) = - X^p (h) Y^k (g) .
  \end{array} \label{eq:anticom2}
\end{equation}
We want to consider some more general setting than the one considered in
equation~{\eqref{eq:examplefinite}} that maintains the important
property~{\eqref{eq:anticom2}}.

\

Consider $F_{\ell} \in \mathcal{L} (\mathcal{S} (\mathbb{R}^d) \otimes
\mathbb{R}^n, \Lambda^{\ell}_{\pi} (\mathcal{C}^{s \asterisk} \otimes
\mathbb{R}^n))$, then $F_{\ell}$ can be seen as an element of $(\mathcal{S}'
(\mathbb{R}^d) \otimes \mathbb{R}^n) \otimes \Lambda^{\ell}_{\pi}
(\mathcal{C}^{s \asterisk} \otimes \mathbb{R}^n)$. Using the same technique as
before we can introduce the Banach space $\mathcal{C}^s_n (\mathbb{R}^d,
\Lambda^{\ell}_{\pi} (\mathcal{C}^{s \asterisk} \otimes \mathbb{R}^n)) \subset
(\mathcal{S}' (\mathbb{R}^d) \otimes \mathbb{R}^n) \otimes
\Lambda^{\ell}_{\pi} (\mathcal{C}^{s \asterisk} \otimes \mathbb{R}^n)$
equipped with its natural norm. In this case we will use the notation
\begin{equation}
  F_{\ell} (X) \assign X \circ F_{\ell}, \label{eq:composition}
\end{equation}
where $(X \circ F_{\ell}) (v) = X (F_{\ell} (v)) \in \mathcal{A}$ for any $v
\in \mathcal{S} (\mathbb{R}^d) \otimes \mathbb{R}^n$.

\begin{definition}
  \label{definition-F1}We say that the linear function $F \in \mathcal{C}^s_n
  \left( \mathbb{R}^d, \bigoplus_{j = 0}^k \Lambda^{2 j + 1}_{\pi}
  (\mathcal{C}^{s \asterisk} \otimes \mathbb{R}^n) \right)$ is a standard odd
  function of degree $2 k + 1$ if there is a sequence of functions $F^r \in
  \mathcal{C}^s_n \left( \mathbb{R}^d, \bigoplus_{j = 0}^k \Lambda^{2 j + 1}
  (\mathcal{C}^{s \asterisk} \otimes \mathbb{R}^n) \right)$ (where $\Lambda^{2
  j + 1} (\mathcal{C}^{s \asterisk} \otimes \mathbb{R}^n) =
  \underbrace{(\mathcal{C}^{s \asterisk} \otimes \mathbb{R}^n) \otimes \cdots
  \otimes (\mathcal{C}^{s \asterisk} \otimes \mathbb{R}^n)}_{\text{$2 j + 1$
  times}}$ and $\otimes$ is the algebraic tensor product) such that for any $g
  \in V$ we have $F^r (g) \rightarrow F (g)$, as $r \rightarrow + \infty$, in
  $\bigoplus_{j = 0}^k \Lambda^{2 j + 1}_{\pi} (\mathcal{C}^{s \asterisk}
  \otimes \mathbb{R}^n)$ equipped with its standard topology.
\end{definition}

\begin{remark}
  \label{rem:discr}When $s > 0$, an example of a function $F = (F^1, \ldots,
  F^n)$ satisfying Definition~\ref{definition-F1} is
  \[ F^k (X) (\delta_x) = \dot{X}^k (x) \sum_{k_1, k_2 = 1}^n
     \int_{\mathbb{R}^d} \mathcal{G}_{k_1, k_2} (x - y) \dot{X}^{k_1} (y)
     \dot{X}^{k_2} (y) \mathd y, \]
  where $\mathcal{G}_{k_1, k_2} \in \mathcal{C}^{s + \epsilon}$ are functions
  that decrease faster than any polynomial at infinity. We note that a finite
  dimensional approximation of $F^k (X) (\delta_x)$ is given by
  \[ F^{k, r} (X) (\delta_x) = \sum_{k_1, k_2 = 1}^n \sum_{D_i \in
     \mathcal{P}_n} \dot{X}^k (x) \dot{X}^{k_1} (x - x_i) \dot{X}^{k_2} (x -
     x_j) \int_{D_i} \mathcal{G}_{k_1, k_2} (y) \mathd y, \]
  where $(\mathcal{P}_r)_r$ is sequence of increasing partitions of
  $\mathbb{R}^d$ for which
  \[ \max_{k_1, k_2 = 1, \ldots, n} \sup_i \left| \int_{D_i} \mathcal{G}_{k_1,
     k_2} (y) \mathd y \right| \rightarrow 0 \]
  as $r \rightarrow \infty$.
\end{remark}

\begin{theorem}
  \label{theorem_F}Let $s \in \mathbb{R}$, let $F \in \mathcal{C}^s_n \left(
  \mathbb{R}^d, \bigoplus_{j = 0}^k \Lambda^{2 j + 1}_{\pi} (\mathcal{C}^{s
  \asterisk} \otimes \mathbb{R}^n) \right)$ be a standard odd function of
  degree $2 k + 1$ (see Definition~\ref{definition-F1}), and consider $X \in
  \mathcal{G}^s (\mathcal{S} (\mathbb{R}^d))$. Then $F (X) \in \mathcal{C}^s_n
  (\mathbb{R}^d, \mathcal{A})$, furthermore for any $v_1, v_2 \in
  \mathcal{C}^{s \asterisk} \otimes \mathbb{R}^n$ we have
  \begin{equation}
    \begin{array}{c}
      F (X) (v_1) F (X) (v_2) = - F (X) (v_2) F (X) (v_1)\\
      F (X) (v_1) X (v_2) = - X (v_2) F (X) (v_1) .
    \end{array} \label{eq:anticommutation}
  \end{equation}
\end{theorem}

\begin{proof}
  The fact that $F (X) \in \mathcal{C}^s_n (\mathbb{R}^d, \mathcal{A})$
  follows from Remark~\ref{remark_regularity}. The
  relations~{\eqref{eq:anticommutation}} are obvious when $F \in
  \mathcal{C}^s_n \left( \mathbb{R}^d, \bigoplus_{j = 0}^k \Lambda^{2 j + 1}
  (\mathcal{C}^{s \asterisk} \otimes \mathbb{R}^n) \right)$. Using a finite
  dimensional approximation $F^r$ converging to $F$, required by Definition
  \ref{definition-F1}, the relations~{\eqref{eq:anticommutation}} are extended
  to the more general case.
\end{proof}

\begin{remark}
  An important consequence of Theorem~\ref{theorem_F} is that if $X \in
  \mathcal{G}^s (V)$ and $F$ satisfies the hypotheses of
  Theorem~\ref{theorem_F}, $F (X)$ can be extended to a Grassmann random
  variable (which in the following we denote by the same symbol $F (X)$) in
  $\mathcal{G}^s (V)$. Furthermore, $F (X)$ and $X$ are compatible (in the
  sense of the corresponding definition in Section~\ref{sec:grassmann-rv}).
\end{remark}

We are now in the position to give a precise meaning to
equation~{\eqref{eq:infinite5}} and to the various hypotheses on its data. Let
the operator $A$ be the (generally unbounded) generator of a strongly
continuous and exponentially contractive semigroup $(e^{A t})_{t \geqslant 0}$
on $L^2 (\mathbb{R}^d)$. Moreover consider the stationary Ornstein-Uhlenbeck
motion
\begin{equation}
  B^A_t (v) = \Xi (\mathbb{I}_{(- \infty, t]} \otimes e^{A (t - \cdot)} v),
  \qquad t \in \mathbb{R}, v \in \mathcal{S} (\mathbb{R}^d) \otimes
  \mathbb{R}^n, \label{eq:ba}
\end{equation}
where $\Xi$ is a Grassmann Gaussian noise on $\mathcal{G} (C^{\infty}
(\mathbb{R}_+) \otimes \mathcal{S} (\mathbb{R}^d) \otimes \mathbb{R}^n)$ with
covariance $\mathbb{I} \otimes C$ (we consider on $C^{\infty}_c (\mathbb{R}_+)
\otimes \mathcal{S} (\mathbb{R}^d) \otimes \mathbb{R}^n$ the pre-Hilbert space
structure given by the natural inclusion of $C^{\infty}_c (\mathbb{R}_+)
\otimes \mathcal{S} (\mathbb{R}^d) \otimes \mathbb{R}^n$ in $L^2 (\mathbb{R}_+
\times \mathbb{R}^d) \otimes \mathbb{R}^n$).

\begin{definition}
  \label{definition_F2}We say that the process $\Psi \in C^0 ([0, T],
  \mathcal{G}^s (V))$ satisfies equation~{\eqref{eq:infinite5}} with initial
  condition $\Psi_0$ if $F$ is \ a standard odd function of finite degree,
  $\Psi_0 \in \mathcal{G}^s (V)$ is independent of $\Xi$, and for any $t \in
  [0, T]$ and any $v \in V$ we have
  \begin{equation}
    \Psi_t (v) = \Psi_0 (e^{A t} v) + \lambda \int_0^t F (\Psi_s) (e^{A (t -
    s)} v) \mathd s + B^A_t (v) - B_0^A (e^{A t} v) \label{eq:infinite2},
  \end{equation}
  where the integral in the variable $s$ is taken in Bochner sense with
  respect the norm of $\mathcal{A}$.
\end{definition}

Hereafter, we introduce the following notations if $F \in \mathcal{C}^s_n
\left( \mathbb{R}^d, \bigoplus_{j = 0}^k \Lambda^{2 j + 1}_{\pi}
(\mathcal{C}^{s \asterisk} \otimes \mathbb{R}^n) \right)$:
\[ \| F \|_{\mathcal{C}^s_n, \pi} \assign \sup_{k = 1, \ldots, n} \| F^k
   \|_{\mathcal{C}^s \left( \mathbb{R}^d, \bigoplus_{j = 0}^k \Lambda^{2 j +
   1}_{\pi} \mathcal{C}^{s \asterisk} \right)}, \]
\[ \deg (F) \assign 2 k + 1. \]
\begin{theorem}
  \label{theorem_drift}Suppose that $F \in \mathcal{C}^s_n \left(
  \mathbb{R}^d, \bigoplus_{j = 0}^k \Lambda^{2 j + 1}_{\pi} (\mathcal{C}^{s
  \asterisk} \otimes \mathbb{R}^n) \right)$ and that $F$ satisfies
  Definition~\ref{definition-F1}, then for any $X, Y \in \mathcal{G}^s (V)$
  compatible we have
  \begin{equation}
    \| F (X) \|_{\mathcal{G}^s (V)} \leqslant \| F \|_{\mathcal{C}^s_n, \pi}
    (1 + \| X \|_{\mathcal{G}^s (V)})^{\deg (F)} \label{eq:propertyF1}
  \end{equation}
  \begin{eqnarray}
    \| F (X) - F (Y) \|_{\mathcal{G}^s (V)} & \leqslant & \deg (F) \| F
    \|_{\mathcal{C}^s_n, \pi} \| X - Y \|_{\mathcal{G}^{s (V)}} \times
    \nonumber\\
    &  & \times (1 + \max (\| X \|_{\mathcal{G}^s (V)}, \| Y
    \|_{\mathcal{G}^s (V)}))^{\deg (F) - 1} .  \label{eq:propertyF2}
  \end{eqnarray}
\end{theorem}

\begin{proof}
  Let $F^{k, r} \in \mathcal{C}^s_n \left( \mathbb{R}^d, \bigoplus_{j = 0}^k
  \Lambda^{2 j + 1} \mathcal{C}^{s \asterisk} \right)$ (where $\bigoplus_{j =
  0}^k \Lambda^{2 j + 1} \mathcal{C}^{s \asterisk}$ is the usual algebraic
  antisymmetric tensor product) be the approximation of $F$ required in
  Definition~\ref{definition-F1}, then using the same reasoning of
  Theorem~\ref{theorem_function1} and Theorem \ref{theorem_function2}, if
  $K_i$ is the function corresponding to the Littlewood--Paley block
  $\Delta_i$, we obtain
  \begin{equation}
    \sup_{k = 1, \ldots, n} \| F^{k, r} (X) (K_i (x - \cdot)) \|_{\mathcal{A}}
    \leqslant \sup_{k = 1, \ldots, n} \| F^{k, r} (K_i (x - \cdot))
    \|_{\Lambda_{\pi} \mathcal{C}^{s \ast}} (1 + \| X \|_{\mathcal{G}^s
    (V)})^{\deg (F)} . \label{eq:propertyF3}
  \end{equation}
  If we multiply both sides of~{\eqref{eq:propertyF3}} by $2^{i s}$, take the
  limit as $r \rightarrow + \infty$ and afterwords the supremum on $x$ and $j$
  we get inequality~{\eqref{eq:propertyF1}}. In a similar way we
  get~{\eqref{eq:propertyF2}} by generalizing the proof of
  Theorem~\ref{theorem_function1} and Theorem~\ref{theorem_function2}.
\end{proof}

\begin{remark}
  \label{remark_drift} Hereafter we use the notations, for $x \geqslant 0$,
  \begin{equation}
    f_F (x) = \| F \|_{\mathcal{C}^s_n, \pi} (1 + x)^{\deg (F)},
  \end{equation}
  \begin{equation}
    g_F (x) = \deg (F) \| F \|_{\mathcal{C}^s_n, \pi} (1 + x)^{\deg (F) - 1} .
  \end{equation}
\end{remark}

We introduce some conditions on $A$ and $\Xi$ in order to prove a result on
existence and uniqueness of solutions to~{\eqref{eq:infinite2}}.

\begin{definition}
  \label{definition_infinitelinear}We say that the unbounded operator $A$ on
  $L^2 (\mathbb{R}^d)$, and the Grassmann Gaussian noise $\Xi$ with covariance
  $C$ are adapted to the space $\mathcal{G}^s (V)$, for a given $s \in
  \mathbb{R}$, if
  \begin{enumerateroman}
    \item The operator $A$ generates a contraction semigroup on $\mathcal{C}^s
    (\mathbb{R}^d)$ such that $\| e^{A t} \|_{\mathcal{L} (\mathcal{C}^s
    (\mathbb{R}^d), \mathcal{C}^s (\mathbb{R}^d))} \lesssim e^{- \lambda_A
    t}$, $t \geqslant 0$, for some strictly positive constant $\lambda_A > 0$.
    
    \item The Gaussian process $B^A$ defined by~{\eqref{eq:ba}} belongs to $C
    (\mathbb{R}, \mathcal{G}^s (V))$ and furthermore we have $\sup_{t \in
    \mathbb{R}} \| B_t^A \|_{\mathcal{C}^s (\mathbb{R}^d, \mathcal{A})} < +
    \infty$.
  \end{enumerateroman}
\end{definition}

\begin{theorem}
  \label{theorem_existence2}Suppose that $F$ satisfies
  Definition~\ref{definition-F1} and $A, \Xi$ satisfy
  Definition~\ref{definition_infinitelinear} and let $\Psi_0$ be a variable in
  $\mathcal{G}^s (V)$ independent of the Grassmann Gaussian noise $\Xi$. Then
  there exists a $T > 0$ such that there is a unique solution $\Psi \in C ([0,
  T], \mathcal{G}^s (V))$ to equation~{\eqref{eq:infinite2}}. Furthermore
  $\Psi_t$ is compatible with $(\Psi_0, \Xi, \Psi_s)$ for every $t, s \in [0,
  T]$.
\end{theorem}

\begin{proof}
  The proof is a straightforward generalization of that of
  Theorem~\ref{theorem_picard} for the case of our infinite dimensional
  setting.
\end{proof}

In order to obtain solutions for all times it is not possible to extend
directly Theorem~\ref{th:global-sol}, since its proof was essentially based on
the hypothesis that $V$ is finite dimensional. On the other hand the proofs of
Theorem~\ref{theorem_stationary}, Lemma~\ref{lemma_bound} and
Theorem~\ref{theorem_approximation} do not depend on the dimension. For this
reason we introduce the space $\mathfrak{K}_K \subset C (\mathbb{R},
\mathcal{G}^s (V))$ defined as the set of functions $\Psi_t \in C (\mathbb{R},
\mathcal{G}^s (V))$ such that
\[ \sup_{t \in \mathbb{R}} \| \Psi_t \|_{\mathcal{C}^s (\mathbb{R}^d,
   \mathcal{A})} < K, \]
for some constant $K > 0$, the function $L$ given by
\[ L (A, C) = \sup_{t \in \mathbb{R}} \| B^A_t \|_{\mathcal{C}^s
   (\mathbb{R}^d, \mathcal{A})}, \]
and the equation for stationary solutions corresponding to
{\eqref{eq:infinite2}} as
\begin{equation}
  \Psi_t^s = \lambda \int_{- \infty}^t F (\Psi_{\tau}^s (e^{A (t - \tau)}))
  \mathd \tau + B^A_t \label{eq:stationary2} \qquad t \in \mathbb{R}.
\end{equation}
We also denote by $\Psi_{- T, t}^X$ the Grassmann process defined by $\Psi_{-
T, t}^X = X$ if $t \leqslant - T$, $T \geqslant 0$, and such that $\Psi_{- T,
t}^X$ is the solution to equation~{\eqref{eq:infinite2}} with initial
condition $X$ at time $t = - T$. \

\begin{theorem}
  \label{theorem_longtime2}Suppose that $F$ satisfies
  Definition~\ref{definition-F1} and $A, \Xi$ satisfy
  Definition~\ref{definition_infinitelinear} and let $K > 2 L (A, C)$. Then
  there exists a function $\lambda_0 (K, L, F)$ such that for any $| \lambda |
  \leqslant 2 \lambda_0 (K, L, F)$ there exists a unique (stationary) solution
  to equation~{\eqref{eq:stationary2}} in $\mathfrak{K}_K$. Furthermore for
  any $| \lambda | \leqslant \lambda_0 (K, L, F)$ and for any random variable
  $X \in \mathcal{G}^s_{\mathcal{}} (V)$, independent of $B^A_t$ and such that
  $\| X \|_{\mathcal{C}^s (\mathbb{R}^d, \mathcal{A})} < \frac{1}{3} K$, we
  have that $\Psi^X_{- T, t}$ converges, as $T \rightarrow + \infty$, in
  $\mathcal{G}^s (V)$ to $\Psi^s_t$ uniformly on compact subsets of
  $\mathbb{R}$.
\end{theorem}

\begin{proof}
  Due to our appropriate infinite dimensional setting the proof is a
  straightforward generalization of Theorem~\ref{theorem_stationary},
  Lemma~\ref{lemma_bound} and Theorem~\ref{theorem_approximation}.
\end{proof}

\begin{remark}
  The function $\lambda_0$ defined in Theorem~\ref{theorem_longtime2} has to
  satisfy the following inequalities
  \begin{equation}
    \frac{2 \lambda_0 (K, L, F) f_F (K)}{\lambda_A} + 2 L < K
    \label{eq:contraction3},
  \end{equation}
  \begin{equation}
    \frac{2 \lambda_0 (K, L, F)}{\lambda_A} g_F (K) < 1
    \label{eq:contraction4}
  \end{equation}
  where $f_F$, $g_F$ and $\lambda_A$ are the quantities introduced in Remark
  \ref{remark_drift} and Definition~\ref{definition_infinitelinear}
  respectively.
\end{remark}

\section{The Yukawa$_2$ model}\label{sec:yukawa}

We are now in position to study the stochastic quantization for an Euclidean
fermionic QFT. We consider, as an example, the Euclidean counterpart of the
Yukawa model describing the interaction of a scalar field with a spin $1 / 2$
field in a $1 + 1$ Minkowski space-time (``Yukawa$_2$
model'')~{\cite{osterwalder_euclidean_1973,cooper_weakly_1977,battle_infinite_1981,lesniewski_effective_1987,summers2012perspective}}.
However we will not attempt here to remove the UV cutoff so the dimension of
the space will not play any fundamental role. The Yukawa$_2$ model contains
already the main features of a large class of Euclidean fermionic theories. It
will be clear from the considerations below that more general results are
easily derivable, including other models involving fermions with ultraviolet
cut-offs, at the price of heavier notation. For the sake of readability we
choose to stick to a definite model.

\subsection{Osterwalder--Schrader fields}\label{sec:os-fields}

We define the Euclidean Yukawa$_2$ quantum field theory following Osterwalder
and Schrader~{\cite{osterwalder_euclidean_1973-2}}. In the language of this
paper, the Euclidean free Fermions introduced by Osterwalder and Schrader
(compare with~{\cite{osterwalder_euclidean_1973-2}}, page~282 and
equation~(3.3), page~284) are given by a quadruplet $\tilde{\psi} = (\psi,
\bar{\psi}) \in \mathcal{G} (\mathcal{S} (\mathbb{R}^2) \otimes \mathbb{C}^4)$
of complex Grassmann Gaussian fields on the complex pre-Hilbert space $V
=\mathcal{S} (\mathbb{R}^2) \otimes \mathbb{C}^4 \subseteq L^2 (\mathbb{R}^2)
\otimes \mathbb{C}^4$ with an Hermitian scalar product denoted $(\cdot,
\cdot)_V$ and satisfying
\begin{equation}
  \omega (\tilde{\psi} (f) \tilde{\psi} (g)) = \langle f, K g \rangle_c,
  \label{eq:free-corrrelation} \qquad f, g \in \mathcal{S} (\mathbb{R}^2)
  \otimes \mathbb{C}^4,
\end{equation}
with the standard real structure $c$ given by the complex conjugation (i.e. $c
f = \bar{f}$) determining the bilinear form $\langle f, g \rangle_c = (c f,
g)_V$, and the correlation $K : L^2 (\mathbb{R}^2) \otimes \mathbb{C}^4
\rightarrow L^2 (\mathbb{R}^2) \otimes \mathbb{C}^4$ given by the bounded
$c$-antisymmetric operator
\begin{equation}
  K = (m_f^2 - \Delta)^{- 1} \left(\begin{array}{cc}
    0 & \left( \slashnabla + m_f \right)\\
    - \left( \slashnabla + m_f \right)^c & 0
  \end{array}\right), \label{eq:matrixK}
\end{equation}
where $\Delta$ is the Laplacian on $\mathbb{R}^2$, $m_f > 0$ is a constant
interpreted as the mass of the fermionic particle associated with the
fermionic fields $\tilde{\psi}$, $\slashnabla \assign \gamma_1 \nabla_1 +
\gamma_2 \nabla_2$ with $\nabla_i$ the partial derivatives with respect to
$x_i$, $i = 1, 2$, and with Euclidean Dirac matrices $\gamma_1, \gamma_2 :
\mathbb{C}^2 \rightarrow \mathbb{C}^2$ given by
\[ \gamma_1 = \left(\begin{array}{cc}
     0 & 1\\
     1 & 0
   \end{array}\right), \quad \gamma_2 = \left(\begin{array}{cc}
     0 & - \mathi\\
     \mathi & 0
   \end{array}\right), \]
and such that
\[ \gamma_a \gamma_b + \gamma_b \gamma_a = 2 \delta_{a b}, \qquad a, b = 1, 2.
\]
Note that $K = (m_f^2 - \Delta)^{- 1 / 2} u$, with $u$ an unitary operator in
$L^2 (\mathbb{R}^2) \otimes \mathbb{C}^4$. Hence the correlation matrix of the
field $\tilde{\psi}$ can also be written as
\[ \omega (\tilde{\psi} (f) \tilde{\psi} (g)) = \langle (m_f^2 - \Delta)^{- 1
   / 4} f, u (m_f^2 - \Delta)^{- 1 / 4} g \rangle_c . \]
Therefore $\omega$ is the state on $\mathcal{A}$ describing the Euclidean free
fermion fields and we can realize $\tilde{\psi}$ as
\[ \tilde{\psi} (f) \assign \tilde{\zeta} ((m_f^2 - \Delta)^{- 1 / 4} f),
   \qquad f \in \mathcal{S} (\mathbb{R}^2) \otimes \mathbb{C}^4, \]
where $\tilde{\zeta} \in \mathcal{G} (\mathcal{S} (\mathbb{R}^2) \otimes
\mathbb{C}^4)$ is the complex Grassmann Gaussian field specified by the
bounded covariance
\[ \omega (\tilde{\zeta} (f) \tilde{\zeta} (g)) = \langle f, u g \rangle_c,
   \qquad f, g \in \mathcal{S} (\mathbb{R}^2) \otimes \mathbb{C}^4 . \]
As a consequence, in this specific realization, we have
\begin{equation}
  \| \tilde{\psi} (f) \| = \| \tilde{\zeta} ((m_f^2 - \Delta)^{- 1 / 4} f) \|
  \lesssim \| (m_f^2 - \Delta)^{- 1 / 4} f \|_{L^2 (\mathbb{R}^2) \otimes
  \mathbb{C}^4} < \infty, \label{eq:dirac-reg}
\end{equation}
where the implicit constant in the inequality depends only on $m_f$.

\

In order to describe the Yukawa$_2$ model we need also to introduce an
Euclidean boson field $\varphi$ of mass $m_b > 0$: it is the centered Gaussian
field on $\mathbb{R}^2$ with covariance $(m^2_b - \Delta)^{- 1}$ where $m_b >
0$ is looked upon as a fixed parameter and $\Delta$ is the Laplacian on
$\mathbb{R}^2$. We can realize $\varphi$ as the canonical multiplication
operator on the Hilbert space $\mathcal{H}_{\varphi} = L^2 \left(
\mathcal{C}^{- \kappa}_{\text{loc}} (\mathbb{R}^2), \nu \right)$ where $\nu$
is the centered Gaussian measure on $\mathcal{C}^{- \kappa}_{\text{loc}}
(\mathbb{R}^2)$ (the local version of the $\mathcal{C}^{- \kappa}$ Besov
space) with covariance $(m^2_b - \Delta)^{- 1}$ and $\kappa > 0$ is any small
positive quantity. We shall denote by $\omega_{\varphi}$ the state associated
with $\nu$ and $\mathcal{C}^{- \kappa}_{\text{loc}} (\mathbb{R}^2)$.

\

In addition, large scale and small scale cutoffs are introduced as follows.
Fix a smooth positive functions $h \in C^{\infty} (\mathbb{R}^2)$ with compact
support and let $\mathfrak{a}$ be a regular enough (see below for the precise
regularity) continuous function with compact support such that $\mathfrak{a}
(0) = 1$. Take $\varepsilon > 0$ and let $\mathfrak{a}_{\varepsilon} (x) =
\mathfrak{a} (x / \varepsilon) \varepsilon^{- 2}$. Then we define new
Grassmann variables
\[ \tilde{\psi}_{\varepsilon} (x) \assign \tilde{\psi}
   (\mathfrak{a}_{\varepsilon} (\cdot - x)), \qquad x \in \mathbb{R}^2 . \]
Then $\tilde{\psi}_{\varepsilon} \in \mathcal{G}^s (\mathcal{S} (\mathbb{R}^2)
\otimes \mathbb{C}^4)$ for some $s > 0$ (in general depending on the
regularity of $\mathfrak{a}$ see below).

\

Following the argument of Section~6 of~{\cite{osterwalder_euclidean_1973-2}}
we define the Euclidean Yukawa$_2$ model as follows.

\begin{definition}
  \label{def:yukawa2}The approximate fermionic Schwinger functions
  $(\rho^n_{\varepsilon, h})_n$ of the Yukawa$_2$ model on $\mathbb{R}^2$ with
  infrared (IR) cutoff $h$ and ultraviolet (UV) fermion cutoff $\varepsilon$
  are given by the Euclidean averages
  \begin{equation}
    \rho^n_{\varepsilon, h} (f_1, \ldots, f_n) = \frac{\omega \otimes
    \omega_{\varphi} (\tilde{\psi}_{\varepsilon} (f_1) \cdots
    \tilde{\psi}_{\varepsilon} (f_n) e^{- V_{\varepsilon, h}
    (\varphi)})}{\omega \otimes \omega_{\varphi} (e^{- V_{\varepsilon}})}
    \label{eq:euclidean-yukawa-main}
  \end{equation}
  with potential $V_{\varepsilon, h}$ given by the unbounded
  (non-self-adjoint) operator on $\mathcal{H} \otimes \mathcal{H}_{\varphi}$,
  where $\mathcal{H}$ is the fermionic Hilbert space,
  \begin{equation}
    V_{\varepsilon, h} (\varphi) = \lambda \int_{\mathbb{R}^2} \mathd x h (x)
    (\bar{\psi}_{\varepsilon} (x) \cdot \psi_{\varepsilon} (x)) \varphi (x),
    \label{eq:yukawa-potential-boson}
  \end{equation}
  $\lambda \in \mathbb{R}$ is a coupling constant and $\omega$, resp.
  $\omega_{\varphi}$ are the fermion, resp. boson, states described above.
\end{definition}

The potential $V_{\varepsilon, h} (\varphi)$ can be expressed in terms of the
fields $\tilde{\psi}_{\varepsilon}$ as follows
\[ V_{\varepsilon, h} (\varphi) = \frac{\lambda}{2} \int_{\mathbb{R}^2} h (x)
   (\tilde{\psi}_{\varepsilon} (x) \cdot J \widetilde{\psi}_{\varepsilon} (x))
   \varphi (x) \mathd x, \]
where $J (\psi, \bar{\psi}) \assign (- \bar{\psi}, \psi)$.

Note that there is no difficulty in making sense
of~{\eqref{eq:euclidean-yukawa-main}}, due to the following observations (see
also~{\cite{osterwalder_euclidean_1973-2}}). The map $x \in \mathbb{R}^2
\mapsto \tilde{\psi}_{\varepsilon} (x) \in \mathcal{A}= \mathcal{L}
(\mathcal{H})$ is $\mathcal{C}^s$ due to the presence of the regularization
$\varepsilon > 0$ and the regularity of the Euclidean field $\tilde{\psi}$
given in~{\eqref{eq:dirac-reg}}. The massive Gaussian free field $\varphi$ has
(local) regularity $\mathcal{C}^{- \kappa}_{\tmop{loc}}$ for arbitrarily small
$\kappa > 0$, therefore, taking $s > \kappa$, by a standard Besov estimates
for Banach valued functions (see e.g.~{\cite{Amann1997,Amann2019}}), we have,
for any \tmtextit{fixed} $\varphi \in \mathcal{C}^{- \kappa}_{\tmop{loc}}
(\mathbb{R}^2)$
\[ \| V_{\varepsilon, h} (\varphi) \|_{\mathcal{A}} \lesssim \lambda \|
   \tilde{\psi}_{\varepsilon} \|^2_{\mathcal{C}^{\kappa} (\mathcal{A})} \| h
   \varphi \|_{B^{- \kappa}_{1, 1} (\mathbb{R}^2)} \lesssim \lambda \| h
   \varphi \|_{\mathcal{C}^{- \kappa} (\mathbb{R}^2)} . \]
It is also not difficult to see that the map $V_{\varepsilon, h} :
\mathcal{C}^{- \kappa} \rightarrow \mathcal{A}$ given by $V_{\varepsilon, h} :
\varphi \mapsto V_{\varepsilon, h} (\varphi)$ is linear and continuous in
$\varphi$, therefore measurable and that
\[ \omega \otimes \omega_{\varphi} (\tilde{\psi} (f_1) \cdots \tilde{\psi}
   (f_n) e^{- V_{\varepsilon, h} (\varphi)}) = \omega (\tilde{\psi} (f_1)
   \cdots \tilde{\psi} (f_n) \omega_{\varphi} (e^{- V_{\varepsilon, h}
   (\varphi)})) \]
e.g. by proving this first for a piecewise approximation of $\varphi$ and then
taking the limit. Moreover we have (by the properties of the Bochner integral)
\[ \| \omega_{\varphi} (e^{- V_{\varepsilon, h} (\varphi)}) \|_{\mathcal{A}}
   \leqslant \omega_{\varphi} (\exp (C | \lambda | \| h \varphi
   \|_{\mathcal{C}^{- \kappa}})) < \infty, \]
for all $\lambda \in \mathbb{R}$. We conclude that the averages
in~{\eqref{eq:euclidean-yukawa-main}} are well defined since
\[ | (\omega \otimes \omega_{\varphi}) (\tilde{\psi} (f_1) \cdots
   \tilde{\psi} (f_n) e^{- V_{\varepsilon, h}}) | \leqslant \| \tilde{\psi}
   (f_1) \cdots \tilde{\psi} (f_n) \omega_{\varphi} (e^{- V_{\varepsilon, h}
   (\varphi)}) \|_{\mathcal{A}} < \infty . \]
By expanding the quantity $\omega_{\varphi} (e^{- V_{\varepsilon, h}
(\varphi)}) \in \mathcal{A}$ in power series and integrating away the Gaussian
field we obtain
\[ \omega_{\varphi} (e^{- V_{\varepsilon, h} (\varphi)}) =
   e^{\mathcal{V}_{\varepsilon, h}}, \]
with a purely fermionic potential given by
\begin{equation}
  \mathcal{V}_{\varepsilon, h} (\tilde{\psi}_{\varepsilon}) = \lambda^2
  \int_{(\mathbb{R}^2)^2} (h \bar{\psi}_{\varepsilon} \cdot
  \psi_{\varepsilon}) (x) \mathcal{G} (x - y) (h \bar{\psi}_{\varepsilon}
  \cdot \psi_{\varepsilon}) (y) \mathd x \mathd y. \label{eq:effectiveV}
\end{equation}
Here $\mathcal{G} = (m^2_b - \Delta)^{- 1}$ is the Green function of the
operator $(m^2_b - \Delta)$ and $m_b$ is the boson mass. This gives us the
purely fermionic representation of the Schwinger functions
{\eqref{eq:euclidean-yukawa-main}} given by:
\begin{equation}
  \rho^n_{\varepsilon, h} (f_1, \ldots, f_n) = \frac{\omega (\tilde{\psi}
  (f_1) \cdots \tilde{\psi} (f_n) e^{\mathcal{V}_{\varepsilon, h}})}{\omega
  (e^{\mathcal{V}_{\varepsilon, h}})} . \label{eq:euclidean-yukawa-fermion}
\end{equation}
Perturbation theory suggests that the boson field requires a mass
renormalization $a_{\varepsilon}$ which depends on the fermion cutoff
$\varepsilon > 0$. In theory one would need to modify the bosonic Gaussian
field covariance to include a mass renormalization and let $\varphi$ have
covariance $(m^2_b - \Delta + a_{\varepsilon})^{- 1}$ instead of the original
$(m^2_b - \Delta)^{- 1}$. However in this paper we will not discuss the
removal of the UV cutoff and keep $\varepsilon > 0$ fixed. Therefore, for
simplicity, we ignore this mass renormalization.

\subsection{Stochastic quantization}\label{sec:yukawa-sq}

The Grassmann SPDEs associated to the measure $\tilde{\omega} (\cdot) = \omega
(\cdot e^{\mathcal{V}_{\varepsilon, h}})$ is formally
\begin{equation}
  \mathd \tilde{\Psi}_{\varepsilon, t}^h = ((\Delta - m_f^2)
  \tilde{\Psi}_{\epsilon, t}^h + \lambda^2 F_{\varepsilon, h, Y}
  (\tilde{\Psi}_{\varepsilon, t}^h)) \mathd t + \mathd \tilde{B}_{\varepsilon,
  t} \label{eq:stochasticquantization}
\end{equation}
where $\tilde{\Psi}_{\varepsilon, t}^h = (\Psi_{\varepsilon, t}^h,
\bar{\Psi}_{\varepsilon, t}^h)$,
\[ F_{\varepsilon, h, Y} (\tilde{\Psi}_{\varepsilon, t}^h) = C_f \left(
   \begin{array}{l}
     \mathfrak{a}_{\varepsilon}^{\asterisk 2} \ast \left( h
     \bar{\Psi}_{\varepsilon, t}^h \int \mathcal{G} (y - \cdot) h (y)
     \Psi_{\varepsilon, t}^h (y) \bar{\Psi}_{\varepsilon, t}^h (y) \mathd y
     \right)\\
     \mathfrak{a}_{\varepsilon}^{\asterisk 2} \ast \left( h \Psi_{\varepsilon,
     t}^h \int \mathcal{G} (y - \cdot) h (y) \Psi_{\varepsilon, t}^h (y)
     \bar{\Psi}_{\varepsilon, t}^h (y) \mathd y \right)
   \end{array} \right), \]
with
\begin{equation}
  C_f = \left(\begin{array}{cc}
    0 & \left( \slashnabla + m_f \right)\\
    - \left( \slashnabla + m_f \right)^{\Tau} & 0
  \end{array}\right), \label{eq:Cf}
\end{equation}
and where $\mathfrak{a}_{\varepsilon} (x) = \varepsilon^{- 2} \mathfrak{a}
\left( \frac{x}{\varepsilon} \right)$, $\varepsilon > 0$, $\mathfrak{a} :
\mathbb{R}^2 \rightarrow \mathbb{R}^2$ is the regular enough function with
compact support (see below for the more precise requirements) considered
before, $\mathfrak{a}_{\varepsilon}^{\asterisk 2} = \mathfrak{a}_{\varepsilon}
\asterisk \mathfrak{a}_{\varepsilon}$, $h$ is a smooth function (in the
beginning with compact support, later on we will consider $h \equiv 1$).

The Brownian motion $\tilde{B}$ is described as follows. Let $\tilde{\Xi}$ be
the white noise on the pre-Hilbert space $L^2 (\mathbb{R}_+) \otimes
\mathcal{S} (\mathbb{R}^2) \otimes \mathbb{C}^4 \subseteq L^2 (\mathbb{R}_+)
\otimes B_{2, 2}^{- 1 / 2} (\mathbb{R}^2) \otimes \mathbb{C}^4$ equipped with
the scalar product
\begin{equation}
  ((f, g)) = ((m_f^2 - \Delta)^{- 1 / 4} f, (m_f^2 - \Delta)^{- 1 / 4} g)_{L^2
  (\mathbb{R}_+ \times \mathbb{R}^2) \otimes \mathbb{C}^4}
  \label{eq:scalarproduct} \qquad f, g \in L^2 (\mathbb{R}_+) \otimes
  \mathcal{S} (\mathbb{R}^2) \otimes \mathbb{C}^4,
\end{equation}
and having bounded covariance
\[ \tilde{C}_f = (m^2_f - \Delta)^{- 1 / 2} C_f . \]
The (two-sided) Brownian motion $\tilde{B}$ on $\mathcal{S} (\mathbb{R}^2)
\otimes \mathbb{C}^4$ is defined by
\[ \tilde{B}_t (v) = \Xi (\mathbb{I}_{[0, t]} \otimes v), \qquad t \in
   \mathbb{R}, v \in \mathcal{S} (\mathbb{R}^2) \otimes \mathbb{C}^4 . \]
where we use the convention
\[ \mathbb{I}_{[0, t]} (s) = \left\{\begin{array}{l}
     1 \tmop{if} \; t \geqslant 0 \text{and} 0 \leqslant s \leqslant t\\
     - 1 \tmop{if} \; t < 0 \text{and} t \leqslant s \leqslant 0\\
     0 \tmop{otherwise}
   \end{array}\right. . \]
We set $\tilde{B}_{\varepsilon} = \mathfrak{a}_{\varepsilon} \asterisk
\tilde{B}$. According to Definition~\ref{definition_F2} the stationary version
of equation~{\eqref{eq:stochasticquantization}} is the solution to the
following integral problem
\begin{equation}
  \tilde{\Psi}^{h, s}_{\varepsilon, t} (v) = \lambda^2 \int_{- \infty}^t
  F_{\varepsilon, h, Y} (\tilde{\Psi}_{\varepsilon, \tau}^{h, s}) (e^{A (t -
  \tau)} v) \mathd \tau + \tilde{B}^A_t (v), \label{eq:yukawas} \qquad t \in
  \mathbb{R}, v \in \mathbb{C}^4,
\end{equation}
where we take $A = \Delta - m_f^2$ and
\[ \tilde{B}^A_t (v) = \Xi (\mathbb{I}_{(- \infty, t]} \otimes e^{A (t -
   \cdot)} v), \qquad t \in \mathbb{R}, v \in \mathcal{S} (\mathbb{R}^2)
   \otimes \mathbb{C}^4 . \]
\begin{remark}
  \label{remark_alpha}From now on we will make the following assumption. We
  fix $\delta > 0$ and $\gamma > 0$, let \ $\alpha = \gamma + 1 / 2 + 2
  \delta$ and take $\mathfrak{a} \in B^{\alpha}_{1, \infty} (\mathbb{R}^2)
  \cap \mathcal{C}^{\alpha} (\mathbb{R}^2)$. Again we recall that $\varepsilon
  > 0$ will be fixed in all our study.
\end{remark}

The main theorem of this section concerns the construction of a stochastic
quantization equation yielding (for $| \lambda |$ small and $\varepsilon > 0$)
the relation~{\eqref{eq:euclidean-yukawa-fermion}} for the Yukawa$_2$ model
with a smooth IR cut-off $h$ (Theorem~\ref{theorem_main1}), resp. after the
removal of $h$~(Theorem~\ref{theorem_main2}).

\begin{theorem}
  \label{theorem_main1}If $| \lambda |$ is small enough and $h$ is a smooth
  function with compact support, then, for any $k \geqslant 1$, $f_1, \ldots,
  f_k \in C^{\infty}_c (\mathbb{R}^2) \otimes \mathbb{C}^4$ and $t \in
  \mathbb{R}$, we have
  \begin{equation}
    \frac{\omega (\tilde{\psi} (f_1) \cdots \tilde{\psi} (f_k)
    e^{\mathcal{V}_{\varepsilon, h}})}{\omega (e^{\mathcal{V}_{\varepsilon,
    h}})} = \omega (\tilde{\Psi}_{\varepsilon, t}^{h, s} (f_1) \cdots
    \tilde{\Psi}_{\varepsilon, t}^{h, s} (f_k)) .
  \end{equation}
\end{theorem}

We are also able to obtain the following result on the removal of the IR
cutoff.

\begin{theorem}
  \label{theorem_main2}If $| \lambda |$ is small enough and if we denote by
  $\tilde{\Psi}_{\varepsilon, t}^s$ the solution to
  equation~{\eqref{eq:yukawas}} with $h \equiv 1$, we have that for any $f_1,
  \ldots, f_k \in C^{\infty}_c (\mathbb{R}^2) \otimes \mathbb{C}^4$ and for
  any $t \in \mathbb{R}$
  \begin{equation}
    \lim_{h_n \rightarrow 1} \frac{\omega \left( \tilde{\psi} (f_1) \cdots
    \tilde{\psi} (f_k) e^{\mathcal{V}_{\varepsilon, h_n}} \right)}{\omega
    \left( e^{\mathcal{V}_{\varepsilon, h_n}} \right)} = \omega
    (\tilde{\Psi}_{\varepsilon, t}^s (f_1) \cdots \tilde{\Psi}_{\varepsilon,
    t}^s (f_k)) .
  \end{equation}
  where $(h_n)_{n \geqslant 1}$ is any sequence of smooth functions with
  compact support converging to $1$ uniformly in $\mathcal{C}^{\gamma +
  \delta}$ (see Remark~\ref{remark_alpha} for the condition on $\gamma,
  \delta$) on compact subsets of $\mathbb{R}^2$.
\end{theorem}

The proof of these theorems will be presented in
Section~\ref{sec:cont-finite-vol-yukawa}, resp.
Section~\ref{eq:yukawa-infinite-vol}. These will be prepared by various
estimates and results in the rest of the present section and in
Section~\ref{sec:finit-dim-yukawa} that might have an interest on their own,
also in view of other applications.

\begin{lemma}
  \label{lemma_noise}We have $\tilde{B}^A_{\varepsilon} \in
  \mathcal{C}^{\delta} (\mathbb{R}, \mathcal{C}^{\gamma} (\mathbb{R}^2,
  \mathcal{A}) \otimes \mathbb{C}^4)$ and
  \[ \sup_{t \in \mathbb{R}} \| \tilde{B}^A_{\varepsilon, t}
     \|_{\mathcal{C}^{\gamma + 2 \delta} (\mathbb{R}^2, \mathcal{A}) \otimes
     \mathbb{C}^4} < + \infty ._{} \]
\end{lemma}

\begin{proof}
  It is sufficient to prove the lemma for $\varepsilon = 0$ and $\gamma = -
  \frac{1}{2} - 2 \delta$, since for $\varepsilon > 0$ we have
  $\mathfrak{a}_{\varepsilon} \asterisk : \mathcal{C}^s \rightarrow
  \mathcal{C}^{s + \alpha}$ by Theorem~\ref{theorem_convolution} in
  Appendix~\ref{appendix_Besov}. By definition of Grassmann Gaussian random
  variables, if $K_i \in \mathcal{S} (\mathbb{R}^2)$ is the function
  corresponding to the Littlewood--Paley block $\Delta_i$ (see Appendix
  \ref{appendix_Besov} for the definition), we have
  \[ \sup_{x \in \mathbb{R}^2} \| \Delta_i \tilde{B}_{0, t} - \Delta_i
     \tilde{B}_{0, s} \|_{\mathcal{A} \otimes \mathbb{C}^4} \lesssim \|
     \mathbb{I}_{(- \infty, t]} (\tau) e^{(\Delta - m_f^2) (t - \tau)} (-
     \Delta + m^2)^{1 / 4} (K_i) + \nobracket \]
  \[ \nobracket \nobracket -\mathbb{I}_{(- \infty, s]} (\tau) e^{(\Delta -
     m_f^2) (s - \tau)} (m^2_f - \Delta)^{1 / 4} (K_i) \| \nobracket_{L^2
     (\mathbb{R}^2 \times \mathbb{R})} \lesssim \]
  \[ \lesssim \| \mathbb{I}_{(s, t]} (\tau) e^{(\Delta + m_f^2) (t - \tau)}
     (m^2_f - \Delta)^{1 / 4} (K_i) \|_{L^2 (\mathbb{R}^2 \times \mathbb{R})}
     + \]
  \[ + \| \mathbb{I}_{(- \infty, s]} (\tau) (e^{(\Delta - m_f^2) (t - s)} - 1)
     e^{- (\Delta - m_f^2) (s - \tau)} (m^2_f - \Delta)^{1 / 4} (K_i) \|_{L^2
     (\mathbb{R}^2 \times \mathbb{R})} \lesssim \]
  \[ \lesssim \left( \int_{\mathbb{R}^2} \int_s^t \phi_i (p)^2 e^{- 2 (| p |^2
     + m_f^2) (t - \tau)} | p | \mathd p \mathd \tau \right)^{1 / 2} + \]
  \[ + \left( \int_{\mathbb{R}^2} \int_{- \infty}^s \left( \int_0^{t - s} | p
     |^2 e^{- (| p |^2 + m_f^2) k} \mathd k^{} \right)^2 \phi_i (p)^2 e^{- 2
     (| p |^2 + m_f^2) (s - \tau)} | p | \mathd p \mathd \tau \right)^{1 / 2}
     \lesssim \]
  \[ \lesssim \left( \int_{2^{i - 1} \lesssim | p | \lesssim 2^i} \left(
     \int_0^{t - s} \frac{1}{\tau^{1 - 2 \delta}} \mathd \tau \right) | p |^{-
     1 + 4 \delta} \mathd p \right)^{1 / 2} + \]
  
  \[ + \left( \int_{2^{i - 1} \lesssim | p | \lesssim 2^i} \left( \int_0^{t -
     s} \frac{1}{k^{1 - \delta}} \mathd k \right)^2 | p |^{- 1 + 4 \delta}
     \mathd p \right)^{1 / 2} \lesssim 2^{i \left( \frac{1}{2} + 2 \delta
     \right)} | t - s |^{\delta} . \]
  From this the thesis follows from the definition of the norm of the Besov
  space $\mathcal{C}^{\gamma} (\mathbb{R}^2, \mathcal{A})$.
\end{proof}

\begin{lemma}
  \label{lemma_drift}For any $s > 0$, the function $F_{\varepsilon, h, Y}$ is
  well defined on $\mathcal{C}^s (\mathbb{R}^d, \mathcal{A})$ into itself, it
  satisfies Definition~\ref{definition-F1}. Furthermore we have $\|
  F_{\varepsilon, h, Y} \|_{\mathcal{C}^s, \pi} \lesssim \| h
  \|_{\mathcal{C}^s}^2$.
\end{lemma}

\begin{proof}
  First we prove that $F \in \mathcal{C}^s (\mathbb{R}^2, \Lambda_{\pi}^3
  (\mathcal{C}^{s \ast} (\mathbb{R}^2) \otimes \mathbb{C}^4))$. We start by
  noticing that, if we denote by $\delta_x$ the Dirac delta with unit mass at
  the point $x \in \mathbb{R}^2$, we have $\delta_x \otimes e_i \in
  \mathcal{C}^{0 \ast} (\mathbb{R}^2)$, for $i = 1, \ldots, 4$, where $e_i$ is
  the projection on the component $i$-th of $\mathbb{C}^4$. Furthermore, since
  for $s \centernot{\in} \mathbb{N}$, $s > 0$, $\mathcal{C}^s (\mathbb{R}^2)$ is the
  space of $s$ continuous H{\"o}lder functions, we have that for any $0 < s' <
  s$ and $s' \centernot{\in} \mathbb{N}$ the map $x \longmapsto \delta_x \otimes
  e_i$ is in $\mathcal{C}^{s'} (\mathbb{R}^2, \mathcal{C}^{s \ast}
  (\mathbb{R}^2) \otimes \mathbb{C}^4)$. This and the fact that, by
  Theorem~\ref{theorem_multiplication}, $\mathcal{C}^{s'} (\mathbb{R}^2,
  \Lambda_{\pi} (\mathcal{C}^{s \ast} (\mathbb{R}^2) \otimes \mathbb{C}^4))$
  is a Banach algebra with respect to the multiplication $\wedge$, imply that
  the map
  \[ \tilde{F} (x) \assign h (x) (\delta_x \otimes e_1) \wedge (\delta_x
     \otimes e_3) + (\delta_x \otimes e_2) \wedge (\delta_x \otimes e_4) \]
  is $s'$-H{\"o}lder as a function from $\mathbb{R}^2$ into $\Lambda_{\pi}^2
  (\mathcal{C}^{s \ast} (\mathbb{R}^2) \otimes \mathbb{C}^4)$, with
  \begin{equation}
    \| \tilde{F} \|_{\mathcal{C}^{s'} (\mathbb{R}^2, \Lambda_{\pi}^2
    (\mathcal{C}^{s \ast} (\mathbb{R}^2) \otimes \mathbb{C}^4))} \lesssim \| h
    \|_{\mathcal{C}^s} \label{eq:inequalityh1}
  \end{equation}
  where the constants in the symbol $\lesssim$ do not depend on $h$. On the
  other hand, by Theorem \ref{theorem_laplacian}, $\mathcal{G} \ast \tilde{F}
  \in \mathcal{C}^{s' + 2} (\mathbb{R}^2, \Lambda_{\pi} (\mathcal{C}^{s \ast}
  (\mathbb{R}^2) \otimes \mathbb{C}^4))$ which, if $s - s' > 2$, is contained
  in $\mathcal{C}^s (\mathbb{R}^2, \Lambda_{\pi} (\mathcal{C}^{s \ast}
  (\mathbb{R}^2) \otimes \mathbb{C}^4))$. Since every component of
  $F_{\varepsilon, h, Y}$ is of the form $\left( \left( \slashnabla + m_f
  \right) \mathfrak{a}_{\varepsilon}^{\asterisk 2} \right) \ast ((h (x)
  \delta_x \otimes e_i) \wedge \mathcal{G} \asterisk \tilde{F})$, using the
  fact that, by Theorem~\ref{theorem_laplacian} and
  Theorem~\ref{theorem_convolution}, $\left( \slashnabla + m_f \right)
  \mathfrak{a}_{\varepsilon}^{\asterisk 2} \asterisk \cdot : \mathcal{C}^s
  (\mathbb{R}^2, \mathcal{A}) \rightarrow \mathcal{C}^{s - 1 + 2 \alpha}
  (\mathbb{R}^2, \mathcal{A}) \subset \mathcal{C}^s (\mathbb{R}^2,
  \mathcal{A})$ (since $\alpha > \frac{1}{2}$ for the conditions in
  Remark~\ref{remark_alpha}) is a continuous operator, and by
  Theorem~\ref{theorem_multiplication} we get the thesis. As byproduct of the
  previous reasoning we obtain also that $\| F_{\varepsilon, h, Y}
  \|_{\mathcal{C}^s, \pi} \lesssim \| h \|_{\mathcal{C}^s}^2$. \ This can be
  seen using the fact that, when $s > 0$, $\tilde{\Psi}_{\varepsilon, t}$ is
  H{\"o}lder continuous in space and that $\mathcal{G}$ decreases
  exponentially at infinity.
  
  In order to finish the proof we have to show that $F_{\varepsilon, h, Y}$
  admits an approximation of the form required by
  Definition~\ref{definition-F1}. This can be obtained by approximating the
  convolution with $\mathcal{G}$ and $\mathfrak{a}_{\varepsilon}$ (the
  regularizing function in the definition of $\tilde{\Psi}_{\varepsilon}$)
  using a finite sum. Indeed, if $\mathcal{P}_n$ is a sequence of increasing
  (finite) partition of the $\mathbb{R}^2$, we can approximate the integral
  $\mathcal{G} \asterisk \tilde{F} (x)$ by the sum
  \[ S_n (x) = \sum_{D_i \in \mathcal{P}_n}  \tilde{F} (x - x_i) \int_{D_i}
     \mathcal{G} (y) \mathd y, \qquad x \in \mathbb{R}^2, \]
  where $x_i \in D_i$ is any point in $D_i$. If we choose the sequence of
  partitions $\mathcal{P}_n$ in such a way that $\lim_{n \rightarrow + \infty}
  \left\{ \sup_i \left( \int_{D_i} \mathcal{G} (y) \mathd y \right) \right\} =
  0$, we have that, as $n \rightarrow \infty$, $S_n \rightarrow \mathcal{G}
  \asterisk \tilde{F}$ in $\mathcal{C}^{s'} (\mathbb{R}^2, \Lambda_{\pi}
  (\mathcal{C}^{s \ast} (\mathbb{R}^2) \otimes \mathbb{C}^4))$, where $s' <
  s$. The convolution with $\left( \left( \slashnabla + m_f \right)
  \mathfrak{a}_{\varepsilon}^{\asterisk 2} \right)$ can be handled in a
  similar way, completing the proof.
\end{proof}

Using the previous two lemmas we are in position to prove the following
theorem.

\begin{theorem}
  \label{theorem_main3}Equation~{\eqref{eq:stochasticquantization}} admits a
  unique (local in time) solution, starting from $0$ in $\mathcal{G}^{\gamma}
  (V)$ ($\gamma$ being defined before Theorem~\ref{theorem_main1} and
  Theorem~\ref{theorem_main2}). Furthermore if $| \lambda |$ is small enough,
  there exists a unique solution to the stationary
  equation~{\eqref{eq:yukawas}} (associated with
  {\eqref{eq:stochasticquantization}}) which is uniformly bounded in
  $\mathcal{C}^{\gamma} (\mathbb{R}^d, \mathcal{A}) \otimes \mathbb{C}^4$ by
  some constant $K > 0$ (that can be chosen independent of $h$ when $\| h
  \|_{\mathcal{C}^{\gamma}} < C$ for some $C > 0$).
\end{theorem}

\begin{proof}
  The proof is based on Theorem~\ref{theorem_existence2} and
  Theorem~\ref{theorem_longtime2}, whose hypotheses are satisfied in this
  setting because of Lemma~\ref{lemma_noise} and Lemma~\ref{lemma_drift}. The
  uniformity on the constant $K$ and the bounds on the admissible $\lambda$
  follows from Lemma~\ref{lemma_drift}, Theorem~\ref{theorem_drift} and the
  inequalities~{\eqref{eq:contraction3}} and~{\eqref{eq:contraction4}}.
\end{proof}

\subsection{Finite dimensional approximations in finite
volume}\label{sec:finit-dim-yukawa}

The next step in order to prove Theorem~\ref{theorem_main1} and
Theorem~\ref{theorem_main2} is to introduce finite dimensional approximations
of the Yukawa model. Consider the torus $\mathbb{T}_R^2 =\mathbb{R}^2
\backslash (2 \pi R\mathbb{Z})^2$ for $R > 0$ and the periodic version
$\tilde{\psi}_R$ of the Fermion field $\tilde{\psi}$ defined by
\[ \tilde{\psi}_R (f) \assign \sum_{n \in \mathbb{Z}^2} \tilde{\psi} (f (\cdot
   - 2 \pi R n)), \]
for all $f \in \mathcal{S} (\mathbb{R}^2) \otimes \mathbb{C}^4$. Then letting
$\tilde{\psi}_{R, \varepsilon} (x) \assign \tilde{\psi}_R
(\mathfrak{a}_{\varepsilon} (\cdot - x))$ we have that the function $x \in
\mathbb{T}^2_R \mapsto \tilde{\psi}_{R, \varepsilon} (x)$ is well defined and
$\mathcal{C}^{\gamma + 2 \delta} (\mathbb{T}^2, \mathcal{A})$. Moreover
$\tilde{\psi}_{R, \varepsilon} (x)$ is uniformly bounded in $\mathcal{A}$ and
\[ \| \tilde{\psi}_{R, \varepsilon} (x) - \tilde{\psi}_{\varepsilon} (x)
   \|_{\mathcal{C}^{\gamma + 2 \delta} (K, \mathcal{A} \otimes \mathbb{C}^4)}
   \rightarrow 0, \]
uniformly on compacts $K$ as $R \rightarrow \infty$. As a consequence we have
\begin{equation}
  \frac{\omega (\tilde{\psi} (f_1) \cdots \tilde{\psi} (f_n)
  e^{\mathcal{V}_{\varepsilon, h}})}{\omega (e^{\mathcal{V}_{\varepsilon,
  h}})} = \lim_{R \rightarrow \infty} \frac{\omega (\tilde{\psi}_R (f_1)
  \cdots \tilde{\psi}_R (f_n) e^{\mathcal{V}_{R, \varepsilon, h}})}{\omega
  (e^{\mathcal{V}_{R, \varepsilon, h}})}
  \label{eq:euclidean-yukawa-fermion-periodic}
\end{equation}
with
\[ \mathcal{V}_{R, \varepsilon, h} \assign \lambda^2 \int_{\mathbb{R}^4}
   \mathd x (h \bar{\psi}_{R, \varepsilon} \cdot \psi_{R, \varepsilon}) (x)
   \mathcal{G} (x - y) (h \bar{\psi}_{R, \varepsilon} \cdot \psi_{R,
   \varepsilon}) (y) \mathd x \mathd y, \]
for all compactly supported $f_1, \ldots, f_n, h \in \mathcal{S}
(\mathbb{R}^2)$.

\

Let now $(e_k)_{k \in \mathbb{Z}^2}$ be the orthonormal Fourier basis of (the
complex) $L^2 (\mathbb{T}_R^2)$ and $P_N : L^2 (\mathbb{T}_R^2) \rightarrow
L^2 (\mathbb{T}_R^2)$ the orthogonal projection to the closed subspace
generated by $\{ e_k : k \in \mathbb{Z}^2, | k | \leqslant N \}$. Let
$\tilde{\psi}_{N, R, \varepsilon} = P_N \tilde{\psi}_{R, \varepsilon}$ and
observe again that $\| \tilde{\psi}_{N, R, \varepsilon} (x) - \tilde{\psi}_{R,
\varepsilon} (x) \|_{\mathcal{C}^{\alpha} (\mathbb{R}^2)} \rightarrow 0$ as $N
\rightarrow \infty$, for any $\alpha > 0$. \ Defining
\[ \mathcal{V}_{N, R, \varepsilon, h} \assign \lambda^2 \int_{\mathbb{R}^4} (h
   \bar{\psi}_{N, R, \varepsilon} \cdot \psi_{N, R, \varepsilon}) (x)
   \mathcal{G}_{\varepsilon} (x - y) (h \bar{\psi}_{N, R, \varepsilon} \cdot
   \psi_{N, R, \varepsilon}) (y) \mathd x \mathd y \]
and we easily conclude that
\begin{equation}
  \frac{\omega (\tilde{\psi} (f_1) \cdots \tilde{\psi} (f_n)
  e^{\mathcal{V}_{\varepsilon, h}})}{\omega (e^{\mathcal{V}_{\varepsilon,
  h}})} = \lim_{R \rightarrow \infty} \lim_{N \rightarrow \infty} \frac{\omega
  (\tilde{\psi}_{N, R} (f_1) \cdots \tilde{\psi}_{N, R} (f_n)
  e^{\mathcal{V}_{N, R, \varepsilon, h}})}{\omega (e^{\mathcal{V}_{N, R,
  \varepsilon, h}})} \label{eq:euclidean-yukawa-fermion-periodic-galerkin} .
\end{equation}
Now consider the finite set $\Lambda_N = \{ k \in \mathbb{Z}^2 : | k |
\leqslant N \}$ and let
\[ \tilde{\theta} (k) \assign \int_{\mathbb{T}_R^2} \overline{e_k (x)}
   \tilde{\psi}_{N, R} (x) \mathd x, \qquad k \in \Lambda_N . \]
We can show easily that $(\tilde{\theta} (k))_{k \in \Lambda_N}$ is a
Grassmann Gaussian vector with correlation given by
\begin{equation}
  \omega (\tilde{\theta}_{\alpha} (k) \tilde{\theta}_{\beta} (\ell)) =
  \delta_{k + \ell = 0} \left(\begin{array}{cc}
    0 & \left( i \slashk + m_f \right)^{- 1}\\
    - \left( i \slashk + m_f \right)^{- 1 \Tau} & 0
  \end{array}\right)_{\alpha, \beta}
  \label{eq:free-corrrelation-fourier-galerkin}
\end{equation}
where $\alpha, \beta = 1, \ldots, 4$ denote the canonical coordinates in
$\mathbb{C}^4$ and $k, \ell \in \Lambda_N$. Moreover
\[ \tilde{\psi}_{N, R} (x) = \sum_{k \in \Lambda_N} e_k (x) \tilde{\theta}
   (k), \qquad x \in \mathbb{T}^2_R \]
and $\tilde{\psi}_{N, R, \varepsilon} = \mathfrak{a}_{\varepsilon}
\ast_{\mathbb{T}_R^2} \tilde{\psi}_{N, R}$ where the convolution is done on
$\mathbb{T}^2_R$ as indicated.

\

At this point we have reduced via two subsequent approximations the number of
degrees of freedom to a finite amount, described by the Grassmann random
variable $\tilde{\theta} \in \mathcal{G} (\Lambda_N \otimes \mathbb{C}^4)$.

\

We introduce the following equation
\begin{equation}
  \tilde{\Psi}^{h_R, s}_{\varepsilon, N, R, t} (v) = \lambda^2 \int_{-
  \infty}^t P_N (F_{\varepsilon, h_R, Y} (P_N (\tilde{\Psi}_{\varepsilon, N,
  R, \tau}^{h_R, s}))) (e^{(\Delta - m^2_f) (t - \tau)} v) \mathd \tau +
  \tilde{B}_{\varepsilon, R, t}^A (v) \label{eq:finitedimensional}
\end{equation}
where $v \in V$, $(\tilde{B}_{\varepsilon, R, t}^A)_{t \in \mathbb{R}}$ is the
Ornstein--Uhlenbeck motion on $\mathbb{T}^2_R$ obtained from
$(\tilde{B}_{\varepsilon, t}^A)_{t \in \mathbb{R}}$ via the restriction from
$\mathbb{R}^2$ to $\mathbb{T}^2$ of the periodicization
\[ \tilde{B}_{\varepsilon, R, t}^A (f) \assign \sum_{n \in \mathbb{Z}^2}
   \tilde{B}_{\varepsilon, t}^A (f (\cdot - 2 \pi R n)), \qquad f \in
   \mathcal{S} (\mathbb{R}^2), \]
and
\[ h_R (z) \assign \sum_{n \in \mathbb{Z}^2} h (z - 2 \pi R n), \qquad z \in
   \mathbb{R}^2 . \]
Obviously we have that
\[ (I - P_N) \tilde{\Psi}_{\varepsilon, N, R, t}^{h_R, s} = (I - P_N)
   \tilde{B}^A_{\varepsilon, R, t} . \]
This implies that $P_N (\tilde{\Psi}_{\varepsilon, N, R, t}^{h_R, s})$ solves
a finite dimensional Grassmann SDE. Furthermore the drifts $P_N
(F_{\varepsilon, h_R, Y} (P_N (\tilde{\Psi}_{\varepsilon, N, R, t}^{h_R,
s})))$ can be looked upon as the differentials of the function
\[ U (P_N (\tilde{\Psi})) = \int_{\mathbb{T}_L^4} h_L (x) h_L (y) (P_N (\Psi)
   P_N (\bar{\Psi})) (x) \mathcal{G} (x - y) (P_N (\Psi) P_N (\bar{\Psi})) (y)
   \mathd x \mathd y \]
times the matrix $P_N \mathfrak{a}_{\varepsilon}^{\asterisk 2} \asterisk C_f$
(that is the covariance of the finite dimensional Brownian motion $P_N
(\tilde{B}_{\varepsilon, R, t})$). This implies the following key lemma.

\begin{lemma}
  \label{lemma_finite1}For $| \lambda |$ small enough (that can be chosen in a
  way independent of the spatial cut-off $h$ of Def.~\ref{def:yukawa2}), $N
  \in \mathbb{N}$, $R > 0$, $\varepsilon > 0$, $t \in \mathbb{R}$, we have
  that for any $k \geqslant 1$ and $f_1, \ldots, f_k$ some linear combination
  of the first $N$ vectors $e_1, \ldots, e_N$ in the ONB $(e_k)_{k \in
  \mathbb{N}}$ of $L^2 (\mathbb{T}^2_R)$,
  \[ \frac{\omega (\tilde{\psi}_{N, R} (f_1) \cdots \tilde{\psi}_{N, R} (f_k)
     e^{\mathcal{V}_{N, R, \varepsilon, h}})}{\omega (e^{\mathcal{V}_{N, R,
     \varepsilon, h}})} = \omega (\tilde{\Psi}_{\varepsilon, N, R, t}^{h_R, s}
     (f_1) \cdots \tilde{\Psi}_{\varepsilon, N, R, t}^{h_R, s} (f_k)) . \]
\end{lemma}

\begin{proof}
  The proof is based on the fact that for any linear combination $f_i$ of
  $e_1, \ldots, e_N$ we have $\tilde{\Psi}_{\varepsilon, N, R, + \infty}^{h_L,
  s} (f_i) = P_N (\tilde{\Psi}_{\varepsilon, N, R, + \infty}^{h_R, s}) (f_i)$.
  Thus the thesis follows by letting $U (\tilde{\psi}) =\mathcal{V}_{N, R,
  \varepsilon, h} (\tilde{\psi}) / 2$ in Theorem~\ref{theorem_invariant2} and
  using Remark \ref{remark_invariant2}.
\end{proof}

\subsection{The infinite dimensional equation in finite
volume}\label{sec:cont-finite-vol-yukawa}

Consider the equation
\begin{equation}
  \tilde{\Psi}_{\varepsilon, R, t}^{h_R, s} (v) = \lambda^2 \int_{- \infty}^t
  F_{\varepsilon, h_R, Y} (\tilde{\Psi}_{\varepsilon, R, \tau}^{h_R, s})
  (e^{(\Delta - m^2_f) \tau} v) \mathd \tau + \tilde{B}^A_{\varepsilon, R, t}
  \label{eq:infinite1} .
\end{equation}
This is a modification of~{\eqref{eq:finitedimensional}} where the finite
dimensional projections $P_N$ are removed from the non-linearity
$F_{\varepsilon, h_R, Y}$.

\begin{remark}
  Lemma~\ref{lemma_drift} holds for the drift of
  equation~{\eqref{eq:infinite1}}, and it is easy to generalize
  Lemma~\ref{lemma_noise} to the OU process $(\tilde{B}_{\varepsilon, R,
  t}^A)_{t \in \mathbb{R}}$. This implies that it is possible to generalize
  Theorem~\ref{theorem_main3} to equation~{\eqref{eq:infinite1}} proving
  existence and uniqueness of its global solution (when $| \lambda |$ is small
  enough). Finally, it is important to note that the bound $K$ on the solution
  $\tilde{\Psi}_{\varepsilon, R, t}^{h_R, s}$ to
  equation~{\eqref{eq:infinite1}}, and the bounds on the constant $| \lambda
  |$ can be chosen independent of $h$ when $\| h \|_{\mathcal{C}^{\gamma}} <
  C$ for some constant $C > 0$.
\end{remark}

Equation~{\eqref{eq:infinite1}} is infinite dimensional and in order to
generalize Lemma~\ref{lemma_finite1} we cannot apply directly the finite
dimensional results contained in Theorem~\ref{theorem_invariant2} and
Remark~\ref{remark_invariant2}. For this reason we need to prove that the
solution to equation~{\eqref{eq:finitedimensional}} converges, as $N
\rightarrow + \infty$, to the solution to equation~{\eqref{eq:infinite1}}.

\begin{theorem}
  \label{theorem_convergence1}For $| \lambda |$ small enough, any $0 < q <
  \gamma$ and $R > 0$, we have that
  \[ \tilde{\Psi}_{\varepsilon, N, R, t}^{h_R, s} \rightarrow
     \tilde{\Psi}_{\varepsilon, R, t}^{h_R, s}, \]
  in $\mathcal{G}^q (C^{\infty} (\mathbb{T}^2) \otimes \mathbb{C}^4)$ when $N
  \rightarrow + \infty$, uniformly in $t \in \mathbb{R}$.
\end{theorem}

\begin{proof}
  We have that the difference $\Psi_{\varepsilon, N, R, t}^{h_R, s} -
  \Psi_{\varepsilon, R, t}^{h_R, s} $ satisfies the following inequality
  \begin{eqnarray}
    \frac{\| \tilde{\Psi}_{\varepsilon, N, R, t}^{h_R, s} -
    \tilde{\Psi}_{\varepsilon, R, t}^{h_R, s} \|_{\mathcal{C}^q}}{\lambda^2} &
    \lesssim & \int_{- \infty}^t e^{- \lambda_A (t - \tau)} \| (I - P_N)
    F_{\varepsilon, h_R, Y} (P_N (\tilde{\Psi}_{\varepsilon, N, R, \tau}^{h_R,
    s})) \|_{\mathcal{C}^q} \mathd \tau + \nonumber\\
    &  & + \int_{- \infty}^t e^{- \lambda_A (t - \tau)} g_F (\max (\|
    \tilde{\Psi}_{\varepsilon, N, R, \tau}^{h_R, s} \|_{\mathcal{C}^q}, \| P_N
    \tilde{\Psi}_{\varepsilon, N, R, \tau}^{h_R, s} \|_{\mathcal{C}^q}))
    \times \nonumber\\
    &  & \times (\| \tilde{\Psi}_{\varepsilon, N, R, \tau}^{h_R, s} - P_N
    \tilde{\Psi}_{\varepsilon, N, R, \tau}^{h_R, s} \|_{\mathcal{C}^q}) \mathd
    \tau + \nonumber\\
    &  & + \int_{- \infty}^t e^{- \lambda_A (t - \tau)} g_F (\max (\|
    \tilde{\Psi}_{\varepsilon, N, R, \tau}^{h_R, s} \|_{\mathcal{C}^q}, \|
    \tilde{\Psi}_{\varepsilon, R, \tau}^{h_R, s} \|_{\mathcal{C}^q})) \times
    \nonumber\\
    &  & \times \| \tilde{\Psi}_{\varepsilon, N, R, \tau}^{h_R, s} -
    \tilde{\Psi}_{\varepsilon, R, \tau}^{h_R, s} \|_{\mathcal{C}^q} \mathd
    \tau . \nonumber
  \end{eqnarray}
  Using Gronwall inequality (see, e.g., Theorem~1.3.3
  in~{\cite{Pachpatte1998}}) we obtain that
  \[ \| \tilde{\Psi}_{\varepsilon, N, R, t}^{h_R, s} -
     \tilde{\Psi}_{\varepsilon, R, t}^{h_R, s} \|_{\mathcal{C}^q} \lesssim (\|
     I - P_N \|_{\mathcal{L} (\mathcal{C}^{\gamma}, \mathcal{C}^q)}) \left( 1
     + \int_{- \infty}^t \exp (- (\lambda_A - \lambda^2 g_F (K)) (t - \tau))
     \mathd \tau \right) . \]
  Since $\lambda$ is small we have that $\lambda_A - \lambda^2 g_F (K) > 0$
  (this is exactly the request~{\eqref{eq:contraction4}} for $| \lambda | <
  \lambda_0 (K, L, F)$). This means that
  \[ \| \tilde{\Psi}_{\varepsilon, N, R, t}^{h_R, s} -
     \tilde{\Psi}_{\varepsilon, R, t}^{h_R, s} \|_{\mathcal{C}^q} \lesssim (\|
     I - P_N \|_{\mathcal{L} (\mathcal{C}^{\gamma}, \mathcal{C}^q)}) . \]
  On the other hand, since, \ $\gamma > q$ by assumption, we have $\| I - P_N
  \|_{\mathcal{L} (\mathcal{C}^{\gamma}, \mathcal{C}^q)} \rightarrow 0$ as $N
  \rightarrow + \infty$. The thesis is proved noting that the convergence
  $\tilde{\Psi}_{\varepsilon, N, R, t}^{h_R, s} \rightarrow
  \tilde{\Psi}_{\varepsilon, R, t}^{h_R, s}$ is uniform in $t$.
\end{proof}

An easy corollary of Theorem~\ref{theorem_convergence1},
Lemma~\ref{lemma_finite1}, and the
convergence~{\eqref{eq:euclidean-yukawa-fermion-periodic-galerkin}} is the
following:

\begin{corollary}
  \label{corollary_formula}For $| \lambda |$ small enough (that can be chosen
  independent of $h$ when $\| h \|_{\mathcal{C}^{\gamma}} < C$ where $C > 0$),
  for any $k \geqslant 1$, $f_1, \ldots, f_k \in C^{\infty} (\mathbb{T}^d_L)
  \otimes \mathbb{C}^4$, $R > 0$, and any $t \in \mathbb{R}$ we have
  \[ \frac{\omega (\tilde{\psi}_R (f_1) \cdots \tilde{\psi}_R (f_k)
     e^{\mathcal{V}_{R, \varepsilon, h}})}{\omega (e^{\mathcal{V}_{R,
     \varepsilon, h}})} = \omega (\tilde{\Psi}_{\varepsilon, R, t}^{h_R, s}
     (f_1) \cdots \tilde{\Psi}_{\varepsilon, R, t}^{h_R, s} (f_k)) . \]
\end{corollary}

\subsection{The infinite volume limit}\label{eq:yukawa-infinite-vol}

In order to tackle the infinite volume limit and prove
Theorem~\ref{theorem_main1} and Theorem~\ref{theorem_main2} we have to
introduce weighted Besov spaces (see the Appendix~\ref{appendix_Besov}
and~{\cite{TriebelIII}} for the details). Let $\ell \in \mathbb{R}$ and
consider $\rho_{\ell} : \mathbb{R}^d \rightarrow \mathbb{R}_+$ defined by
$\rho_{\ell} (x) = (1 + | x |^2)^{(- \ell / 2)}$. Consider the following norm
on $\mathcal{S}' (\mathbb{R}^d) \otimes \mathcal{A}$ (where as before
$\mathcal{A} = \mathcal{L} (\mathcal{H})$ and $\mathcal{H}$ is the fermionic
Hilbert space)
\[ \| u \|_{\mathcal{C}^s_{\ell} (\mathbb{R}^d, \mathcal{A})} = \sup_{j
   \geqslant - 1} \sup_{x \in \mathbb{R}^d} 2^{j s} \| \Delta_j u (x)
   \rho_{\ell} (x) \|_{\mathcal{A}} . \]
Define the space $\mathcal{C}^s_{\ell} (\mathbb{R}^d, \mathcal{A})$ as the
subspace of $\mathcal{S}' (\mathbb{R}^d) \otimes \mathcal{A}$ where $\| \cdot
\|_{\mathcal{C}^s_{\ell} (\mathbb{R}^d, \mathcal{A})}$ is finite. We define
also $\mathcal{G}^s_{\ell} (V) = \mathcal{C}^s_{\ell} (\mathbb{R}^d,
\mathcal{A}) \cap \mathcal{G}^{- \infty} (V)$ (recall that $V =\mathcal{S}
(\mathbb{R}^d) \otimes \mathbb{R}^n$ or $\mathcal{S} (\mathbb{R}^d) \otimes
\mathbb{C}^n$ for some $n \geqslant 1$) and the distance
$d_{\mathcal{G}^s_{\ell} (V)} (X - Y) \assign \| X - Y
\|_{\mathcal{C}^s_{\ell} (\mathbb{R}^d, \mathcal{A})}$.

\begin{remark}
  \label{remark_weight}It is important to note that, by
  Theorem~\ref{theorem_weight} below,
  \[ \sup_{x \in \mathbb{R}^d} \| \Delta_j u (x) \rho_{\ell} (x)
     \|_{\mathcal{A}} \sim \sup_{x \in \mathbb{R}^d} \| \Delta_j (\rho_{\ell}
     u) (x) \|_{\mathcal{A}}, \]
  where $\sim$ means that both quantities considered can be bounded from above
  and from below by some positive constant times the other quantity.
\end{remark}

\begin{lemma}
  \label{lemma_convergenceB}Let $\gamma, \delta$ be as stated in
  Remark~\ref{remark_alpha}, for any $s < \gamma + 2 \delta$ and any $\ell >
  0$ we have that
  \[ \sup_{t \in \mathbb{R}} d_{\mathcal{G}^s_{\ell} (V)}
     (\tilde{B}_{\varepsilon, t}^A, \tilde{B}^A_{\varepsilon, R, t})
     \rightarrow 0 \]
  when $R \rightarrow + \infty$.
\end{lemma}

\begin{proof}
  Using a method similar to the one in the proof of Lemma \ref{lemma_noise},
  it is simple to see that
  \[ \sup_{R > 1} (\sup_{t \in \mathbb{R}} \| (\Delta + 1)
     \tilde{B}^A_{\varepsilon, R, t} \|_{\mathcal{C}^{\gamma + 2 \delta - 2}
     (\mathbb{R}^2, \mathcal{A}) \otimes \mathbb{C}^4}) < + \infty,_{} \]
  which implies that, for any $s < \gamma + \delta$, we have
  \begin{equation}
    \| \Delta_i (\tilde{B}^A_{\varepsilon, R, t} - \tilde{B}^A_{\varepsilon,
    t}) \|_{\infty} \leqslant C_1 2^{- i (\gamma + 2 \delta)}
    \label{eq:weightB1}
  \end{equation}
  where the constant $C_1 \in \mathbb{R}_+$ does not depend on $i$ and $R >
  1$. This means that, fixing $s < s' < \gamma + 2 \delta$, we have
  \begin{equation}
    \| \tilde{B}^A_{\varepsilon, R, t} - \tilde{B}^A_{\varepsilon, t}
    \|_{\mathcal{C}^s_{\ell} (\mathbb{R}^2, \mathcal{A}) \otimes \mathbb{C}^4}
    \lesssim \sum_{i \geqslant - 1} 2^{i s'} \| \Delta_i
    (\tilde{B}^A_{\varepsilon, R, t} - \tilde{B}^A_{\varepsilon, t})
    \rho_{\ell} \|_{\infty} \label{eq:weightB5}
  \end{equation}
  and, thus by the uniform estimate {\eqref{eq:weightB1}} and Lebesgue's
  dominated convergence theorem, if, for any $i \geqslant - 1$, $\| \Delta_i
  (\tilde{B}^A_{\varepsilon, R, t} - \tilde{B}^A_{\varepsilon, t}) \rho_{\ell}
  \|_{\infty} \rightarrow 0$ as $R \rightarrow + \infty$, then $\|
  \tilde{B}^A_{\varepsilon, R, t} - \tilde{B}^A_{\varepsilon, t}
  \|_{\mathcal{C}^s_{\ell} (\mathbb{R}^2, \mathcal{A}) \otimes \mathbb{C}^4}
  \rightarrow 0$ as $R \rightarrow + \infty$, and thus, by inequality
  {\eqref{eq:weightB5}}, $\| \tilde{B}^A_{\varepsilon, R, t} -
  \tilde{B}^A_{\varepsilon, t} \|_{\mathcal{C}^s_{\ell} (\mathbb{R}^2,
  \mathcal{A}) \otimes \mathbb{C}^4} \rightarrow 0$ as $R \rightarrow +
  \infty$.
  
  Fix $i \geqslant 1$, then we have
  \begin{equation}
    \begin{array}{c}
      \| \Delta_i (\tilde{B}^A_{\varepsilon, R, t} - \tilde{B}^A_{\varepsilon,
      t}) (x) \rho_{\ell} (x) \|_{\mathcal{A} \otimes \mathbb{C}^4}
      \leqslant\\
      \leqslant \| \rho_{\ell} (x) \Xi_R (\tilde{K}_i (x - \cdot)
      \mathcal{G}_t (\cdot) \mathbb{I}_{[- \kappa R, \kappa R]^2}) -
      \rho_{\ell} (x)_{} \Xi (\tilde{K}_i (x - \cdot) \mathcal{G}_t (\cdot)
      \mathbb{I}_{[- \kappa R, \kappa R]^2}) \|_{\mathcal{A} \otimes
      \mathbb{C}^4} +\\
      + \| \rho_{\ell} (x) \Xi_R (\tilde{K}_i (x - \cdot) \mathcal{G}_t
      (\cdot) \mathbb{I}_{([- \kappa R, \kappa R]^2)^c}) - \rho_{\ell} (x)_{}
      \Xi (\tilde{K}_i (x - \cdot) \mathcal{G}_t (\cdot) \mathbb{I}_{([-
      \kappa R, \kappa R]^2)^c}) \|_{\mathcal{A} \otimes \mathbb{C}^4}
      \leqslant\\
      \leqslant \| \rho_{\ell} (x) \Xi_R (\tilde{K}_i (x - \cdot)
      \mathcal{G}_t (\cdot) \mathbb{I}_{([- \kappa R, \kappa R]^2)^c})
      \|_{\mathcal{A} \otimes \mathbb{C}^4} +\\
      + \| \rho_{\ell} (x) \Xi (\tilde{K}_i (x - \cdot) \mathcal{G}_t (\cdot)
      \mathbb{I}_{([- \kappa R, \kappa R]^2)^c}) \|_{\mathcal{A} \otimes
      \mathbb{C}^4},
    \end{array}
  \end{equation}
  where $\tilde{K}_i$ is given by
  \[ \tilde{K}_i (x) \assign (\Delta + m_f)^{1 / 4} (K_i), \]
  $K_i$ being the function related to the block $i$-th Littlewood-Paley, i.e.
  $\Delta_i f = K_i \asterisk f$, $\mathcal{G}_t$ is the Green function of the
  operator $e^{(\Delta + m^2) t}$, and where we use that $\Xi_R (\tilde{K}_i
  (x - \cdot) \mathcal{G}_t (\cdot) \mathbb{I}_{[- \kappa R, \kappa R]^2})
  =_{} \Xi (\tilde{K}_i (x - \cdot) \mathcal{G}_t (\cdot) \mathbb{I}_{[-
  \kappa R, \kappa R]^2})_{}$ since they coincide on $[- \kappa R, \kappa
  R]^2$ for any $0 \leqslant \kappa < \frac{1}{2}$. We focus on the term $J
  (x) \assign \| \rho_{\ell} (x) \Xi (\tilde{K}_i (x - \cdot) \mathcal{G}_t
  (\cdot) \mathbb{I}_{([- \kappa R, \kappa R]^2)^c} (\cdot)) \|_{\mathcal{A}
  \otimes \mathbb{C}^4}$. In this case we can write for $| x | > \kappa_1 R$
  \begin{eqnarray}
    J (x) & \lesssim & \frac{1}{(\kappa_1 R)^{\ell}} \sup_{x \in \mathbb{R}^2}
    \left. \left| \int_0^{+ \infty} \int_{\mathbb{R}^4} \tilde{K}_i (x - y_1)
    \tilde{K}_i (x - y_2) \right. \times \right. \nonumber\\
    &  & \left. \left. \left( \int_{([- \kappa R, \kappa R]^2)^c}
    \mathcal{G}_t (y_1 - z) \mathcal{G}_t (y_2 - z) \mathd z \right) \mathd
    y_1 \mathd y_2 \mathd t \right| \right. \nonumber\\
    & \lesssim & \frac{1}{(\kappa_1 R)^{\ell}} \left. \left|
    \int_{\mathbb{R}_+ \times \mathbb{R}^4} \exp (- \beta_i | y_1 |^{\alpha /
    2} - \beta_i | y_2 |^{\alpha / 2}) \times \right. \right. \nonumber\\
    &  & \left. \left. \times \left( \int_{\mathbb{R}^2} \mathcal{G}_t (y_1 -
    z) \mathcal{G}_t (y_2 - z) \mathd z \right) \mathd y_1 \mathd y_2 \mathd t
    \right| \right. \nonumber\\
    & \lesssim & \frac{1}{(\kappa_1 R)^{\ell}} \int_{\mathbb{R}^4} e^{(-
    \beta_i | y_1 |^{\alpha / 2} - \beta_i | y_2 |^{\alpha / 2})} \left|
    \int_0^{+ \infty} \exp \left( - \frac{| y_1 - y_2 |^2}{2 t} - 2 m^2 t
    \right) \frac{\mathd t}{t} \right| \mathd y_1 \mathd y_2 \nonumber\\
    & \lesssim & \frac{1}{(\kappa_1 R)^{\ell}} \int_{\mathbb{R}^4} e^{(-
    \beta_i | y_1 |^{\alpha / 2} - \beta_i | y_2 |^{\alpha / 2})} (| \log (|
    y_1 - y_2 |) | + 1) \mathd y_1 \mathd y_2 \nonumber\\
    & \lesssim & \frac{1}{(\kappa_1 R)^{\ell}}, \nonumber
  \end{eqnarray}
  where we used the fact that there are $\beta_i$ (depending on $i \geqslant -
  1$) and $\alpha \in (0, 1)$ such that $| \tilde{K}_i (y) | \lesssim e^{-
  \beta_i | y |^{\alpha}}$ (see Remark \ref{remark_decay} below). Consider now
  $| x | \leqslant \kappa_1 R$, then we have
  \begin{eqnarray}
    J (x) & \lesssim & \left| \int_0^{+ \infty} \int_{| y_1 |, | y_2 | >
    \kappa_2 R} \tilde{K}_i (x - y_1) \tilde{K}_i (x - y_2)_{} P (y_1, y_2, t)
    \mathd y_1 \mathd y_2 \mathd t \right| + \nonumber\\
    &  & \left| \int_0^{+ \infty} \int_{| y_1 |, | y_2 | < \kappa_2 R}
    \tilde{K}_i (x - y_1) \tilde{K}_i (x - y_2)_{} P (y_1, y_2, t) \mathd y_1
    \mathd y_2 \mathd t \right| + \nonumber\\
    &  & 2 \left| \int_0^{+ \infty} \int_{| y_1 | > R, | y_2 | < \kappa_2 R}
    \tilde{K}_i (x - y_1) \tilde{K}_i (x - y_2)_{} P (y_1, y_2, t) \mathd y_1
    \mathd y_2 \mathd t \right| \nonumber\\
    & \lesssim & J_1 (x) + J_2 (x) + J_3 (x), \nonumber
  \end{eqnarray}
  where $\kappa_1 < \kappa_2 < \kappa < \frac{1}{2}$ and
  \[ P (y_1, y_2, t) = \left( \int_{([- \kappa R, \kappa R]^2)^c}
     \mathcal{G}_t (y_1 - z) \mathcal{G}_t (y_2 - z) \mathd z \right) \]
  Then we get
  \begin{eqnarray}
    J_1 (x) & \lesssim & \int_0^{+ \infty} \int_{| y_1 |, | y_2 | > \kappa_2
    R} \exp (- \beta_i | x - y_1 |^{\alpha} - \beta_i | x - y_1 |^{\alpha}) P
    (y_1, y_2, t) \mathd y_1 \mathd y_2 \mathd t \nonumber\\
    & \lesssim & \int_0^{+ \infty} \int_{| y_1 |, | y_2 | > \kappa_2 R} \exp
    (- \beta_i (| y_1 | - | x |)^{\alpha} - \beta_i (| y_2 | - | x
    |)^{\alpha}) \times \nonumber\\
    &  & \left( \int_{\mathbb{R}^2} \mathcal{G}_t (y_1 - z) \mathcal{G}_t
    (y_2 - z) \mathd z \right) \mathd y_1 \mathd y_2 \mathd t \nonumber\\
    & \lesssim & (e^{- \beta_i (\kappa_2 - \kappa_1)^{\alpha} R^{\alpha}} (1
    + R^2)) \int_{\mathbb{R}^2} e^{(- \beta_i | y_1' |^{\alpha} - \beta_i |
    y_2' |^{\alpha})} \left| \int_0^{+ \infty} \exp \left( - \frac{| y_1 - y_2
    |^2}{2 t} - 2 m^2 t \right) \frac{\mathd t}{t} \right| \mathd y_1 \mathd
    y_2 \nonumber\\
    & \lesssim & (e^{- \beta_i (\kappa_2 - \kappa_1)^{\alpha} R^{\alpha}} (1
    + R^2)) \int_{\mathbb{R}^2} e^{(- \beta_i | y_1' |^{\alpha} - \beta_i |
    y_2' |^{\alpha})} (| \log (| y_1 - y_2 |) | + 1) \mathd y_1 \mathd y_2
    \nonumber\\
    & \lesssim & (e^{- \beta_i (\kappa_2 - \kappa_1)^{\alpha} R^{\alpha}} (1
    + R^2)) \nonumber
  \end{eqnarray}
  For $J_2 (x)$ we have
  \begin{eqnarray}
    J_2 (x) & \lesssim & (\kappa_2 R^2)^2 \sup_{| y_1 | < \kappa_2 R, | y_1 |
    < \kappa_2 R} \left| \int_0^{+ \infty} \left( \int_{([- \kappa R, \kappa
    R]^2)^c} \mathcal{G}_t (y_1 - z) \mathcal{G}_t (y_2 - z) \mathd z \right)
    \mathd t \right| \nonumber\\
    & \lesssim & (\kappa_2 R^2)^2 \int_0^{+ \infty} e^{\frac{(\kappa -
    \kappa_2)^2 R^2}{2 t} - 2 m^2 t} \frac{\mathd t}{t} \nonumber\\
    & \lesssim & \frac{(\kappa - \kappa_2)^{2 n}}{R^{2 n}} (\kappa_2 R^2)^2
    \int_0^{+ \infty} t^{n - 1} e^{- 2 m^2 t} \mathd t \lesssim \frac{(\kappa
    - \kappa_2)^2 \kappa_2^2}{R^{2 n - 4}} . \nonumber
  \end{eqnarray}
  where we used that $\exp (- x) \lesssim x^{- n}$ for any $n \in \mathbb{N}$.
  The term $J_3$ can be estimated in the same way combining the techniques
  used for $J_1$ and $J_2$, In conclusion
  \[ \sup_{x \in \mathbb{R}^2} J (x) \lesssim \left( \frac{(\kappa -
     \kappa_2)^2 \kappa_2^2}{R^{2 n - 4}} + \frac{1}{(\kappa_1 R)^{\ell}} +
     e^{- \beta_i (\kappa_2 - \kappa_1)^{\alpha} R^{\alpha}} (1 + R^2) \right)
     \rightarrow 0 \]
  when $n > 2$ and $\ell > 0$.
  
  For $\| \Xi_R (\tilde{K}_i (x - \cdot) \mathcal{G}_t (\cdot) \mathbb{I}_{([-
  \kappa R, \kappa R]^2)^c}) \|_{\mathcal{A} \otimes \mathbb{C}^4}$ we have
  \begin{equation}
    \begin{array}{c}
      \| \Xi_R (\tilde{K}_i (x - \cdot) \mathcal{G}_t (\cdot) \mathbb{I}_{([-
      \kappa R, \kappa R]^2)^c}) \|_{\mathcal{A} \otimes \mathbb{C}^4}
      \lesssim\\
      \lesssim \left| \int_0^{+ \infty} \int_{\mathbb{R}^4} \tilde{K}_i (x -
      y_1) \tilde{K}_i (x - y_2) \nobracket \nobracket P_R (y_1, y_2, t)
      \mathd y_1 \mathd y_2 \mathd t \right|,
    \end{array} \label{eq:weightB7}
  \end{equation}
  where
  \[ P_R (y_1, y_2, t) \assign \int_{\mathbb{T}_R^2 \backslash ([- \kappa R,
     \kappa R]^2)^c} \mathcal{G}_t (y_1 - z) \mathcal{G}_t (y_2 - z) \mathd z
     + \sum_{z' \in R\mathbb{Z}^2} \int_{\mathbb{T}_R^2} \mathcal{G}_t (y_1 -
     z) \mathcal{G}_t (y_2 - z) \mathd z. \]
  By noting that for $| y_1 |, | y_2 | > \kappa_2 R$
  \[ \int_0^{+ \infty} P_R (y_1, y_2, t) \mathd t \leqslant C_2 (| \log (| y_1
     - y_2 |) | + 1) \]
  for some constant $C_2 \in \mathbb{R}_+$ independent of $R > 1$ and, if one
  of the inequalities $| y_1 | < \kappa_2 R$ or $| y_2 | < \kappa_2 R$ holds,
  then
  \[ P_R (y_1, y_2, t) \leqslant \frac{\exp \left( - \beta' \frac{(\kappa -
     \kappa_2)^2 R^2}{t} + 2 m^2 t \right)}{t} . \]
  Then using inequality {\eqref{eq:weightB7}}, it is possible to prove that
  $\sup_{x \in \mathbb{R}^2} \| \rho_{\ell} (x) \Xi_R (\tilde{K}_i (x - \cdot)
  \mathcal{G}_t (\cdot) \mathbb{I}_{([- \kappa R, \kappa R]^2)^c})
  \|_{\mathcal{A} \otimes \mathbb{C}^4} \rightarrow 0$ as $R \rightarrow +
  \infty$ by exploiting an argument analogous to the one mentioned obove for
  proving $\sup_{x \in \mathbb{R}^2} \| \rho_{\ell} (x) \Xi (\tilde{K}_i (x -
  \cdot) \mathcal{G}_t (\cdot) \mathbb{I}_{([- \kappa R, \kappa R]^2)^c})
  \|_{\mathcal{A} \otimes \mathbb{C}^4} \rightarrow 0$.
\end{proof}

We now prove the analogue of Theorem~\ref{theorem_convergence1} for the case
where the space cut-off $R$ is removed (infinite volume case):

\begin{theorem}
  \label{theorem_torusplane}Let $| \lambda |$ be sufficiently small and $0 < q
  < \gamma$, then $\tilde{\Psi}_{\varepsilon, R, t}^{h_R, s} \rightarrow
  \tilde{\Psi}_{\varepsilon, t}^{h, s}$ in $\mathcal{G}^q_{\ell} (\mathcal{S}
  (\mathbb{R}^2) \otimes \mathbb{C}^4)$ uniformly in $t \in \mathbb{R}$ when
  $R \rightarrow + \infty$.
\end{theorem}

\begin{proof}
  The proof is based on the fact that for any $\tilde{\Psi}_1, \tilde{\Psi}_2
  \in \mathcal{} \mathcal{G}^{\gamma}_0 (\mathcal{S} (\mathbb{R}^2) \otimes
  \mathbb{C}^4)$ (it is important to note that in these spaces we take the
  constant weight) we have, by Theorem~\ref{theorem_multiplication} in
  Appendix \ref{appendix_Besov},
  \[ \| F_{\varepsilon, h_R, Y} (\tilde{\Psi}_1) - F_{\varepsilon, h, Y}
     (\tilde{\Psi}_2) \|_{\mathcal{C}^q_{\ell} (\mathbb{R}^2, \mathcal{A}
     \otimes \mathbb{C}^4)} \lesssim g_F (\max (\| h_R \tilde{\Psi}_1
     \|_{\mathcal{C}^q}, \| h \tilde{\Psi}_1 \|_{\mathcal{C}^q}, \| h
     \tilde{\Psi}_2 \|_{\mathcal{C}^q})) \times \]
  \begin{equation}
    \times (\| h \|_{\mathcal{C}^q} \| \tilde{\Psi}_1 - \tilde{\Psi}_2
    \|_{\mathcal{C}^q_{\ell} (\mathbb{R}^2, \mathcal{A} \otimes \mathbb{C}^4)}
    + \| \Psi_1 \|_{\mathcal{C}^q_0 (\mathbb{R}^2, \mathcal{A} \otimes
    \mathbb{C}^4)} \| h_R - h \|_{\mathcal{C}^q_{\ell}})
    \label{eq:inequalityh},
  \end{equation}
  where $g_F$ is the function in Remark~\ref{remark_drift} and $h$ the IR
  cut-off in Definition~\ref{def:yukawa2}. Using a reasoning similar to the
  one exploited in the proof of Theorem~\ref{theorem_convergence1} we get, for
  any $t \in \mathbb{R}$,
  \[ \| \tilde{\Psi}_{\varepsilon, R, t}^{h_R, s} -
     \tilde{\Psi}_{\varepsilon, t}^{h, s} \|_{\mathcal{C}^q_{\ell}
     (\mathbb{R}^2, \mathcal{A} \otimes \mathbb{C}^4)} \lesssim (\| h_R - h
     \|_{\mathcal{C}^q_{\ell}} + \sup_{\tau \in \mathbb{R}} \|
     \tilde{B}_{\varepsilon, \tau}^A - \tilde{B}^A_{\varepsilon, R, \tau}
     \|_{\mathcal{C}^q_{\ell} (\mathbb{R}^2, \mathcal{A} \otimes
     \mathbb{C}^4)}) \times \]
  \[ \times \left( 1 + \int_{- \infty}^t \exp (- (\lambda_A - \lambda^2 g_F
     (K)) (t - \tau)) \mathd \tau \right) . \]
  By $| \lambda |$ small enough, the bound converges to $0$ as $R \rightarrow
  + \infty$ uniformly in $t \in \mathbb{R}$, since, by
  Lemma~\ref{lemma_convergenceB},
  \[ \lim_{R \rightarrow + \infty} \sup_{\tau \in \mathbb{R}} \|
     \tilde{B}_{\varepsilon, \tau}^A - \tilde{B}^A_{\varepsilon, R, \tau}
     \|_{\mathcal{C}^q_{\ell} (\mathbb{R}^2, \mathcal{A} \otimes
     \mathbb{C}^4)} \rightarrow 0 \]
  and $\| h_R - h \|_{\mathcal{C}^q_{\ell}} \rightarrow 0$, as $R \rightarrow
  + \infty$, by the definition of $h_R$.
\end{proof}

Hereafter we denote by $\tilde{\Psi}_{\varepsilon, t}^s$ the solution to
equation~{\eqref{eq:stochasticquantization}} with $h \equiv 1$ (so with the
spatial cut-off $h$ removed).

\begin{theorem}
  \label{theorem_cutoffh}Using the previous notations and hypotheses, let
  $(h_n : \mathbb{R}^2 \rightarrow \mathbb{R}^{}_{})_{n \geqslant 1}$ be a
  sequence of smooth functions with compact supports converging to $1$
  uniformly on compacts in the norm of $\mathcal{C}^{\gamma}$ and with a
  uniform bound in $\mathcal{C}^{\gamma}$. For any $0 < q < \gamma$ we have
  \[ \tilde{\Psi}_{\varepsilon, t}^{h_n, s} \rightarrow
     \tilde{\Psi}_{\varepsilon, t}^s \]
  in $\mathcal{G}^q_{\ell} (\mathcal{S} (\mathbb{R}^2) \otimes \mathbb{C}^4)$,
  for any $\ell > 0$ and uniformly in $t \in \mathbb{R}$ when $n \rightarrow +
  \infty$.
\end{theorem}

\begin{proof}
  The proof is analogous to the one of Theorem~\ref{theorem_torusplane} where
  in inequality~{\eqref{eq:inequalityh}} we have $\| h_n - 1
  \|_{\mathcal{C}^q_{\ell}}$ (instead of $\| h_R - h
  \|_{\mathcal{C}^q_{\ell}}$) which converges to $0$, as $n \rightarrow +
  \infty$, by hypothesis, when $\ell > 0$.
\end{proof}

\begin{proof*}{Proof of Theorem~\ref{theorem_main1} and
Theorem~\ref{theorem_main2}}
  The proof is analogous to the one of Corollary~\ref{corollary_formula} where
  Theorem~\ref{theorem_convergence1} is replaced by
  Theorem~\ref{theorem_torusplane} and Theorem~\ref{theorem_cutoffh}
  respectively.
\end{proof*}

These theorems provide an alternative construction of the finite, resp.
infinite, volume Euclidean Yukawa model with UV cut-off as an invariant state
given by the solution of an SPDE for a Grassmann valued process. In this sense
also a major step in the program of the stochastic quantization for the Yukawa
model has been achieved.

\appendix\section{Remarks on the Grassmann Gaussians}\label{app:GP}

In this appendix we provide some standard facts on real structures on complex
(pre-) Hilbert spaces and then discuss an alternative realization of Grassmann
Gaussian variables on a tracial state.

\

There are various ways to constructs Grassmann Gaussians and while in the
main body of the paper we used the approach of Osterwalder and Schrader, based
on the CAR algebra of creation and annihilation operators on Fock space and
the state generated by the Fock vacuum, here we would like to discuss other
approaches based on self-dual Clifford algebras and their associated tracial
state. This goes back to the work of Gross~{\cite{Gross_Euclidean_1976}} and
Palmer~{\cite{palmer_euclidean_1980}} and to the ``real-wave'' representation
of the Clifford algebra introduced by
Segal~{\cite{segal_tensor_1956,Segalbook}} and nicely described by Gross
in~{\cite{gross_existence_1972}}. For more background information, aside from
the already cited works~{\cite{Gross_Euclidean_1976}}
and~{\cite{palmer_euclidean_1980}}, we refer to
Araki~{\cite{araki_representations_1964,araki_diagonalization_1968,araki_quasifree_1970}}.

\subsection{Gross-Palmer formulation}

Let $V$ be a complex Hilbert space with Hermitian scalar product $(\cdot,
\cdot)_V$ anti-linear in the left variable. We denote $A^{\ast}$ the usual
Hermitian adjoint of the linear operator $A : V \rightarrow V$.

\begin{definition}
  \label{GPdef:real-str}A real structure $\varkappa : V \rightarrow V$
  compatible with $(\cdummy, \cdummy)_V$ is an anti-unitary involution, i.e a
  map such that, for all $\alpha, \beta \in \mathbb{C}$, $v, w \in V$,
  \[ \varkappa (\alpha v + \beta w) = \bar{\alpha} \varkappa v + \bar{\beta}
     \varkappa w, \quad (\varkappa v, \varkappa w)_V = \overline{(v, w)}_V,
     \quad \varkappa \varkappa = 1. \]
\end{definition}

\begin{remark*}
  Since we will deal with Hilbert spaces we will always assume a real
  structure to be compatible with the scalar product and we will not
  explicitly say that it is compatible. 
\end{remark*}

For a real structure $\varkappa$ the following properties hold:
\begin{itemizedot}
  \item Let $\tmop{Re}_{\varkappa} v \assign (v + \varkappa v) / 2$, \
  $\tmop{Im}_{\varkappa} v \assign (v - \varkappa v) / (2 i)$, then $v =
  \tmop{Re}_{\varkappa} v + i \tmop{Im}_{\varkappa} v$, $\varkappa
  \tmop{Re}_{\varkappa} = \tmop{Re}_{\varkappa}$, \ $\varkappa
  \tmop{Im}_{\varkappa} = - \tmop{Im}_{\varkappa}$, $\tmop{Re}_{\varkappa} =
  \tmop{Im}_{\varkappa} i$, \ $\tmop{Im}_{\varkappa} = - \tmop{Re}_{\varkappa}
  i$, so
  \[ V = \tmop{Re}_{\varkappa} V + i \tmop{Im}_{\varkappa} V =
     \tmop{Re}_{\varkappa} V + i \tmop{Re}_{\varkappa} i V, \]
  which is a direct sum in the sense of vector spaces but it is not orthogonal
  with respect to $(\cdummy, \cdummy)_V$.
  
  \item The form
  \[ \llangle v, w \rrangle_{\varkappa} \assign (\varkappa v, w)_V \]
  is bilinear, symmetric and non-degenerate.
  
  \item For any linear operator $A : V \rightarrow V$ we define its
  $\varkappa$-transpose
  \[ A^{\varkappa} \assign \varkappa A^{\ast} \varkappa \]
  which satisfies $\llangle v, A w \rrangle_{\varkappa} = \llangle
  A^{\varkappa} v, w \rrangle_{\varkappa}$.
  
  \item If $\kappa$ is another real structure on $V$, then letting $U \assign
  \kappa \varkappa$ we have that $U$ is a linear operator which is invertible
  and $(U v, U w)_V = (v, w)_V$, that is $U$ is unitary.
\end{itemizedot}


Note that $V$ is a space of complex functions, then the map $\mathrm{c}: V
\rightarrow V$ given by taking the complex conjugate is a real structure.

\

Let $\Gamma_a V$ be the antisymmetric Fock space over the complex pre-Hilbert
space $V$, i.e. the completion of the vector space $\Lambda V$ with respect to
the scalar product
\[ (v_1 \cdots v_n, w_1 \cdots w_m)_{\Gamma_a V} = \delta_{n, m} \underset{1
   \leqslant j, \ell \leqslant n}{\det} (v_j, w_{\ell})_V, \qquad v_1, \ldots,
   w_m \in V, n, m \geqslant 0. \]
Let $\tmop{CAR} (V)$ denote the algebra of canonical anti-commutation
relations (CAR) over $V$, that is the unital, associative, complex,
$\ast$-algebra generated by the identity and the generators $c (v), a (v), v
\in V$ which satisfy the following relations, for $\alpha, \beta \in
\mathbb{C}$, $v, w \in V$: $c (\alpha v + \beta w) = \alpha c (v) + \beta c
(w)$, $c (v) = a (v)^{\ast}$, and
\[ 0 = \{ a (v), a (w) \} = \{ c (v), c (w) \}, \qquad \{ a (v), c (w) \} =
   (v, w)_V 1, \]
where $\{ a, b \}$ stands for the anticommutator between $a$ and b i.e. $\{ a,
b \} = a b + b a$. Note in particular that the creation operators $c (v)$ are
linear in $v \in V$, whereas the annihilation operators $a (v)$ are
anti-linear in $v \in V$.

The algebra $\tmop{CAR} (V)$ has the following standard representation on
$\Gamma_a V$. Let $(C_v)_{v \in V}$ be the linear bounded operators in
$\Gamma_a V$ acting on $\Gamma_a V$ as $C_v F = v F$ for any $F \in \Lambda V$
and let $A_v = C_v^{\ast}$ be their adjoint. Then $v \mapsto C_v$ is a linear
map, $v \mapsto A_v$ is anti-linear and
\[ \{ C_v, C_w \} = \{ A_v, A_w \} = 0, \qquad \{ A_v, C_w \} = (v, w)_V,
   \qquad v, w \in V. \]
So $(c (v), (v)) \mapsto (C_v, A_v)$, $v \in V$, is a representation of the
CAR algebra $\tmop{CAR} (V)$ on $\Gamma_a V$. Let $\Omega$ be the vacuum
vector of $\Gamma_a V$, i.e the unit element of $\Lambda V$ and let $\omega :
\mathcal{L} (\Gamma_a V) \rightarrow \mathbb{C}$ the vacuum state defined as
$\omega (F) = (\Omega, F \Omega)_{\Gamma_a V}$ for $F \in \mathcal{L}
(\Gamma_a V)$.

\

Recall that we constructed the centered complex Grassmann Gaussian $X \in
\mathcal{G} (V)$ associated to a complex Hilbert space $V$ with real structure
$\varkappa$ and covariance $S : V \rightarrow V$ via the CAR algebra $(C_v,
A_v)_{v \in V}$ as
\[ X (v) = C_{S v} + A_{\varkappa v}, \qquad v \in V \]
on the algebraic probability space $(\mathcal{L} (\Gamma_a V), \omega)$. In
particular let us recall that we have (cfr.~{\eqref{eq:pfaff}})
\begin{equation}
  \omega (X (f_1) \cdots X (f_n)) = \underset{1 \leqslant i, j \leqslant
  n}{\tmop{Pf}} \llangle f_i, S f_j \rrangle_{\varkappa},
  \label{eq:complex-gauss-GP}
\end{equation}


We shall now collect some standard facts on Clifford $C^{\ast}$-algebras, and
in particular their realization on the Fock space and the relation between the
Fock vacuum and the unique tracial state on the Clifford algebra.

\begin{definition}
  \label{GPdef:Clifford}Let $E$ be a complex vector space and $\llangle
  \cdummy, \cdummy \rrangle$ a non-degenerate symmetric bilinear form on $E.$
  We denote by $\mathbb{C}l (E, \llangle \cdummy, \cdummy \rrangle)$ the
  complex Clifford algebra on $(E, \llangle \cdummy, \cdummy \rrangle)$, that
  is the unital, associative algebra, over $\mathbb{C}$, generated by the
  identity $1$ and the symbols $(\gamma (v))_{v \in E}$ which satisfy $\{
  \gamma (v), \gamma (w) \} = 2 \llangle v, w \rrangle 1$, $\gamma (\alpha v +
  \beta w) = \alpha \gamma (v) + \beta \gamma (w)$, for all $\alpha, \beta \in
  \mathbb{C}, v, w \in E$.
\end{definition}

\begin{definition}
  \label{GPdef:self-dual}Let $V$ be a complex pre-Hilbert space with a real
  structure $\varkappa$. We define a self-dual CAR algebra
  $\mathbb{C}l_{\tmop{SD}} (V, \varkappa)$ to be a complex Clifford algebra
  with respect to the symmetric bilinear form $\llangle \cdummy, \cdummy
  \rrangle_{\varkappa}$, where in addition the generators satisfy $\gamma
  (v)^{\ast} = \gamma (\varkappa v)$ for $v \in V$.
\end{definition}

\begin{proposition}
  \label{GPpro:trace}If $\| \cdummy \|$ denotes the (unique) $C^{\ast}$-norm
  of $\mathbb{C}l_{\tmop{SD}} (V, \varkappa)$, then $\| B (v) \| = (v, v)_V^{1
  / 2} $, where $(B (v))_{v \in V}$ are the generators of
  $\mathbb{C}l_{\tmop{SD}} (V, \mathrm{U})$. The $C^{\ast}$-completion
  $\overline{\mathbb{C}l}_{\tmop{SD}} (V, \varkappa)$ admits a unique trace,
  that is, a state $\tau : \overline{\mathbb{C}l}_{\tmop{SD}} \rightarrow
  \mathbb{C}$ with the properties of being even (it is zero on the odd part of
  $\mathbb{C}l_{\tmop{SD}}$) and central ($\tau (A B) = \tau (B A)$, $A, B \in
  \mathbb{C}l_{\tmop{SD}}$). Moreover we have, $k \in \mathbb{N}$,
  \begin{equation}
    \tau (B (v_1) \cdots B (v_{2 k})) = \underset{1 \leqslant \ell, m
    \leqslant 2 k}{\tmop{Pf}} [\tmop{sgn} (\ell - m) \llangle v_{\ell}, v_m
    \rrangle_{\varkappa}], \label{eq:trace-GP}
  \end{equation}
  and $\tau (B (v_1) \cdots B (v_{2 k + 1})) = 0$, where $\tmop{Pf}$ denotes
  the Pfaffian and $\tmop{sgn}$ is the sign function.
\end{proposition}

\begin{proof}
  The relation $\| B (v) \| = (v, v)^{1 / 2}, v \in V$, is a consequence of
  the commutation relation $\{ B (v)^{\ast}, B (v) \} = \{ B (\varkappa v), B
  (v) \} = (v, v)$ and the $C^{\ast}$-algebra property of the norm. It follows
  at once from the tracial property of $\tau$ that $\tau (B (v) B (w)) =
  \tfrac{1}{2} \tau (\{ B (v), B (w_{}) \}) = \llangle v, w
  \rrangle_{\varkappa}$. The formula for the trace follows by induction.
  Indeed, by writing
  \[ \tau (B (v_1) \cdots B (v_{2 k})) = \frac{1}{2} \tau (B (v_1) \cdots B
     (v_{2 k}) + B (v_2) \cdots B (v_{2 k}) B (v_1)), \]
  and noting that
  \[ B (v_1) \cdots B (v_{2 k}) + B (v_2) \cdots B (v_{2 k}) B (v_1) = \sum_{j
     = 2}^{2 k} (- 1)^{j - 1} \{ B (v_1), B (v_j) \} B (v_2) \cdots \cancel{B
     (v_j)} \cdots B (v_{2 k}), \]
  we get, using the induction hypothesis for $\tau \left( B (v_2) \cdots
  \cancel{B (v_j)} \cdots B (v_n) \right)$,
  \begin{eqnarray*}
    \tau (B (v_1) \cdots B (v_{2 k})) & = & \sum^{2 k}_{j = 2} (- 1)^{j - 1}
    \{ B (v_1), B (v_j) \} \tau \left( B (v_2) \cdots \cancel{B (v_j)} \cdots B
    (v_{2 k}) \right)\\
    & = & \sum_{j = 2}^{2 k} (- 1)^{j - 1} \llangle v_1, v_j
    \rrangle_{\varkappa}  \underset{\ell, m \in \left\{ 2, \ldots, \centernot{j},
    \ldots, 2 k \right\}}{\tmop{Pf}} (\tmop{sgn} (m - \ell) \llangle v_{\ell},
    v_m \rrangle_{\varkappa})\\
    & = & \underset{1 \leqslant m, \ell \leqslant 2 k}{\tmop{Pf}} (\tmop{sgn}
    (m - \ell) \llangle v_{\ell}, v_m \rrangle_{\varkappa}),
  \end{eqnarray*}
  where in the last line we used the extension rule for Pfaffians
  (cf.~{\cite{caianiello_combinatorics_1973}} formula~(1.5)).
  
  Hence the trace $\tau$ is now uniquely defined for $n = 2 k$ even. By
  hypothesis the trace is normalized, hence $\tau (1) = 1$, and even, i.e.
  $\tau (B (v_1) \cdots B (v_{2 k + 1})) = 0$, $k \in \mathbb{N}$ (actually
  the fact of being even is implied by centrality when $V$ is even
  dimensional). Therefore we have proved the formula in the statement and the
  uniqueness of $\tau$. Finally, $\tau$ can be extended to the
  $C^{\ast}$-completion by continuity
  (cf.~{\cite{plymen_spinors_1994}}~Theorem 1.2.8).
\end{proof}

\begin{definition}
  \label{def:clifford-car}Assume that we have a complex pre-Hilbert space $V$
  with a real structure $\varkappa$. Let $\tmop{CAR} (V)$ be the CAR algebra
  over $V$ with generators $C_v, A_v, v \in V$. We define the Gross-Palmer
  (G-P) fields $B_{\tmop{GP}} (v), v \in V$, to be
  \[ B_{\tmop{GP}} (v) \assign C_v + A_{\varkappa v}, \qquad v \in V. \]
\end{definition}

\begin{remark}
  \label{GPrem:GP-CAR}Let $(B (v))_{v \in V}$, be the generators of
  $\mathbb{C}l_{\tmop{SD}} (V, \varkappa)$ then the $\ast$-homomorphism which
  extends $B (v) \mapsto B_{\tmop{GP}} (v)$ to the whole
  $\mathbb{C}l_{\tmop{SD}} (V, \varkappa)$ is a realization of
  $\mathbb{C}l_{\tmop{SD}} (V, \varkappa)$ inside $\tmop{CAR} (V)$. Indeed,
  since $\varkappa^2 = 1$, we have
  \[ B_{\tmop{GP}} (\varkappa f) = C_{\varkappa f} + A_{\varkappa^2 f} =
     C_{\varkappa f} + A_f = B_{\tmop{GP}} (f)^{\ast} . \]
  Moreover, since $(\varkappa f, \varkappa g) = \overline{(f, g)}$,
  \begin{eqnarray*}
    \{ B_{\tmop{GP}} (f), B_{\tmop{GP}} (g) \} & = & \{ C_f, A_{\varkappa g}
    \} + \{ A_{\varkappa f}, C_g \}\\
    & = & (\varkappa g, f)_V + (\varkappa f, g)_V = \llangle g, f
    \rrangle_{\varkappa} + \llangle f, g \rrangle_{\varkappa}\\
    & = & 2 \llangle g, f \rrangle_{\varkappa}
  \end{eqnarray*}
  Hence $(B_{\tmop{GP}} (f))_{f \in V}$ give a representation of
  $\mathbb{C}l_{\tmop{SD}} (V, \varkappa)$ within $\tmop{CAR} (V)$.
\end{remark}

\begin{proposition}
  Let $(B_{\tmop{GP}} (v))_{v \in V}$ be as above. Then $(B_{\tmop{GP}}
  (v))_{v \in V}$ is a $\ast$--representation of $\mathbb{C}l_{\tmop{SD}} (V,
  \varkappa)$ on $\Gamma_a V$ and the Fock vacuum state $\omega$ restricted to
  this representation coincides with the tracial state $\tau$ of
  $\mathbb{C}l_{\tmop{SD}} (V, \varkappa)$, namely
  \begin{equation}
    \tau (B (v_1) \cdots B (v_n)) = \omega (B_{\tmop{GP}} (v_1) \cdots
    B_{\tmop{GP}} (v_n)), \qquad v_1, \ldots, v_n \in V, n \in \mathbb{N},
    \label{eq:equiv-gp}
  \end{equation}
  where $(B (v))_{v \in V}$ are the generators of $\mathbb{C}l_{\tmop{SD}} (V,
  \varkappa)$. 
\end{proposition}

\begin{proof}
  Since $\varkappa^2 = 1$, we have
  \[ B_{\tmop{GP}} (v)^{\ast} = C_v^{\ast} + A_{\varkappa v}^{\ast} =
     A_{\varkappa \varkappa v} + C_{\varkappa v} = B_{\tmop{GP}} (\varkappa
     v), \qquad v \in V. \]
  Moreover,
  \begin{eqnarray*}
    \{ B_{\tmop{GP}} (v), B_{\tmop{GP}} (w) \} & = & \{ C_v, A_{\varkappa w}
    \} + \{ A_{\varkappa v}, C_w \}\\
    & = & (\varkappa w, v)_V + (\varkappa v, w)_V = \llangle w, v
    \rrangle_{\varkappa} + \llangle v, w \rrangle_{\varkappa}\\
    & = & 2 \llangle v, w \rrangle_{\varkappa}
  \end{eqnarray*}
  Hence $(B_{\tmop{GP}} (v))_{v \in V}$ give a representation of
  $\mathbb{C}l_{\tmop{SD}} (V, \varkappa)$ inside $\tmop{CAR} (V)$. A
  straightforward computation shows that the statement~{\eqref{eq:equiv-gp}}
  holds for $n = 2$. Then the result follows by noting that both states,
  $\tau$ and $\omega$, are quasi-free (in the sense of
  Araki~{\cite{araki_quasifree_1970}}), hence they are uniquely determined by
  their evaluation for $n = 2$.
\end{proof}

\subsection{A tracial formulation of Grassmann variables}

In this subsection we describe one way to relate the expectation of a complex
Gaussian Grassmann random variable to the trace on a self-dual Clifford
algebra. In the approach given in this subsection the construction of
isotropic subspaces for certain bilinear forms reveals itself as crucial. One
remarkable aspect of this approach is its ``dual'' nature with respect to the
Osterwalder-Schrader approach. Perhaps its biggest disadvantage is that the
random variables will be realized as elements of a Clifford algebra at the
expenses of their anti-commutativity.

\

Assume we are given an orthogonal decomposition $V = V_+ \oplus V_-$ such
that $\varkappa V_{\pm} = V_{\mp}$. Then the subspaces $V_{\pm}$ are isotropic
subspaces for the bilinear form $\llangle \cdot, \cdot \rrangle_{\varkappa}$,
i.e. $0 = \llangle v, w \rrangle_{\varkappa}$ if $v, w \in V_+$ or $v, w \in
V_-$. That it is always possible to find such an isotropic decomposition is
the subject of next proposition.

\begin{proposition}
  \label{GPpro:orthogonal-dec}Let $\varkappa$ be a real structure on $V$ and
  assume $V$ has even or infinite dimension. Then $V$ admits an orthogonal
  decomposition $V = V_+ \oplus V_-$ such that $\varkappa V_{\pm} = V_{\mp}$.
  
  Moreover, if $\mathrm{c}$ is another real structure on $V$ then $\varkappa =
  U\mathrm{c}$ with $U$ a unitary operator on $V$ such that $U^{\mathrm{c}} =
  U$, where, as before, $U^{\mathrm{c}} \assign \mathrm{c}U^{\ast}
  \mathrm{c}$, with $U^{\ast}$ the adjoint of $U$.
  
  Finally, if $\mathrm{c}$ maps each $V_+, V_-$ to itself, then, with respect
  to the decomposition $V = W \oplus W$ we have the block decomposition
  \begin{equation}
    U = \left( \begin{array}{cc}
      0 & u\\
      u^{\mathrm{c}} & 0
    \end{array} \right) \label{GPeq:UsDecomposition},
  \end{equation}
  for some unitary operator $u : W \rightarrow W$. Conversely if $U$ and
  $\mathrm{c}$ are as above for any (complex linear) unitary $u : W
  \rightarrow W$, then $\varkappa \assign U\mathrm{c}$ is a real structure on
  $V$ compatible with $(\cdummy, \cdummy)_V$.
\end{proposition}

\begin{proof}
  Let $(f_j)_{j \in \mathbb{N}}$ be an orthonormal basis of $V$. If $f_j +
  \varkappa f_j \neq 0$, let $e_j \assign f_j + \varkappa f_j$, otherwise let
  $e_j \assign \mathi (f_j - \varkappa f_j)$, then $(e_j)_{j \in \mathbb{N}}$,
  is an orthogonal basis of $V$ invariant under $\varkappa$. Let $g^+_j
  \assign e_{2 j - 1} + \mathi e_{2 j}, g^-_j \assign e_{2 j - 1} - \mathi
  e_{2 j}$, $j \in \mathbb{N}$. Then $(g_j^+)_{j \in \mathbb{N}}, (g_j^-)_{j
  \in \mathbb{N}}$ generate orthogonal isomorphic subspaces $V_+ \cong V_-$.
  Let $\pi_+, \pi_-$ be the respective projections. Then, since $\varkappa :
  g^+_j \mapsto g_j^-$, we have $\varkappa \pi_+ = \pi_- \varkappa$. Hence
  $V_+, V_-$ are isotropic.
  
  Let $\mathrm{c}$ be another real structure on $V$. Let $U \assign \varkappa
  \mathrm{c}$, then $U$ is linear and unitary since $(U v, U w)_V = (\varkappa
  \mathrm{c}v, \kappa \mathrm{c}w)_V = (v, w)_V$. Moreover $U U^{\ast} = 1 =
  \varkappa^2 = U\mathrm{c}U\mathrm{c}$ gives that \ $\mathrm{c}U\mathrm{c}=
  U^{\ast}$, that is $U = U^{\mathrm{c}}$.
  
  If $\mathrm{c}V_{\pm} = V_{\pm}$ then both $\pi_+, \pi_-$ commute with
  $\mathrm{c}$. Because of this and the fact that $V_+, V_-$ are isotropic we
  have $U = \pi_+ A \pi_- + \pi_- B \pi_-$ for some (complex linear) operators
  $A, B$. Let us identify $V_+ \cong V_-$ and denote them both by $W$, then
  $A, B : W \rightarrow W$. Because $U$ is unitary $A$ and $B$ are unitary.
  Finally, from $U = U^{\mathrm{c}}$, we get $A\mathrm{c}B\mathrm{c}= 1$.
  Hence if we let $u \assign A$ we get the
  decomposition~{\eqref{GPeq:UsDecomposition}}.
  
  The last statement follows by direct verification that, for any unitary
  operator $u : W \rightarrow W$, $\varkappa = U\mathrm{c}$, $U, \mathrm{c}$
  as stated, we have $\varkappa^2 =\mathbb{I}, \varkappa \mathi = - \mathi
  \varkappa$, and $(\varkappa v, \varkappa w)_V = (w, v)_V$ so $\varkappa$ is
  another real structure.
\end{proof}

Following Gross (cf. Lemma~3.1 in~{\cite{Gross_Euclidean_1976}}) we give the
following definition.

\begin{definition}
  \label{GPdef:theta-map}Let $V = V_+ \oplus V_-$ be a complex vector space
  which decomposes into a direct sum of isotropic subspaces with respect to a
  symmetric, non-degenerate, bilinear form $\llangle \cdummy, \cdummy \rrangle
  : V \times V \rightarrow \mathbb{C}$. We define the vector space isomorphism
  $\theta : \Lambda V \rightarrow \mathbb{C}l (V, \llangle \cdummy, \cdummy
  \rrangle)$ by
  \begin{enumerateroman}
    \item $\theta (1) = 1, \quad \theta (v) = B (v)$, $v \in V$;
    
    \item $\theta (v_1 \wedge \cdots \wedge v_n \wedge w_1 \wedge \cdots
    \wedge w_m) = B (v_1) \cdots B (v_n) B (w_1) \cdots B (w_m)$,
    
    if $v_j \in V_+$, $w_{\ell} \in V$, or $w_{\ell} \in V_-$, $v_j \in V$,
    $j = 1, \ldots, .n, \ell = 1, \ldots, m$.
  \end{enumerateroman}
\end{definition}

\begin{remark}
  \label{GPrem:theta-map}The map $\theta$ order the factors so that the
  generators belonging to $V_+$ are all on the left, in essence is a Wick
  ordering map with respect to the decomposition $V = V_+ \oplus V_-$. Note
  that $\theta : \Lambda V \rightarrow \mathbb{C}l (V, \llangle \cdummy,
  \cdummy \rrangle)$ cannot be extended to an algebra isomorphism since the
  exterior algebra $\Lambda V$ and the complex Clifford algebra $\mathbb{C}l
  (V, \llangle \cdummy, \cdummy \rrangle)$ are not isomorphic. By the isotropy
  of the symmetric bilinear form $\llangle \cdummy, \cdummy \rrangle$, the map
  $\theta$ is an algebra homomorphism when restricted to $V_+$ or $V_-$.
  Finally note that $\theta$ has the following nice interpretation in the
  language of Hopf algebras (cf. Section~\ref{sec:top}). Since the
  restrictions $\theta_+$, respectively $\theta_-$, of $\theta$ restricted to
  $\Lambda V_+$, respectively $\Lambda V_-$, are algebra isomorphisms, we let
  $\vartheta : \Lambda V_+ \otimes \Lambda V_- \rightarrow \mathbb{C}l (V,
  \llangle \cdummy, \cdummy \rrangle) \otimes \mathbb{C}l (V, \llangle
  \cdummy, \cdummy \rrangle)$ be the algebra isomorphism (with its image)
  defined by $\vartheta \assign \theta_+ \otimes \theta_-$. Moreover let
  $m_{\mathbb{C}l} : \mathbb{C}l (V, \llangle \cdummy, \cdummy \rrangle)
  \otimes \mathbb{C}l (V, \llangle \cdummy, \cdummy \rrangle) \rightarrow
  \mathbb{C}l (V, \llangle \cdummy, \cdummy \rrangle)$ denote the
  multiplication in $\mathbb{C}l (V, \llangle \cdummy, \cdummy \rrangle)$.
  Then we have
  \begin{equation}
    \theta (F) = m_{\mathbb{C}l} [\vartheta (\Pi_{\Lambda V_+} \otimes
    \Pi_{\Lambda V_-}) \Delta F], \qquad F \in \Lambda V,
    \label{GPeq:theta-map}
  \end{equation}
  where $\Delta$ denotes the co-product of $\Lambda V$. To see
  {\eqref{GPeq:theta-map}} we first note that $(\Pi_{\Lambda V_+} \otimes
  \Pi_{\Lambda V_-}) \Delta$ it is nothing but the isomorphism $\Lambda (V_+
  \oplus V_-) \cong \Lambda V_+ \otimes \Lambda V_-$. Now, this isomorphism
  sends $F \rightarrow F_+ \otimes F_-$, $F_+ \in \Lambda V_+, F_- \in \Lambda
  V_-$. Finally by the definition of $\theta$, we see that $\theta$ applied to
  an element $F = F_+ F_- \in \Lambda V$, gives $\theta (F_+ F_-) = \theta
  (F_+) \theta (F_-) = m_{\mathbb{C}l} \vartheta (F_+ \otimes F_-)$. Hence
  {\eqref{GPeq:theta-map}} is verified.
\end{remark}

\begin{remark}
  Let $\varkappa$ be a real structure on $V$ and let $V = V_+ \oplus V_-$ be
  an orthogonal, $\llangle \cdummy, \cdummy \rrangle_{\varkappa}$-isotropic
  decomposition. With respect to this decomposition, let $\mathrm{c} \assign
  \sigma \varkappa$ with
  \[ \sigma = \left( \begin{array}{ll}
       0 & 1\\
       1 & 0
     \end{array} \right) . \]
  Then $\mathrm{c}$ is a real structure on $V$ and it preserves the
  decomposition $V_+ \oplus V_-$ . If $\kappa$ is another real structure on
  $V$ which with respect to the decomposition $V_+ \oplus V_-$ has the form
  \[ \kappa = \left(\begin{array}{cc}
       0 & u\\
       u^{\varkappa} & 0
     \end{array}\right) \mathrm{c} \]
  for some unitary $u : V_+ \rightarrow V_+$, then
  \[ \kappa = \left(\begin{array}{cc}
       u & 0\\
       0 & u^{\varkappa}
     \end{array}\right) \varkappa . \]
  If $S : V \rightarrow V$ is unitary, $\varkappa$-antisymmetric, and
  preserves the decomposition $V_+ \oplus V_-$ then $S$ has the form
  \[ S = \left(\begin{array}{cc}
       v & 0\\
       0 & - v^{\varkappa}
     \end{array}\right) \]
  for some unitary $v : V_+ \rightarrow V_+$.
\end{remark}

\begin{theorem}
  \label{GPth:relation}Let $X \in \mathcal{G} (V)$ be a complex Gaussian
  Grassmann variable with covariance associated to a triple $(V, \varkappa,
  S)$, where $\kappa$ is a real structure on $V$ and $S = - S^{\varkappa}$ is
  the covariance of $X$. Assume $S$ is unitary and assume there exists an
  orthogonal, $\llangle \cdummy, \cdummy \rrangle_{\varkappa}$-isotropic
  decomposition $V = V_+ \oplus V_-$ which is preserved by $S : V_{\pm}
  \rightarrow V_{\pm}$. Let $\pi_+, \pi_-$ denote the projections onto
  respectively $V_+$ and $V_-$. Then
  \[ \kappa \assign S (\pi_+ - \pi_-) \varkappa \]
  is a real structure on $V$ with respect to which, if $v_{\pm}, w_{\pm} \in
  V_{\pm}$ and $v, w \in V$ we have,
  \begin{equation}
    \llangle v_{\pm}, w \rrangle_{\kappa} = (\kappa v_{\pm}, w)_V = \pm
    (\varkappa S v_{\pm}, w)_V = \pm \llangle S v_{\pm}, w
    \rrangle_{\varkappa} \label{eq:cov-pm}
  \end{equation}
  \begin{equation}
    \llangle v, w_- \rrangle_{\kappa} = \llangle w_-, v \rrangle_{\kappa} = -
    \llangle S w_-, v \rrangle_{\varkappa} = \llangle w_-, S v
    \rrangle_{\varkappa} = \llangle S v, w_- \rrangle_{\varkappa} .
    \label{eq:cov-pm-2}
  \end{equation}
  As a consequence the map $\theta : \Lambda V \rightarrow
  \mathbb{C}l_{\tmop{SD}} (W, \kappa)$ in Definition~\ref{GPdef:theta-map}
  satisfies
  \begin{equation}
    \omega (X (F)) = \tau (\theta (F)), \qquad F \in \Lambda V.
    \label{eq:relation-OS-GP}
  \end{equation}
\end{theorem}

\begin{proof}
  By Proposition~\ref{GPpro:orthogonal-dec} we see that $\kappa$ is a real
  structure on $V$ and that equations~{\eqref{eq:cov-pm}}
  and~{\eqref{eq:cov-pm-2}} are satisfied. Hence, from the definition of
  $\theta$, $\tau,$ and $\omega$, we get
  \[ \tau (\theta (B (v) B (w))) = \omega (X (v) X (w)), \qquad v, w \in V, \]
  which proves~{\eqref{eq:relation-OS-GP}} for $n = 2$. To prove the general
  case, let us first note that without loss of generality we can take
  \[ S = - \pi_+ + \pi_- . \]
  Moreover, from~{\eqref{GPeq:theta-map}} in Remark~\ref{GPrem:theta-map}, it
  is enough to consider elements of the form $F = F_+ F_- \in \Lambda V$, $F_+
  \in \Lambda V_+, F_- \in \Lambda V_-$. Finally, by linearity we can restrict
  our attention to $F = F_+ F_- = v_1^+ \wedge \cdots \wedge v_k^+ \wedge v_{k
  + 1}^- \wedge \cdots \wedge v_n^-$. Now, by the definition of $\theta$, we
  have
  \[ \theta (v_1^+ \wedge \cdots \wedge v_k^+ \wedge v_{k + 1}^- \wedge \cdots
     \wedge v_n^-) = B (v_1^+) \cdots B (v_k^+) B (v^-_{k + 1}) \cdots B
     (v_n^-) . \]
  Hence, by comparing~{\eqref{eq:trace-GP}}
  with~{\eqref{eq:complex-gauss-GP}}, we see that to
  prove~{\eqref{eq:relation-OS-GP}} it will suffice to show that
  \[ \underset{1 \leqslant \ell, m \leqslant n}{\tmop{Pf}} [\tmop{sgn} (\ell
     - m) \llangle v^{\#}_{\ell}, v^{\#}_m \rrangle_{\varkappa}] = \underset{1
     \leqslant \ell, m \leqslant n}{\tmop{Pf}} \llangle v^{\#}_{\ell}, S
     v^{\#}_m \rrangle_{\varkappa}, \]
  where $v_j^{\#} = v_j^+$ for $j = 1, \ldots, k$ and $v^{\#}_j = v_j^-$ for
  $j > k$. Note first that if $n \neq 2 k$ then both terms are zero because
  the matrices inside the Pfaffians would be degenerate. Finally observe that,
  since $S = - \pi_+ + \pi_-$ and since $V_+, V_-$ are isotropic, we have
  \[ 0 = \llangle v^+_{\ell}, v^+_m \rrangle = \llangle v^+_{\ell}, S v^+_m
     \rrangle, \quad \llangle v^+_{\ell}, v^-_{k + m} \rrangle = \llangle
     v^+_{\ell}, S v^-_{k + m} \rrangle, \qquad 1 \leqslant \ell, m \leqslant
     k, \]
  and
  \[ \llangle v^-_{k + \ell}, v^+_m \rrangle = - \llangle v^-_{k + \ell}, S
     v^+_m \rrangle, \quad 0 = \llangle v^-_{k + \ell}, v^-_{k + m} \rrangle =
     \llangle v^-_{k + \ell}, S v^-_{k + m} \rrangle, \qquad 1 \leqslant \ell,
     m \leqslant k. \]
  Hence the matrices inside the Pfaffians are the same. This concludes the
  proof.
\end{proof}

With notation as in Section~\ref{sec:os-fields}, let now $W = L^2
(\mathbb{R}^2) \otimes \mathbb{C}^2$, $u \assign T^{- 1 / 2} \left(
\slashnabla + m_f \right)$ and $T \assign m_f^2 - \Delta$, with $T$ as
self-adjoint operator on $L^2 (\mathbb{R}^2) \otimes \mathbb{C}^2$ with domain
$W \assign \mathcal{S} (\mathbb{R}^2) \otimes \mathbb{C}^2$. Note that $T$ is
positive essentially self-adjoint and $u$ is unitary. Take $V \assign W \oplus
W$ with $\pi_+, \pi_-$ the projections onto the two copies of $W$ and consider
the scalar product
\[ (v, w)_V \assign (\pi_+ v, T^{- 1 / 2} \pi_+ w)_{L^2 (\mathbb{R}^2) \otimes
   \mathbb{C}^2} + (\pi_- v, T^{- 1 / 2} \pi_- w)_{L^2 (\mathbb{R}^2) \otimes
   \mathbb{C}^2}, \qquad v, w \in V. \]
We denote by $\mathrm{c}$ the operation of complex conjugation of a function
in $L^2 (\mathbb{R}^2) \otimes \mathbb{C}^4$. Note that $T$ commutes with
$\mathrm{c}$, hence $\mathrm{c}$ is a real structure on $V$. If we let
$\varkappa \assign U\mathrm{c}$, with
\[ U \assign \left( \begin{array}{cc}
     0 & u^{\ast}\\
     \mathrm{c}u\mathrm{c} & 0
   \end{array} \right), \]
by Proposition~\ref{GPpro:orthogonal-dec} and the fact that both $U$ and
$\mathrm{c}$ commute with $T$, we have that $\mathrm{U}$ is another real
structure on $V$. If we take $S : V \rightarrow V$ given by
\begin{equation}
  S = \left(\begin{array}{cc}
    1 & 0\\
    0 & - 1
  \end{array}\right), \label{eq:S-diag}
\end{equation}
in the decomposition $V = W \oplus W$, then we have for all $v, w \in V$,
\[ (\varkappa v, S w)_V = \left( \left( \begin{array}{cc}
     0 & u^{\ast}\\
     \mathrm{c}u\mathrm{c} & 0
   \end{array} \right) \mathrm{c}v, T^{- 1 / 2} S w \right)_{L^2
   (\mathbb{R}^2) \otimes \mathbb{C}^2} \]
\[ = \left( \mathrm{c}v, T^{- 1 / 2} \left( \begin{array}{cc}
     0 & u\\
     -\mathrm{c}u^{\ast} \mathrm{c} & 0
   \end{array} \right) w \right)_{L^2 (\mathbb{R}^2) \otimes \mathbb{C}^2} \]
\[ = (\mathrm{c}v, K w)_{L^2 (\mathbb{R}^2) \otimes \mathbb{C}^2} = \omega
   (\tilde{\psi} (v) \tilde{\psi} (w)) \]
with $K : L^2 (\mathbb{R}^2) \otimes \mathbb{C}^2 \rightarrow L^2
(\mathbb{R}^2) \otimes \mathbb{C}^2$ as given in Section~\ref{sec:os-fields}
and $(\tilde{\psi} (v))_{v \in V}$ the complex Grassmann Gaussian introduced
there to represent the Euclidean Dirac field. As a consequence we have:

\begin{corollary}
  \label{GPcor:schwinger}On the pre-Hilbert space $V =\mathcal{S}
  (\mathbb{R}^2) \otimes \mathbb{C}^4$ there exists a real structure
  $\varkappa$ and orthogonal projectors $\pi_+, \pi_-$ such that $\pi_+ V
  \cong \pi_- V$ and, for $f_j, g_j \in V$ and $n \in \mathbb{N}$,
  \begin{equation}
    \omega (\tilde{\psi} (\pi_+ f_1) \cdots \tilde{\psi} (\pi_+ f_n)
    \tilde{\psi} (\pi_- g_1) \cdots \tilde{\psi} (\pi_- g_n)) = \tau (B (\pi_+
    f_1) \cdots B (\pi_+ f_n) B (\pi_- g_1) \cdots B (\pi_- g_n)),
    \label{eq:Schwinger-GP}
  \end{equation}
  where $(B (v))_{v \in V}$, are the generators of $\mathbb{C}l_{\tmop{SD}}
  (V, \mathrm{U})$ and $\tau$ is its trace and where $(\tilde{\psi} (v))_{v
  \in V}$ is the Gaussian Grassmann field of Section~\ref{sec:os-fields}.
\end{corollary}

\begin{proof}
  This is a consequence of Theorem \ref{GPth:relation} once we
  show~{\eqref{eq:Schwinger-GP}} for $n = 2$. By direct computation, for all
  $v, w \in V$ we have indeed
  \[ \tau (B (\pi_+ v) B (\pi_- w)) = (\varkappa \pi_+ v, \pi_- w)_V =
     (\mathrm{c} \pi_+ v, T^{- 1 / 2} u \pi_- w)_{L^2 (\mathbb{R}^2) \otimes
     \mathbb{C}^2} \]
  \[ = (\mathrm{c} \pi_+ v, K \pi_- w)_{L^2 (\mathbb{R}^2) \otimes
     \mathbb{C}^2} = \omega (\tilde{\psi} (\pi_+ v) \tilde{\psi} (\pi_- w)) .
  \]
\end{proof}

\subsection{The real-wave representation of Grassmann Gaussians}

Theorem~\ref{GPth:relation} gives a realization of the Grassmann Gaussian
variables in terms of the trace state on a Clifford algebra. This realization
is obtained through the vector-space map $\theta$, which can be interpreted as
a sort of ``Wick-ordering'' with respect to a suitable decomposition $V = V_+
\oplus V_-$. Since the map $\theta$ is not an algebra isomorphism that
realization is not a representation. Here we are going to give an actual
representation of Grassmann Gaussian variables as multiplication operators on
a non-commutative $L^2$ space in full analogy to the case of the usual
(commutative) Gaussian variables. Moreover the representation we give here
will not depend on any decomposition $V = V_+ \oplus V_-$, hence making the
construction intrinsically more natural.

\

Let $\mathcal{K}= \Gamma_a V$ be the antisymmetric Fock space over the
complex Hilbert space $V$ with a real structure $\varkappa$. Following the
exposition of Gross~{\cite{gross_existence_1972}} one considers the Clifford
algebra $\mathfrak{C}_0 \subseteq \mathcal{L} (\mathcal{K})$ generated
algebraically by the operators $(B_{\tmop{GP}} (v))_{v \in V}$. Let
$\mathfrak{C}$ be the von Neumann algebra generated by $\mathfrak{C}_0$ and
consider on $\mathfrak{C}$ the state $\tmop{tr} (u) = \omega (u) = \langle
\Omega, u \Omega \rangle_{\mathcal{K}}$ for $u \in \mathfrak{C}$. This state
is tracial (on $\mathfrak{C}$), that is $\tmop{tr} (C B) = \tmop{tr} (B C)$
for all $B, C \in \mathfrak{C}$ and faithful, i.e. $\tmop{tr} (A^{\ast} A) = 0
\Rightarrow A = 0$. So the scalar product $(A, B) \mapsto \tmop{tr} (A^{\ast}
B)$ is positive definite and the closure of $\mathfrak{C}$ with respect to the
corresponding norm defines the Hilbert space $L^2 (\mathfrak{C})$. By a
theorem of Segal ({\cite{segal_tensor_1956}}, but see also Theorem 5
in~{\cite{gross_existence_1972}}) there is an isomorphism $D : L^2
(\mathfrak{C}) \rightarrow \mathcal{K}$ such that $u \in \mathfrak{C}
\subseteq L^2 (\mathfrak{C})$ is sent to $u \Omega \in \mathcal{K}$. Let
$\varkappa$ be, as before, the real structure on $V$. By the proof of
Proposition~\ref{GPpro:orthogonal-dec} we know that there exists a
$\varkappa$-invariant orthonormal basis $(e_j)_j$ of $V$. Then the inverse
$D^{- 1} : \mathcal{K} \rightarrow \mathfrak{C}$ satisfies $D^{- 1} (e_{i_1}
\wedge \cdots \wedge e_{i_n}) = B_{\tmop{GP}} (e_{i_1}) \cdots B_{\tmop{GP}}
(e_{i_n})$. Note however that $D$ is a canonical isomorphism and in particular
(unlike the map $\theta$ of the previous section) does not depend on any
particular decomposition $V = V_+ \oplus V_-$. By Corollary 5.2
of~{\cite{gross_existence_1972}}, writing $L_a$ (resp. $R_a$) for the left
(resp. right) multiplication of $a \in \mathfrak{C}$ on $L^2 (\mathfrak{C})$,
one has
\begin{equation}
  D L_{B_v} D^{- 1} = C_v + A_{\varkappa v} \beta, \qquad D R_{B_v} D^{- 1} =
  (C_v - A_{\varkappa v}) \beta, \qquad v \in V \label{eq:maps}
\end{equation}
where $\beta = \Gamma (- 1) : \mathcal{K} \rightarrow \mathcal{K}$ is the
unitary map corresponding to the second quantization of the unitary operator
$v \mapsto - v$ on $V$. Let us denote by $\hat{\beta} = D^{- 1} \beta D$ the
pullback on $L^2 (\mathfrak{C})$ of this unitary.

\begin{remark}
  Recall that the standard centered Gaussian measure $\mu$ corresponding to an
  Hilbert space $V$ supports the It\^{o}--Wiener isomorphism $L^2 (\mu)
  \approx \Gamma_s V$, where $\Gamma_s V$ is the symmetric Fock space over
  $V$. The gaussian random variable $X$ with law $\mu$ is then realized as an
  unbounded self-adjoint operator acting in $L^2 (\mu)$ given by
  multiplication. The space $L^2 (\mathfrak{C})$ can be considered to be the
  fermionic equivalent of this construction, in that it provides the
  representation of the antisymmetric Fock space as a space of ``$L^2$
  functions'' over a non-commutative probability space $(\mathfrak{C},
  \tmop{tr})$. The von-Neumann algebra $\mathfrak{C}$ can be interpreted as
  $L^{\infty} (\mathfrak{C})$ acting on $L^2 (\mathfrak{C})$ and fermionic
  random variables can be constructed, according to~{\eqref{eq:maps}}, as
  (left and right) multiplication operators on $L^2 (\mathfrak{C})$.
\end{remark}

We want now to understand this construction in the context of our Grassmann
Gaussian variables. This will allow us to replace the ``particle'' Fock space
$\mathcal{K}$ by the space $L^2 (\mathfrak{C})$ of ``wave-functions'' for the
von Neumann Clifford algebra $\mathfrak{C}$. While this is not important in
the theory exposed in our work, it could play a role as soon as we have to
deal with unbounded operators, for which the $L^p$ spaces of unbounded
operators associated to the non-commutative probability space $(\tmop{tr},
\mathfrak{C})$ could become useful~{\cite{gross_existence_1972}}.

\

As before we consider an operator $S : V \rightarrow V$ such that
$S^{\varkappa} = - S$. The triple $(V, \varkappa, S)$ determines a complex
Grassmann Gaussian variable $X \in \mathcal{G} (V)$ with covariance $\llangle
\cdot, S \cdot \rrangle_{\varkappa}$ (cf. Definition
\ref{def:complex-gaussians}). By the Osterwalder--Schrader construction, this
Gaussian has been obtained via the concrete representation as sum of creation
and destruction operators over the antisymmetric Fock space $\Gamma_a V$ with
the state $\omega$ given by the Fock vacuum and as in the previous subsection
we have
\[ X (v) = C_{S v} + A_{\varkappa v}, \qquad v \in V. \]
By Segal's duality $D$ there will be multiplication operators on $L^2
(\mathfrak{C})$ corresponding to our Grassmann Gaussian. Let us work out their
precise form: by~{\eqref{eq:maps}} we have
\[ D (L_{B_v} + R_{B_v} \hat{\beta}) D^{- 1} = 2 C_v, \qquad D (L_{B_v} -
   R_{B_v} \hat{\beta}) D^{- 1} = 2 A_{\varkappa v}, \qquad v \in V. \]
Therefore for any $v \in V$,
\[ \hat{X} (v) \assign D^{- 1} X (v) D = \frac{1}{2} [L_{B_{S v}} + L_{B_v}] +
   \frac{1}{2} [R_{B_{S v}} - R_{B_v}] \hat{\beta} \]
\[ = L_{B_{Q_+ v}} - R_{B_{Q_- v}} \hat{\beta}, \]
where $Q_{\pm} = (1 \pm S) / 2$. By construction, $(\hat{X} (v))_{v \in V}$
are the generators of a Grassmann algebra and we can extend them to an
homomorphism $\hat{X} : \Lambda V \rightarrow L^{\infty} (\mathfrak{C})$ such
that
\[ \tmop{tr} (\hat{X} (F) 1) = \omega (X (F)), \]
for all $F \in \Lambda V$. This shows that, if needed, we can realize
Grassmann Gaussians over the $L^2$-space over the non-commutative measure
space $(\mathfrak{C}, \tmop{tr})$ in full analogy with the bosonic setting.
Note that in the particular case of the Euclidean Dirac field the operator $S$
has the diagonal form~{\eqref{eq:S-diag}} and therefore the operators $Q_{\pm}
= \pi_{\pm}$ are the projections onto the isotropic decomposition $V = V_+
\oplus V_-$ corresponding to the two fields $\psi$ and $\bar{\psi}$ which
constitute the Euclidean Dirac field $\tilde{\psi}$ (cf. Section
\ref{sec:os-fields}).

\

\section{Convergence of the perturbative series: finite dimensional
case}\label{app:series}

In this appendix we sketch some implications of the Pauli exclusion principle
for the existence of global solutions to finite-dimensional Grassmann SDEs. In
principle some of these considerations also apply to some more realistic
models like the Yukawa$_2$ model however, to our surprise, stochastic
quantization, at least in the perturbative regime, can be carried on without
establishing them, as we demonstrated in Section~\ref{sec:yukawa}.

\

The solutions to the non-linear equation~{\eqref{eq:sde-lambda}} can be
represented via a series. We investigate now the properties of this series and
its absolute convergence. The explicit control of the series is not needed for
the development of the theory of stochastic quantization as we have
shown.\quad For simplicity let us assume that $(v_{\alpha})_{\alpha = 1,
\ldots, N}$ is a finite basis of $V$ and that $F \in \tmop{Hom} (V, \Lambda
V)$ is given by a sum of cubic monomials
\[ F (v_{\alpha}) = \sum_{\alpha_1, \alpha_2, \alpha_3}
   \lambda^{\alpha}_{\alpha_1 \alpha_2 \alpha_3} \psi (v_{\alpha_1}) \psi
   (v_{\alpha_2}) \psi (v_{\alpha_2}), \]
where $\psi \in \tmop{Hom} (V, \Lambda V)$ is the canonical injection of $V$
into $\Lambda V$ and $(\lambda^{\alpha}_{\alpha_1 \alpha_2 \alpha_3})$ a
family of coefficients in $\mathbb{R}$. Moreover we take $A = -\mathbb{I}$.
Equation~{\eqref{eq:sde-lambda}} has then the integral formulation given by
\begin{equation}
  \Psi_t (v) = \Phi_t (v) + \int_0^t \Psi_s (e^{- (t - s)} F (v)) \mathd s,
  \qquad t \geqslant 0, v \in V, \label{eq:mild}
\end{equation}
where
\[ \Phi_t (v) \assign \Psi_0 (e^{- t} v) + B^A_t (v) - B^A_0 (e^{- t} v),
   \qquad t \geqslant 0, v \in V. \]
By iteratively expanding $\Psi_s$ on the right hand side of~{\eqref{eq:mild}}
we obtain a series expansion for $\Psi_t$ of the form
\begin{equation}
  \Psi_t = \sum_{\tau} J_{\tau} (\Phi) (t) \label{eq:sums}
\end{equation}
\[ = J_{\bullet} (\Phi) (t) + J_{[\bullet \bullet \bullet]} (\Phi) (t) +
   J_{[[\bullet \bullet \bullet] \bullet \bullet]} (\Phi) (t) + \cdots +
   J_{[[\bullet \bullet \bullet] [\bullet [\bullet \bullet [\bullet \bullet
   \bullet]] \bullet] \bullet]} (\Phi) (t) + \cdots \]
The series is indexed by (planar) trees $\tau$ which have branches of order 3
and where $J$ is a multilinear integral operator such that
\[ J_{\bullet} (\Phi) (t)^{\alpha} = \Phi_t (v_{\alpha}) \]
\[ J_{[\tau_1 \tau_2 \tau_3]} (\Phi) (t)^{\alpha} = \sum_{\alpha_1, \alpha_2,
   \alpha_3} \int_0^t e^{- (t - s)} \lambda^{\alpha}_{\alpha_1, \alpha_2,
   \alpha_3} J_{\tau_1} (\Phi) (s)^{\alpha_1} J_{\tau_2} (\Phi) (s)^{\alpha_2}
   J_{\tau_3} (\Phi) (s)^{\alpha_3} \mathd s \]
where $\bullet$ denotes the simple tree and $[\tau_1, \ldots, \tau_3]$ the
tree with branches $\tau_1, \ldots, \tau_3$. Our goal is to prove that the
above series converges for all times and derive estimates for the norm $\|
\Psi_t \|$ of the solution $\Psi$. Expanding the expression for $J_{\tau}
(\Phi)$ we have
\[ J_{\tau} (\Phi) = \int \left( \prod_{p \in I (\tau)}
   \lambda^{\alpha_p}_{\alpha_{p_1} \alpha_{p_2} \alpha_{p_3}} \right) \left(
   \prod_{p \in L (\tau)} \Phi^{\alpha_p} (s_{p_-}) \right) \left( \prod_{p
   \in I (\tau)} \mathd s_p \right) \]
where $I (\tau)$ denotes the internal nodes, $L (\tau)$ the leaves, $p$ the
parent node and $p_1, p_2, p_3$ the three children of every internal node and
$\Phi^{\alpha} (t) = \Phi_t (v_{\alpha})$. We have $| \tau | = | L (\tau) | +
| I (\tau) |$, $| \tau | = 1 + 3 | I (\tau) |$, $| L (\tau) | = 1 + 3 | I
(\tau) | - | I (\tau) | = 1 + 2 | I (\tau) |$.

\

Note that $(\Phi^{\alpha}_t)^2 = 0$ and that the increments $\Phi^{\alpha}
(t, s) \assign \Phi^{\alpha}_t - \Phi^{\alpha}_s$ have the norm bound
\[ \| \Phi^{\alpha} (t, s) \|^2 \leqslant \int_s^t e^{- (t - r)} \mathd r + |
   e^{- (t - s)} - 1 |^2 \leqslant 2 (1 - e^{- (t - s)}) \backassign H (t - s)
\]
which can be estimated as
\[ \| \Phi^{\alpha} (t, s) \| \leqslant | t - s |^{1 / 2} . \]
The following Lemma~uses $(\Phi^{\alpha}_t)^2 = 0$ (Pauli exclusion principle)
to derive good estimates for products of fields.

\begin{lemma}
  Assume $m \geqslant 1$. For any $n \geqslant 1$ and any $t_1, \ldots, t_n
  \in [0, T]$, $\alpha_1, \ldots, \alpha_n \in \{ 1, \ldots, N \}$ we have
  \[ \| \Phi^{\alpha_1}_{t_1} \cdots \Phi^{\alpha_n}_{t_n} \| \leqslant
     \frac{C^{n + 1} T^{n / 2}}{(n!)^{1 / 2}} \]
  where $C$ is a universal constant depending only on $N$.
\end{lemma}

\begin{proof}
  Partition $[0, T]$ in $r = n / (2 N)$ equal intervals $(I_k)_k$ and let
  $(s_k)_k$ be the centers of those intervals. Now in the product replace each
  $\Phi^{\alpha_i}_{t_1}$ by $\Phi^{\alpha_i}_{t_1} = \Phi^{\alpha_i}_{s_k} +
  \Phi^{\alpha_i} (t_i, s_k)$ where $s_k$ is the nearest to $t_i$ of the
  centers. By doing so we rewrite $\Phi^{\alpha_1}_{t_1} \cdots
  \Phi^{\alpha_n}_{t_n}$ as a sum $S$ of $2^n$ products of fields which are
  either $\Phi^{\alpha_i}_{s_k}$ or \ $\Phi^{\alpha_i} (t_i, s_k)$, moreover
  for a given $I_k$ there cannot be more than $N$ factors
  $\Phi^{\alpha_i}_{s_k}$ by the exclusion principle. Denote by $n_k$ the
  number of $(t_i)_i$ in $I_k$. Then $n = \sum_k n_k$ and using the estimate
  $\| \Phi^{\alpha_i} (t_i, s_k) \| \leqslant | t_i - s_k | \leqslant | I_k |
  = (2 N T) / n$ and the fact that in each cell we have at most $n_k - N$
  increments, we have
  \[ \| S \| \leqslant 2^n \prod_{k = 1}^r | I_k |^{n_k - N} \leqslant 2^n
     \left( \frac{2 N T}{n} \right)^{n - N r} = 2^n \left( \frac{2 N T}{n}
     \right)^{n / 2} \leqslant \frac{C^{n + 1} T^{n / 2}}{(n!)^{1 / 2}} . \]
\end{proof}

Using this lemma above we can estimate
\[ \| J_{\tau} (\Phi) (t) \| \leqslant \int \left( \prod_{p \in I (\tau)}
   \left| \lambda^{\alpha_p}_{\alpha_{p_1} \alpha_{p_2} \alpha_{p_3}} \right|
   \right) \left\| \prod_{p \in L (\tau)} \Phi^{\alpha_p}_{s_{p_-}} \right\|
   \left( \prod_{p \in I (\tau)} \mathd s_p \right) \]
\[ \lesssim | \lambda |^{| I (\tau) |} N^{| L (\tau) |} \frac{C^{| L (\tau) |
   + 1} t^{| L (\tau) | / 2}}{(| L (\tau) | !)^{1 / 2}} \]
where we estimated all the integrals by constants uniformly in $t$ due to the
presence of the exponential factors. The number of trees $\tau$ with a given
number of nodes $| \tau | = n$ is no more than $D^n$ for some $n$ and $D > 0$
and $| L (\tau) | = 1 + 2 | I (\tau) | = 1 + (2 / 3) (| \tau | - 1) = 1 / 3 +
2 | \tau | / 3$ so we obtain at the end for the solution $\Psi_t$
of~{\eqref{eq:mild}}
\begin{equation}
  \| \Psi_t \| \leqslant \sum_{\tau} \| J_{\tau} (\Phi) (t) \| \lesssim
  \sum_{n \geqslant 0} D^n | \lambda |^{(2 / 3) (n - 1)} N^{2 n / 3}
  \frac{C^{2 n / 3} t^{n / 3 + 1 / 6}}{(n / 3 + 1 / 6) !}
  \label{eq:estimateexpansion}
\end{equation}
which is a series which converges for all $t$ (with some stretched exponential
behavior).

\begin{theorem}
  There exists an increasing function $E (t)$ depending on $N, | \lambda |, m$
  such that for all $m > 0$,
  \[ \sup_{s \leqslant t} \| \Psi_s \| \leqslant E (t), \]
  where recall that $(\Psi_t)_{t \geqslant 0}$ is the unique solution
  of~{\eqref{eq:mild}}.
\end{theorem}

In particular this shows again that explosion is not possible for this kind of
SDEs, unlike their bosonic analogs.

\

Our goal now is to remove the time dependence on our estimate for the
solution exploiting the exponential decay. Let us partition the interval $[0,
\infty]$ into intervals $(I_h)_{h \geqslant 0}$ of size $1$. In each of these
intervals we will use the finer partition of the previous lemma to estimate
products, which will then not depend on $T$ anymore. Of course now the problem
is that we have a bound of the form
\[ \| \Phi^{\alpha_1} (t_1) \cdots \Phi^{\alpha_n} (t_n) \| \leqslant \prod_{h
   \geqslant 0} \frac{C^{n_h + 1}}{(n_h !)^{1 / 2}} \]
where $n_h$ is the number of time variables in $I_h$ and $\sum_{h \geqslant 0}
n_h = n$, and so far there is no useful upper bound on this (since we observe
that we can have the situation where $n_h = 1$ for all $h$ and so we loose the
factorial contribution). Let $Q$ be the number of intervals with $n_h > 0$. In
the expression for $J_{\tau} (\Phi) (t)$ there are at least $Q$ factors of the
form $e^{- (s - s')}$, $s' < s$, for which the sum of the quantities $s - s'$
is at least $Q$. So we have the bounds
\[ \| J_{\tau} (\Phi) (t) \| \leqslant | \lambda |^{(2 / 3) (n - 1)} N^n
   \sum_{Q \geqslant 1} e^{- Q / 2} \prod_{h : n_h > 0} \frac{C^{n_h +
   1}}{(n_h !)^{1 / 2}} \leqslant \frac{C^n | \lambda |^{(2 / 3) (n -
   1)}}{(n!)^{1 / 2}} \sum_{Q \geqslant 1} e^{- Q / 2} Q^{n / 2} \]
where we used that
\[ \prod_{h : n_h > 0} \frac{1}{n_h !} \leqslant \frac{1}{n!}
   \sum_{\tmscript{\begin{array}{c}
     k_1, \ldots, k_Q \geqslant 1\\
     k_1 + \cdots + k_Q = n
   \end{array}}} \frac{n!}{k_1 ! \cdots k_Q !} = \frac{1}{n!} Q^n . \]
Now, using the bound $\sum_{Q \geqslant 1} e^{- Q / 2} Q^{n / 2} \lesssim c^n
(n!)^{1 / 2}$ we obtain that
\begin{equation}
  \| J_{\tau} (\phi) (t) \| \lesssim | \lambda |^{(2 / 3) (n - 1)} C^n
  \label{eq:estimateexpansion2}
\end{equation}
so provided $| \lambda |$ is small enough (depending on our choice of $N$) we
conclude, from the representation~{\eqref{eq:sums}}, that the uniform bound
$\sup_{t \geqslant 0} \| \Psi_t \| \lesssim 1$ holds.

\section{Besov spaces of Banach algebras}\label{appendix_Besov}

In this section we want to recall some results about Besov spaces of functions
(or distributions) from $\mathbb{R}^d$ taking values in a Banach algebra
$\mathcal{A}$. All the results of this section can be found
in~{\cite{Amann1997,Amann2019}} for the theory of Besov spaces taking values
in a Banach space and~{\cite{TriebelIII}} for weighted Besov spaces. We
present here only the $\mathbb{R}^d$ case. The definition of Besov spaces on
$\mathbb{T}^d$ is similar, and all the results announced here also hold on
$\mathbb{T}^d$.

\

We denote by $\mathcal{S} (\mathbb{R}^d)$ the set of smooth functions $f \in
\mathcal{S} (\mathbb{R}^d)$ such that
\[ \| f \|_{\ell, \alpha} \assign \| (1 + | x |)^{\ell} | D^{\alpha} f |
   \|_{L^{\infty} (\mathbb{R}^d)} < + \infty \]
where $\ell \in \mathbb{R}_+$ and $\alpha \in \mathbb{N}^d$. We denote by
$\mathcal{S}' (\mathbb{R}^d)$ the strong dual of $\mathcal{S} (\mathbb{R}^d)$
(equipped with the topology of the semi-norms $\| \cdot \|_{\ell, \alpha}$)
and by $\mathcal{O}_M (\mathbb{R}^d)$ the set of smooth functions $f \in
\mathcal{O}_M (\mathbb{R}^d)$ such that for any $\alpha \in \mathbb{N}^d$
there exists $\ell_{\alpha} \in \mathbb{R}_+$ for which
\[ \| (1 + | x |)^{- \ell_{\alpha}} | D^{\alpha} f | \|_{L^{\infty}
   (\mathbb{R}^d)} < + \infty . \]
We use the notations
\[ \mathcal{S} (\mathbb{R}^d, \mathcal{A}) = \mathcal{S} (\mathbb{R}^d)
   \hat{\otimes} \mathcal{A}, \quad \mathcal{S}' (\mathbb{R}^d, \mathcal{A}) =
   \mathcal{S}' (\mathbb{R}^d) \hat{\otimes} \mathcal{A}, \quad \mathcal{O}_M
   (\mathbb{R}^d, \mathcal{A}) = \mathcal{O}_M (\mathbb{R}^d) \hat{\otimes}
   \mathcal{A}, \]
where $\mathcal{A}$ is a Banach algebra and $A \hat{\otimes} B$ is the
completion with respect to the natural metric of the algebraic tensor product
$A \otimes B$ when $A$ is a Fr{\'e}chet nuclear space and $B$ is a Banach
space (see Remark ). Note that these tensor products are well defined since
the spaces $\mathcal{S} (\mathbb{R}^d), S' (\mathbb{R}^d), \mathcal{O}
(\mathbb{R}^d)$ are nuclear. It is important to note that
\[ \mathcal{S} (\mathbb{R}^d, \mathcal{A}) \overset{d}{\hookrightarrow}
   \mathcal{O}_M (\mathbb{R}^d, \mathcal{A}) \overset{d}{\hookrightarrow}
   \mathcal{S}' (\mathbb{R}^d, \mathcal{A}) \]
where the arrows mean that each space is continuously embedded and dense in
the following one.

\begin{theorem}
  It is possible to define uniquely $\langle \cdot, \cdot \rangle :
  \mathcal{S} (\mathbb{R}^d, \mathcal{A}) \times \mathcal{S}' (\mathbb{R}^d,
  \mathcal{A}) \rightarrow \mathcal{A}$, $\cdot : \mathcal{S} (\mathbb{R}^d,
  \mathcal{A}) \times \mathcal{S}' (\mathbb{R}^d, \mathcal{A}) \rightarrow
  \mathcal{S}' (\mathbb{R}^d, \mathcal{A})$, $\cdot : \mathcal{S}
  (\mathbb{R}^d) \times \mathcal{S}' (\mathbb{R}^d, \mathcal{A}) \rightarrow
  \mathcal{S}' (\mathbb{R}^d, \mathcal{A})$, $\asterisk : \mathcal{S}
  (\mathbb{R}^d) \times \mathcal{S}' (\mathbb{R}^d, \mathcal{A}) \rightarrow
  \mathcal{O}_M (\mathbb{R}^d, \mathcal{A})$ and $D^{\alpha} : \mathcal{S}'
  (\mathbb{R}^d, \mathcal{A}) \rightarrow \mathcal{S}' (\mathbb{R}^d,
  \mathcal{A})$ (where $\alpha \in \mathbb{N}^d$) which extend in an
  epicontinuous way the following operations: any $f \in \mathcal{S}
  (\mathbb{R}^d)$, $u \in \mathcal{S}' (\mathbb{R}^d)$ and $a_1, a_2 \in
  \mathcal{A}$ we have
  \[ \begin{array}{rll}
       \langle f \otimes a_1, u \otimes a_2 \rangle & = & \langle f, u \rangle
       a_1 a_2\\
       (f \otimes a_1) \cdot (u \otimes a_2) & = & (f u) \otimes (a_1 a_2)\\
       f \cdot (u \otimes a_2) & = & (f u) \otimes a_1\\
       f \asterisk (u \otimes a_2) & = & (f \asterisk u) \otimes a_1\\
       D^{\alpha} (u \otimes a_1) & = & (D^{\alpha} u) \otimes a_1
     \end{array} \]
  where $\langle f, u \rangle$ is the normal pairing in $\mathcal{S}
  (\mathbb{R}^d) \times \mathcal{S}' (\mathbb{R}^d)$, $(f u)$ is the product
  in $\mathcal{S} (\mathbb{R}^d) \times \mathcal{S}' (\mathbb{R}^d)$, $(f
  \asterisk u)$ is the convolution in $\mathcal{S} (\mathbb{R}^d) \times
  \mathcal{S}' (\mathbb{R}^d)$ and $D^{\alpha}$ is the $\alpha$ derivative in
  $\mathcal{S}' (\mathbb{R}^d)$.
\end{theorem}

\begin{proof}
  The proof of this theorem can be found in~{\cite{Amann2019}} Appendix~1.
\end{proof}

We recall the definition of Littlewood--Paley blocks: let $\chi, \varphi$ be
smooth non-negative functions from $\mathbb{R}^n$ into $\mathbb{R}$ such that
\begin{itemize}
  \item $\tmop{supp} (\chi) \subset B_{\frac{4}{3}} (0)$ and $\tmop{supp}
  (\varphi) \subset B_{\frac{8}{3}} (0) \setminus B_{\frac{3}{4}} (0)$,
  
  \item $\chi, \varphi \leqslant 1$ and $\chi (y) + \sum_{j \geq 0} \varphi
  (2^{- j} y) = 1$ for any $y \in \mathbb{R}^n$,
  
  \item $\tmop{supp} (\chi) \cap \tmop{supp} (\varphi (2^{- i} \cdot)) =
  \emptyset$ for $i \geqslant 1$,
  
  \item $\tmop{supp} (\varphi (2^{- j} \cdot)) \cap \tmop{supp} (\varphi (2^{-
  i} \cdot)) = \emptyset$ if $| i - j | > 1$,
\end{itemize}
where by $B_r (x)$ we denote the ball of center $x \in \mathbb{R}^n$ and
radius $r > 0$. We introduce the following notations: $\varphi_{- 1} = \chi$,
$\varphi_j (\cdot) = \varphi (2^{- j} \cdot)$, $K_j = \hat{\varphi}_j$ and,
for any $\ell > 0$, we define $\rho_{\ell} (x) \assign (1 + | x |^2)^{- \ell /
2}$.

\begin{remark}
  \label{remark_decay}It is always possible to choose $\chi$ and $\varphi$ in
  such a way that there exists some constants $\gamma_{- 1}, \gamma > 0$ and
  $0 < \theta_{- 1}, \theta < 1$ such that
  \[ | K_{- 1} (y) | \lesssim \exp (- \gamma_{- 1} | y |^{\theta_{- 1}})
     \quad \tmop{and} \quad | K_0 (y) | \lesssim \exp (- \gamma | y
     |^{\theta}), \]
  (see, i.e., {\cite{SchmeisserTriebel}} Section 1.2.2 Proposition 1 or
  {\cite{mourrat_plane_2015}}).
\end{remark}

If $v \in \mathcal{S}' (\mathbb{R}^d, \mathcal{A})$ and if $i \in \mathbb{Z}$,
$i \geqslant - 1$ we define $i$-th Littlewood--Paley block as follows
\[ \Delta_i v = K_i \asterisk v \in \mathcal{O}_M (\mathbb{R}^d, \mathcal{A})
   . \]
Then, if $s \in \mathbb{R}$, $p, q \in [1, + \infty]$ and $\ell \in
\mathbb{R}$ we define the function
\[ \| v \|_{B^s_{p, q, \ell} (\mathbb{R}^d, \mathcal{A})} = \left( \sum_{j = -
   1}^{+ \infty} 2^{j s q} \| \Delta_j v \|_{L^p_{\ell} (\mathbb{R}^d,
   \mathcal{A})}^q \right)^{1 / q}, \]
when $q \in [1, + \infty)$ and $\| v \|_{B^s_{p, + \infty, \ell}
(\mathbb{R}^d, \mathcal{A})} = \sup_j (2^{j s} \| \Delta_j v \|_{L^p_{\ell}
(\mathbb{R}^d, \mathcal{A})})$, where $\| \cdot \|_{L^p_{\ell} (\mathbb{R}^d,
\mathcal{A})}$ is the norm in the space $L_{\ell}^p (\mathbb{R}^d,
\mathcal{A})$ that is
\[ \| f \|_{L^p_{\ell} (\mathbb{R}^d, \mathcal{A})} = \left(
   \int_{\mathbb{R}^n} (\| f (y) \|_{\mathcal{A}} \rho_{\ell} (y))^p \mathd y
   \right)^{1 / p}, \]
if $p \in [1, + \infty)$ and
\[ \| f \|_{L^{\infty}_{\ell} (\mathbb{R}^d, \mathcal{A})} = \sup_{y \in
   \mathbb{R}^d} (\| f (y) \|_{\mathcal{A}} \rho_{\ell} (y)), \]
if $p = + \infty$. For any $v \in \mathcal{S} (\mathbb{R}^d, \mathcal{A})$ the
norm $\| v \|_{B^s_{p, q, \ell} (\mathbb{R}^d, \mathcal{A})} < + \infty$ is
finite. Then we look at $B^s_{p, q, \ell} (\mathbb{R}^d, \mathcal{A})$ as the
closure of $\mathcal{S} (\mathbb{R}^d, \mathcal{A})$ in $\mathcal{S}'
(\mathbb{R}^d, \mathcal{A})$ with respect to the norm $\| \cdot \|_{B^s_{p, q,
\ell} (\mathbb{R}^d, \mathcal{A})}$. Hereafter, if $s \in \mathbb{R}$, $p, q
\in [1, + \infty]$ and $\ell \in \mathbb{R}$, we use the following notations
$\mathcal{C}^s_{\ell} (\mathbb{R}^d, \mathcal{A}) \assign B^s_{\infty, \infty,
\ell} (\mathbb{R}^d, \mathcal{A})$, $B^s_{p, q} (\mathbb{R}^d, \mathcal{A}) =
B^s_{p, q, 0} (\mathbb{R}^d, \mathcal{A})$, $\mathcal{C}^s (\mathbb{R}^d,
\mathcal{A}) \assign B^s_{\infty, \infty, 0} (\mathbb{R}^d, \mathcal{A})$,
$B^s_{p, q, \ell} \assign B^s_{p, q, \ell} (\mathbb{R}^d, \mathbb{R})$ etc.

\

In this paper we need only the next three results.

\begin{theorem}
  \label{theorem_weight}For any $s, \ell \in \mathbb{R}$ and $v \in
  \mathcal{C}^s_{\ell} (\mathbb{R}^d, \mathcal{A})$ we have that $\rho_{\ell}
  v \in \mathcal{C}^s (\mathbb{R}^d, \mathcal{A})$ and
  \[ \| v \|_{\mathcal{C}^s_{\ell} (\mathbb{R}^d, \mathcal{A})} \sim \|
     \rho_{\ell} v \|_{\mathcal{C}^s (\mathbb{R}^d, \mathcal{A})} . \]
\end{theorem}

\begin{proof}
  The estimate holds by Theorem~6.5 in~{\cite{TriebelIII}} for weighted Besov
  spaces on $\mathbb{R}^d$ with polynomial-like weights. The theorem for a
  general Banach algebra $\mathcal{A}$ is a straightforward generalization.
\end{proof}

\begin{theorem}
  \label{theorem_laplacian}Consider $m > 0$, $\alpha, s, \ell \in \mathbb{R}$
  such that $s, s + \alpha \centernot{\in} \mathbb{N}$ then we have that $(- \Delta
  + m)^{- \alpha}$, where $\Delta$ is the standard Laplacian on
  $\mathbb{R}^d$, is a continuous linear map from $\mathcal{C}^s_{\ell}
  (\mathbb{R}^d, \mathcal{A})$ into $\mathcal{C}^{s + \alpha}_{\ell}
  (\mathbb{R}^d, \mathcal{A})$
\end{theorem}

\begin{proof}
  This is exactly Theorem 5.3.2 of~{\cite{Amann2019}} for the unweighted case
  $\ell = 0$. The theorem can be easily extended using the techniques of
  Chapter 6 of~{\cite{TriebelIII}} to Besov spaces with weights $\rho_{\ell}$.
\end{proof}

\begin{theorem}
  \label{theorem_convolution}Let $\alpha > 0$ and consider $\mathfrak{a} \in
  B^{\alpha}_{1, \infty}$, then we have that the convolution $\mathfrak{a}
  \asterisk$ can be extended in a unique continuous way from an operator from
  $\mathcal{S} (\mathbb{R}^d, \mathcal{A})$ into itself, into an operator from
  $\mathcal{C}^s (\mathbb{R}^d, \mathcal{A})$ into $\mathcal{C}^{s + \alpha}
  (\mathbb{R}^d, \mathcal{A})$.
\end{theorem}

\begin{proof}
  Using the notations of~{\cite{Amann1997}}, if $\mathfrak{a \in}
  B^{\alpha}_{1, \infty}$ then $((- \Delta + 1)^{\alpha / 2} (\mathfrak{a}))
  \in \mathcal{F} L_1$, then the thesis follows from Corollary~6.4
  in~{\cite{Amann1997}}.
\end{proof}

\begin{theorem}
  \label{theorem_multiplication}Let $s > 0$ and $\ell_1, \ell_2 \in
  \mathbb{R}$ we have that the product $\cdot : \mathcal{S} (\mathbb{R}^d,
  \mathcal{A}) \times \mathcal{S} (\mathbb{R}^d, \mathcal{A}) \rightarrow
  \mathcal{S} (\mathbb{R}^d, \mathcal{A})$ can be (uniquely) extended in a
  continuous way from $\mathcal{C}^s_{\ell_1} (\mathbb{R}^d, \mathcal{A})
  \times \mathcal{C}^s (\mathbb{R}^d, \mathcal{A})$ into
  $\mathcal{C}^s_{\ell_1 + \ell_2} (\mathbb{R}^d, \mathcal{A})$, furthermore
  for any $v_1 \in \mathcal{C}^s_{\ell_1} (\mathbb{R}^d, \mathcal{A})$ and
  $v_2 \in \mathcal{C}^s_{\ell_2} (\mathbb{R}^d, \mathcal{A})$ we have
  \[ \| v_1 \cdot v_2 \|_{\mathcal{C}^s_{\ell_1 + \ell_2} (\mathbb{R}^d,
     \mathcal{A})} \lesssim \| v_1 \|_{\mathcal{C}^s_{\ell_1} (\mathbb{R}^d,
     \mathcal{A})}  \| v_2 \|_{\mathcal{C}^s_{\ell_2} (\mathbb{R}^d,
     \mathcal{A})}, \]
\end{theorem}

\begin{proof}
  The proof can be found in~{\cite{Amann2019}} in the unweighted case. The
  proof in the Besov spaces with polynomial weights $\rho_{\ell}$ is similar
  and can be done using the techniques of Chapter~6 of~{\cite{TriebelIII}}. 
\end{proof}

\section{Convergence of the perturbative series: infinite dimensional
case}\label{appendix:convergenceinfinitedim}

In this section we want to consider the stationary solution to the equation


\begin{equation}
  \tilde{\Psi}^s_{\varepsilon, \lambda^2, t} (x) = \lambda^2 \int_{- \infty}^t
  e^{(\Delta - m_f^2) (t - \tau)} (F_{\varepsilon, Y} (\tilde{\Psi}_{\tau}
  (\cdot))) \mathd \tau + \tilde{B}^A_{\varepsilon, t} \label{eq:real1}
\end{equation}
as a function of $\lambda^2$, i.e. we want to study the regularity of the
function $\tilde{\lambda} \longmapsto Z_{F_{\varepsilon, Y},
\tilde{B}^A_{\varepsilon, t}} (\tilde{\lambda})$ from a neighborhood $U
\subset \mathbb{R}$ of 0 into $C^0 (\mathbb{R}, \mathcal{C}^q (\mathbb{R}^2,
\mathcal{A} \otimes \mathbb{C}^2))$ (here $0 < q < \gamma$ where $\gamma$ is
the constant chosen in Remark \ref{remark_alpha}), where
\[ Z_{F_{\varepsilon, Y}, \tilde{B}^A_{\varepsilon, t}} (\tilde{\lambda})
   \assign \tilde{\Psi}^s_{\varepsilon, \tilde{\lambda}, t} (x) \in C^0
   (\mathbb{R}, \mathcal{C}^q (\mathbb{R}^2, \mathcal{A} \otimes
   \mathbb{C}^4)) . \]
In particular we want to prove the following theorem.

\begin{theorem}
  \label{theorem:realanalyticity}Under the Hypotheses of Theorem
  \ref{theorem_main2}, the function $Z_{F_{\varepsilon, Y}, \tilde{B}^A_t}
  (\tilde{\lambda})$ is real analytic (see Definition
  \ref{definition:realanalytic}) in a neighborhood of $0$.
\end{theorem}

An easy consequence of Theorem \ref{theorem:realanalyticity} is the following
corollary on the Schwinger functions of the (regularized) Yukawa$_2$ model.

\begin{corollary}
  \label{coroallry:realanalytic}Under the Hypotheses of Theorem
  \ref{theorem_main2}, we have that the Schwinger functions of the
  (regularized) Yukawa$_2$ model, i.e. for any $f_1, \ldots, f_n \in
  C^{\infty}_c (\mathbb{R}^2) \otimes \mathbb{C}^4$
  \begin{equation}
    \mathcal{S}_n (f_1, \ldots, f_n | \lambda) = \lim_{h_n \rightarrow 1}
    \frac{\omega \left( \tilde{\psi} (f_1) \cdots \tilde{\psi} (f_k)
    e^{\mathcal{V}_{\varepsilon, h_n}} \right)}{\omega \left(
    e^{\mathcal{V}_{\varepsilon, h_n}} \right)} = \omega
    (\tilde{\Psi}_{\varepsilon, \lambda^2, t}^s (f_1) \cdots
    \tilde{\Psi}_{\varepsilon, \lambda^2, t}^s (f_k))
  \end{equation}
  (where $h_n$ is a sequence of spatial cut-off as defined in Section
  \ref{sec:yukawa}) is a real analytic function of the coupling constant
  $\lambda$.
\end{corollary}

\begin{proof}
  We have that the functional $(Y_1, \ldots, Y_n) \mapsto \omega (Y_1 (f_1)
  \cdots Y_n (f_n))$ (where $(Y_1, \ldots, Y_n \in \mathcal{C}^q
  (\mathbb{R}^2, \mathcal{A} \otimes \mathbb{C}^4))$ is a multilinear
  continuous functional (and so real analytic), from $\mathcal{C}^q
  (\mathbb{R}^2, \mathcal{A} \otimes \mathbb{C}^4)$ into $\mathbb{C}$ (since
  $\omega$ is continuous and linear from $\mathcal{A}$ into $\mathbb{C}$ and
  the product in $\mathcal{A}$ is also continuous and multilinear). Since the
  composition of real analytic functionals is real analytic the thesis
  follows.
\end{proof}

\begin{remark}
  Theorem \ref{theorem:realanalyticity} and Corollary
  \ref{coroallry:realanalytic} do not depend on the fact that the
  non-linearity $F_{\varepsilon, h, Y}$ has the specific form of the case of
  Yukawa interaction (i.e., for example, that is a third degree polynomial or
  on the specific non local form). In particular Theorem
  \ref{theorem:realanalyticity} and Corollary \ref{coroallry:realanalytic}
  extend easily to the (regularized) quantum field fermionic models mentioned
  in Appendix \ref{appendix_models}.
\end{remark}

The rest of this appendix contains the proof of Theorem
\ref{theorem:realanalyticity}. Theorem \ref{theorem:realanalyticity} can be
seen as an infinite dimensional generalization of the results in Appendix
\ref{app:series}, where, between the other things, it is proved that the real
analyticity of Theorem \ref{theorem:realanalyticity} holds for finite
dimensional systems. In our case we can consider the finite dimensional
approximation
\begin{equation}
  \tilde{\Psi}_{\varepsilon, \tilde{\lambda}, N, R, t}^{h_R, s} =
  \tilde{\lambda} \int_{- \infty}^t e^{(\Delta - m_f^2) (t - \tau)} P_N
  (F_{\varepsilon, h_R, Y} (\tilde{\Psi}_{\varepsilon, R, \tau}^{h_R, s}))
  \mathd \tau + \tilde{B}^A_{\varepsilon, R, t} (v) \label{eq:real2}
\end{equation}
and the relative function $Z_{P_N (F_{\varepsilon, h_R, Y}),
\tilde{B}^A_{\varepsilon, R, t}} (\tilde{\lambda}) \assign
\tilde{\Psi}_{\varepsilon, \tilde{\lambda}, N, R, t}^{h_R, s}$.

\begin{lemma}
  \label{lemma:realanalytic}The function $\tilde{\lambda} \mapsto Z_{P_N
  (F_{\varepsilon, h_R, Y}), \tilde{B}^A_{\varepsilon, R, t}}
  (\tilde{\lambda})$ is real analytic in a neighborhood of $0$
\end{lemma}

\begin{proof}
  The result follows from the proof of convergence of perturbation series
  given in Appendix \ref{app:series} (in particular from inequalities
  {\eqref{eq:estimateexpansion}} and {\eqref{eq:estimateexpansion2}}) for
  finite dimensional Grassmannian equation and from the fact that the topology
  on $\mathcal{G} (P_N (C^{\infty} (\mathbb{T}^2_R, \mathcal{A} \otimes
  \mathbb{C}^4)))$ does depend on the norm chosen on $P_N (C^{\infty}
  (\mathbb{T}^2_R, \mathcal{A} \otimes \mathbb{C}^4))$ since all the norms on
  $P_N (C^{\infty} (\mathbb{T}^2_R, \mathcal{A} \otimes \mathbb{C}^4))$ are
  equivalent being, this space, finete dimensional. 
\end{proof}

Instead of proving the Theorem \ref{theorem:realanalyticity} by trying to
generalize the computation done in Appendix \ref{app:series} we present a
proof using the approximation method described in Section \ref{sec:yukawa},
since it is closer to the SPDE approach followed in the main part of the
present paper. For this reason we want to introduce some sort of apriori
estimate characterizing real analytic functions. This is done exploiting the
notion of formal power series and majorants method (see, e.g.,
{\cite{Van2003}} for an introduction to the subject). We call a formal power
series around $0$ taking values in the Banach space $X$ the following object:
\begin{equation}
  f = \sum_{n = 0}^{+ \infty} f^{(n)} \lambda^n, \quad \lambda \in \mathbb{C},
  \label{eq:formalseries}
\end{equation}
where $f^{(n)} \in X$. We write $f (0) = f^{(0)}$ and also
\[ \frac{\mathd^n f}{\mathd \lambda^n} = \sum_{k = n}^{+ \infty} f^{(n)}
   \lambda^{n - k} \prod_{i = 0}^{n - 1} (k - i) . \]
When the power series {\eqref{eq:formalseries}} is absolutely convergent for
$0 < | \lambda | < \varepsilon$ (for some $\varepsilon > 0$) we identify the
formal series {\eqref{eq:formalseries}} with the real analytic function
$\lambda \mapsto f (\lambda) \assign \sum_{n = 0}^{+ \infty} f^{(n)} \lambda^n
\in \mathbb{C}$ defined for $| \lambda | < \varepsilon$.

\begin{definition}
  Let $f_X$ be a formal power series around $0$ taking values in the Banach
  space $X$ and let $K$ be a formal power series around $0$ taking values in
  $\mathbb{C}$. We say that $f$ is majored by $K$, and we write $f_X
  \trianglelefteqslant K$, if for any $n \in \mathbb{N}$ we have
  \[ \left\| \frac{\mathd^n f_X}{\mathd \lambda^n} (0) \right\|_X = n! \|
     f^{(n)} \|_X \leqslant \frac{\mathd^n K}{\mathd \lambda^n} (0) =
     n!K^{(n)}, \]
  where $K (\lambda) = \sum_{n = 0}^{+ \infty} K^{(n)} \lambda^n$.
\end{definition}

\begin{remark}
  \label{remark:seriesboundedness}If $f_{\mathbb{C}}$ is a formal power series
  in $\mathbb{C}$ and $K$ is an analytic function such that $f_{\mathbb{C}}
  \trianglelefteqslant K$ then also $f_{\mathbb{C}}$ is an analytic function.
  Furthermore if the power series defining $g$ (in a neighborhood of $0$)
  converges absolutely for $\lambda = R > 0$, then the power series defining
  $f_{\mathbb{C}}$ converges absolutely for $\lambda \in \mathbb{C}$ with $|
  \lambda | \leqslant R$ and furthermore
  \[ \sup_{\lambda \in \mathbb{C}, | \lambda | \leqslant R} | f_{\mathbb{C}}
     (\lambda) | \leqslant K (R) . \]
\end{remark}

In general it is possible to define the sum and product of formal power
series. Furthermore if $g : \mathbb{C} \rightarrow \mathbb{C}$ is an entire
function and $f_{\mathbb{C}}$ is a formal power series taking values in
$\mathbb{C}$ it is possible to define the composition $g \circ f_{\mathbb{C}}$
as a formal power series taking values in $0$.

Let $g$ be an analytic function such that $0 \trianglelefteqslant g$, i.e. $g
(\lambda) \assign \sum_{n = 0}^{+ \infty} g^{(n)} \lambda^n$ such that
$g^{(n)} \geqslant 0$, consider $c \geqslant 0$ and define
\[ G_{g, c} (a, \lambda) = \lambda g (a) + c. \]
\begin{lemma}
  \label{lemma:majorant}For any entire function $g$ such that $0
  \trianglelefteqslant g$ and $c \geqslant 0$ there exists a unique analytic
  function $K : U \rightarrow \mathbb{R}$ (where $U$ is a suitable
  neighborhood of $0$) such that
  \begin{equation}
    K (\lambda) = G_{g, c} (K (\lambda), \lambda) \label{eq:equationK1}
  \end{equation}
  for any $\lambda \in U$. Furthermore if $f_{\mathbb{C}}$ is a formal power
  series such that $f_{\mathbb{C}} \trianglelefteqslant G_{g, c}
  (f_{\mathbb{C}} (\cdot), \cdot)$ then $f_{\mathbb{C}} \trianglelefteqslant
  K$.
\end{lemma}

\begin{proof}
  Equation {\eqref{eq:equationK1}} admits a unique real analytic solution $K
  (\lambda)$ defined in a neighborhood $U$ of $0$ and taking values in a
  neighborhood $V$ of $c \in \mathbb{R}$. Indeed we have that the point $(c,
  0) = (k, \lambda)$ is solution to the equation
  \begin{equation}
    k - G_{g, c} (k, \lambda) = 0. \label{eq:equationklambda}
  \end{equation}
  Furthermore we have that
  \[ \partial_k (k - G_{g, c} (k, \lambda)) |_{(k, \lambda) = (c, 0)} = (1 -
     \lambda g' (k)) |_{(k, \lambda) = (c, 0)} = 1_{} . \]
  Thus by the implicit function theorem for holomorphic functions (see, e.g.,
  {\cite{Markushevich1965}} Chapter 3, Theorem 3.11), there is are two sets
  $U, V \subset \mathbb{R}$ (such that $V \times U$ is a neigborhood of $(c,
  0)$) and a unique real analytic function $K : U \rightarrow V$ solution to
  equation {\eqref{eq:equationK1}}. Furthermore this $K$ is the unique
  solution to equation {\eqref{eq:equationK1}} of the form $K : U \rightarrow
  \mathbb{R}$ since $(c, 0)$ is the unique solution to equation
  {\eqref{eq:equationklambda}} of the form $(\cdot, 0) \in \mathbb{R}^2$.
  
  The fact that $K (\lambda) = \sum_{n = 0}^{+ \infty} K^{(n)} \lambda^n$ is
  the unique solution to equation {\eqref{eq:equationK1}} in a neighborhood of
  $0$ implies that
  \[ K^{(1)} = \left. \left. \frac{\mathd}{\mathd \lambda} (\lambda g (K
     (\lambda))) \right| \right._{\lambda = 0} = g (K (0)) = \sum_{k = 0}^{+
     \infty} g^{(k)} (K^{(0)})^k \]
  and also, for any $n \geqslant 2$,
  \begin{eqnarray}
    K^{(n)} & = & \frac{1}{n!} \left. \left. \frac{\mathd^n}{\mathd \lambda^n}
    (\lambda g (K (\lambda))) \right| \right._{\lambda = 0} \nonumber\\
    & = & \frac{n}{n!} \left. \left. \frac{\mathd^{n - 1}}{\mathd \lambda^{n
    - 1}} g (K (\lambda)) \right| \right._{\lambda = 0} \nonumber\\
    & = & \sum_{k = 1}^{+ \infty} \sum_{\ell_1 + \cdots + \ell_k = n - 1}
    g^{(k)} K^{(\ell_1)} \cdots K^{(\ell_k)} .  \label{eq:serieK}
  \end{eqnarray}
  Suppose that $f_{\mathbb{C}} \trianglelefteqslant G_{g, c} (f_{\mathbb{C}}
  (\cdot), \cdot)$, with $f_{\mathbb{C}} (\lambda) = \sum_{n = 0}^{+ \infty}
  f^{(n)} \lambda^n$, then we want to prove that $| f^{(n)} | \leqslant
  K^{(n)}$ by induction on $n$. For $n = 0$ we have
  \[ | f^{(0)} | \leqslant G_{g, c} (f (0), 0) = c = K^{(0)} . \]
  For $n = 1$ we have
  \[ | f^{(1)} | \leqslant \frac{\mathd}{\mathd \lambda} (G_{g, c}
     (f_{\mathbb{C}} (\lambda), \lambda)) = g (f_{\mathbb{C}} (0)) = \sum_{k =
     0}^{+ \infty} g^{(k)} (f^{(0)})^k \leqslant \sum_{k = 0}^{+ \infty}
     g^{(k)} (K^{(0)})^k = K^{(1)} . \]
  Supposing that $| f^{(k)} | \leqslant K^{(k)}$ for any $k < n$, where $n
  \geqslant 2$, then we will have
  \begin{eqnarray}
    | f^{(n)} | & \leqslant & \frac{1}{n!} \left. \frac{\mathd^n}{\mathd
    \lambda^n} G_{g, c} (f_{\mathbb{C}} (\lambda), \lambda) \right|_{\lambda =
    0} \nonumber\\
    & \leqslant & \frac{n}{n!} \left. \frac{\mathd^{n - 1}}{\mathd \lambda^{n
    - 1}} g (f_{\mathbb{C}} (\lambda)) \right|_{\lambda = 0} \nonumber\\
    & \leqslant & \sum_{k = 1}^{+ \infty} \sum_{\ell_1 + \cdots + \ell_k = n
    - 1} g^{(k)} f^{(\ell_1)} \cdots f^{(\ell_k)} \nonumber\\
    & \leqslant & \sum_{k = 1}^{+ \infty} \sum_{\ell_1 + \cdots + \ell_k = n
    - 1} g^{(k)} K^{(\ell_1)} \cdots K^{(\ell_k)} \nonumber\\
    & \leqslant & K^{(n)} \nonumber
  \end{eqnarray}
  where we used equality {\eqref{eq:serieK}} and the induction hypothesis.
  This proves that $f_{\mathbb{C}} \trianglelefteqslant K$.
\end{proof}

\begin{remark}
  \label{remark:monotonicity}From formula {\eqref{eq:serieK}} it is clear that
  the map $(g, c) \mapsto K$ is monotone, i.e. if $K_1 (\lambda) = G_{g_1,
  c_1} (K_1 (\lambda), \lambda)$, $K_2 (\lambda) = G_{g_2, c_2} (K_2
  (\lambda), \lambda)$ and $g_1 \trianglelefteqslant g_2$ and $c_1 \leqslant
  c_2$ then $K_1 \trianglelefteqslant K_2$.
\end{remark}

We want to use Lemma \ref{lemma:majorant} for giving an apriori bound on the
norm of $\frac{\partial^n \tilde{\Psi}_{\varepsilon, \lambda^2, t}^s}{\partial
\lambda^n}$ (with $\tilde{\Psi}_{\varepsilon, \lambda^2, t}^s$ as in equation
{\eqref{eq:real1}}) when $\lambda = 0$. This would imply that the
$\tilde{\Psi}_{\varepsilon, \lambda^2, t}^s$ is real analytic in $\lambda$
when $\lambda$ is small enough.

\

First of all we introduce the notion of real-analyticity for functions taking
values in a Banach space.

\begin{definition}
  \label{definition:realanalytic}Let $X$ be a Banach space and let $Z : U
  \rightarrow X$ be an infinitly differentiable function (where $U \subset
  \mathbb{R}$ is an open set). We say that $Z$ is real analytic in $U$ if for
  any $\lambda_0 \in U$ we have
  \begin{equation}
    Z (\lambda) = \sum_{k = 0}^{+ \infty} \frac{1}{n!} \frac{\mathd^n
    Z}{\mathd \lambda^n} (\lambda_0) (\lambda - \lambda_0)^n
    \label{eq:seriesZ}
  \end{equation}
  for $| \lambda - \lambda_0 | > 0$ small enough (where the series on the
  right hand side of {\eqref{eq:seriesZ}} is supposed to converge absolutely
  in $X$).
\end{definition}

If $Z : U \rightarrow X$ is differentiable an infinite number of times at $0$
we can define the formal power series
\begin{equation}
  \tilde{Z} (\lambda) \assign \sum_{n = 0}^{+ \infty} Z^{(n)} \lambda^n
  \label{eq:expansionZ1}
\end{equation}
where $Z^{(n)} \in X$ is given by
\[ Z^{(n)} \assign \frac{1}{n!} \frac{\mathd^n Z}{\mathd \lambda^n} (0) \in X.
\]
We can consider also the formal power series $S_Z (\lambda)$ defined as
\begin{equation}
  S_Z (\lambda) \assign \sum_{n = 0}^{+ \infty} \| Z^{(n)} \| \lambda^n
  \label{eq:definitionS} .
\end{equation}
It is useful to recall an equivalent characterization of real analytic
functions in Banach spaces.

\begin{remark}
  We have that $\tilde{Z}$ (as given by {\eqref{eq:expansionZ1}}) is
  convergent in a neighborhood of $0$ (and so it defines an analytic function
  in a neighborhood of $0$) if and only if $S_Z$ (as given by
  {\eqref{eq:definitionS}}) is real analytic. Furthermore for any $0
  \trianglelefteqslant g$, we have $\tilde{Z} \trianglelefteqslant g$ if and
  only if $S_Z \trianglelefteqslant g$.
\end{remark}

It is important to note that if $\tilde{Z}$ is convergent (and so it defines
an analytic function in a neighborhood of $0$) this is not enough for proving
that $Z$ is also real analytic in a neighborhood of $0$. Indeed, in general,
if $\tilde{Z}$ is real analytic we could have $\tilde{Z} (\lambda) \centernot{=} Z
(\lambda)$ for $\lambda \centernot{=} 0$.

In order to overcome this problem we introduce the following useful theorems.

\begin{theorem}
  \label{theorem:propertiesanalytic}Let $X$ be a Banach space and let $Z : U
  \subset \mathbb{R} \rightarrow X$ be a function defined in a neighborhood of
  $0$. Then $Z$ is real-analytic in $U$ if and only if, for any $\ell^{\ast}
  \in X^{\ast}$, the function $\lambda \mapsto \ell^{\ast} (Z (\lambda))$ is
  real analytic in $U$.
\end{theorem}

\begin{proof}
  The proof can be found in {\cite{Browder1962}} Proposition 1 and Remark 2.
\end{proof}

\begin{theorem}
  \label{theorem:convergencerealanalytic}Suppose that $Z_n (\lambda)$ is a
  sequence of functions from an open neighborhood $U \subset \mathbb{R}$ of
  $0$ into a Banach space $X$, converging pointwise $Z^n (\lambda) \rightarrow
  Z (\lambda)$ to some function from $U$ into $X$. Suppose furthermore that
  each $Z_n (\lambda)$ is real analytic and that there is a real analytic
  function $0 \trianglelefteqslant K : U \rightarrow \mathbb{R}$ such that
  $Z_n \trianglelefteqslant K$, then $Z (\lambda)$ is real analytic in a
  neighborhood of $0$ and $Z \trianglelefteq K$.
\end{theorem}

\begin{proof}
  Consider $\ell^{\ast} \in X^{\ast}$ such that $\| \ell^{\ast} \|_{X^{\ast}}
  \leqslant 1$ and define $\ell^{\ast}_{Z_n} (\lambda) \assign \ell^{\ast}
  (Z_n (\lambda))$. Since $Z_n$ are real analytic also $\ell^{\ast}_{Z_n}$ is
  real analytic, and also $\ell^{\ast}_{Z_n} (\lambda) \rightarrow
  \ell^{\ast}_Z (\lambda) \assign \ell^{\ast} (Z (\lambda))$ for any $\lambda
  \in U$. Furthermore $\ell^{\ast}_{Z_n} \trianglelefteqslant S_{Z_n}
  \trianglelefteqslant K$ which implies that, by Remark
  \ref{remark:seriesboundedness}, for any $R$ such that $K (R)$ is well
  defined, we have
  \[ \sup_{\lambda \in \mathbb{C}, | \lambda | \leqslant R} |
     \ell^{\ast}_{Z_n} (\lambda) | \leqslant K (R), \]
  i.e. $\ell^{\ast}_{Z_n}$ are uniformly bounded in a (complex)-neighborhood
  of $0$. This means by Vitali-Porter theorem (see, i.e. {\cite{Schiff1993}}
  Chapter 2, Section 2.4) that also $\ell^{\ast}_Z (\lambda)$ is real analytic
  for $\lambda \in U \cap \{ | \lambda | < R \}$. Since $\ell^{\ast}_Z$ is
  real analytic for any $\ell^{\ast} \in X^{\ast}$ such that $\| \ell^{\ast}
  \|_{X^{\ast}} = 1$, the function $\tilde{\ell} (Z (\lambda))$ is real
  analytic in $U \cap \{ | \lambda | < R \}$ which implies, by Theorem
  \ref{theorem:propertiesanalytic}, that $Z$ is real analytic in $U \cap \{ |
  \lambda | < R \}$ which is a non-empty neighborhood of $0$.
  
  Finally we have
  \[ \left| \frac{\mathd^k S_Z}{\mathd \lambda^k} (0) \right| =
     \sup_{\ell^{\ast} \in X^{\ast}, \| \ell^{\ast} \|_{X^{\ast}} = 1} \left|
     \frac{\mathd^k \ell^{\ast}_Z}{\mathd \lambda^k} (0) \right| \leqslant
     \sup_{\ell^{\ast} \in X^{\ast}, \| \ell^{\ast} \|_{X^{\ast}} = 1, n \in
     \mathbb{N}} \left| \frac{\mathd^k \ell^{\ast}_{Z_n}}{\mathd \lambda^k}
     (0) \right| \leqslant K^{(k)} \]
  where we used that, from the Vitali-Porter theorem, $\frac{\mathd^k
  \ell^{\ast}_{Z_n}}{\mathd \lambda^k}$ converges to $\frac{\mathd^k
  \ell^{\ast}_Z}{\mathd \lambda^k}$ uniformly.
\end{proof}

\begin{remark}
  \label{remark:extensionconvergence}It is easy to generalize the proof of
  Theorem \ref{theorem:convergencerealanalytic} replacing the convergence of
  $Z_n (\lambda)$ to $Z (\lambda) \in X$ in the norm of $X$, with the
  convergence of $Z_n (\lambda)$ to $Z (\lambda) \in X$ with respect to a
  (weaker) norm of a Banach space $X' \supset X$ in which $X$ is densely
  contained.
\end{remark}

Let $P : X \rightarrow X$ be a polynomial map, i.e. there are some $P_0, P_1,
\ldots, P_n$ such that $P_k : X^k \rightarrow X$ are continuous $k$-linear
symmetric maps such that
\[ P (x) \assign \sum_{k = 0}^n P_k (x, \ldots, x) . \]
We also suppose that there is $\mathfrak{P} : \mathbb{R} \rightarrow
\mathbb{R}$ defined as
\[ \mathfrak{P} (\lambda) \assign \sum_{k = 0}^n \| P_k \|_{\mathcal{L} (X^k,
   X)} \lambda^k \]
where $\| \cdot \|_{\mathcal{L} (X^k, X)}$ is the natural operator norm of
$\mathcal{L} (X^k, X)$ that is the space of $k$-linear maps between $X^k$ into
$X$. Suppose that $Z_{P, C} : U \rightarrow X$ (where $U \subset \mathbb{R}$
is a neighborhood of 0) is the unique solution to the equation
\begin{equation}
  Z_{P, C} (\lambda) = \lambda P (Z (\lambda)) + C \label{eq:fixpointX}
\end{equation}
where $\lambda \in U$ and $C \in X$. We are now ready to prove a general
apriori estimate for equations of the form {\eqref{eq:fixpointX}}.

\begin{theorem}
  \label{theorem:aprioriexpansion}Suppose that a solution $Z_{P, C}$ to
  equation {\eqref{eq:fixpointX}} is differentiable an infinite number of time
  in $\lambda = 0$, then we have $Z_{P, C} \trianglelefteqslant K_{P, C}$
  where \ $K_{P, C}$ is the unique analytic solution (in a neighborhood of
  $0$) to the equation
  \begin{equation}
    K_{P, C} (\lambda) = \lambda \mathfrak{P} (K_{P, C} (\lambda)) + \| C \|_X
    . \label{eq:equationflambda}
  \end{equation}
\end{theorem}

\begin{proof}
  If we take the derivatives in the variable $\lambda$ in the point $\lambda =
  0$ and write
  \[ Z_{P, C}^{(n)} \assign \frac{1}{n!} \frac{\mathd^n Z_{P, C}}{\mathd
     \lambda^n} (0) \in X, \]
  we obtain from the formal series {\eqref{eq:expansionZ1}} and equation
  {\eqref{eq:fixpointX}}
  \[ Z^{(0)}_{P, C} = C, \quad Z^{(1)}_{P, C} = P (Z^{(0)}) = P (C) \]
  and, for $r \geqslant 2$,
  \[ Z^{(r)}_{P, C} = \sum_{k = 0}^n \sum_{j_1 + \cdots + j_k = r - 1} P_k
     (Z^{(j_1)}_{P, C}, \ldots, Z^{(j_k)}_{P, C}) . \]
  This implies that
  \[ \| Z^{(0)}_{P, C} \|_X = \| C \|_X, \quad \| Z^{(1)}_{P, C} \| \leqslant
     \mathfrak{P} (\| C \|_X) \]
  and, for $r \geqslant 2$,
  \begin{equation}
    \| Z^{(r)}_{P, C} \| \leqslant \sum_{k = 0}^n \sum_{j_1 + \cdots + j_{2 k
    + 1} = r - 1} \| P_{2 k + 1} \|_{\mathcal{L} (X^{2 k + 1}, X)} \|
    Z^{(j_1)}_{P, C} \|_X \cdots \| Z^{(j_{2 k + 1})}_{P, C} \|_X .
    \label{eq:inequalityZ1}
  \end{equation}
  \[ \  \]
  If we set $S_{Z_{P, C}} = \sum_{n = 0}^{+ \infty} \| Z^{(n)}_{P, C} \|
  \lambda^n$, then we have
  \begin{eqnarray}
    S_{Z_{P, C}} (\lambda) & = & \| C \|_X + \sum_{r = 1}^{+ \infty} \|
    Z^{(k)}_{P, C} \| \lambda^r \nonumber\\
    & \trianglelefteqslant & \| C \|_X + \sum_{r = 1}^{+ \infty} \lambda^r
    \sum_{k = 0}^n \sum_{j_1 + \cdots + j_k = r - 1} \| P_k \|_{\mathcal{L}
    (X^k, X)} \| Z^{(j_1)}_{P, C} \|_X \cdots \| Z^{(j_k)}_{P, C} \|_X
    \nonumber\\
    & \trianglelefteqslant & \| C \|_X + \lambda \sum_{k = 0}^n \sum_{r =
    0}^{+ \infty} \lambda^r \sum_{j_1 + \cdots + j_k = r} \| P_k
    \|_{\mathcal{L} (X^k, X)} \| Z^{(j_1)}_{P, C} \|_X \cdots \| Z^{(j_k)}_{P,
    C} \|_X \nonumber\\
    & \trianglelefteqslant & \| C \|_X + \lambda \mathfrak{P} (S_{Z_{P, C}}
    (\lambda)) \nonumber
  \end{eqnarray}
  where we used inequality {\eqref{eq:inequalityZ1}} and the fact that, by
  Theorem \ref{theorem:propertiesanalytic}, the function $S_{P, C}$ is real
  analytic in a neighborhood of $0$ with positive derivatives in $0$. By Lemma
  \ref{lemma:majorant} the thesis follows by taking $G_{g, c} (a, \lambda) =
  \lambda \mathfrak{P} (a) + \| C \|_X$.
\end{proof}

\begin{proof*}{Proof of Theorem \ref{theorem:realanalyticity}}
  We want to apply the previous theorems for proving Theorem
  \ref{theorem:realanalyticity}. Here the Banach space $X$ is $X = C^0
  (\mathbb{R}, \mathcal{C}^q (\mathbb{R}^2, \mathcal{A} \otimes
  \mathbb{C}^4))$ with the norm given by
  \[ \| \tilde{\Psi} \|_X = \sup_{t \in \mathbb{R}} \| \tilde{\Psi}
     \|_{\mathcal{C}^q (\mathbb{R}^2, \mathcal{A} \otimes \mathbb{C}^4)} . \]
  The function $Z_{P_N (F_{\varepsilon, h_R, Y}), \tilde{B}^A_{\varepsilon, R,
  t}} (\tilde{\lambda})_{}$ is real analytic in $\tilde{\lambda}$ in its
  radius of convergence. Furthermore it is a solution to equation
  {\eqref{eq:real2}}, which is of the form {\eqref{eq:fixpointX}} with
  \begin{equation}
    P_{3, h_R, N} (\tilde{\Psi}) (t, \cdot) = \int_{- \infty}^t e^{(\Delta -
    m_f^2) (t - \tau)} P_N (F_{\varepsilon, h_R, Y} (\tilde{\Psi} (\tau,
    \cdot))) \mathd \tau, \label{eq:P3etc}
  \end{equation}
  $C = \tilde{B}^A_{\varepsilon, R, t}$, and $\mathfrak{P}_{h_R, N}
  (\tilde{\lambda}) = \| P_{3, h_R, N} \|_{\mathcal{L} (X^3, X)}
  \tilde{\lambda}^3$. This means that, by Theorem
  \ref{theorem:aprioriexpansion}, we have \[Z_{P_N (F_{\varepsilon, h_R, Y}),
  \tilde{B}^A_{\varepsilon, R, t}} \trianglelefteqslant K_{P_{3, h_R, N},
  \tilde{B}^A_{\varepsilon, R}},\] where $K_{P_{3, h, N},
  \tilde{B}^A_{\varepsilon, R}}$ is the solution to the equation \[K_{P_{3,
  h_R, N}, \tilde{B}^A_{\varepsilon, R}} (\tilde{\lambda}) = \tilde{\lambda}
  \| P_{3, h_R, N} \|_{\mathcal{L} (X^3, X)} \left( K_{P_{3, h_R, N},
  \tilde{B}^A_{\varepsilon, R X}} (\tilde{\lambda}) \right)^3 + \|
  \tilde{B}^A_{\varepsilon, R} \|_X .\] 
  Since $\| P_{3, h_R, N} \|_{\mathcal{L}
  (X^3, X)} \leqslant \| P_{3, h_R} \|_{\mathcal{L} (X^3, X)}$ (i.e. the
  polynomial $P_{3, h_R, N}$ without the finite dimensional projection), by
  Remark \ref{remark:monotonicity}, we have $Z_{P_N (F_{\varepsilon, h_R, Y}),
  \tilde{B}^A_{\varepsilon, R, t}} \trianglelefteqslant K_{P_{3, h_R},
  \tilde{B}^A_{\varepsilon, R}}$. Since, by Theorem
  \ref{theorem_convergence1}, $Z_{P_N (F_{\varepsilon, h_R, Y}),
  \tilde{B}^A_{\varepsilon, R, t}}$ converges pointwise to $Z_{F_{\varepsilon,
  h_R, Y}, \tilde{B}^A_{\varepsilon, R, t}}$ (i.e. the solution to equation
  {\eqref{eq:real2}} without the finite dimensional projection in the
  non-linearity) we have that, by Theorem
  \ref{theorem:convergencerealanalytic}, $Z_{F_{\varepsilon, h_R, Y},
  \tilde{B}^A_{\varepsilon, R, t}}$ is real analytic and $Z_{F_{\varepsilon,
  h_R, Y}, \tilde{B}^A_{\varepsilon, R, t}} \trianglelefteqslant K_{P_{3,
  h_R}, \tilde{B}^A_{\varepsilon, R}}$.
  
  Using again Remark \ref{remark:monotonicity} and the fact that $\|
  \tilde{B}^A_{\varepsilon, R, t} \| \leqslant \| \tilde{B}^A_{\varepsilon, t}
  \|_X$ (where, we recall, that $\tilde{B}^A_{\varepsilon, R, t}$ is the
  periodized version of $\tilde{B}^A_{\varepsilon, t}$), we get
  $Z_{F_{\varepsilon, h_R, Y}, \tilde{B}^A_{\varepsilon, R, t}}
  \trianglelefteqslant K_{P_{3, h_R}, \tilde{B}^A_{\varepsilon, R}}
  \trianglelefteqslant K_{P_{3, h}, \tilde{B}^A_{\varepsilon}}$. Furthermore,
  since, by Theorem \ref{theorem_torusplane}, $Z_{F_{\varepsilon, h_R, Y},
  \tilde{B}^A_{\varepsilon, R, t}}$ converges to $Z_{F_{\varepsilon, h, Y},
  \tilde{B}^A_{\varepsilon, t}}$ (where $Z_{F_{\varepsilon, h, Y},
  \tilde{B}^A_{\varepsilon, t}}$ is the solution to equation
  {\eqref{eq:real2}} where $\tilde{B}^A_{\varepsilon, R, t}$ is replaced by
  the non-periodized version $\tilde{B}^A_{\varepsilon, t}$ of
  $\tilde{B}^A_{\varepsilon, R, t}$) in $X' = C^0 (\mathbb{R},
  \mathcal{C}^q_{\ell} (\mathbb{R}^2, \mathcal{A} \otimes \mathbb{C}^4))
  \supset X$ (where we consider the weighted space $\mathcal{C}^q_{\ell}$
  instead of the unweighted one), by Theorem
  \ref{theorem:convergencerealanalytic} and Remark
  \ref{remark:extensionconvergence}, we obtain that $Z_{F_{\varepsilon, h, Y},
  \tilde{B}^A_{\varepsilon, t}}$ is real analytic with respect to
  $\tilde{\lambda}$ and also $Z_{F_{\varepsilon, h, Y},
  \tilde{B}^A_{\varepsilon, t}} \trianglelefteqslant K_{P_{3, h},
  \tilde{B}^A_{\varepsilon}}$. The thesis follows in a similar way, using the
  fact that, by Remark \ref{remark:monotonicity} and the fact that $\| P_{3,
  h} \|_{\mathcal{L} (X^3, X)} \leqslant \| P_3 \|_{\mathcal{L} (X^3, X)}$
  (where $P_3$ is the nonlinearity defined in equation {\eqref{eq:P3etc}}
  where all the cut-offs have been removed), we have $K_{P_{3, h},
  \tilde{B}^A_{\varepsilon}} \trianglelefteqslant K_{P_3,
  \tilde{B}^A_{\varepsilon}}$, and the convergence $Z_{F_{\varepsilon, h, Y},
  \tilde{B}^A_{\varepsilon, t}}$ to $Z_{F_{\varepsilon, Y},
  \tilde{B}^A_{\varepsilon, t}}$ (i.e. the solution to equation
  {\eqref{eq:real1}}).
\end{proof*}

\section{Other interacting models}\label{appendix_models}

In this last appendix we want to give an idea of other Euclidean quantum
fields models which can be approached within the theory developed in this
paper.

\

\tmtextit{Non-local, Yukawa$_2$-like models.} The discussion in
Section~\ref{sec:yukawa} can be extended in a straightforward way to other
fermionic quantum field theories. One of them is two dimensional quantum
electrodynamics (QED$_2$), see,
e.g.~{\cite{mastropietro_non-perturbative_2008}} Chapter~3. In this case we
consider a fermionic (free) field $\tilde{\psi} \in \mathcal{S} (\mathbb{R}^2,
\mathcal{A} \otimes \mathbb{C}^4)$ with the same covariance as
in~{\eqref{eq:free-corrrelation}} and~{\eqref{eq:matrixK}}, and two
(massive-free-scalar) bosons $A = (A^1, A^2) \in \mathcal{S} (\mathbb{R}^2,
\mathcal{A} \otimes \mathbb{R}^2)$ independent of each other with covariance
\[ \langle A^i (x) A^j (y) \rangle = \delta_{i, j} (- \Delta + M^2)^{- 1} (x
   - y), \]
where $i, j \in \{ 1, 2 \}$, and $M \in \mathbb{R}_+$ is their mass. The
interaction potential is given (formally) by
\[ V_{\tmop{QED}} (A, \tilde{\psi}) = \lambda \sum_{i = 1, 2}
   \int_{\mathbb{R}^2}  (\bar{\psi} (x) \cdot \gamma_i \psi (x)) A^i (x)
   \mathd x, \]
where $\lambda \in \mathbb{R}$, which, when the bosons $A$ are integrated out,
gives rise to the effective potential
\begin{equation}
  \mathcal{V} (\tilde{\psi}) = \lambda^2 \sum_{i = 1, 2} \int_{\mathbb{R}^4}
  (\bar{\psi} (x) \cdot \gamma_i \psi (x)) \mathcal{G}_M (x - y) (\bar{\psi}
  (y) \cdot \gamma_i \psi (y)) \mathd x \mathd y \label{eq:QED1}
\end{equation}
where $\mathcal{G}_M (x - y) = (- \Delta + M^2)^{- 1} (x - y)$. If we consider
a regularized version $\tilde{\psi}_{\varepsilon} (x) = \tilde{\psi}
(\mathfrak{a}_{\varepsilon} (\cdot - x))$, $\varepsilon > 0$, as it was done
in Section \ref{sec:yukawa}, of the fermion field $\tilde{\psi}$, with the
purely fermionic interaction $\mathcal{V}$, the stochastic quantization
equation is of the form~{\eqref{eq:stochasticquantization}}, where the
non-linearity $F_{\varepsilon, Y} = F_{\varepsilon, 1, Y}$ is replaced by the
following term
\[ F_{\varepsilon, \tmop{QED}} (\tilde{\Psi}_{\varepsilon, t}) = C_f \left(
   \begin{array}{l}
     \sum_{i = 1, 2} \mathfrak{a}_{\varepsilon}^{\asterisk 2} \ast \left( -
     \bar{\Psi}_{\varepsilon, t} \gamma_i \int_{\mathbb{R}^2} \mathcal{G}_M (y
     - \cdot) \bar{\Psi}_{\varepsilon, t} (y) \cdot \gamma_i
     \Psi_{\varepsilon, t} (y) \mathd y \right)\\
     \sum_{i = 1, 2} \mathfrak{a}_{\varepsilon}^{\asterisk 2} \ast \left(
     \gamma_i \Psi_{\varepsilon, t} \int_{\mathbb{R}^2} \mathcal{G}_M (y -
     \cdot) \bar{\Psi}_{\varepsilon, t} (y) \cdot \gamma_i \Psi_{\varepsilon,
     t} (y) \mathd y \right)
   \end{array} \right), \]
and where $C_f$ is defined in eq.~{\eqref{eq:Cf}}. It is simple to generalize
Lemma~\ref{lemma_drift} for the nonlinearity $F_{\varepsilon, \tmop{QED}}$ and
this permits to prove an analogous version of Theorem~\ref{theorem_main2}
using the same methods as in Section~\ref{sec:yukawa}.

An analogous reasoning can be extended also to the Yukawa model with
pseudo-scalar current having an effective potential of the form
\[ \mathcal{V} (\tilde{\psi}) = \lambda^2 \int_{\mathbb{R}^4} (\bar{\psi} (x)
   \cdot \gamma_5 \psi (x)) \mathcal{G} (x - y) (\bar{\psi} (y) \cdot \gamma_5
   \psi (y)) \mathd x \mathd y, \]
where $\gamma_5 = \mathi \gamma_1 \gamma_2$ and $\mathcal{G} (x - y) = (-
\Delta + m_b^2)^{- 1} (x - y)$ as in eq.~{\eqref{eq:effectiveV}}.

\

\tmtextit{Local models.} The treatment proposed for the regularized Yukawa
model can extended also to purely fermionic models with quartic interaction.
We describe as an example the Gross--Neveu model
(see~{\cite{mastropietro_non-perturbative_2008}} Chapter~6 for this model) in
two dimensions. In this case we consider a (free) fermionic Dirac field
$\tilde{\psi}$ with interaction given by the potential
\[ \mathcal{V} (\tilde{\psi}) = - \lambda^2 \int_{\mathbb{R}^2} (\bar{\psi}
   (x) \cdot \psi (x))^2 \mathd x. \]
The main difference in the corresponding stochastic quantization equation,
compared with that of the QED$_2$ model above, is that $F_{\varepsilon,
\tmop{QED}}$ is replaced by
\[ F_{\varepsilon, \tmop{GN}} (\tilde{\Psi}_{\varepsilon, t}) = C_f \left(
   \begin{array}{l}
     \mathfrak{a}_{\varepsilon}^{\asterisk 2} \ast (- \bar{\Psi}_{\varepsilon,
     t} (\bar{\Psi}_{\varepsilon, t} \cdot \Psi_{\varepsilon, t}))\\
     \mathfrak{a}_{\varepsilon}^{\asterisk 2} \ast (\Psi_{\varepsilon, t}
     (\bar{\Psi}_{\varepsilon, t} \cdot \Psi_{\varepsilon, t}))
   \end{array} \right) \]
where $C_f$ is defined, as before, in eq.~{\eqref{eq:Cf}}. Also for
$F_{\varepsilon, \tmop{GN}}$ a generalization of Lemma~\ref{lemma_drift}
holds. This allows us to prove an analogous version of
Theorem~\ref{theorem_main2} using the same methods as in
Section~\ref{sec:yukawa}.

\

\tmtextit{Supersymmetric models.} Supersymmetric quantum field models
describe quantum fields where the bosonic and fermionic fields are related,
one to the other, by a special transformation called supersymmetry. It is not
possible to reduce a supersymmetric non free (i.e. interacting) model to a
purely fermionic one since the presence of interaction and supersymmetry
forces the potential to have a part of bosonic self interaction. Nevertheless
from an analytical point of view it can be interesting to study supersymmetric
models considering free bosons and integrating them out as was done in
Section~\ref{sec:yukawa}.

One of the most important examples of super-symmetric models is Wess--Zumino
model composed by two scalar bosons interacting with a Dirac fermionic field
through a Yukawa scalar and pseudo-scalar types interaction
(see~{\cite{Jaffe_supersymmetry,Arai_supersymmetry,haba_supersymmetry_1995}}).
This model can also be handled as in Section~\ref{sec:yukawa}. Another source
of supersymmetric quantum field models is the supersymmetric formulation of
stochastic differential equations (see~{\cite{ZinnJustin}} Chapter~15). In
particular here we discuss the model related with elliptic stochastic
differential equations
(see~{\cite{klein_supersymmetry_1983,klein_supersymmetry_1985,albeverio_elliptic_2018}}
see also~{\cite{de_vecchi_note_2019}}).

In this case we have a fermionic field $\tilde{\psi} = (\psi, \bar{\psi}) \in
\mathcal{S} (\mathbb{R}^2, \mathcal{A} \otimes \mathbb{R}^2)$ with covariance
\[ \langle \bar{\psi} (x) \psi (y) \rangle = (- \Delta + m^2)^{- 1} (x - y),
\]
where $m \in \mathbb{R}_+$ is a positive constant. The corresponding effective
potential $\mathcal{V}$, obtained after integrating over the boson field, is
\[ \mathcal{V} (\tilde{\psi}) = \lambda^2 \int_{\mathbb{R}^4} (\bar{\psi} (x)
   \psi (x)) \mathcal{K}_m (x - y) (\bar{\psi} (y) \psi (y)) \mathd x \mathd
   y, \]
where $\mathcal{K}_m (x - y) \assign (- \Delta + m^2)^{- 2} (x - y)$. The
stochastic quantization equation reads
\begin{equation}
  \mathd \tilde{\Psi}_{\varepsilon, t} = ((- \Delta + m^2)
  \tilde{\Psi}_{\varepsilon, t} + \lambda^2 F_{\varepsilon, E}
  (\tilde{\Psi}_{\varepsilon, t})) \mathd t + \mathd \hat{B}_{\varepsilon, t}
  \label{eq:elliptic}
\end{equation}
where $\hat{B}_{\varepsilon, t} = \mathfrak{a}_{\varepsilon} \asterisk (B_t,
\bar{B}_t)$, with $\mathfrak{a}_{\varepsilon}$ as in Section
\ref{sec:os-fields}, $(B_t, \bar{B}_t)$ is a couple of Brownian motion, i.e.
they are the integral of a Gaussian noise with covariance $J = \left(
\begin{array}{ll}
  0 & 1\\
  - 1 & 0
\end{array} \right)$, and $F_{\varepsilon, E}$ is defined by
\begin{equation}
  F_{\varepsilon, E} (\tilde{\Psi}_{\varepsilon, t}) = - J \cdot \left(
  \begin{array}{l}
    \mathfrak{a}_{\varepsilon}^{\asterisk 2} \ast \left( -
    \bar{\Psi}_{\varepsilon, t} \int_{\mathbb{R}^2} \mathcal{K}_m (y - \cdot)
    \bar{\Psi}_{\varepsilon, t} (y) \Psi_{\varepsilon, t} (y) \mathd y
    \right)\\
    \mathfrak{a}_{\varepsilon}^{\asterisk 2} \ast \left( \Psi_{\varepsilon, t}
    \int_{\mathbb{R}^2} \mathcal{K}_m (y - \cdot) \bar{\Psi}_{\varepsilon, t}
    (y) \Psi_{\varepsilon, t} (y) \mathd y \right)
  \end{array} \right) . \label{eq:elliptic2}
\end{equation}
Also in this case a generalization of Lemma~\ref{lemma_drift} and of
Theorem~\ref{theorem_main2} can be proven.

\

\tmtextit{Comments.} As it was shown, the theory presented in this article is
able to construct the regularized version of many models of fermionic quantum
field theories. The mathematical differences between these models can only be
seen when the regularization is removed. More precisely the solution to the
supersymmetric model~{\eqref{eq:elliptic}} and~{\eqref{eq:elliptic2}} is
expected to be a distribution of regularity $\mathcal{C}^{- \delta}
(\mathbb{R}^2, \mathcal{A} \otimes \mathbb{R}^2)$ (for any $\delta > 0$) as in
$\Phi^4_2$ bosonic case. The non-local interactions (resulting after the
bosons are integrated out) such as in the Yukawa model (in
Section~\ref{sec:yukawa}) and in the QED model (defined by
eqs.~{\eqref{eq:QED1}}) have solutions in $\mathcal{C}^{- 1 / 2 - \delta}
(\mathbb{R}^2, \mathcal{A} \otimes \mathbb{R}^2)$ (for any $\delta > 0$) and
they are (formally) super-renormalizable much like the $\Phi^4_3$ bosonic
model. The main difference between Yukawa and QED being that in the Yukawa
case the mass $m_b$ of the boson $\varphi$ used for defining $\mathcal{G}$
goes to $+ \infty$ as the regularisation is removed and have to be
renormalized, while in the QED setting the mass of the photon $M$ can be kept
constant. Finally, let us also mention that the removal of the regularization
$\varepsilon$ in the study of the Gross--Neveu model in our setting would be
particularly interesting (the model being only renormalizable and not
superrenormalizable as are the other ones, it remains to see whether it is
``asymptotic free'' character, according to physical terminology, will help
this construction).

\begin{thebibliography}{100}
  \bibitem[1]{accardi_outline_2015}Luigi Accardi. {\newblock}An outline of
  quantum probability. {\newblock}Preprint
  \href{https://doi.org/10.13140/rg.2.1.3078.3844}{}10.13140/rg.2.1.3078.3844,
  2015.{\newblock}
  
  \bibitem[2]{accardi_quantum_1982}Luigi Accardi, Alberto Frigerio, and 
  John~T.~Lewis. {\newblock}Quantum stochastic processes.
  {\newblock}\tmtextit{Kyoto University. Research Institute for Mathematical
  Sciences. Publications}, 18(1):97--133, 1982.
  {\newblock}\href{Https://doi.org/10.2977/prims/1195184017}{}10.2977/prims/1195184017.{\newblock}
  
  \bibitem[3]{albeverio_elliptic_2019}Sergio Albeverio, Francesco~C.~De
  Vecchi, and  Massimiliano Gubinelli. {\newblock}The elliptic stochastic
  quantization of some two dimensional Euclidean QFTs.
  {\newblock}\tmtextit{ArXiv:1906.11187 [math-ph]}, jun 2019.
  {\newblock}ArXiv: 1906.11187, to appear in {\textit{Annales de l'Institut
  Henri Poincar{\'e}, Probabilit{\'e}s et Statistiques}}.{\newblock}
  
  \bibitem[4]{albeverio_elliptic_2018}Sergio Albeverio, Francesco~C.~De
  Vecchi, and  Massimiliano Gubinelli. {\newblock}Elliptic stochastic
  quantization. {\newblock}\tmtextit{Ann. Probab.}, 48(4):1693--1741,
  2020.{\newblock}
  
  \bibitem[5]{albeverio_invariant_2017}Sergio Albeverio  and  Seiichiro
  Kusuoka. {\newblock}The invariant measure and the flow associated to the
  $\Phi^4_3$-quantum field model. {\newblock}\tmtextit{Annali della Scuola
  Normale di Pisa - Classe di Scienze}, 2018.
  {\newblock}\href{Https://doi.org/10.2422/2036-2145.201809_008}{}10.2422/2036-2145.201809\_008.{\newblock}
  
  \bibitem[6]{alpay_distribution_2019}Daniel Alpay, Ismael~L.~Paiva, and 
  Daniele~C.~Struppa. {\newblock}Distribution spaces and a new construction of
  stochastic processes associated with the Grassmann algebra.
  {\newblock}\tmtextit{Journal of Mathematical Physics}, 60(1):13508--21,
  2019.
  {\newblock}\href{Https://doi.org/10.1063/1.5052010}{}10.1063/1.5052010.{\newblock}
  
  \bibitem[7]{Amann1997}Herbert Amann. {\newblock}Operator-valued Fourier
  multipliers, vector-valued Besov spaces, and applications.
  {\newblock}\tmtextit{Math. Nachr.}, 186:5--56, 1997.
  {\newblock}\href{Https://doi.org/10.1002/mana.3211860102}{}10.1002/mana.3211860102.{\newblock}
  
  \bibitem[8]{Amann2019}Herbert Amann. {\newblock}\tmtextit{Linear and
  quasilinear parabolic problems. Vol. II},  volume  106  of
  \tmtextit{Monographs in Mathematics}. {\newblock}Birkh{\"a}user/Springer,
  Cham, 2019.
  {\newblock}\href{Https://doi.org/10.1007/978-3-030-11763-4}{}10.1007/978-3-030-11763-4.{\newblock}
  
  \bibitem[9]{angelis_berezin_1998}G.~F.~De Angelis, G.~Jona-Lasinio, and 
  V.~Sidoravicius. {\newblock}Berezin integrals and Poisson processes.
  {\newblock}\tmtextit{Journal of Physics A: Mathematical and General},
  31(1):289--308, jan 1998.
  {\newblock}\href{Https://doi.org/10.1088/0305-4470/31/1/026}{}10.1088/0305-4470/31/1/026.{\newblock}
  
  \bibitem[10]{Applebaum2}D.~B.~Applebaum  and  R.~L.~Hudson.
  {\newblock}Fermion It{\^o}'s formula and stochastic evolutions.
  {\newblock}\tmtextit{Comm. Math. Phys.}, 96(4):473--496, 1984.{\newblock}
  
  \bibitem[11]{Applebaum1}D.~Applebaum  and  R.~L.~Hudson. {\newblock}Fermion
  diffusions. {\newblock}\tmtextit{J. Math. Phys.}, 25(4):858--861, 1984.
  {\newblock}\href{Https://doi.org/10.1063/1.526236}{}10.1063/1.526236.{\newblock}
  
  \bibitem[12]{Arai_supersymmetry}Asao Arai. {\newblock}A general class of
  infinite-dimensional Dirac operators and path integral representation of
  their index. {\newblock}\tmtextit{J. Funct. Anal.}, 105(2):342--408,
  1992.{\newblock}
  
  \bibitem[13]{araki_diagonalization_1968}Huzihiro Araki. {\newblock}On the
  Diagonalization of a Bilinear Hamiltonian by a Bogoliubov Transformation.
  {\newblock}\tmtextit{Publications of the Research Institute for Mathematical
  Sciences, Kyoto University. Ser. A}, 4(2):387--412, 1968.
  {\newblock}\href{Https://doi.org/10.2977/prims/1195194882}{}10.2977/prims/1195194882.{\newblock}
  
  \bibitem[14]{araki_quasifree_1970}Huzihiro Araki. {\newblock}On Quasifree
  States of CAR and Bogoliubov Automorphisms.
  {\newblock}\tmtextit{Publications of the Research Institute for Mathematical
  Sciences}, 6(3):385--442, 1970.
  {\newblock}\href{Https://doi.org/10.2977/prims/1195193913}{}10.2977/prims/1195193913.{\newblock}
  
  \bibitem[15]{MR1799198}Huzihiro Araki. {\newblock}\tmtextit{Mathematical
  theory of quantum fields},  volume  101  of \tmtextit{International Series
  of Monographs on Physics}. {\newblock}Oxford University Press, New York,
  1999. {\newblock}Translated from the 1993 Japanese original by Ursula
  Carow-Watamura.{\newblock}
  
  \bibitem[16]{araki_representations_1964}Huzihiro Araki  and  Walter Wyss.
  {\newblock}Representations of canonical anticommutation relations.
  {\newblock}\tmtextit{Helvetica Physica Acta}, 37(2):136, 1964.{\newblock}
  
  \bibitem[17]{Segalbook}John~C.~Baez, Irving~E.~Segal, and  Zheng-Fang Zhou.
  {\newblock}\tmtextit{Introduction to algebraic and constructive quantum
  field theory}. {\newblock}Princeton Series in Physics. Princeton University
  Press, Princeton, NJ, 1992.
  {\newblock}\href{Https://doi.org/10.1515/9781400862504}{}10.1515/9781400862504.{\newblock}
  
  \bibitem[18]{Doplicher}Dorothea Bahns, Sergio Doplicher, Gerardo Morsella,
  and  Gherardo Piacitelli. {\newblock}Quantum spacetime and algebraic quantum
  field theory. {\newblock}In \tmtextit{Advances in algebraic quantum field
  theory}, Math. Phys. Stud.,  pages  289--329. Springer, Cham, 2015.
  {\newblock}\href{Https://doi.org/10.1007/978-3-319-21353-8_7}{}10.1007/978-3-319-21353-8\tmrsub{7}.{\newblock}
  
  \bibitem[19]{Streater3}C.~Barnett, R.~F.~Streater, and  I.~F.~Wilde.
  {\newblock}The It{\^o}-Clifford integral. {\newblock}\tmtextit{J. Functional
  Analysis}, 48(2):172--212, 1982.
  {\newblock}\href{Https://doi.org/10.1016/0022-1236(82)90066-0}{}10.1016/0022-1236(82)90066-0.{\newblock}
  
  \bibitem[20]{Streater2}C.~Barnett, R.~F.~Streater, and  I.~F.~Wilde.
  {\newblock}The It{\^o}-Clifford integral. II. Stochastic differential
  equations. {\newblock}\tmtextit{J. London Math. Soc. (2)}, 27(2):373--384,
  1983.
  {\newblock}\href{Https://doi.org/10.1112/jlms/s2-27.2.373}{}10.1112/jlms/s2-27.2.373.{\newblock}
  
  \bibitem[21]{Streater1}C.~Barnett, R.~F.~Streater, and  I.~F.~Wilde.
  {\newblock}The It{\^o}-Clifford integral. III. The Markov property of
  solutions to stochastic differential equations. {\newblock}\tmtextit{Comm.
  Math. Phys.}, 89(1):13--17, 1983.{\newblock}
  
  \bibitem[22]{battle_infinite_1981}Guy~A.~Battle  and  Lon Rosen.
  {\newblock}On the infinite volume limit of the strongly coupled
  Yukawa\tmrsub{$2$} model. {\newblock}\tmtextit{Journal of Mathematical
  Physics}, 22(4):770--776, apr 1981.
  {\newblock}\href{Https://doi.org/10.1063/1.524982}{}10.1063/1.524982.{\newblock}
  
  \bibitem[23]{Baumgartel1}Hellmut Baumg{\"a}rtel.
  {\newblock}\tmtextit{Operator algebraic methods in quantum field theory}.
  {\newblock}Akademie Verlag, Berlin, 1995. {\newblock}A series of
  lectures.{\newblock}
  
  \bibitem[24]{Baumgartel2}Hellmut Baumg{\"a}rtel  and  Manfred Wollenberg.
  {\newblock}\tmtextit{Causal nets of operator algebras},  volume~80  of
  \tmtextit{Mathematische Lehrb{\"u}cher und Monographien, II. Abteilung:
  Mathematische Monographien [Mathematical Textbooks and Monographs, Part II:
  Mathematical Monographs]}. {\newblock}Akademie-Verlag, Berlin, 1992.
  {\newblock}Mathematical aspects of algebraic quantum field
  theory.{\newblock}
  
  \bibitem[25]{Belavkin}Viacheslav~P.~Belavkin. {\newblock}Quantum stochastic
  calculus and quantum nonlinear filtering. {\newblock}\tmtextit{J.
  Multivariate Anal.}, 42(2):171--201, 1992.
  {\newblock}\href{Https://doi.org/10.1016/0047-259X(92)90042-E}{}10.1016/0047-259X(92)90042-E.{\newblock}
  
  \bibitem[26]{benfatto_functional_2007}G.~Benfatto, P.~Falco, and 
  V.~Mastropietro. {\newblock}Functional Integral Construction of the Massive
  Thirring model: Verification of Axioms and Massless Limit.
  {\newblock}\tmtextit{Communications in Mathematical Physics},
  273(1):67--118, jul 2007.
  {\newblock}\href{Https://doi.org/10.1007/s00220-007-0254-y}{}10.1007/s00220-007-0254-y.{\newblock}
  
  \bibitem[27]{MR1340626}Y.~M.~Berezansky  and  Y.~G.~Kondratiev.
  {\newblock}\tmtextit{Spectral methods in infinite-dimensional analysis. Vol.
  1},  volume  12/1  of \tmtextit{Mathematical Physics and Applied
  Mathematics}. {\newblock}Kluwer Academic Publishers, Dordrecht, 1995.
  {\newblock}Translated from the 1988 Russian original by P. V. Malyshev and
  D. V. Malyshev and revised by the authors.
  \href{https://doi.org/10.1007/978-94-011-0509-5}{}10.1007/978-94-011-0509-5.{\newblock}
  
  \bibitem[28]{Berezin1}F.~A.~Berezin. {\newblock}\tmtextit{The method of
  second quantization}. {\newblock}Translated from the Russian by Nobumichi
  Mugibayashi and Alan Jeffrey. Pure and Applied Physics, Vol. 24. Academic
  Press, New York-London, 1966.{\newblock}
  
  \bibitem[29]{Berezin2}Felix~Alexandrovich Berezin.
  {\newblock}\tmtextit{Introduction to superanalysis},  volume~9  of
  \tmtextit{Mathematical Physics and Applied Mathematics}. {\newblock}D.
  Reidel Publishing Co., Dordrecht, 1987.
  {\newblock}\href{Https://doi.org/10.1007/978-94-017-1963-6}{}10.1007/978-94-017-1963-6.{\newblock}
  
  \bibitem[30]{bertini_stochastic_1993}Lorenzo Bertini, Giovanni Jona-Lasinio,
  and  Claudio Parrinello. {\newblock}Stochastic Quantization, Stochastic
  Calculus and Path Integrals: Selected Topics. {\newblock}\tmtextit{Progress
  of Theoretical Physics Supplement}, 111:83--113, jan 1993.
  {\newblock}\href{Https://doi.org/10.1143/PTPS.111.83}{}10.1143/PTPS.111.83.{\newblock}
  
  \bibitem[31]{biane_calcul_1995}Philippe Biane. {\newblock}Calcul
  stochastique non-commutatif. {\newblock}In  Pierre Bernard, editor,
  \tmtextit{Lectures on Probability Theory},  volume  1608,  pages  1--96.
  Springer Berlin Heidelberg, Berlin, Heidelberg, 1995.
  {\newblock}\href{Https://doi.org/10.1007/BFb0095746}{}10.1007/BFb0095746.{\newblock}
  
  \bibitem[32]{bogolubov_general_1990}N.~N.~Bogolubov, Anatoly~A.~Logunov,
  A.~I.~Oksak, and  I.~Todorov. {\newblock}\tmtextit{General Principles of
  Quantum Field Theory}. {\newblock}Mathematical Physics and Applied
  Mathematics. Springer Netherlands, 1990.{\newblock}
  
  \bibitem[33]{bonsall_complete_1973}Frank~F.~Bonsall  and  John Duncan.
  {\newblock}\tmtextit{Complete normed algebras}. {\newblock}Springer-Verlag,
  New York-Heidelberg, 1973.{\newblock}
  
  \bibitem[34]{borasi}Luigi~Marcello Borasi.
  {\newblock}\tmtextit{Probabilistic and differential geometric methods for
  relativistic and Euclidean Dirac and radiation fields}. {\newblock}PhD
  thesis, University of Bonn, Bonn, Germany, jul 2019.{\newblock}
  
  \bibitem[35]{bozejko_interpolations_1996}Marek Bo{\.z}ejko  and  Roland
  Speicher. {\newblock}Interpolations between bosonic and fermionic relations
  given by generalized brownian motions. {\newblock}\tmtextit{Mathematische
  Zeitschrift}, 222(1):135--160, may 1996.
  {\newblock}\href{Https://doi.org/10.1007/BF02621861}{}10.1007/BF02621861.{\newblock}
  
  \bibitem[36]{Bratteli1}Ola Bratteli  and  Derek~W.~Robinson.
  {\newblock}\tmtextit{Operator algebras and quantum statistical mechanics.
  1}. {\newblock}Texts and Monographs in Physics. Springer-Verlag, New York,
  Second  edition, 1987.
  {\newblock}\href{Https://doi.org/10.1007/978-3-662-02520-8}{}10.1007/978-3-662-02520-8.{\newblock}
  
  \bibitem[37]{Bratteli2}Ola Bratteli  and  Derek~W.~Robinson.
  {\newblock}\tmtextit{Operator algebras and quantum statistical mechanics.
  2}. {\newblock}Texts and Monographs in Physics. Springer-Verlag, Berlin,
  Second  edition, 1997.
  {\newblock}\href{Https://doi.org/10.1007/978-3-662-03444-6}{}10.1007/978-3-662-03444-6.{\newblock}
  
  \bibitem[38]{Browder1962}Felix~E.~Browder. {\newblock}Analyticity and
  partial differential equations. I. {\newblock}\tmtextit{Amer. J. Math.},
  84:666--710, 1962.{\newblock}
  
  \bibitem[39]{caianiello_combinatorics_1973}Eduardo~R.~Caianiello.
  {\newblock}\tmtextit{Combinatorics and Renormalization in Quantum Field
  Theory},  volume~38. {\newblock}Benjamin, 1973.{\newblock}
  
  \bibitem[40]{Carlen1}Eric Carlen  and  Paul Kr{\'e}e. {\newblock}$L^p$
  estimates on iterated stochastic integrals. {\newblock}\tmtextit{Ann.
  Probab.}, 19(1):354--368, 1991.{\newblock}
  
  \bibitem[41]{catellier_paracontrolled_2013}R{\'e}mi Catellier  and  Khalil
  Chouk. {\newblock}Paracontrolled distributions and the 3-dimensional
  stochastic quantization equation. {\newblock}\tmtextit{The Annals of
  Probability}, 46(5):2621--2679, 2018.
  {\newblock}\href{Https://doi.org/10.1214/17-AOP1235}{}10.1214/17-AOP1235.{\newblock}
  
  \bibitem[42]{Cipriani}Fabio Cipriani. {\newblock}Noncommutative potential
  theory: a survey. {\newblock}\tmtextit{J. Geom. Phys.}, 105:25--59, 2016.
  {\newblock}\href{Https://doi.org/10.1016/j.geomphys.2016.03.016}{}10.1016/j.geomphys.2016.03.016.{\newblock}
  
  \bibitem[43]{connes_noncommutative_1994}Alain Connes.
  {\newblock}\tmtextit{Noncommutative geometry}. {\newblock}Academic Press,
  San Diego, 1994.{\newblock}
  
  \bibitem[44]{cooper_weakly_1977}Alan Cooper  and  Lon Rosen. {\newblock}The
  Weakly Coupled Yukawa\$\_2\$ Field Theory: Cluster Expansion and Wightman
  Axioms. {\newblock}\tmtextit{Transactions of the American Mathematical
  Society}, 234(1):1--88, 1977.
  {\newblock}\href{Https://doi.org/10.2307/1997994}{}10.2307/1997994.{\newblock}
  
  \bibitem[45]{corry_modern_2004}Leo Corry. {\newblock}\tmtextit{Modern
  Algebra and the Rise of Mathematical Structures}. {\newblock}Birkh{\"a}user
  Basel, 2  edition, 2004.
  {\newblock}\href{Https://doi.org/10.1007/978-3-0348-7917-0}{}10.1007/978-3-0348-7917-0.{\newblock}
  
  \bibitem[46]{damgaard_stochastic_1984}P.H.~Damgaard  and  K.~Tsokos.
  {\newblock}Stochastic quantization with fermions.
  {\newblock}\tmtextit{Nuclear Physics B}, 235(1):75--92, may 1984.
  {\newblock}\href{Https://doi.org/10.1016/0550-3213(84)90149-4}{}10.1016/0550-3213(84)90149-4.{\newblock}
  
  \bibitem[47]{damgaard_stochastic_1988}Poul~Henrik Damgaard  and  Helmuth
  H{\"u}ffel. {\newblock}\tmtextit{Stochastic Quantization}. {\newblock}World
  Scientific, 1988.{\newblock}
  
  \bibitem[48]{de_vecchi_note_2019}Francesco~C.~De Vecchi  and  Massimiliano
  Gubinelli. {\newblock}A note on supersymmetry and stochastic differential
  equations. {\newblock}\tmtextit{ArXiv:1912.04830 [math-ph]}, dec 2019.
  {\newblock}ArXiv: 1912.04830.{\newblock}
  
  \bibitem[49]{DeWitt}Bryce DeWitt. {\newblock}\tmtextit{Supermanifolds}.
  {\newblock}Cambridge Monographs on Mathematical Physics. Cambridge
  University Press, Cambridge, Second  edition, 1992.
  {\newblock}\href{Https://doi.org/10.1017/CBO9780511564000}{}10.1017/CBO9780511564000.{\newblock}
  
  \bibitem[50]{dimock_estimates_1972}Jonathan Dimock. {\newblock}Estimates,
  renormalized currents, and field equations for the Yukawa\tmrsub{$2$} field
  theory. {\newblock}\tmtextit{Annals of Physics}, 72(1):177--242, jul 1972.
  {\newblock}\href{Https://doi.org/10.1016/0003-4916(72)90241-2}{}10.1016/0003-4916(72)90241-2.{\newblock}
  
  \bibitem[51]{dirksen_noncommutative_2010}Sjoerd Dirksen.
  {\newblock}Noncommutative stochastic integration through decoupling.
  {\newblock}\tmtextit{Journal of Mathematical Analysis and Applications},
  370(1):200--223, oct 2010.
  {\newblock}\href{Https://doi.org/10.1016/j.jmaa.2010.04.062}{}10.1016/j.jmaa.2010.04.062.{\newblock}
  
  \bibitem[52]{disertori_continuous_1998}M.~Disertori  and  V.~Rivasseau.
  {\newblock}Continuous Constructive Fermionic Renormalization.
  {\newblock}\tmtextit{Annales Henri Poincar{\'e}}, 1(1):1--57, feb 2000.
  {\newblock}\href{Https://doi.org/10.1007/PL00000998}{}10.1007/PL00000998.{\newblock}
  
  \bibitem[53]{efremov_stochastic_2019}A.~N.~Efremov. {\newblock}Stochastic
  Quantization of Massive Fermions. {\newblock}\tmtextit{International Journal
  of Theoretical Physics}, 58(4):1150--1156, apr 2019.
  {\newblock}\href{Https://doi.org/10.1007/s10773-019-04006-w}{}10.1007/s10773-019-04006-w.{\newblock}
  
  \bibitem[54]{Emch}G{\'e}rard~G.~Emch. {\newblock}\tmtextit{Mathematical and
  conceptual foundations of 20th-century physics},  volume  100  of
  \tmtextit{North-Holland Mathematics Studies}. {\newblock}North-Holland
  Publishing Co., Amsterdam, 1984. {\newblock}Notas de Matem{\'a}tica
  [Mathematical Notes], 100.{\newblock}
  
  \bibitem[55]{Federbush1971}Paul Federbush  and  Basilis Gidas.
  {\newblock}Renormalization of the one-space dimensional Yukawa model by
  unitary transformation. {\newblock}\tmtextit{Ann. Physics}, 68:98--101,
  1971.{\newblock}
  
  \bibitem[56]{MR1905424}Joel Feldman, Horst Kn{\"o}rrer, and  Eugene
  Trubowitz. {\newblock}\tmtextit{Fermionic functional integrals and the
  renormalization group},  volume~16  of \tmtextit{CRM Monograph Series}.
  {\newblock}American Mathematical Society, Providence, RI, 2002.{\newblock}
  
  \bibitem[57]{feldman_renormalizable_1986}Joel Feldman, J.~Magnen,
  V.~Rivasseau, and  R.~S{\'e}n{\'e}or. {\newblock}A renormalizable field
  theory: The massive Gross-Neveu model in two dimensions.
  {\newblock}\tmtextit{Commun.Math. Phys.}, 103(1):67--103, mar 1986.
  {\newblock}\href{Https://doi.org/10.1007/BF01464282}{}10.1007/BF01464282.{\newblock}
  
  \bibitem[58]{friedrichs_mathematical_1953}Kurt~Otto Friedrichs.
  {\newblock}\tmtextit{Mathematical Aspects of the Quantum Theory of Fields}.
  {\newblock}Interscience Publishers, 1953.{\newblock}
  
  \bibitem[59]{frohlich_is_1975}J{\"u}rg Fr{\"o}hlich  and  Konrad
  Osterwalder. {\newblock}Is There a Euclidean Field Theory for Fermions.
  {\newblock}\tmtextit{Helv. Phys. Acta}, 47:781, 1975.{\newblock}
  
  \bibitem[60]{fukai_stochastic_1983}T.~Fukai, H.~Nakazato, I.~Ohba, K.~Okano,
  and  Y.~Yamanaka. {\newblock}Stochastic Quantization Method of Fermion
  Fields. {\newblock}\tmtextit{Progress of Theoretical Physics},
  69(5):1600--1616, may 1983.
  {\newblock}\href{Https://doi.org/10.1143/PTP.69.1600}{}10.1143/PTP.69.1600.{\newblock}
  
  \bibitem[61]{gawedzki_gross_neveu_1985}K.~Gaw{\k e}dzki  and  A.~Kupiainen.
  {\newblock}Gross-Neveu model through convergent perturbation expansions.
  {\newblock}\tmtextit{Communications in Mathematical Physics}, 102(1):1--30,
  1985.{\newblock}
  
  \bibitem[62]{Gelfand}I.~Gelfand  and  M.~Neumark. {\newblock}On the
  imbedding of normed rings into the ring of operators in Hilbert space.
  {\newblock}\tmtextit{Rec. Math. [Mat. Sbornik] N.S.}, 12(54):197--213,
  1943.{\newblock}
  
  \bibitem[63]{Jaffe1971}James Glimm  and  Arthur Jaffe. {\newblock}The
  \text{Yukawa}\tmrsub{$2$} quantum field theory without cutoffs.
  {\newblock}\tmtextit{J. Functional Analysis}, 7:323--357, 1971.{\newblock}
  
  \bibitem[64]{glimm_quantum_1985}James Glimm  and  Arthur Jaffe.
  {\newblock}Quantum Field Theory Models. {\newblock}In  James Glimm  and 
  Arthur Jaffe, editors, \tmtextit{Quantum Field Theory and Statistical
  Mechanics: Expositions},  pages  11--121. Birkh{\"a}user, Boston, MA, 1985.
  {\newblock}\href{Https://doi.org/10.1007/978-1-4612-5158-3_4}{}10.1007/978-1-4612-5158-3\tmrsub{4}.{\newblock}
  
  \bibitem[65]{glimm_quantum_1987}James Glimm  and  Arthur Jaffe.
  {\newblock}\tmtextit{Quantum physics}. {\newblock}Springer-Verlag, New York,
  Second  edition, 1987. {\newblock}A functional integral point of
  view.{\newblock}
  
  \bibitem[66]{Gordina}Maria Gordina. {\newblock}Stochastic differential
  equations on noncommutative $L^2$. {\newblock}In \tmtextit{Finite and
  infinite dimensional analysis in honor of Leonard Gross (New Orleans, LA,
  2001)},  volume  317  of \tmtextit{Contemp. Math.},  pages  87--98. Amer.
  Math. Soc., Providence, RI, 2003.
  {\newblock}\href{Https://doi.org/10.1090/conm/317/05521}{}10.1090/conm/317/05521.{\newblock}
  
  \bibitem[67]{Gross1967}Leonard Gross. {\newblock}Abstract Wiener spaces.
  {\newblock}In \tmtextit{Proc. Fifth Berkeley Sympos. Math. Statist. and
  Probability (Berkeley, Calif., 1965/66), Vol. II: Contributions to
  Probability Theory, Part 1},  pages  31--42. Univ. California Press,
  Berkeley, Calif., 1967.{\newblock}
  
  \bibitem[68]{gross_existence_1972}Leonard Gross. {\newblock}Existence and
  uniqueness of physical ground states. {\newblock}\tmtextit{Journal of
  Functional Analysis}, 10(1):52--109, may 1972.
  {\newblock}\href{Https://doi.org/10.1016/0022-1236(72)90057-2}{}10.1016/0022-1236(72)90057-2.{\newblock}
  
  \bibitem[69]{gross_hypercontractivity_1975}Leonard Gross.
  {\newblock}Hypercontractivity and logarithmic Sobolev inequalities for the
  Clifford-Dirichlet form. {\newblock}\tmtextit{Duke Mathematical Journal},
  42(3):383--396, sep 1975.
  {\newblock}\href{Https://doi.org/10.1215/S0012-7094-75-04237-4}{}10.1215/S0012-7094-75-04237-4.{\newblock}
  
  \bibitem[70]{Gross_Euclidean_1976}Leonard Gross. {\newblock}Euclidean
  fermion fields. {\newblock}In \tmtextit{Les m{\'e}thodes math{\'e}matiques
  de la th{\'e}orie quantique des champs (Colloq. Internat. CNRS, No. 248,
  Marseille, 1975)},  pages  131--146. 1976.{\newblock}
  
  \bibitem[71]{gross_formula_1977}Leonard Gross. {\newblock}On the formula of
  Mathews and Salam. {\newblock}\tmtextit{Journal of Functional Analysis},
  25(2):162--209, jun 1977.
  {\newblock}\href{Https://doi.org/10.1016/0022-1236(77)90039-8}{}10.1016/0022-1236(77)90039-8.{\newblock}
  
  \bibitem[72]{gross_noncommutative_1971}Lonard Gross. {\newblock}A
  noncommutative extension of the Perron-Frobenius theorem.
  {\newblock}\tmtextit{Bulletin of the American Mathematical Society},
  77(3):343--348, may 1971.
  {\newblock}\href{Https://doi.org/10.1090/S0002-9904-1971-12686-1}{}10.1090/S0002-9904-1971-12686-1.{\newblock}
  
  \bibitem[73]{gubinelli_pde_2018}Massimiliano Gubinelli  and  Martina
  Hofmanov{\'a}. {\newblock}A PDE construction of the Euclidean $\Phi^4_3$
  quantum field theory. {\newblock}\tmtextit{ArXiv:1810.01700 [math-ph]}, oct
  2018. {\newblock}ArXiv: 1810.01700.{\newblock}
  
  \bibitem[74]{gubinelli_global_2019}Massimiliano Gubinelli  and  Martina
  Hofmanov{\'a}. {\newblock}Global Solutions to Elliptic and Parabolic
  $\Phi^4$ Models in Euclidean Space. {\newblock}\tmtextit{Communications in
  Mathematical Physics}, 368(3):1201--1266, 2019.{\newblock}
  
  \bibitem[75]{gubinelli_paracontrolled_2015}Massimiliano Gubinelli, Peter
  Imkeller, and  Nicolas Perkowski. {\newblock}Paracontrolled distributions
  and singular PDEs. {\newblock}\tmtextit{Forum of Mathematics. Pi}, 3:0,
  2015.
  {\newblock}\href{Https://doi.org/10.1017/fmp.2015.2}{}10.1017/fmp.2015.2.{\newblock}
  
  \bibitem[76]{gubinelli_paracontrolled_2018}Massimiliano Gubinelli, Herbert
  Koch, and  Tadahiro Oh. {\newblock}Paracontrolled approach to the
  three-dimensional stochastic nonlinear wave equation with quadratic
  nonlinearity. {\newblock}\tmtextit{ArXiv:1811.07808 [math]}, nov 2018.
  {\newblock}ArXiv: 1811.07808.{\newblock}
  
  \bibitem[77]{gubinelli_renormalization_2018}Massimiliano Gubinelli, Herbert
  Koch, and  Tadahiro Oh. {\newblock}Renormalization of the two-dimensional
  stochastic nonlinear wave equations. {\newblock}\tmtextit{Transactions of
  the American Mathematical Society},  page  1, 2018.
  {\newblock}\href{Https://doi.org/10.1090/tran/7452}{}10.1090/tran/7452.{\newblock}
  
  \bibitem[78]{MR0452281}Francesco Guerra. {\newblock}Local algebras in
  Euclidean quantum field theory. {\newblock}In \tmtextit{Symposia
  Mathematica, Vol. XX (Convegno sulle Algebre C* e loro Applicazioni in
  Fisica Teorica, Convegno sulla Teoria degli Operatori Indice e Teoria K,
  INDAM, Roma, 1974)},  pages  13--26. 1976.{\newblock}
  
  \bibitem[79]{Haag1}Rudolf Haag. {\newblock}\tmtextit{Local quantum physics}.
  {\newblock}Texts and Monographs in Physics. Springer-Verlag, Berlin, Second 
  edition, 1996.
  {\newblock}\href{Https://doi.org/10.1007/978-3-642-61458-3}{}10.1007/978-3-642-61458-3.{\newblock}
  
  \bibitem[80]{haba_supersymmetry_1995}Z.~Haba  and  J.~Kupsch.
  {\newblock}Supersymmetry in Euclidean Quantum Field Theory.
  {\newblock}\tmtextit{Fortschritte der Physik/Progress of Physics},
  43(1):41--66, 1995.
  {\newblock}\href{Https://doi.org/10.1002/prop.2190430103}{}10.1002/prop.2190430103.{\newblock}
  
  \bibitem[81]{hairer_theory_2014}M.~Hairer. {\newblock}A theory of regularity
  structures. {\newblock}\tmtextit{Inventiones mathematicae}, 198(2):269--504,
  2014.
  {\newblock}\href{Https://doi.org/10.1007/s00222-014-0505-4}{}10.1007/s00222-014-0505-4.{\newblock}
  
  \bibitem[82]{MR0277208}K.~Hepp. {\newblock}\tmtextit{Th{\'e}orie de la
  renormalisation}. {\newblock}Cours donn{\'e} {\`a} l'{\'E}cole
  Polytechnique, Paris. Lecture Notes in Physics, Vol. 2. Springer-Verlag,
  Berlin-New York, 1969.{\newblock}
  
  \bibitem[83]{hida_white_2013}Takeyuki Hida, Hui-Hsiung Kuo, J{\"u}rgen
  Potthoff, and  Walter Streit. {\newblock}\tmtextit{White Noise: An Infinite
  Dimensional Calculus}. {\newblock}Springer Science \& Business Media, jun
  2013.{\newblock}
  
  \bibitem[84]{Holevo}Alexander Holevo. {\newblock}\tmtextit{Probabilistic and
  statistical aspects of quantum theory},  volume~1  of
  \tmtextit{Quaderni/Monographs}. {\newblock}Edizioni della Normale, Pisa,
  Second  edition, 2011.
  {\newblock}\href{Https://doi.org/10.1007/978-88-7642-378-9}{}10.1007/978-88-7642-378-9.{\newblock}
  
  \bibitem[85]{hsu_stochastic_2002}Elton~P.~Hsu.
  {\newblock}\tmtextit{Stochastic Analysis on Manifolds}. {\newblock}Graduate
  Studies in Mathematics 38. American Mathematical Society, 2002.{\newblock}
  
  \bibitem[86]{hudson_unification_1986}R.~L.~Hudson  and  K.~R.~Parthasarathy.
  {\newblock}Unification of fermion and Boson stochastic calculus.
  {\newblock}\tmtextit{Communications In Mathematical Physics},
  104(3):457--470, sep 1986.
  {\newblock}\href{Https://doi.org/10.1007/BF01210951}{}10.1007/BF01210951.{\newblock}
  
  \bibitem[87]{ignatyuk_convergence_1993}I.~A.~Ignatyuk, V.~A.~Malyshev, and 
  V.~Sidoravi{\v c}ius. {\newblock}Convergence of a Method of the Stochastic
  Quantization II. {\newblock}\tmtextit{Theory of Probability \& Its
  Applications}, 37(4):599--620, jan 1993.
  {\newblock}\href{Https://doi.org/10.1137/1137117}{}10.1137/1137117.{\newblock}
  
  \bibitem[88]{Itzykson}Claude Itzykson  and  Jean~Bernard Zuber.
  {\newblock}\tmtextit{Quantum field theory}. {\newblock}McGraw-Hill
  International Book Co., New York, 1980. {\newblock}International Series in
  Pure and Applied Physics.{\newblock}
  
  \bibitem[89]{ivashchuk_infinite_dimensional_2000}V.~D.~Ivashchuk.
  {\newblock}Infinite-dimensional Grassmann-Banach algebras.
  {\newblock}\tmtextit{ArXiv:math-ph/0009006}, sep 2000. {\newblock}ArXiv:
  math-ph/0009006.{\newblock}
  
  \bibitem[90]{Jaffe_supersymmetry}Arthur Jaffe, Andrzej Lesniewski, and 
  Jonathan Weitsman. {\newblock}The two-dimensional, $N = 2$ Wess-Zumino model
  on a cylinder. {\newblock}\tmtextit{Comm. Math. Phys.}, 114(1):147--165,
  1988.{\newblock}
  
  \bibitem[91]{janson_gaussian_1997}Svante Janson.
  {\newblock}\tmtextit{Gaussian Hilbert Spaces}. {\newblock}Cambridge
  University Press, jun 1997.{\newblock}
  
  \bibitem[92]{jona_lasinio_stochastic_1987}G.~Jona-Lasinio.
  {\newblock}Stochastic quantization: a new domain for stochastic analysis.
  {\newblock}In \tmtextit{Proceedings of the 1st World Congress of the
  Bernoulli Society, Vol. 1 (Tashkent, 1986)},  pages  535--546. VNU Sci.
  Press, Utrecht, 1987.{\newblock}
  
  \bibitem[93]{jona_lasinio_stochastic_1985}G.~Jona-Lasinio  and 
  P.~K.~Mitter. {\newblock}On the stochastic quantization of field theory.
  {\newblock}\tmtextit{Communications in Mathematical Physics (1965-1997)},
  101(3):409--436, 1985.{\newblock}
  
  \bibitem[94]{jona_lasinio_large_1990}G.~Jona-Lasinio  and  P.~K.~Mitter.
  {\newblock}Large deviation estimates in the stochastic quantization of
  $\phi^4_2$. {\newblock}\tmtextit{Communications in Mathematical Physics},
  130(1):111--121, may 1990.
  {\newblock}\href{Https://doi.org/10.1007/BF02099877}{}10.1007/BF02099877.{\newblock}
  
  \bibitem[95]{jost_general_1965}Res Jost. {\newblock}\tmtextit{The general
  theory of quantized fields},  volume  1960  of \tmtextit{Mark Kac, editor.
  Lectures in Applied Mathematics (Proceedings of the Summer Seminar, Boulder,
  Colorado}. {\newblock}American Mathematical Society, Providence, R.I.,
  1965.{\newblock}
  
  \bibitem[96]{Kadison2}Richard~V.~Kadison  and  John~R.~Ringrose.
  {\newblock}\tmtextit{Fundamentals of the theory of operator algebras. Vol.
  I},  volume~15  of \tmtextit{Graduate Studies in Mathematics}.
  {\newblock}American Mathematical Society, Providence, RI, 1997.
  {\newblock}Elementary theory, Reprint of the 1983 original.{\newblock}
  
  \bibitem[97]{Kadison1}Richard~V.~Kadison  and  John~R.~Ringrose.
  {\newblock}\tmtextit{Fundamentals of the theory of operator algebras. Vol.
  II},  volume~16  of \tmtextit{Graduate Studies in Mathematics}.
  {\newblock}American Mathematical Society, Providence, RI, 1997.
  {\newblock}Advanced theory, Corrected reprint of the 1986 original.
  \href{https://doi.org/10.1090/gsm/016/01}{}10.1090/gsm/016/01.{\newblock}
  
  \bibitem[98]{kakudo1983gauge}Yujiro Kakudo, Yukio Taguchi, Azuma Tanaka, and
  Kunio Yamamoto. {\newblock}Gauge-independent calculation of S-matrix
  elements in quantum electrodynamics. {\newblock}\tmtextit{Progress of
  Theoretical Physics (Kyoto)}, 69(4):1225--1235, 1983.{\newblock}
  
  \bibitem[99]{khrennikov_equations_1991}A.~Yu Khrennikov.
  {\newblock}Equations on a superspace. {\newblock}\tmtextit{Mathematics of
  the USSR-Izvestiya}, 36(3):597, 1991. {\newblock}Publisher: IOP Publishing.
  \href{https://doi.org/10.1070/IM1991v036n03ABEH002036}{}10.1070/IM1991v036n03ABEH002036.{\newblock}
  
  \bibitem[100]{klein_supersymmetry_1985}Abel Klein. {\newblock}Supersymmetry
  and a two-dimensional reduction in random phenomena. {\newblock}In
  \tmtextit{Quantum probability and applications, II (Heidelberg, 1984)}, 
  volume  1136  of \tmtextit{Lecture Notes in Math.},  pages  306--317.
  Springer, Berlin, 1985.
  {\newblock}\href{Https://doi.org/10.1007/BFb0074481}{}10.1007/BFb0074481.{\newblock}
  
  \bibitem[101]{klein_supersymmetry_1983}Abel Klein  and  J.~Fernando Perez.
  {\newblock}Supersymmetry and dimensional reduction: A non-perturbative
  proof. {\newblock}\tmtextit{Physics Letters B}, 125(6):473--475, jun 1983.
  {\newblock}\href{Https://doi.org/10.1016/0370-2693(83)91329-1}{}10.1016/0370-2693(83)91329-1.{\newblock}
  
  \bibitem[102]{Kummer}Burkhard K{\"u}mmerer. {\newblock}Survey on a theory of
  noncommutative stationary Markov processes. {\newblock}In \tmtextit{Quantum
  probability and applications, III (Oberwolfach, 1987)},  volume  1303  of
  \tmtextit{Lecture Notes in Math.},  pages  154--182. Springer, Berlin, 1988.
  {\newblock}\href{Https://doi.org/10.1007/BFb0078061}{}10.1007/BFb0078061.{\newblock}
  
  \bibitem[103]{kupiainen_renormalization_2016}A.~Kupiainen.
  {\newblock}Renormalization Group and Stochastic PDEs.
  {\newblock}\tmtextit{Annales Henri Poincar{\'e}}, 17(3):497--535, 2016.
  {\newblock}\href{Https://doi.org/10.1007/s00023-015-0408-y}{}10.1007/s00023-015-0408-y.{\newblock}
  
  \bibitem[104]{kupsch_fermionic_1993}J.~Kupsch. {\newblock}Fermionic and
  supersymmetric stochastic processes. {\newblock}\tmtextit{Journal of
  Geometry and Physics}, 11(1):507--516, jun 1993.
  {\newblock}\href{Https://doi.org/10.1016/0393-0440(93)90074-O}{}10.1016/0393-0440(93)90074-O.{\newblock}
  
  \bibitem[105]{lanford_iii_construction_1966}O.~E.~Lanford III.
  {\newblock}\tmtextit{Construction of quantum fields interacting by a cutoff
  Yukawa coupling.} {\newblock}PhD thesis, Princeton university, jan
  1966.{\newblock}
  
  \bibitem[106]{le_jan_temps_1987}Y.~Le Jan. {\newblock}Temps local et
  superchamp. {\newblock}In \tmtextit{S{\'e}minaire de Probabilit{\'e}s, XXI},
  volume  1247  of \tmtextit{Lecture Notes in Math.},  pages  176--190.
  Springer, Berlin, 1987.
  {\newblock}\href{Https://doi.org/10.1007/BFb0077633}{}10.1007/BFb0077633.{\newblock}
  
  \bibitem[107]{lehmann_probabilistic_1991}Detlef Lehmann. {\newblock}A
  probabilistic approach to Euclidean Dirac fields.
  {\newblock}\tmtextit{Journal of Mathematical Physics}, 32(8):2158--2166, aug
  1991. {\newblock}Publisher: American Institute of Physics.
  \href{https://doi.org/10.1063/1.529189}{}10.1063/1.529189.{\newblock}
  
  \bibitem[108]{Leites}D.~A.~Leites. {\newblock}Introduction to the theory of
  supermanifolds. {\newblock}\tmtextit{Uspekhi Mat. Nauk}, 35(1(211)):3--57,
  1980.{\newblock}
  
  \bibitem[109]{leppard_feynman_kac_2001}Steven Leppard  and  Alice Rogers.
  {\newblock}A Feynman-Kac formula for anticommuting Brownian motion.
  {\newblock}\tmtextit{Journal of Physics A: Mathematical and General},
  34(3):555--568, jan 2001.
  {\newblock}\href{Https://doi.org/10.1088/0305-4470/34/3/315}{}10.1088/0305-4470/34/3/315.{\newblock}
  
  \bibitem[110]{lesniewski_effective_1987}A.~Lesniewski. {\newblock}Effective
  action for the Yukawa\tmrsub{$2$} quantum field theory.
  {\newblock}\tmtextit{Communications in Mathematical Physics},
  108(3):437--467, sep 1987.
  {\newblock}\href{Https://doi.org/10.1007/BF01212319}{}10.1007/BF01212319.{\newblock}
  
  \bibitem[111]{Magnen1976}J.~Magnen  and  R.~S{\'e}n{\'e}or. {\newblock}The
  Wightman axioms for the weakly coupled Yukawa model in two dimensions.
  {\newblock}\tmtextit{Comm. Math. Phys.}, 51(3):297--313, 1976.{\newblock}
  
  \bibitem[112]{Magnen1978}J.~Magnen  and  R.~S{\'e}n{\'e}or.
  {\newblock}Yukawa quantum field theory in three dimensions $(Y_3)$.
  {\newblock}In \tmtextit{Third International Conference on Collective
  Phenomena (Moscow, 1978)},  volume  337  of \tmtextit{Ann. New York Acad.
  Sci.},  pages  13--43. New York Acad. Sci., New York, 1980.{\newblock}
  
  \bibitem[113]{Markushevich1965}A.~I.~Markushevich.
  {\newblock}\tmtextit{Theory of functions of a complex variable. Vol. II}.
  {\newblock}Revised English edition translated and edited by Richard A.
  Silverman. Prentice-Hall, Inc., Englewood Cliffs, N.J., 1965.{\newblock}
  
  \bibitem[114]{mastropietro_non-perturbative_2008}Vieri Mastropietro.
  {\newblock}\tmtextit{Non-Perturbative Renormalization}. {\newblock}World
  Scientific Publishing Co Pte Ltd, Hackensack, NJ, 2008.{\newblock}
  
  \bibitem[115]{matthews_propagators_1955}P.~T.~Matthews  and  A.~Salam.
  {\newblock}Propagators of quantized field. {\newblock}\tmtextit{Il Nuovo
  Cimento (1955-1965)}, 2(1):120--134, jul 1955.
  {\newblock}\href{Https://doi.org/10.1007/BF02856011}{}10.1007/BF02856011.{\newblock}
  
  \bibitem[116]{mcbryan_volume_1975}Oliver~A.~McBryan. {\newblock}Volume
  dependence of Schwinger function in the \text{Yukawa}\tmrsub{$2$} quantum
  field theory. {\newblock}\tmtextit{Communications in Mathematical Physics},
  45(3):279--294, 1975. {\newblock}Publisher: Springer-Verlag.{\newblock}
  
  \bibitem[117]{McBryan1975}Oliver~A.~McBryan  and  Yong~Moon Park.
  {\newblock}Lorentz covariance of the Yukawa\tmrsub{$2$} quantum field
  theory. {\newblock}\tmtextit{J. Mathematical Phys.}, 16:104--110,
  1975.{\newblock}
  
  \bibitem[118]{Meyerbook}Paul-Andr{\'e} Meyer. {\newblock}\tmtextit{Quantum
  probability for probabilists},  volume  1538  of \tmtextit{Lecture Notes in
  Mathematics}. {\newblock}Springer-Verlag, Berlin, 1993.
  {\newblock}\href{Https://doi.org/10.1007/978-3-662-21558-6}{}10.1007/978-3-662-21558-6.{\newblock}
  
  \bibitem[119]{moinat_space_time_2018}A.~Moinat  and  H.~Weber.
  {\newblock}Space-time localisation for the dynamic $\Phi^4_3$ model.
  {\newblock}\tmtextit{ArXiv:1811.05764}, nov 2018. {\newblock}ArXiv:
  1811.05764.{\newblock}
  
  \bibitem[120]{mourrat_plane_2015}J.-C.~Mourrat  and  H.~Weber.
  {\newblock}Global well-posedness of the dynamic $\Phi^4$ model in the plane.
  {\newblock}\tmtextit{The Annals of Probability}, 45(4):2398--2476, 2017.
  {\newblock}\href{Https://doi.org/10.1214/16-AOP1116}{}10.1214/16-AOP1116.{\newblock}
  
  \bibitem[121]{MWcomedown}J.-C.~Mourrat  and  H.~Weber. {\newblock}The
  dynamic $\Phi^4_3$ model comes down from infinity.
  {\newblock}\tmtextit{Comm. Math. Phys.}, 356(3):673--753, 2017.
  {\newblock}\href{Https://doi.org/10.1007/s00220-017-2997-4}{}10.1007/s00220-017-2997-4.{\newblock}
  
  \bibitem[122]{Naimark}M.~A.~Naimark. {\newblock}\tmtextit{Normed algebras}.
  {\newblock}Wolters-Noordhoff Publishing, Groningen, Third  edition, 1972.
  {\newblock}Translated from the second Russian edition by Leo F. Boron,
  Wolters-Noordhoff Series of Monographs and Textbooks on Pure and Applied
  Mathematics.{\newblock}
  
  \bibitem[123]{nakano_quantum_1959}Tadao Nakano. {\newblock}Quantum Field
  Theory in Terms of Euclidean Parameters. {\newblock}\tmtextit{Progress of
  Theoretical Physics}, 21(2):241--259, feb 1959. {\newblock}Publisher: Oxford
  Academic.
  \href{https://doi.org/10.1143/PTP.21.241}{}10.1143/PTP.21.241.{\newblock}
  
  \bibitem[124]{nelson_construction_1973}Edward Nelson.
  {\newblock}Construction of quantum fields from Markoff fields.
  {\newblock}\tmtextit{Journal of Functional Analysis}, 12(1):97--112, jan
  1973.
  {\newblock}\href{Https://doi.org/10.1016/0022-1236(73)90091-8}{}10.1016/0022-1236(73)90091-8.{\newblock}
  
  \bibitem[125]{nelson_notes_1974}Edward Nelson. {\newblock}Notes on
  non-commutative integration. {\newblock}\tmtextit{Journal of Functional
  Analysis}, 15(2):103--116, feb 1974.
  {\newblock}\href{Https://doi.org/10.1016/0022-1236(74)90014-7}{}10.1016/0022-1236(74)90014-7.{\newblock}
  
  \bibitem[126]{osipov_quantum_2000}Edward~P.~Osipov. {\newblock}Quantum
  interaction $\phi^4_4$, the construction of quantum field defined as a
  bilinear form. {\newblock}\tmtextit{Journal of Mathematical Physics},
  41(2):759--786, 2000.
  {\newblock}\href{Https://doi.org/10.1063/1.533162}{}10.1063/1.533162.{\newblock}
  
  \bibitem[127]{osterwalder_euclidean_1973}Konrad Osterwalder.
  {\newblock}Euclidean fermi fields. {\newblock}In \tmtextit{Constructive
  Quantum Field Theory}, Lecture Notes in Physics,  pages  326--331. Springer,
  Berlin, Heidelberg, 1973.
  {\newblock}\href{Https://doi.org/10.1007/BFb0113094}{}10.1007/BFb0113094.{\newblock}
  
  \bibitem[128]{osterwalder_euclidean_1973-2}Konrad Osterwalder  and  Robert
  Schrader. {\newblock}Euclidean Fermi fields and a Feynman-Kac formula for
  Boson-Fermions models. {\newblock}\tmtextit{Helvetica Physica Acta},
  46:277--302, 1973.{\newblock}
  
  \bibitem[129]{Pachpatte1998}B.~G.~Pachpatte.
  {\newblock}\tmtextit{Inequalities for differential and integral equations}, 
  volume  197  of \tmtextit{Mathematics in Science and Engineering}.
  {\newblock}Academic Press, Inc., San Diego, CA, 1998.{\newblock}
  
  \bibitem[130]{palmer_euclidean_1980}John Palmer. {\newblock}Euclidean Fermi
  fields. {\newblock}\tmtextit{Journal of Functional Analysis},
  36(3):287--312, may 1980.
  {\newblock}\href{Https://doi.org/10.1016/0022-1236(80)90092-0}{}10.1016/0022-1236(80)90092-0.{\newblock}
  
  \bibitem[131]{parisi_perturbation_1981}G.~Parisi  and  Yong~Shi Wu.
  {\newblock}Perturbation theory without gauge fixing.
  {\newblock}\tmtextit{Scientia Sinica. Zhongguo Kexue}, 24(4):483--496,
  1981.{\newblock}
  
  \bibitem[132]{parthasarathy_introduction_1992}K.~R.~Parthasarathy.
  {\newblock}\tmtextit{An Introduction to Quantum Stochastic Calculus}.
  {\newblock}Springer Science \& Business Media, 1992.{\newblock}
  
  \bibitem[133]{pestov_analysis_1994}Vladimir~G.~Pestov. {\newblock}Analysis
  on superspace: an overview. {\newblock}\tmtextit{Bulletin of the Australian
  Mathematical Society}, 50(1):135--165, aug 1994. {\newblock}Publisher:
  Cambridge University Press.
  \href{https://doi.org/10.1017/S0004972700009643}{}10.1017/S0004972700009643.{\newblock}
  
  \bibitem[134]{plymen_spinors_1994}Roger Plymen  and  Paul Robinson.
  {\newblock}\tmtextit{Spinors in Hilbert Space}. {\newblock}Cambridge
  University Press, 1994.{\newblock}
  
  \bibitem[135]{prastaro_geometry_1994}Agostino Prastaro  and 
  Themistocles~M.~Rassias. {\newblock}\tmtextit{Geometry in Partial
  Differential Equations}. {\newblock}World Scientific, 1994.{\newblock}
  
  \bibitem[136]{rivasseau_perturbative_1991}V.~Rivasseau.
  {\newblock}\tmtextit{From Perturbative to Constructive Renormalization}.
  {\newblock}Princeton University Press, Princeton, N.J, 2 edition  edition,
  may 1991.{\newblock}
  
  \bibitem[137]{rogers_stochastic_1992_I}A.~Rogers. {\newblock}Stochastic
  calculus in superspace. I. Supersymmetric Hamiltonians.
  {\newblock}\tmtextit{Journal of Physics A: Mathematical and General},
  25(2):447--468, jan 1992.
  {\newblock}\href{Https://doi.org/10.1088/0305-4470/25/2/024}{}10.1088/0305-4470/25/2/024.{\newblock}
  
  \bibitem[138]{rogers_stochastic_1992_II}A.~Rogers. {\newblock}Stochastic
  calculus in superspace. II. Differential forms, supermanifolds and the
  Atiyah-Singer index theorem. {\newblock}\tmtextit{Journal of Physics A:
  Mathematical and General}, 25(22):6043--6062, nov 1992.
  {\newblock}\href{Https://doi.org/10.1088/0305-4470/25/22/027}{}10.1088/0305-4470/25/22/027.{\newblock}
  
  \bibitem[139]{rogers_graded_1986}Alice Rogers. {\newblock}Graded manifolds,
  supermanifolds and infinite-dimensional Grassmann algebras.
  {\newblock}\tmtextit{Communications in Mathematical Physics},
  105(3):375--384, sep 1986.
  {\newblock}\href{Https://doi.org/10.1007/BF01205932}{}10.1007/BF01205932.{\newblock}
  
  \bibitem[140]{rogers_path_1995}Alice Rogers. {\newblock}Path integration,
  anticommuting variables, and supersymmetry. {\newblock}\tmtextit{Journal of
  Mathematical Physics}, 36(5):2531--2545, may 1995.
  {\newblock}\href{Https://doi.org/10.1063/1.531049}{}10.1063/1.531049.{\newblock}
  
  \bibitem[141]{Rogers2007}Alice Rogers. {\newblock}\tmtextit{Supermanifolds}.
  {\newblock}World Scientific Publishing Co. Pte. Ltd., Hackensack, NJ, 2007.
  {\newblock}\href{Https://doi.org/10.1142/9789812708854}{}10.1142/9789812708854.{\newblock}
  
  \bibitem[142]{rosenberg_noncommutative_2008}Jonathan Rosenberg.
  {\newblock}Noncommutative variations on Laplace's equation.
  {\newblock}\tmtextit{Analysis \& PDE}, 1(1):95--114, oct 2008.
  {\newblock}Publisher: Mathematical Sciences Publishers.
  \href{https://doi.org/10.2140/apde.2008.1.95}{}10.2140/apde.2008.1.95.{\newblock}
  
  \bibitem[143]{Ryan2002}Raymond~A.~Ryan. {\newblock}\tmtextit{Introduction to
  tensor products of Banach spaces}. {\newblock}Springer Monographs in
  Mathematics. Springer-Verlag London, Ltd., London, 2002.{\newblock}
  
  \bibitem[144]{salmhofer_renormalization}Manfred Salmhofer.
  {\newblock}\tmtextit{Renormalization: An Introduction}. {\newblock}Springer,
  Berlin ; New York, 1st Corrected ed. 1999, Corr. 2nd printing 2007 edition 
  edition, sep 2007.{\newblock}
  
  \bibitem[145]{savin_noncommutative_2010}A.~Yu.~Savin  and  B.~Yu.~Sternin.
  {\newblock}Noncommutative elliptic theory. Examples.
  {\newblock}\tmtextit{Proceedings of the Steklov Institute of Mathematics},
  271(1):193--211, dec 2010.
  {\newblock}\href{Https://doi.org/10.1134/S0081543810040152}{}10.1134/S0081543810040152.{\newblock}
  
  \bibitem[146]{Schiff1993}Joel~L.~Schiff. {\newblock}\tmtextit{Normal
  families}. {\newblock}Universitext. Springer-Verlag, New York,
  1993.{\newblock}
  
  \bibitem[147]{SchmeisserTriebel}Hans-J{\"u}rgen Schmeisser  and  Hans
  Triebel. {\newblock}\tmtextit{Topics in Fourier analysis and function
  spaces}. {\newblock}A Wiley-Interscience Publication. John Wiley \& Sons,
  Ltd., Chichester, 1987.{\newblock}
  
  \bibitem[148]{schmudgen_unbounded_1990}Konrad Schm{\"u}dgen.
  {\newblock}\tmtextit{Unbounded operator algebras and representation theory},
  volume~37  of \tmtextit{Operator Theory: Advances and Applications}.
  {\newblock}Birkh{\"a}user Verlag, Basel, 1990.
  {\newblock}\href{Https://doi.org/10.1007/978-3-0348-7469-4}{}10.1007/978-3-0348-7469-4.{\newblock}
  
  \bibitem[149]{schurman_quantum_1991}Michael Sch{\"u}rman. {\newblock}Quantum
  q-white noise and a q-central limit theorem.
  {\newblock}\tmtextit{Communications in Mathematical Physics},
  140(3):589--615, oct 1991.
  {\newblock}\href{Https://doi.org/10.1007/BF02099136}{}10.1007/BF02099136.{\newblock}
  
  \bibitem[150]{schwinger_euclidean_1958}Julian Schwinger. {\newblock}On the
  euclidean structure of relativistic field theory.
  {\newblock}\tmtextit{Proceedings of the National Academy of Sciences of the
  United States of America}, 44(9):956--965, sep 1958.{\newblock}
  
  \bibitem[151]{schwinger_euclidean_1959}Julian Schwinger.
  {\newblock}Euclidean Quantum Electrodynamics. {\newblock}\tmtextit{Physical
  Review}, 115(3):721--731, aug 1959. {\newblock}Publisher: American Physical
  Society.
  \href{https://doi.org/10.1103/PhysRev.115.721}{}10.1103/PhysRev.115.721.{\newblock}
  
  \bibitem[152]{segal_non_commutative_1953}I.~E.~Segal. {\newblock}A
  Non-Commutative Extension of Abstract Integration.
  {\newblock}\tmtextit{Annals of Mathematics}, 57(3):401--457, 1953.
  {\newblock}\href{Https://doi.org/10.2307/1969729}{}10.2307/1969729.{\newblock}
  
  \bibitem[153]{Segal1}I.~E.~Segal. {\newblock}A non-commutative extension of
  abstract integration. {\newblock}\tmtextit{Ann. of Math. (2)}, 57:401--457,
  1953.
  {\newblock}\href{Https://doi.org/10.2307/1969729}{}10.2307/1969729.{\newblock}
  
  \bibitem[154]{segal_tensor_1956}I.~E.~Segal. {\newblock}Tensor Algebras Over
  Hilbert Spaces. II. {\newblock}\tmtextit{Annals of Mathematics},
  63(1):160--175, 1956.
  {\newblock}\href{Https://doi.org/10.2307/1969994}{}10.2307/1969994.{\newblock}
  
  \bibitem[155]{Segalbook2}Irving~E.~Segal. {\newblock}\tmtextit{Mathematical
  problems of relativistic physics},  volume  1960  of \tmtextit{With an
  appendix by George W. Mackey. Lectures in Applied Mathematics (proceedings
  of the Summer Seminar, Boulder, Colorado}. {\newblock}American Mathematical
  Society, Providence, R.I., 1963.{\newblock}
  
  \bibitem[156]{seiler_schwinger_1975}Erhard Seiler. {\newblock}Schwinger
  functions for the Yukawa model in two dimensions with space-time cutoff.
  {\newblock}\tmtextit{Communications in Mathematical Physics},
  42(2):163--182, 1975. {\newblock}Publisher: Springer-Verlag.{\newblock}
  
  \bibitem[157]{seiler_bounds_1975}Erhard Seiler  and  Barry Simon.
  {\newblock}Bounds in the \text{Yukawa}\tmrsub{$2$} quantum field theory:
  upper bound on the pressure, Hamiltonian bound and linear lower bound.
  {\newblock}\tmtextit{Communications in Mathematical Physics}, 45(2):99--114,
  1975. {\newblock}Publisher: Springer-Verlag.{\newblock}
  
  \bibitem[158]{shcherbakov_elements_1993-2}V.~V.~Shcherbakov.
  {\newblock}Elements of stochastic analysis for the case of Grassmann
  variables. I. Grassmann stochastic integrals and random processes.
  {\newblock}\tmtextit{Theoretical and Mathematical Physics}, 96(1):792--800,
  jul 1993.
  {\newblock}\href{Https://doi.org/10.1007/BF01074107}{}10.1007/BF01074107.{\newblock}
  
  \bibitem[159]{shcherbakov_elements_1993-1}V.~V.~Shcherbakov.
  {\newblock}Elements of stochastic analysis for the case of Grassmann
  variables. II. Stochastic partial differential equations for Grassmann
  processes. {\newblock}\tmtextit{Theoretical and Mathematical Physics},
  97(2):1229--1235, nov 1993.
  {\newblock}\href{Https://doi.org/10.1007/BF01016868}{}10.1007/BF01016868.{\newblock}
  
  \bibitem[160]{shcherbakov_elements_1993}V.~V.~Shcherbakov.
  {\newblock}Elements of stochastic analysis for the case of Grassmann
  variables. III. Correlation functions. {\newblock}\tmtextit{Theoretical and
  Mathematical Physics}, 97(3):1323--1332, dec 1993.
  {\newblock}\href{Https://doi.org/10.1007/BF01015761}{}10.1007/BF01015761.{\newblock}
  
  \bibitem[161]{simon_po2_1974}Barry Simon. {\newblock}\tmtextit{$P (\phi)_2$
  Euclidean (Quantum) Field Theory}. {\newblock}Princeton University Press,
  Princeton, N.J, 1974.{\newblock}
  
  \bibitem[162]{Sinha}Kalyan~B.~Sinha  and  Debashish Goswami.
  {\newblock}\tmtextit{Quantum stochastic processes and noncommutative
  geometry},  volume  169  of \tmtextit{Cambridge Tracts in Mathematics}.
  {\newblock}Cambridge University Press, Cambridge, 2007.
  {\newblock}\href{Https://doi.org/10.1017/CBO9780511618529}{}10.1017/CBO9780511618529.{\newblock}
  
  \bibitem[163]{streater_classical_2000}R.~F.~Streater. {\newblock}Classical
  and quantum probability. {\newblock}\tmtextit{Journal of Mathematical
  Physics}, 41(6):3556--3603, jun 2000.
  {\newblock}\href{Https://doi.org/10.1063/1.533322}{}10.1063/1.533322.{\newblock}
  
  \bibitem[164]{streater_pct_2000}R.~F.~Streater  and  A.~S.~Wightman.
  {\newblock}\tmtextit{PCT, spin and statistics, and all that}.
  {\newblock}Princeton Landmarks in Physics. Princeton University Press,
  Princeton, NJ, 2000.{\newblock}
  
  \bibitem[165]{strocchi_introduction_2013}Franco Strocchi.
  {\newblock}\tmtextit{An Introduction to Non-Perturbative Foundations of
  Quantum Field Theory}. {\newblock}OUP Oxford, Oxford, feb 2013.{\newblock}
  
  \bibitem[166]{summers2012perspective}Stephen~J Summers. {\newblock}A
  perspective on constructive quantum field theory. {\newblock}\tmtextit{ArXiv
  preprint arXiv:1203.3991}, 2012.{\newblock}
  
  \bibitem[167]{symanzik}K.~Symanzik. {\newblock}Euclidean quantum field
  theory. {\newblock}In  R.~Jost, editor, \tmtextit{Local quantum field theory
  (Varenna lectures)},  pages  152--226. Academic Press, New York,
  1969.{\newblock}
  
  \bibitem[168]{Treves}Fran{\c c}ois Tr{\`e}ves.
  {\newblock}\tmtextit{Topological vector spaces, distributions and kernels}.
  {\newblock}Academic Press, New York-London, 1967.{\newblock}
  
  \bibitem[169]{TriebelIII}Hans Triebel. {\newblock}\tmtextit{Theory of
  function spaces. III},  volume  100  of \tmtextit{Monographs in
  Mathematics}. {\newblock}Birkh{\"a}user Verlag, Basel, 2006.{\newblock}
  
  \bibitem[170]{Van2003}Joris van der~Hoeven. {\newblock}\tmtextit{Majorants
  for formal power series}. {\newblock}Technical Report 2003-15.
  Universit{\'e} de Paris-Sud. D{\'e}partement de Math{\'e}matique, Orsay,
  2003.{\newblock}
  
  \bibitem[171]{van2003algebra1}B.L.~van der~Waerden.
  {\newblock}\tmtextit{Algebra: Volume I}. {\newblock}Algebra : based in part
  on lectures by E. Artin and E. Noether. Springer New York, 2003.{\newblock}
  
  \bibitem[172]{van2003algebra2}B.L.~van der~Waerden.
  {\newblock}\tmtextit{Algebra: Volume II}. {\newblock}Algebra : based in part
  on lectures by E. Artin and E. Noether. Springer New York, 2003.{\newblock}
  
  \bibitem[173]{van_nieuwenhuizen_euclidean_1996}Peter van~Nieuwenhuizen  and 
  Andrew Waldron. {\newblock}On Euclidean spinors and Wick rotations.
  {\newblock}\tmtextit{Physics Letters B}, 389(1):29--36, dec 1996.
  {\newblock}\href{Https://doi.org/10.1016/S0370-2693(96)01251-8}{}10.1016/S0370-2693(96)01251-8.{\newblock}
  
  \bibitem[174]{VonNeumann}J.~von~Neumann. {\newblock}On rings of operators.
  III. {\newblock}\tmtextit{Ann. of Math. (2)}, 41:94--161, 1940.
  {\newblock}\href{Https://doi.org/10.2307/1968823}{}10.2307/1968823.{\newblock}
  
  \bibitem[175]{von_waldenfels_non_commutative_1986}Wilhelm von~Waldenfels.
  {\newblock}Non-commutative algebraic central limit theorems. {\newblock}In 
  Herbert Heyer, editor, \tmtextit{Probability Measures on Groups VIII},
  Lecture Notes in Mathematics,  pages  174--202. Berlin, Heidelberg, 1986.
  Springer.
  {\newblock}\href{Https://doi.org/10.1007/BFb0077184}{}10.1007/BFb0077184.{\newblock}
  
  \bibitem[176]{williams_euclidean_1974}David~N.~Williams.
  {\newblock}Euclidean Fermi fields with a hermitean Feynman-Kac-Nelson
  formula. I. {\newblock}\tmtextit{Communications in Mathematical Physics},
  38(1):65--80, mar 1974.
  {\newblock}\href{Https://doi.org/10.1007/BF01651549}{}10.1007/BF01651549.{\newblock}
  
  \bibitem[177]{xiong_sobolev_2018}Xiao Xiong, Quanhua Xu, and  Zhi Yin.
  {\newblock}\tmtextit{Sobolev, Besov and Triebel-Lizorkin Spaces on Quantum
  Tori},  volume  252  of \tmtextit{Memoirs of the American Mathematical
  Society}. {\newblock}American Mathematical Society, mar 2018.
  {\newblock}\href{Https://doi.org/10.1090/memo/1203}{}10.1090/memo/1203.{\newblock}
  
  \bibitem[178]{she-sheng_stochastic_1985}She-Sheng Xue  and  Ting-chang
  Hsien. {\newblock}Stochastic quantization of fermions on lattice.
  {\newblock}\tmtextit{Chinese Physics Letters}, 2(10):474--476, oct 1985.
  {\newblock}\href{Https://doi.org/10.1088/0256-307X/2/10/012}{}10.1088/0256-307X/2/10/012.{\newblock}
  
  \bibitem[179]{MR0448079}Wies{\l}aw {\.Z}elazko. {\newblock}\tmtextit{Banach
  algebras}. {\newblock}Elsevier Publishing Co., Amsterdam-London-New York;
  PWN--Polish Scientific Publishers, Warsaw, 1973. {\newblock}Translated from
  the Polish by Marcin E. Kuczma.{\newblock}
  
  \bibitem[180]{ZinnJustin}J.~Zinn-Justin. {\newblock}\tmtextit{Quantum field
  theory and critical phenomena},  volume~85  of \tmtextit{International
  Series of Monographs on Physics}. {\newblock}The Clarendon Press, Oxford
  University Press, New York, Second  edition, 1993. {\newblock}Oxford Science
  Publications.{\newblock}
  
  \bibitem[181]{zinoviev_equivalence_1995}Yury~M.~Zinoviev.
  {\newblock}Equivalence of Euclidean and Wightman field theories.
  {\newblock}\tmtextit{Communications in Mathematical Physics}, 174(1):1--27,
  1995. {\newblock}Publisher: Springer-Verlag.{\newblock}
\end{thebibliography}

\end{document}
